\documentclass[aps,pra,onecolumn,superscriptaddress,nofootinbib,10pt]{revtex4-2}
\usepackage{amsmath,amssymb,amsthm,mathtools}
\usepackage{booktabs}
\usepackage{enumitem}
\usepackage[colorlinks=true,linkcolor=blue,citecolor=blue,urlcolor=blue]{hyperref}

\newtheorem{theorem}{Theorem}
\newtheorem{proposition}[theorem]{Proposition}
\newtheorem{lemma}[theorem]{Lemma}
\newtheorem{corollary}[theorem]{Corollary}
\theoremstyle{definition}
\newtheorem{definition}[theorem]{Definition}
\newtheorem{postulate}{Postulate}
\renewcommand{\thepostulate}{C\arabic{postulate}}
\newtheorem{problem}{Open problem}
\newtheorem{conjecture}[theorem]{Conjecture}
\theoremstyle{remark}
\newtheorem{remark}[theorem]{Remark}
\newtheorem{example}[theorem]{Example}

\newcommand{\R}{\mathbb{R}}
\newcommand{\C}{\mathbb{C}}
\newcommand{\HH}{\mathbb{H}}
\newcommand{\Z}{\mathbb{Z}}
\newcommand{\g}{\mathfrak{g}}
\newcommand{\su}{\mathfrak{su}}
\newcommand{\cL}{\mathcal{L}}
\newcommand{\cR}{\mathcal{R}}
\newcommand{\cD}{\mathcal{D}}
\newcommand{\Aut}{\operatorname{Aut}}
\newcommand{\Ad}{\operatorname{Ad}}
\newcommand{\ad}{\operatorname{ad}}
\newcommand{\rk}{\operatorname{rank}}
\newcommand{\re}{\operatorname{Re}}
\newcommand{\im}{\operatorname{Im}}
\newcommand{\tr}{\operatorname{Tr}}
\newcommand{\Var}{\operatorname{Var}}
\newcommand{\Cov}{\operatorname{Cov}}
\newcommand{\Herm}{\operatorname{Herm}}
\newcommand{\ip}[2]{\langle #1,#2\rangle}
\newcommand{\hp}[3]{[#1,#2,#3]}
\newcommand{\SU}{\mathrm{SU}}
\newcommand{\SO}{\mathrm{SO}}
\newcommand{\OO}{\mathrm{O}}
\newcommand{\SL}{\mathrm{SL}}
\newcommand{\U}{\mathrm{U}}
\newcommand{\PI}{\cite{PaperI}}
\newcommand{\st}[1]{\textsc{#1}}
\newcommand{\paperpart}[1]{\par\bigskip\begin{center}\large\bfseries #1\end{center}\par\medskip}

\begin{document}

\title{Axiom $\Omega$ and the origin of quantum mechanics:\\ an origin-free comparison, its clock, and its space of events}
% AUTHORS: to be supplied by the owner
\author{[byline to be confirmed]}
\date{19 November 2026}

\begin{abstract}
We study a connected manifold $M$ with a ternary transport $\hp{x}{y}{z}$ that has no distinguished
origin (a heap) and a continuous transport-invariant comparison $\kappa\colon M\times M\to[0,1]$ that is symmetric,
faithful ($\kappa(x,y)=1$ only for $x=y$) and has a unique antipode (for every $x$ exactly one $x^\perp$ with
$\kappa(x,x^\perp)=0$), tied together by the round-trip law $\kappa(w,\hp{x}{w}{x})=(2\kappa(w,x)-1)^2$.  Given the other axioms,
this law is equivalent to two-readout determinacy: along a uniform transport the comparison of any
reference point with the next point is an affine function of its comparisons with the two previous
points.  By the classification of a companion paper the only noncommutative solution is the
three-sphere of unit quaternions with $\kappa(x,y)=\lVert(x+y)/2\rVert^2$.  Its orthogonal complex
structures commuting with the left transports are the quarter turns of clocks, all equivalent.  We
classify the readings of states: a closed group of automorphisms, normalised by the left transports,
declared unobservable, and a weight that is a function of the supremum of the comparison over that
group.  The real, complex and quaternionic structures of the same $\R^4$ give four-, two- and one-level
systems.  If the elementary alternative is binary at the level of states, the reading is through a clock;
equivalently, the unobservable group contains a clock (the phase of a closed system's own clock is
unobservable) and some pair of classes is orthogonal.  In either formulation, if the weight is a function of the class correlation and multiplies on
independent systems, normalisation on complete tests makes it the square of the class correlation, the
Born weight; for one system normalisation alone does not.  Symmetries are then unitary or antiunitary with
Kramers' sign, and for one elementary system Schr\"odinger evolution is the basepoint-free recurrence
$x_{n+1}=\hp{x_n}{x_{n-1}}{x_n}$.  The reading in which points are states (a four-dimensional ball
theory) is consistent for one system; it does not compose if the composite weight vanishes only at
antipodes.  All inputs are listed in one place.  For composites of clock-read factors we assume that product
configurations are given by a unit-valued multilinear map, equivariant for the clocks and the left
transports of the factors, whose values span the composite, that a composite hides nothing beyond the
clocks of its parts, closure of states under round-trip reflections, a weight that is symmetric, as the
comparison is, equals one on each class against itself, vanishes exactly where the class correlation
does and is normalised over all maximal orthogonal families, that each factor keeps its orientation, and that
local automorphisms and pair exchanges are operations.  We prove that a composite then has no
superselection, hence a single clock, and is the complex tensor product: the real tensor product with
two independent clock phases, complex quantum theory with two superselection sectors, multiplies all
class correlations but has maximal orthogonal families of sizes $8$ and $4$, on which no symmetric
weight is normalised; reflection closure is essential to this.  We further prove that probabilities of independent tests multiply
(on a second route that assumes the single clock and drops bilinearity, this multiplicativity gives the
tensor product up to four types without continuity assumptions; the orientation clause selects the
complex one, and for three or more systems reflection closure and spanning are also used), that
entangled states follow, that
the Born weights are the only nonnegative functions of the transition probability that are normalised,
certain on the state's own ray and consistent with adjoining ancillas, and that ancillas give all effects, Kraus instruments and conditional states; complex Hilbert spaces
of every finite dimension occur as round-trip-closed, clock-invariant sub-structures of these registers,
and are systems under an explicit restriction of preparations, tests and operations.  No
outcome rule covariant at ties exists; for a stochastic record process with one terminal record, the Born rule
at every stage is equivalent to the weights being martingales, and a hidden-threshold or simplex
selection is Born exactly when its auxiliary law is the invariant one; which dynamics is realised is
not decided.  For finite memories, a Gram window weakly narrows the corridor of the next comparison, a
memory reset costs at least $k_BT$ times the drop in conditional entropy while index restriction can cost
nothing, an encoded phase is exactly recoverable iff the inaccessible conditional states agree, and
local windows of several observers glue iff a positive semidefinite completion exists (on chordal
patterns iff every clique is positive), which still gives neither a joint state nor joint outcomes.  In a closed register of $2n$ systems in a global singlet, an
exchange constraint gives conditional Schr\"odinger dynamics relative to an internal clock, with an
explicit source from exchanges across the clock cut, and the unobservable clock [Cl] follows from the
invisibility of the total phase under a stated bridge hypothesis.  The dichotomies of configurations form four-dimensional
de~Sitter space with inclusion as causal order, a classical fact; every invariant metric on it solves
the vacuum Einstein equations with $\Lambda\ell^2=3$.  That physical events are dichotomies is a hypothesis; under
it, observers correspond to binary alternatives whose static patches are the dichotomies separating them.  With an ungauged
back-reacting texture source the invariant metric is not a solution, and no maximal solution carries the
full order of the dichotomies.  Under labelled hypotheses, fields of
structures give $\U(2)$ with $m_W=m_Z\cos\theta_W$ at tree level, in an electroweak sector that contains no
scalar and is contradicted by the observed $125$\,GeV scalar; with an optional variable radius of the
comparison sphere it becomes the custodial Standard Model scalar sector, compatible at tree level with
the measured couplings, with the scalar mass fitted and the potential not determined;
with an optional round four-sphere target from a quaternionic structure on a pair it becomes the minimal
composite-Higgs geometry, with $\kappa_V=\sqrt{1-\xi}$ and $\kappa_{2V}=1-2\xi$, $\xi=v^2/f^2$, which the quoted modifiers bound by
$\xi\le0.039$ in a separate-Gaussian treatment;
with an optional Skyrme term the field carries localised textures whose energy is bounded by
$6\pi^2v/e_{\mathrm S}$ times the degree and which admit, without forcing, spin $\tfrac12$, and a closed pair with a third structure gives the
Pati--Salam group broken to the Standard Model group with the hypercharges of one generation;
$\sin^2\theta_W=3/8$ requires equal couplings, which one-loop running does not realise.  Fourteen
physical examples are tagged.  Exact instances of theorems: Ramsey fringes, stroboscopic Larmor
precession, Grover search, isotropic rotational diffusion of a spin, the geometric phase, Kramers
degeneracy of one spin, partial polarisers, the complementarity identity, and the state-independent
fringe $\kappa(\psi,U\psi)=\cos^2(\alpha/4)$ of neutron $4\pi$ interferometry.  Instances under named postulates or
hypotheses: the reading of interferometric fringes and phases as $\kappa$ and as holonomy, Malus's law,
Kramers degeneracy of many spins, the realisation of Grover search, exchange-driven $\sqrt{\mathrm{SWAP}}$
gates, real-versus-complex network tests, the reading of complementarity, and the electroweak
relations.  Analogy: $\Lambda>0$.  The kinematic statements reproduce standard quantum mechanics, and none
of the examples distinguishes the inputs of this paper from it; the $3:8$ ratio of spinor contrast to
polarisation decay under isotropic noise is a standard consequence of spinor rotational diffusion that
has not been compared directly.  We compare with a bottom-up reconstruction that
meets this one at the Hopf fibration.
\end{abstract}

\maketitle

%=====================================================================================
\section{Introduction}\label{sec:intro}

\subsection{The physical questions}
Textbook quantum mechanics postulates a complex Hilbert space, the transition probability
$\lvert\ip{\psi}{\phi}\rvert^2$, unitary or antiunitary symmetries, the tensor product for composites
and a Hamiltonian evolution.  Special relativity adds a limiting velocity and a Lorentzian manifold
of events, cosmology a positive cosmological constant, and the electroweak theory the gauge group
$\U(2)$ with the relation $m_W=m_Z\cos\theta_W$.  Each of these is an input.  This paper asks which of
them follow from one structural axiom, and states for each item whether it is a theorem, a theorem
under a stated postulate or hypothesis, or open.  The questions are:
\begin{enumerate}[label=(Q\arabic*),leftmargin=3em]
\item Why a complex Hilbert space, rather than a real or quaternionic one?
\item Why is the probability of a transition $\lvert\ip{\psi}{\phi}\rvert^2$?
\item Why are symmetries unitary or antiunitary, and why does time reversal of a two-level system
square to $-1$?
\item Where does entanglement come from?
\item Why is there a limiting velocity, and why does the Lorentz group act on states?
\item Why is the manifold of events four-dimensional and Lorentzian, with $\Lambda>0$?
\item What fixes the electroweak gauge group and the relation between the vector masses?
\end{enumerate}

\subsection{The axiom}
The axiom, called \textsc{Axiom}-$\Omega$ here and in the companion paper \PI, has
two primitives: a transport and a comparison.  The transport $\hp{x}{y}{z}$ reads ``apply to $z$ the
displacement that carries $y$ to $x$''; no point plays the role of an origin.  The comparison
$\kappa(x,y)\in[0,1]$ is invariant under transport and equals $1$ only on the diagonal.  The single law
connecting the two is the round-trip law (Definition~\ref{def:cs}).  It has an equivalent form
(Theorem~\ref{thm:trd}): for every reference point $y$, the comparison of $y$ with the reflection
$s_x(w)=\hp{x}{w}{x}$ is an affine function of its comparisons with $x$ and with $w$, with coefficients
that depend only on $(w,x)$; along a uniform transport $x_{n+1}=\hp{x_n}{x_{n-1}}{x_n}$ the comparison of
every reference point with each step is fixed by its comparisons with the two previous steps.  We call
this \emph{two-readout determinacy}.  When a reading of $\kappa$ is needed we use the one that requires
no further structure: $\kappa(x,y)$ is the acceptance probability of an interference (Hadamard) test
between two branches carrying $x$ and $y$ with a shared phase reference
(Proposition~\ref{prop:linear}).  The
classification of solutions is proved in the companion paper \PI; Section~\ref{sec:axiom} states it
and sketches the proof.  Everything after Section~\ref{sec:axiom} concerns the noncommutative solution
$M\cong S^3$.

\subsection{Two routes to the qubit}
An unpublished manuscript \cite{CRQO} reconstructs the qubit from below: from a positive quadratic
self-relation and a phase circle it obtains the Bloch ball and its rank-one boundary $S^2$; an
oriented three-point clause then lifts $S^2$ to $S^3$ through a $\U(1)$ fibre of minimal Chern class,
and further clauses (positive kernel, reversible continuation, complete readout, composition) yield
the Hilbert-space formalism.  The route of the present paper runs from above: the axiom yields $S^3$
with its comparison first; the choice of a clock, a circle subgroup of transports, yields a complex
structure; the quotient by the clock is $S^2$ with the Born kernel.  Both routes pass through the
Hopf fibration $S^3\to S^2$, in opposite directions.  Section~\ref{sec:crqo} makes the comparison
exact clause by clause.

\subsection{Assumptions}\label{sec:assumptions}
Every input of the paper is listed here, once.  Each later statement names the inputs it uses.
\begin{itemize}[leftmargin=1.6em]
\item[$(\Omega)$] Axiom $\Omega$: $(\Omega1)$, $(\Omega2\mathrm{a{-}d})$, $(\Omega3)$ (Definition~\ref{def:cs}).
\item[{[NC]}] The noncommutative solution $M\cong S^3$ is the one studied (the commutative one is the
real qubit, Proposition~\ref{prop:rebit}).
\item[{[L]}] A reading of states declares a closed group $K$ of automorphisms unobservable, and the left
transports remain symmetries (they normalise $K$) (Definition~\ref{def:reading}).
\item[{[W]}] \emph{Weight form.}  The weight of a reading is $\mathsf F\circ s_K$, $s_K=\sup_Kc$, with $\mathsf F$ continuous and
$K\subseteq\Aut(M)$; for a clock this is automatic for continuous $\cL$-invariant weights.
\item[{[B]}] \emph{Binary states.}  Every class of states has exactly one orthogonal class: the form of
$(\Omega2\mathrm d)$ at the level of states.  Orthogonality is $s_K=0$ here and in [NT]; for a weight $\mathsf F\circ s_K$
it agrees with $w=0$ iff $\mathsf F^{-1}(0)=\{0\}$.  Under [L] it is equivalent to [Cl] with [NT] (Proposition~\ref{prop:clockB}).
\item[{[Cl]}] \emph{Unobservable clock.}  The unobservable group of an elementary system contains a clock:
for a closed system the phase of its own clock is not observable.  An input with this physical reading,
not derived from the axiom; it is a theorem under [Ph] (Corollary~\ref{cor:clbridge}), which is itself an input.
\item[{[NT]}] \emph{Non-triviality.}  Some pair of classes is orthogonal.
\item[{[T]}] Convention~\ref{conv:T}: a complete test is a maximal family of pairwise orthogonal classes,
weights are numbers in $[0,1]$ read as probabilities, and weights of disjoint blocks add.
\item[{[N]}] \emph{Normalisation.}  For every maximal family of pairwise orthogonal classes, the weights of
any class against the family sum to one (for one system in Theorem~\ref{thm:readings}, for a composite in
Theorem~\ref{thm:oneclock}).
\item[{[Comp]}] \emph{Composability of readings.}  Class-correlations and weights of independent systems
multiply (Definition~\ref{def:reading}).
\item[] That states are clock orbits is then a theorem under [L], [B] (Theorem~\ref{thm:readings}), equivalently
under [L], [Cl], [NT] (Proposition~\ref{prop:clockB}); that the weight is $\beta$ needs [W], [Comp] and [N] in
addition.
\item[{[Bil]}] Product configurations of $N$ systems are a unit-valued real-multilinear map into a real
inner-product space $\mathcal W$, equivariant for the clock torus; [Bil$_\cL$]: also for the left transports
(Section~\ref{sec:oneclock}).
\item[{[Span]}] $\mathcal W$, and $V_N$, are by definition the span of the product configurations.
\item[{[Loc]}] The unobservable group of a composite is the group generated by the clocks of its factors.
\item[{[O]}] The weight of a composite is symmetric, as $\kappa$ is by $(\Omega2\mathrm a)$, equals one on the diagonal,
and vanishes exactly where $s_K=0$; $\mathsf F\circ s_K$ inherited from the factors is a special case.
\item[{[C2]}] The states of a composite contain the products and are closed under round-trip reflections
(Postulate~\ref{post:C2}).
\item[{[C1]}] (b) Each factor keeps its orientation; (a) probabilities of independent ray tests multiply
(Postulate~\ref{post:C1}).  For clock readings (a) is [Comp].
\item[{[C0]}] \emph{One clock.}  A composite is read in the same way as its factors, through one clock, as the
complex tensor product.  For clock-read factors this is a theorem: one clock under [Bil], [Loc], [O],
[C2] and [N], the complex tensor product and C1(a) with [Bil$_\cL$] and [Span] in addition
(Theorem~\ref{thm:oneclock}, Corollary~\ref{post:C0}); C1(b) selects the orientation.  C2 is essential
(Proposition~\ref{prop:productsalone}).  The content is: if every class of unit vectors reachable by
round-trip reflections is a state, with a symmetric transition probability, and nothing is hidden beyond
the local clocks, there is no superselection (Proposition~\ref{prop:superselection}) and hence one clock.  On a
second route C0 and C1(a) are assumed instead of [Bil], [Bil$_\cL$], [Loc], [O] and [N] for the composite, and
Theorem~\ref{thm:product} gives the tensor product without linearity or continuity (with C2 for $N\ge3$,
Corollary~\ref{cor:ampmult}).  ``C0--C2'' and ``C0--C3'' denote C0 with
Postulates C1--C2 or C1--C3.
\item[{[C3]}] Local clock-preserving automorphisms and the pair exchanges are operations
(Postulate~\ref{post:C3}).
\item[{[Sub]}] A clock-invariant round-trip-closed sub-structure $S(U)$ of a register is treated as a system:
only its configurations are prepared, complete tests are maximal orthogonal families within it, and
operations are those of the host that preserve $U$ (Section~\ref{sec:otherdim}).
\item[{[E]}] Physical events are the dichotomies of an elementary structure (Hypothesis~E).
\item[{[F]}] Fields of structures with a clock field, clock-preserving connections and an
$\Aut(M)$-invariant kinetic form (Hypothesis~F, clauses (i)--(iv)).
\item[{[F$_{\mathrm g}$]}] (a) A dynamical connection with a positive Yang--Mills term; (b) the cap geometry with the
centre field as its flat covariantly constant vacuum (Hypothesis~F$_{\mathrm g}$, Section~\ref{sec:centre}).
\item[{[F$_{\mathrm r}$]}] The configuration of F takes values in $V\setminus\{0\}$ (or $V$), with the unit transports and the
canonical renormalisable action (Hypothesis~F$_{\mathrm r}$, Section~\ref{sec:radial}).
\item[{[F$_{\mathrm H}$]}] The scalar target is the round $\HH P^1\cong S^4=\SO(5)/\SO(4)$ of radius $f$ from a quaternionic structure on a
pair, gauged in the $\SO(4)$ fixing a pole, truncated at two derivatives, with a broken vacuum; [F$_{\mathrm{Hf}}$]: the
fermion embeddings (Section~\ref{sec:hopfbase}).
\item[{[F$_{\mathrm S}$]}] A Skyrme (area-square) term with coefficient $1/(32e_{\mathrm S}^2)$; [F$_{\mathrm{Sg}}$]: its covariant form with a
dynamical connection; [F$_{\mathrm{Sr}}$]: its polynomial extension with the radial mode (Section~\ref{sec:skyrme}).
\item[{[P]}] A localised degree texture is read as a baryon-like, a clock-quotient configuration as a
lepton-like excitation (Hypothesis~P, Section~\ref{sec:skyrme}).
\item[{[Ph]}] \emph{Phase bridge.}  Composite statistics of elementary configurations are those of the complex
composite reading built on one clock's complex structure, and a right transport leaving all internal
statistics unchanged is unobservable (Section~\ref{sec:closedreg}).
\item[{[U], [A]}] The physical specification of a closed-universe state, and of an observer's accessible
algebra (Hypotheses~U and~A, Section~\ref{sec:closedreg}).
\item[{[Gw]}] A finite Gram window of stored comparisons; [Ue]: the entropy drop of a reset equals the log
area ratio of the cap corridors; [Ca]: causal record access in tensor-factor form (Section~\ref{sec:finitemem}).
\item[{[R]}] A record process: one terminal vertex and absorbing record faces in the observer's filtration
(Definition~\ref{def:soprocess}).  [Tr]: a linear number of terminal transcripts in the tick count.  [Ap]: a
preparation-independent, unitary-invariant apparatus sampling law (Section~\ref{sec:outcomes}).
\item[{[G]}] The group, multiplets, breaking configuration and equal couplings of
Section~\ref{sec:sm} (Hypothesis~G, clauses (i)--(v)).
\item[{[GR]}] Einstein's equations with a source (Hypothesis~GR).
\end{itemize}

\subsection{Answers and their status}
Write $c:=2\kappa-1$ (the \emph{correlation}) and $s_x(w):=\hp{x}{w}{x}$ (the \emph{reflection of $w$ in
$x$}).  Table~\ref{tab:answers} gives the answers with the assumptions they use;
Table~\ref{tab:scope} lists every statement.

\begin{table}[h]
\caption{Answers to (Q1)--(Q7), the assumptions they use, and the physical examples of Section~\ref{sec:examples}
with their tags: (a) exact instance of a theorem, (b) instance under a named postulate or hypothesis,
(c) analogy; (a,b): the identity is (a), its physical reading (b).  The first row collects the examples
that concern the comparison and the axiom itself.}
\label{tab:answers}
\footnotesize
\begin{tabular}{@{}p{0.04\textwidth}p{0.5\textwidth}p{0.24\textwidth}p{0.16\textwidth}@{}}
\toprule
& Answer & Assumptions & Examples\\
\midrule
$\kappa$ & The comparison, the round-trip law and two-readout determinacy; the clock holonomy; Brownian motion on $S^3$ (Sections~\ref{sec:axiom}, \ref{sec:one}, \PI). & $(\Omega)$, [NC]; readings: C0, C1, [Span], [T] & \ref{ex:neutron}(a,b), \ref{ex:ramsey}(a), \ref{ex:larmor}(a), \ref{ex:diffusion}(a), \ref{ex:pancharatnam}(a,b)\\
Q1 & The orthogonal complex structures commuting with $\cL$ are the quarter turns of clocks, all conjugate (theorem).  Among all readings of states (closed $K$ normalised by $\cL$, weight a function of the class-correlation), the binary ones are exactly the clock readings, equivalently those whose $K$ contains a clock and has orthogonal classes (Proposition~\ref{prop:clockB}); composable ones have $w=s_T^p$; $\R$, $\C$, $\HH$ on the same $\R^4$ give $4$, $2$, $1$ distinguishable states (Theorem~\ref{thm:readings}).  That a composite has one clock is Theorem~\ref{thm:oneclock}; every finite dimension occurs as a sub-structure of a register, a system under [Sub] (Section~\ref{sec:otherdim}). & $(\Omega)$, [NC], [L], [B], or [L], [Cl], [NT]; composite of clock-read factors: [Bil], [Span], [Loc], [O], [C2], [N]; other dimensions: [Sub] & \ref{ex:real}(b)\\
Q2 & The weight of a composable normalised clock reading is $s_T^2=\beta$: the square follows from multiplicativity and normalisation (Theorem~\ref{thm:readings}); it is the clock-invariant second moment of $\kappa$; for composites it is the only nonnegative function of $\beta_N$ with the properties of Theorem~\ref{thm:bornN} (Gleason). & [L], [W], [Comp], [B], [N]; probability: [T]; composites: [C0], [C1], [Span], [C2] and (a)--(c) of Theorem~\ref{thm:bornN}(iii), including marginal consistency (c) & \ref{ex:malus}(b)\\
Q3 & Clock-preserving symmetries are unitary, clock-reversing ones antiunitary; clock reversals commuting with $\cL$ square to $-1$ (one system). & $(\Omega)$, [NC], with the clock reading & \ref{ex:kramers}(a,b)\\
Q4 & Entangled states follow from reflection closure in the tensor product; effects, Kraus instruments and conditional states follow by dilation. & [C0], [C1], [Span], [C2], [C3], [T] & \ref{ex:exchange}(b), \ref{ex:grover}(a,b), \ref{ex:pvc}(a,b)\\
U & Closed register of $2n$ systems: singlet dimension $\frac1{n+1}\binom{2n}{n}$, which marginals can be flat; conditional Schr\"odinger equation with a cross-cut source relative to an internal clock (exact eight-site example); [Cl] from the invisible total phase under [Ph]; finite modular flow equals exchange time iff Gibbs; local equilibration in the singlet sector: conjecture. & [C0]--[C3], [Span], [T]; [Ph]; [U], [A] & ---\\
Q5 & Not answered: only the geometry of the state space (hyperbolic interior, filters, $\gamma^{-1}=$ concurrence). & --- & \ref{ex:boosts}(a), \ref{ex:pvc}(a,b)\\
Q6 & The dichotomies of $M$ form $dS_4$ with inclusion as causal order (classical); every invariant metric solves the vacuum equations; that events are dichotomies is [E]; with an ungauged back-reacting centre field no maximal solution carries the full cap order (Props.~\ref{prop:nullobstruction}, \ref{prop:friedmann}); under [F$_{\mathrm g}$] the centre field is pure gauge. & $(\Omega)$, [NC]; [E]; [F], [GR], [F$_{\mathrm g}$] & \ref{ex:desitter}(c)\\
H & Under F$_{\mathrm H}$: one scalar, $m_W=gv/2$, $\rho=1$, $\kappa_V=\sqrt{1-\xi}$, $\kappa_{2V}=1-2\xi$ ($\xi=v^2/f^2$); bound $\xi\le0.039$, $f\ge1.24$\,TeV (separate Gaussians); scalar mass not fixed by the geometry; the pair does not force this target. & [F], [F$_{\mathrm g}$](a), [F$_{\mathrm H}$]; Yukawas: [F$_{\mathrm{Hf}}$] & ---\\
P & Under F$_{\mathrm S}$: degree-$\mathcal Q$ textures with $E\ge6\pi^2(v/e_{\mathrm S})\lvert\mathcal Q\rvert$, minimisers in degree $\pm1$ (cited), Hopf charge with the Vakulenko--Kapitanski\u{\i} bound; spin $\tfrac12$ allowed, not forced; a dynamical connection or radius removes the protection; particle readings are a hypothesis. & [F], [F$_{\mathrm S}$]; [F$_{\mathrm{Sg}}$], [F$_{\mathrm{Sr}}$]; [P] & ---\\
Q7 & $\U(2)$ and the ratio $g^2:(g^2+g'^2)$ of the mass-form eigenvalues; with [F$_{\mathrm g}$](a), $m_W=m_Z\cos\theta_W$ as tree-level masses; under [F], [F$_{\mathrm g}$](a) alone this sector contains no scalar and is contradicted by the observed one; with [F$_{\mathrm r}$] it is the custodial Standard Model scalar sector, with the potential free (Section~\ref{sec:radial}); the Pati--Salam group, its breaking to $(\SU(3)\times\SU(2)\times\U(1))/\Z_6$ and the hypercharges of one generation; $\sin^2\theta_W=\tfrac38$ under equal couplings, which one-loop running with Standard Model content does not realise at any scale. & [F]; masses: [F$_{\mathrm g}$](a); scalar: [F$_{\mathrm r}$]; [C0], [C1], [Span], [C2], [C3], [G] & \ref{ex:ew}(b)\\
\bottomrule
\end{tabular}
\end{table}

\subsection{What is not claimed}
The Hamiltonian, the coupling constants and the direction of any external field are free.  A
mechanism for single outcomes is not proposed: Section~\ref{sec:outcomes} characterises the statistics of any
single-outcome record process (Theorem~\ref{thm:sorigidity}) and does not select a dynamics.  The probabilistic reading of weights is [T].  That
states are clock orbits is a theorem under [L], [B], equivalently [L], [Cl], [NT]
(Proposition~\ref{prop:clockB}), and the square needs [W], [Comp] and [N] (Theorem~\ref{thm:readings}); without [B] the
real four-level reading ($K=\{\pm1\}$, weight $c^2$) is also composable.  The single clock of a composite (C0) is derived for clock-read factors under [Bil], [Loc], [O], [C2] and
[N] (Theorem~\ref{thm:oneclock}); [Loc], [O] and C2 are not derived, and C2 is essential.  That physical
events are dichotomies is [E]; dichotomies are threshold comparisons of configurations, not tests in
the sense of [T].  With an ungauged back-reacting centre field no maximal solution carries the full order of the
dichotomies (Propositions~\ref{prop:nullobstruction}, \ref{prop:friedmann}).  The registers of $N$ two-level systems have dimension $2^N$; every finite dimension occurs as a
clock-invariant round-trip-closed sub-structure of a register, and is a system under [Sub]
(Section~\ref{sec:otherdim}).  No elementary comparison structure of another dimension exists: the models of
the axiom are $S^1$ and $S^3$.  The Pati--Salam group, the multiplets, the breaking
configuration and equal couplings are [G]; the number of structures, three generations, the
breaking scale, masses and Yukawa couplings are not derived.

\subsection{Physical examples}
Section~\ref{sec:examples} gives fourteen examples, each tagged (a) exact instance of a theorem, (b)
instance under a named postulate or hypothesis of Section~\ref{sec:assumptions}, or (c) analogy.  Exact
instances (a): Ramsey fringes (Example~\ref{ex:ramsey}), stroboscopic Larmor precession
(Example~\ref{ex:larmor}), Grover search (Example~\ref{ex:grover}), isotropic rotational diffusion
(Example~\ref{ex:diffusion}), the geometric phase (Example~\ref{ex:pancharatnam}), Kramers degeneracy of one spin
(Example~\ref{ex:kramers}), partial polarisers (Example~\ref{ex:boosts}), the complementarity identity
(Example~\ref{ex:pvc}) and the state-independent fringe $\kappa(\psi,U\psi)=\cos^2(\alpha/4)$ of neutron $4\pi$
interferometry (Example~\ref{ex:neutron}).  Instances under postulates or hypotheses (b): the reading of
interferometric fringes and phases as $\kappa$ and as holonomy, Malus's law (Example~\ref{ex:malus}), Kramers
degeneracy of many spins, the realisation of Grover search, exchange-driven $\sqrt{\mathrm{SWAP}}$
(Example~\ref{ex:exchange}), real-versus-complex network tests (Example~\ref{ex:real}), the reading of
complementarity, and the electroweak relations (Example~\ref{ex:ew}).  Analogy (c): $\Lambda>0$
(Example~\ref{ex:desitter}).  None of the examples distinguishes the inputs of this paper from standard
quantum mechanics; Section~\ref{sec:closed} states what a distinguishing consequence would have to come
from.  The ratio $3:8$ of spinor contrast to polarisation decay under isotropic noise is a standard
consequence of spinor rotational diffusion that has not been compared directly.

\subsection{Related work and what is new}
\emph{Reconstructions.}  Operational reconstructions \cite{Hardy,CDP,MM,MMAP,DB} start from convex state
spaces with tomography, purification or reversibility; here the single-system convex structure
(Section~\ref{sec:cone}) follows from the transport, and the complex Hilbert space of a composite of clock-read factors follows from [Bil$_\cL$], [Span],
[Loc], [O], C2 and [N] (Theorem~\ref{thm:oneclock}): a symmetric transition probability on all reachable
unit vectors excludes superselection; no tomography assumption is made.  Masanes, M\"uller, Augusiak and P\'erez-Garc\'{\i}a \cite{MMAP,MMPA14} single out the
three-dimensional Bloch ball among ball state spaces by properties of composites;
Theorem~\ref{thm:readings} is a counterpart in the present setting.  M\"uller and Masanes \cite{MuM} derive the dimension three of
space and two-qubit quantum theory from postulates on observers who exchange direction information by
rotating devices, with continuous reversible interactions; here the Bloch ball is three-dimensional
because $M$ is, and nothing is asserted about space.  Von Weizs\"acker's ur hypothesis \cite{vW}
identifies $\SU(2)\cong S^3$, the symmetry group of a binary alternative, with physical space; we prove no
such identification.
\emph{Geometry of states.}  That a complex structure on a real Hilbert space is a choice of $\U(1)$, and that
the transition probability is determined by the real metric and that choice, is the content of geometric
quantum mechanics \cite{Mielnik68,Kibble79,AshtekarSchilling,BrodyHughston}; Stueckelberg \cite{Stueckelberg}
described complex quantum theory as real quantum theory with an operator $J$ commuting with all
observables.  Propositions~\ref{prop:clock} and Theorem~\ref{thm:born} restate these facts for the
heap $S^3$.  Renou \emph{et al.}\ \cite{Renou} showed that real quantum theory with the real tensor product
predicts network correlations that differ from complex quantum theory.
\emph{Composites.}  Derivations of the tensor product from conditions on product states exist in
lattice-theoretic and operational form \cite{AertsDaubechies,Matolcsi75,CMA21}.
\emph{Events.}  The description of the spheres of $S^n$ by points of de~Sitter space, with nesting as
causal order, is classical in M\"obius geometry \cite{HJ} and reappears as the kinematic space of balls in
holography \cite{CLMS15,dBHMN16,dBHHM16}.
\emph{Fields.}  The custodial symmetry behind $m_W=m_Z\cos\theta_W$ is due to \cite{Veltman77,SSVZ}; the
fixed-length doublet is the gauged nonlinear sigma model \cite{AB80}.  Section~\ref{sec:sm} reproduces the
Pati--Salam embedding \cite{PatiSalam}, the left--right symmetric structure
\cite{MohapatraPati,SenjanovicMohapatra}, the value of \cite{GQW} and the global form of the Standard Model
group inside the Pati--Salam group \cite{BaezHuerta}.  Heaps go back to Pr\"ufer and Baer
\cite{Prufer24,Baer29}.  The Skyrme term, its degree bound and the spinorial quantisation of its
textures are classical \cite{n6:skyrme,n6:faddeev,n6:fr-original,n6:giulini}; Section~\ref{sec:skyrme}
places them in the comparison field, where the clock quotient carries the Hopf charge of the same
configuration and a dynamical connection or radius removes the topological protection.
\emph{What is new}, in the scope of the literature we searched: the axiom and its classification
(Paper~I), including the two-readout form of the round-trip law, proved there in a stronger form
\cite[Thm.~4.5]{PaperI}, of which Theorem~\ref{thm:trd} is a second proof; the product-rule theorem
for maps on products of ray spaces without linearity or continuity (Theorem~\ref{thm:product});
entanglement from closure of states under round-trip reflections (Theorem~\ref{thm:comp}); the
classification of readings over all closed subgroups of $\SU(2)$ acting on the comparison sphere,
by which binarity selects the clock and composability (under [W]) with normalisation gives the square
(Theorem~\ref{thm:readings}); the single clock of a composite as a theorem under [Bil], [Span], [Loc], C2, [N] and [O], whose
symmetry of the weight does the work (Theorem~\ref{thm:oneclock}), mathematically a reformulation of the
absence of superselection in terms of what is unobservable; and the composition no-go for the comparison reading under its stated
hypothesis (Proposition~\ref{prop:nogo}); the description of an
observer's static patch as the dichotomies that separate a binary alternative
(Proposition~\ref{prop:static}); and the identification of the Pati--Salam ingredients with objects of
the structure (Section~\ref{sec:sm}).


\subsection{Roadmap}
Part~I (Sections~\ref{sec:axiom}--\ref{sec:geometry}) is the kinematics: the axiom and its solutions, one
system, many systems, sub-structures of every dimension, and the geometry of states.  Part~II
(Sections~\ref{sec:closedreg}--\ref{sec:outcomes}) concerns observers inside the structure: a closed universe with internal
clocks; finite memory (Gram windows, the arrow of forgetting, erasure and consensus); and single outcomes as
bounded martingales with finite-tick certification; a section on computation is to follow.  Part~III will collect bounds and budgets: speed limits, the Tsirelson bound and an
entanglement budget.  Part~IV (Sections~\ref{sec:eventspace}--\ref{sec:sm}) contains the space of events,
fields of structures, the variable radius, the round four-sphere target from a pair, topological
particles with the Skyrme term, and the
three-structure model; sections on unification scales, gravity and the origin of the three structures,
of three generations and of the coupling ratios are to follow.  Part~V (Sections~\ref{sec:examples}--\ref{sec:lean}) collects the
physical examples, the status of every statement, what could fail, the open problems and the
numerical checks.  Statements that are not proved keep their status as hypothesis, conjecture or
programme.

%=====================================================================================
\paperpart{Part I.  Kinematics}

%=====================================================================================
\section{The axiom and its solutions}\label{sec:axiom}

\begin{definition}[Axiom $\Omega$]\label{def:cs}
A \emph{comparison structure} is a triple $(M,\hp{\cdot}{\cdot}{\cdot},\kappa)$ satisfying:
\begin{enumerate}[label=$(\Omega\arabic*)$,leftmargin=3.2em]
\item \emph{Origin-free transport.}  $M$ is a connected topological manifold of finite dimension
$n\ge1$ and $\hp{\cdot}{\cdot}{\cdot}\colon M^3\to M$ is continuous with
$\hp{x}{x}{z}=z$, $\hp{x}{z}{z}=x$, $\hp{\hp{x}{y}{z}}{u}{v}=\hp{x}{y}{\hp{z}{u}{v}}$.
\item \emph{Faithful invariant comparison.}  $\kappa\colon M\times M\to[0,1]$ is continuous, and
(a) $\kappa(x,y)=\kappa(y,x)$; (b) $\kappa(x,y)=1$ iff $x=y$;
(c) $\kappa(\hp{x}{u}{v},\hp{y}{u}{v})=\kappa(\hp{u}{v}{x},\hp{u}{v}{y})=\kappa(x,y)$;
(d) for every $x$ there is exactly one $x^\perp$ with $\kappa(x,x^\perp)=0$.
\item \emph{Round-trip law.}  $\kappa\bigl(w,\hp{x}{w}{x}\bigr)=\bigl(2\kappa(w,x)-1\bigr)^2$ for all
$w,x\in M$.
\end{enumerate}
\end{definition}
Law $(\Omega2\mathrm d)$ says that an elementary alternative is binary.  Earlier versions of the axiom
added a clause $(\Omega2\mathrm e)$, that $\kappa$ is $C^2$ near the diagonal; it is redundant
\cite[Thm.~3.11]{PaperI}.  By $(\Omega2\mathrm c)$,
$\kappa(x,s_xw)=\kappa(w,x)$, so $(\Omega3)$ says $\kappa(w,s_xw)=c(w,x)\,c(x,s_xw)$: the comparison of
a point with its mirror image factorises through the mirror.

\begin{example}\label{ex:models}
Let $A=\C$ or $\HH$, $M=\{x\in A:\lvert x\rvert=1\}$, $\hp{x}{y}{z}=x\bar yz$ and
$\kappa(x,y)=\tfrac12(1+\re\bar xy)=\lVert\tfrac12(x+y)\rVert^2$.  Then $x^\perp=-x$, and since
$x\bar wx=2\ip{x}{w}x-w$ for the Euclidean inner product, $\kappa(w,s_xw)=\ip{x}{w}^2=c(w,x)^2$.
\end{example}

\begin{theorem}[classification {\cite[Thms.~3.1, 3.11]{PaperI}}]\label{thm:main}
Every comparison structure is isomorphic to exactly one of the two structures of
Example~\ref{ex:models}.  In particular $\dim M\in\{1,3\}$, $M$ is compact, $\kappa$ is positive
semidefinite of rank $\dim M+2$, and the transport is noncommutative iff $M\cong S^3$.
\end{theorem}

\begin{proof}[Sketch, following {\cite[Thm.~3.11]{PaperI}}]
Fix $o\in M$.  By Certaine \cite{Certaine}, $x\cdot y:=\hp{x}{o}{y}$ is a group law with unit $o$ and
$\hp{x}{y}{z}=xy^{-1}z$; changing $o$ gives an isomorphic group.  A locally Euclidean topological group
is a Lie group (Gleason--Montgomery--Zippin; see \cite[Lemma~2.2]{PaperI}).  Invariance under both transports makes
$\psi:=2\kappa(o,\cdot)-1$ a continuous symmetric class function with $\psi(g)=1$ only at $g=o$, and with
$g=w^{-1}x$, $(\Omega3)$ reads $\psi(g^2)=2\psi(g)^2-1$.  (1) On a one-parameter subgroup, a continuous
solution of the duplication formula with $\psi<1$ off $0$ cannot live on a line, so every one-parameter
subgroup closes into a circle.  (2) On a circle such a solution is rigid: $\psi=\cos$ of the angle.
(3) An abelian subalgebra of dimension $\ge2$ would have a kernel of $\exp$ meeting every line through
$0$; that kernel is a discrete subgroup, hence countable, and meets only countably many of the
uncountably many lines; so abelian subalgebras are lines.  (4) Every $\ad X$ is then
diagonalisable with imaginary eigenvalues, and a Lie algebra with these two properties is $\R$ or
$\su(2)$, so $G\in\{S^1,\SU(2),\SO(3)\}$.  For $\SO(3)$, $\psi=\cos\alpha$ and $\kappa=\cos^2(\alpha/2)$ vanishes on a
real projective plane, contradicting $(\Omega2\mathrm d)$; for $S^1$ and $\SU(2)$, $\psi=\re$, the comparison of
Example~\ref{ex:models}.  Positive semidefiniteness follows from $\kappa(x,y)=\ip{\Phi(x)}{\Phi(y)}$ with
$\Phi(x)=(1,x)/\sqrt2$; no regularity beyond continuity is used \cite[Thms.~3.1, 3.11]{PaperI}.  For $C^2$
pre-comparison structures with nondegenerate second-order form $h$ (and, for the first, $\kappa$ positive
semidefinite with $\inf\kappa=0$), Paper~I also shows that two other closing laws (minimal dimension and
rank; completeness of comparison) select $\su(2)$ \cite[Thms.~4.1, 4.2]{PaperI}.
\end{proof}
The geometric conclusion needs no associativity.  By the roundness theorem of Paper~I
\cite[Thm.~5.18, Cor.~5.19]{PaperI}, a connected manifold with a ternary operation satisfying the Mal'cev
identities $\hp aax=x$, $\hp axx=a$, continuous in its outer arguments, and a comparison with
$(\Omega2\mathrm{a{-}d})$ and $(\Omega3)$, is the round $S^1$, $S^3$ or $S^7$ with $\kappa=\cos^2(d/2)$; no form of the third
identity of $(\Omega1)$ is used.  In particular $c=2\kappa-1$ is positive semidefinite of rank $\dim M+1$, which
gives the linear structure of Proposition~\ref{prop:linear} without the classification.  The third identity
excludes $S^7$, since it makes the base loops groups (Certaine) and no isotope of the octonions is
associative \cite[Rem.~5.25]{PaperI}, and it fixes the operation on $S^1$ and $S^3$.  Invariance under the right
transports cannot be dropped even if invariance under the left transports and one clock is kept: the
example of \cite[Prop.~5.26]{PaperI} has a comparison invariant under the maps $x\mapsto qxe^{i\alpha}$, the left
transports and right multiplication by a clock, that is, under the clock-preserving group $\cL T\cong\U(2)$ of
Proposition~\ref{prop:clock} (and not under right multiplication by $e^{j\alpha}$); it is not round.  Its comparison
is continuous; along every geometric sequence $\epsilon=2^{-k}\epsilon_0$ the ratio $\theta(e^{\epsilon u})/\epsilon$ converges to a
function periodic in $\log_2\epsilon$, which would be identically $1$ if $\kappa$ were twice differentiable on the
diagonal, and according to a numerical computation reported in the companion (not a certified one) it
is not, so that $\kappa$ is not twice differentiable on the diagonal \cite[Prop.~5.26(iv), Rem.~5.29]{PaperI}.  If so, it
has no second-order form there, which Hypothesis~F(iv) needs.

\begin{theorem}[two-readout determinacy]\label{thm:trd}
Let $(M,\hp{\cdot}{\cdot}{\cdot},\kappa)$ satisfy $(\Omega1)$ and $(\Omega2)$.  Then $(\Omega3)$ holds if and
only if
\begin{enumerate}[label=(TRD),leftmargin=3.2em]
\item for all $w,x\in M$ there are $a,b,d\in\R$ such that $\kappa(y,s_xw)=a\,\kappa(y,x)+b\,\kappa(y,w)+d$
for every $y\in M$.
\end{enumerate}
When $(\Omega3)$ holds one may take $a=2c(w,x)$, $b=-1$, $d=1-c(w,x)$; for $w\notin\{x,x^\perp\}$ these are the
only coefficients, while for $w\in\{x,x^\perp\}$ other choices work (e.g.\ $w=x$: any $a+b=1$, $d=0$).
\end{theorem}
\begin{proof}
($\Rightarrow$) By Theorem~\ref{thm:main}, $M$ is one of the models, where $x\bar wx=2\ip{x}{w}x-w$
gives $c(y,s_xw)=2c(w,x)c(y,x)-c(y,w)$; substitute $c=2\kappa-1$.

($\Leftarrow$) Fix $x$ and use the group $G=G_x$ with unit $x$ (Certaine's lemma), a Lie group; then
$s_x(w)=w^{-1}$ and $\kappa(y,z)=\varphi(y^{-1}z)$ with $\varphi:=\kappa(x,\cdot)$ a class function,
$\varphi(g^{-1})=\varphi(g)$, $\varphi(g)=1$ iff $g=x$.  By $(\Omega2\mathrm d)$, $\varphi$ has exactly one
zero $z_0$; uniqueness and invariance under conjugation and inversion make $z_0$ central with
$z_0^2=x$.  Write $e:=x$.  Using $\varphi(y^{-1}w^{-1})=\varphi(wy)$ and $\varphi(y^{-1}w)=\varphi(w^{-1}y)$,
(TRD) reads
\begin{equation}\label{eq:trd1}
\varphi(wy)=a\varphi(y)+b\varphi(w^{-1}y)+d\qquad(y\in G).
\end{equation}
Replacing $y$ by $y^{-1}$ and using the same identities gives
$\varphi(w^{-1}y)=a\varphi(y)+b\varphi(wy)+d$; subtracting, $(1+b)(\varphi(wy)-\varphi(w^{-1}y))=0$.  If
$b\ne-1$, then $y=w$ gives $\varphi(w^2)=1$, so $w^2=e$.  Hence for $w^2\ne e$,
\begin{equation}\label{eq:trd2}
\varphi(wy)+\varphi(w^{-1}y)=a(w)\varphi(y)+d(w)\qquad(y\in G),
\end{equation}
and $y=e$, $y=z_0$ give $d(w)=2\varphi(wz_0)$, $a(w)=2\varphi(w)-2\varphi(wz_0)$.

\emph{Step 1: some $w$ with $w^2\ne e$ has $a(w)\ne0$.}  If $a(w)=0$, then \eqref{eq:trd2} with
$y=w^k$ gives $u_{k+1}+u_{k-1}=2u_1$ for $u_k:=\varphi(w^k)$, $u_0=1$; hence $u_2=2u_1-1$, $u_3=u_1$,
$u_4=1$, so $w^4=e$.  But $w=\exp X$ with $X\ne0$ small has $w^4\ne e$.

\emph{Step 2: $\varphi(z_0y)=1-\varphi(y)$.}  Take $w$ as in Step~1.  Then $(wz_0)^2=w^2\ne e$ and
$(wz_0)^{-1}=w^{-1}z_0$, so \eqref{eq:trd2} for $wz_0$ and for $w$ (at $z_0y$) give
$a(w)\varphi(z_0y)+d(w)=a(wz_0)\varphi(y)+d(wz_0)$.  Thus $\varphi(z_0y)=A\varphi(y)+B$, and $y=e$,
$y=z_0$ give $A=-1$, $B=1$.

\emph{Step 3.}  By Step~2, $d(w)=2-2\varphi(w)$ and $a(w)=4\varphi(w)-2$ whenever $w^2\ne e$, and
\eqref{eq:trd2} at $y=w$ gives $\varphi(w^2)=(2\varphi(w)-1)^2$.  The set $\{w:w^2=e\}$ is an analytic
subset of the connected analytic group $G$ that is not all of $G$, so its complement is dense, and the
identity extends by continuity to all $g\in G$.  For arbitrary $w,x'\in M$ put $g=w^{-1}x'$; then
$\kappa(w,\hp{x'}{w}{x'})=\varphi(g^2)$ and $2\kappa(w,x')-1=2\varphi(g)-1$, which is $(\Omega3)$.  (Only
the instances of (TRD) with the fixed point $x$ were used.)
\end{proof}

\begin{remark}\label{rem:trd}
Paper~I \cite[Thm.~4.5]{PaperI} proves more: under $(\Omega1)$, $(\Omega2\mathrm{a,b,c})$ and (TRD),
$2\kappa(o,\cdot)-1=m+(1-m)\re$ for a constant $m\in[0,1)$, and $m=0$ iff $\kappa$ has a zero.  The proof
above is different: it uses the zero of $(\Omega2\mathrm d)$ directly and needs (TRD) only at one point
$x$.
(i) Law $(\Omega2\mathrm d)$ cannot be dropped: on $S^3$, $\kappa'(x,y)=\tfrac34+\tfrac14\ip{x}{y}$
satisfies $(\Omega1)$, $(\Omega2\mathrm{a,b,c})$ and (TRD), with
$\kappa'(y,s_xw)=2c\,\kappa'(y,x)-\kappa'(y,w)+\tfrac32(1-c)$, $c=\ip{w}{x}$, but has no zero and
violates $(\Omega3)$.
(ii) Along a uniform transport $x_{n+1}=s_{x_n}(x_{n-1})$ one has $c(x_n,x_{n+1})=c(x_0,x_1)=:c$ for all
$n$ by $(\Omega2\mathrm c)$, and $r_n:=\kappa(y,x_n)$ obeys $r_{n+1}=2c\,r_n-r_{n-1}+1-c$ for every test
$y$.  The statistics of all tests along the orbit are fixed by those at two consecutive steps and by
$c$, itself the statistic of the test $x_0$ at step one.  In operational reconstructions the state of
a system is fixed by the statistics of finitely many fiducial tests \cite{Hardy,MM}; (TRD) is of a
different kind: it fixes the statistics of the next point of a transport from those of the two
previous ones, test by test, and presupposes no linear state space.  Linearity is derived afterwards
(Proposition~\ref{prop:linear}).
\end{remark}

\begin{remark}[d'Alembert's equation]\label{rem:dalembert}
For $\psi=2\varphi-1$ the interference law of Proposition~\ref{prop:linear}(iv) is d'Alembert's
equation $\psi(ab)+\psi(ab^{-1})=2\psi(a)\psi(b)$.  Davison \cite{Davison08} solved it on compact groups:
the solutions come from continuous homomorphisms into $\SU(2)$, and only compact groups isomorphic to
compact subgroups of $\SU(2)$ carry a basic solution ($f(xy)=f(y)$ for all $y$ only if $x=e$; faithfulness in
our sense implies this); pre-d'Alembert functions on topological monoids are of the form
$\tfrac12\tr\circ h$ with $h$ a continuous homomorphism into $\mathrm{Mat}_2(\C)$ \cite[Thm.~4.12]{Davison09}; see also
\cite{Stetkaer}.  The law $(\Omega3)$ is the diagonal case $a=b$.  Paper~I shows that the two-variable law
alone, with $(\Omega1)$ and $(\Omega2\mathrm{a,b,c})$, already gives the two models without $(\Omega2\mathrm d)$
\cite[Cor.~3.12]{PaperI}, whereas the diagonal case needs $(\Omega2\mathrm d)$: on $\SO(3)$ the function
$\psi=\cos\alpha$ of the rotation angle is faithful, smooth and satisfies the diagonal case, but for quarter
turns $a,b$ about orthogonal axes $\psi(ab)+\psi(ab^{-1})=-1\neq0=2\psi(a)\psi(b)$.
\end{remark}

\begin{proposition}[the commutative solution is the rebit]\label{prop:rebit}
The map $e^{i\phi}\mapsto[(\cos\tfrac\phi2,\sin\tfrac\phi2)]$ is a bijection from the solution $S^1$
onto the pure states $\R P^1$ of a real two-level system, under which $\kappa$ becomes the real
transition probability $\ip{\psi}{\phi}^2$ and the transports become the rotations of the Bloch circle.
\end{proposition}
\begin{proof}
The real ray of $(\cos\alpha,\sin\alpha)$ determines $\alpha$ modulo $\pi$, that is $\phi=2\alpha$
modulo $2\pi$.  Then $\ip{\psi}{\phi}^2=\cos^2\tfrac{\phi-\phi'}2=\lvert(e^{i\phi}+e^{i\phi'})/2\rvert^2$.
A rotation of $\R^2$ by $\alpha$ shifts $\phi$ by $2\alpha$, which is a transport
$z\mapsto e^{2i\alpha}z$ of $S^1$.
\end{proof}
So a real qubit occurs as the commutative solution, and noncommutativity of transport is what
excludes it.  From now on $M$ is the noncommutative solution $S^3$.

\begin{theorem}[obstruction {\cite[Thm.~5.5]{PaperI}}]\label{thm:obstruction}
A sphere $S^n$, $n\ge1$, carries a continuous ternary operation satisfying $(\Omega1)$ only for
$n\in\{1,3\}$.  In particular the unit sphere $S^{2^{N+1}-1}$ of $(\C^2)^{\otimes N}$, $N\ge2$, carries
no origin-free transport.
\end{theorem}
The proof (Paper~I) combines Certaine's lemma with the fact that a compact connected Lie group is a
torus or has nonzero third cohomology.  Paper~I also shows that the transport is equivalent to a pair
of flat metric connections with opposite torsion whose mean is the round Levi-Civita connection, and
that $s_x$ is the geodesic symmetry at $x$ \cite[Prop.~5.3]{PaperI}; we use only Theorem~\ref{thm:obstruction}.

%=====================================================================================
\section{One system}\label{sec:one}
We compute in the quaternion model; every statement is formulated without a base point.

\subsection{Linear structure}
\begin{proposition}\label{prop:linear}
Let $V$ be the reproducing-kernel Hilbert space of $c=2\kappa-1$.
\begin{enumerate}[label=(\roman*)]
\item $\dim V=4$, $x\mapsto c(x,\cdot)$ embeds $M$ as the unit sphere of $V$, and $c(x,y)=\ip{x}{y}_V$.
\item $\kappa(x,y)=\lVert\tfrac12(x+y)\rVert_V^2=1-\tfrac14\lVert x-y\rVert_V^2$.
\item $s_x(w)=2\ip{x}{w}x-w$.
\item (interference law) $c(y,s_xw)+c(y,w)=2\,c(w,x)\,c(y,x)$ for all $w,x,y\in M$.
\end{enumerate}
\end{proposition}
\begin{proof}
In the model $c(x,y)=\re(\bar xy)$ is the Euclidean inner product on $\HH=\R^4$; (ii), (iii) are
Example~\ref{ex:models}, and (iv) is (iii) paired with $y$.
\end{proof}
Sums and real multiples of points of $M$ are therefore defined by $\kappa$ alone.  Statement (ii)
has a reading inside quantum theory, which no proof below uses: if a control qubit is prepared in an
equal superposition of two branches carrying $x$ and $y$ and then measured in the conjugate basis, the
probability of the constructive outcome is $\lVert\tfrac12(x+y)\rVert^2=\kappa(x,y)$ (the Hadamard test
\cite{AJL}).  It depends on the relative phase of $x$ and $y$ and requires a phase reference shared by
the two branches.

\subsection{Automorphisms and clocks}
For $a,b\in M$ put $\ell_{a,b}(z):=\hp{a}{b}{z}$ and $r_{a,b}(z):=\hp{z}{a}{b}$, and let $\cL$, $\cR$ be
the sets of left and right transports.  In the model these are $z\mapsto a\bar bz$ and $z\mapsto z\bar ab$.

\begin{proposition}\label{prop:aut}
(i) $\cL$ and $\cR$ are commuting groups, each acting simply transitively on $M$ and isomorphic to
$\SU(2)$, with $\cL\cap\cR=\{\mathrm{id},x\mapsto x^\perp\}$.
(ii) $\Aut(M)=\cL\cR\cong\SO(4)$; the group $\Aut^\pm(M)$ of maps preserving $\kappa$ and preserving or
reversing the transport is $\OO(4)$; all act orthogonally on $V$.
(iii) $\Aut(M)$ is transitive on pairs with equal $\kappa$.
(iv) The commutant of $\cL$ in $\mathrm{End}_\R(V)$ is the real span of $\cR$, a copy of $\HH$.
\end{proposition}
\begin{proof}
(i) and (iv) are standard facts on left and right multiplication in $\HH$.  (ii) If $f\in\Aut(M)$,
some $\ell\in\cL$ gives $\ell f(o)=o$, and $\ell f$ is a continuous automorphism of $\SU(2)$, hence
inner; every element of $\SO(4)$ is $z\mapsto azb$, and $z\mapsto\bar z$ reverses the transport.
(iii) $\SO(4)$ is transitive on pairs of unit vectors with given inner product.
\end{proof}

\begin{definition}
A \emph{clock} is a one-parameter subgroup $T=\{\rho_\chi\}_{\chi\in\R/2\pi\Z}$ of $\cR$ with
$\rho_{2\pi}=\mathrm{id}$ as first return.  Its \emph{quarter turn} is $J_T:=\rho_{\pi/2}$.
\end{definition}
Since $\rho_\pi$ is the antipodal map, $J_T^2=-1$ on $V$.

\begin{proposition}\label{prop:clock}
(i) All clocks are conjugate under $\cR$; transport reversal exchanges clocks with one-parameter
subgroups of $\cL$.
(ii) $(V,J_T)$ is a complex vector space of dimension $2$, and
$\ip{x}{y}_T:=c(x,y)+i\,c(J_Tx,y)$
is a Hermitian inner product, conjugate-linear in $x$; $\rho_\chi$ acts as the scalar $e^{i\chi}$
and $\cL$ acts as $\SU(2)$.
(iii) The orthogonal complex structures on $V$ commuting with $\cL$ are exactly the quarter turns of clocks.
\end{proposition}
\begin{proof}
(i) One-parameter subgroups of $\SU(2)$ are conjugate.  (ii) $J_T$ is orthogonal with $J_T^2=-1$, so
$c(J_Tx,y)=-c(x,J_Ty)$, which gives Hermitian symmetry and $\ip{J_Tx}{y}_T=-i\ip{x}{y}_T$; $\cL$
commutes with $J_T$ and preserves $c$.  (iii) By Proposition~\ref{prop:aut}(iv) such a structure is
right multiplication by a unit imaginary quaternion $u$, the quarter turn of $\chi\mapsto e^{\chi u}$.
\end{proof}
A complex structure is thus a choice of clock; the choices form a two-sphere and are permuted
transitively by symmetries.  For computations take $\rho_\chi(z)=ze^{i\chi}$ and identify
$(V,J_T)$ with $\C^2$ by $z_0+jz_1\mapsto(z_0,z_1)$; then $\ip{x}{y}_T$ is the standard inner product,
and we write $r_x:=(\ip{x}{\sigma_a x}_T)_{a=1,2,3}$ for the Bloch vector of $x$.

\begin{remark}[why complex]\label{rem:complex}
Without a clock, $V\cong\HH$ carries two commuting quaternionic structures, by $\cL$ and by $\cR$, and
neither a complex nor a quaternionic structure is distinguished; as a quaternionic space $V$ is one-dimensional, a
quaternionic line, so the structure does not contain a quaternionic qubit (whose pure states form
$\HH P^1\cong S^4$).  The cone of states without a clock (Theorem~\ref{thm:cone}(i)) is the Lorentz
cone of $\R^{1,4}$, the cone of squares of the spin factor $\R\oplus\R^4$, whose normalised slice is
a four-ball with boundary $M$.  The state spaces of real, complex, quaternionic and octonionic qubits
are the spin factors $\Herm_2(\mathbb K)$ of dimensions $3,4,6,10$; the clock-free state space has
dimension $5$ and is none of them.  A clock selects $\C\subset\HH$ acting on the right, and the clocked
system is the complex qubit.  The real qubit occurs only as the commutative solution
(Proposition~\ref{prop:rebit}).
\end{remark}

\subsection{The Born rule}
Fix a clock $T$, let $P_T:=M/T$ be the set of its orbits (\emph{rays}) and $[x]$ the ray of $x$.

\begin{theorem}\label{thm:born}
For $x,y\in M$ the following numbers coincide and depend only on $[x],[y]$:
\[
\beta([x],[y]):=\max_{\chi}c(\rho_\chi x,y)^2=\frac1\pi\int_0^{2\pi}\!c(\rho_\chi x,y)^2d\chi
=\max_\chi\kappa\bigl(y,s_{\rho_\chi x}(y)\bigr)=\lvert\ip{x}{y}_T\rvert^2=\tfrac12\bigl(1+\ip{r_x}{r_y}\bigr).
\]
Moreover: (i) $x\mapsto r_x$ identifies $P_T$ with $S^2$, and $M\to P_T$ is the Hopf fibration \cite{Hopf};
(ii) $\frac1{2\pi}\int_0^{2\pi}\kappa(\rho_\chi x,y)\,d\chi=\tfrac12$ and $\beta=8\Var_\chi\kappa(\rho_\chi x,y)$;
(iii) $\beta$ is a positive semidefinite kernel of rank $4$ on $P_T$, $\beta(p,q)=1$ iff $p=q$, every
ray $p$ has exactly one $p^\perp$ with $\beta(p,p^\perp)=0$, and $\beta(q,p)+\beta(q,p^\perp)=1$;
(iv) as a function on $M\times M$, $\beta$ is invariant under $\cL$ and the normaliser of $T$ in $\cR$,
and under no larger subgroup of $\Aut(M)$.
\end{theorem}
\begin{proof}
With $z=\ip{x}{y}_T$, Proposition~\ref{prop:clock}(ii) gives $c(\rho_\chi x,y)=\re(e^{-i\chi}z)
=\lvert z\rvert\cos(\chi-\arg z)$, whose square has maximum $\lvert z\rvert^2$ and mean $\lvert z\rvert^2/2$, and
whose mean is $0$; this gives the first two expressions and (ii).  The third is $(\Omega3)$; the last
is the Bloch identity.  (iii) $\beta$ is the Gram kernel of $[x]\mapsto(1,r_x)/\sqrt2$, and $p^\perp$
has $r_{p^\perp}=-r_p$.  (iv) $\beta(x,y)=1$ iff $y\in Tx$, so an $f\in\cL\cR$ preserving $\beta$ maps
$T$-orbits to $T$-orbits, i.e.\ lies in the group $N$ of Theorem~\ref{thm:wigner}, generated by $\cL$
and $N_\cR(T)$; these preserve $\beta$.
\end{proof}

The theorem identifies $\beta$ among functions of $\kappa$; that $\beta$ is the probability of an
outcome is the reading fixed by Convention~\ref{conv:T}, which Theorem~\ref{thm:bornN} supports: it
shows that, once composites are admitted, $\beta$ is the only nonnegative function of the transition
probability with the properties a probability of outcomes needs.  For one system alone these properties do not single $\beta$ out (the weight
$3\beta^2-2\beta^3$ also sums to one on every complete test).  The comparison $\kappa$ and the Born
kernel $\beta$ are different objects: $\kappa$ lives on $S^3$, has
rank $5$ and is the statistic of a phase-referenced interference test; $\beta$ lives on $S^2$, has
rank $4$ and is what remains when the clock phase is not shared.  The first clock moment of $\kappa$
carries no information, and $\beta$ is its second moment.  In terms of distances,
$\kappa=\cos^2(d_{S^3}/2)$, while $\beta=\cos^2(\vartheta/2)$ with $\vartheta$ the Bloch angle.

\begin{remark}[the ball reading]\label{rem:ball}
Read uniformly with the commutative solution (Proposition~\ref{prop:rebit}), the points of $M=S^3$
would be the pure states and $\kappa(y,x)=\frac12(1+\ip yx)$ the probability that $x$ passes the test $y$;
complete tests are $\{y,-y\}$.  This is the state space of the four-dimensional ball, a consistent
single-system theory that is not quantum \cite{MMAP,MMPA14}.  For one system it cannot be distinguished
from the clock reading by any statement of this section; the selection comes from composition and
binarity (Theorem~\ref{thm:readings}, Proposition~\ref{prop:nogo}), proved in Section~\ref{sec:readings}.
The results of this section are stated for the clock reading.
\end{remark}

\begin{example}[a qubit computed]\label{ex:qubit}
Take $x=1\mapsto(1,0)$ and $y=\cos\tfrac\vartheta2+j\sin\tfrac\vartheta2\,e^{i\varphi}\mapsto
(\cos\tfrac\vartheta2,\sin\tfrac\vartheta2e^{i\varphi})$.  Then $r_y=(\sin\vartheta\cos\varphi,
\sin\vartheta\sin\varphi,\cos\vartheta)$, $\kappa(x,y)=\tfrac12(1+\cos\tfrac\vartheta2)$, and
$c(\rho_\chi x,y)=\cos\chi\cos\tfrac\vartheta2$.  The comparison depends on the global phase of $y$;
its clock maximum $\cos^2\frac\vartheta2$ is the Born probability of finding $y$ in the state $x$.
\end{example}

\subsection{Symmetries}
\begin{theorem}\label{thm:wigner}
Let $N$ be the group of $f\in\Aut(M)$ mapping $T$-orbits to $T$-orbits.
(i) $N=\cL\cdot N_\cR(T)$ with $N_\cR(T)=T\cup T\theta$ and $\theta\rho_\chi\theta^{-1}=\rho_{-\chi}$.
(ii) $N$ acts on $P_T\cong S^2$ through a surjection onto $\OO(3)=\Aut(P_T,\beta)$ with kernel $T$;
$\cL$ maps onto $\SO(3)$ and $T\theta$ onto $r\mapsto-r$.
(iii) $f\in N$ is complex-linear, hence unitary, on $(V,J_T)$ iff it commutes with the clock, and
conjugate-linear, hence antiunitary, iff it reverses the clock.
(iv) Every clock-reversing $\Theta\in\cR$ satisfies $\Theta^2=-1$, commutes with $\cL$ and acts on
$P_T$ by $r\mapsto-r$.  For $\Theta(z)=zj$ one has $\Theta(z_0,z_1)=(-\bar z_1,\bar z_0)$, that is
$\Theta=-i\sigma_yK$ with $K$ complex conjugation.  The clock-reversing elements of $N$ that commute with
all of $\cL$ are exactly those of $T\theta\subset\cR$.
\end{theorem}
\begin{proof}
$f(z)=azb$ maps $zT$ to $azb(b^{-1}Tb)$, so $f\in N$ iff $b\in T\cup jT$.  For $b=je^{i\chi}$:
$b^2=-1$ and $J_Tf=-fJ_T$.  The formula for $\Theta$ follows from $zj=j\bar z$ for $z\in\C$.  If
$z\mapsto azb$ commutes with all left multiplications then $a$ is central, $a=\pm1$.
\end{proof}
Statement (iv) is the time reversal of a spin-$\frac12$ system with Kramers' sign; Wigner's
alternative \cite{Wigner32,Bargmann} appears as the alternative between preserving and reversing the
clock.

\subsection{Dynamics}
\begin{theorem}\label{thm:dyn}
(i) A sequence $(x_n)_{n\in\Z}$ in $M$ satisfies $x_{n+1}=\hp{x_n}{x_{n-1}}{x_n}$ for all $n$ iff
$x_n=\ell^n(x_0)$ for some $\ell\in\cL$, iff $x_n=r^n(x_0)$ for some $r\in\cR$.
(ii) A continuous curve satisfies $x(t+s)=\hp{x(t)}{x(0)}{x(s)}$ for all $s,t$ iff $x(t)=e^{tA}x(0)$
with $A\in\su(2)$ acting on the left; these are the constant-speed geodesics.
(iii) In $(V,J_T)$, $A=-J_TH$ with $H$ Hermitian and traceless, and $i\dot x=Hx$; on $P_T$ the motion is
$\dot r=\omega\times r$ for $H=\tfrac12\omega\cdot\sigma$.
(iv) The continuous one-parameter groups in $\Aut(M)$ commuting with the clock are
$e^{tA}\rho_{-\lambda t}$, with generator $H+\lambda\mathbf 1$.
\end{theorem}
\begin{proof}
(i) With $g:=x_1x_0^{-1}$, induction on $x_{n+1}=x_nx_{n-1}^{-1}x_n$ gives $x_n=g^nx_0$, and conversely.
(ii) $g(t):=x(t)x(0)^{-1}$ is a continuous homomorphism.  (iii), (iv): $\cL$ acts on $(V,J_T)$ as
$\SU(2)$, $\rho_{-\lambda t}$ as $e^{-i\lambda t}$, and the centraliser of $T$ in $\cL\cR$ is $\cL T$.
\end{proof}
Unitary evolution is the transport that repeats at each step the displacement of the previous step,
stated without a base point.  The generator $H$ is not determined by the axiom.

\subsection{Observables, effects and uncertainty}
Let $W_T$ be the reproducing-kernel space of $\beta$: the functions $A(r)=a_0+\ip{a}{r}$ on
$P_T\cong S^2$.  A nonconstant $A$ has two critical rays $p_\pm=\pm a/\lvert a\rvert$, orthogonal, and
$A(q)=\sum_\pm(a_0\pm\lvert a\rvert)\beta(q,p_\pm)$; so $A$ is the expectation function of a two-valued
observable with outcome probabilities given by $\beta$, and $A\mapsto a_0\mathbf1+a\cdot\sigma$
identifies $W_T$ with $\Herm(2)$.

\begin{proposition}[effects]\label{prop:effects}
An element $A\in W_T$ takes values in $[0,1]$ iff $E_A:=a_0\mathbf1+a\cdot\sigma$ satisfies
$0\le E_A\le\mathbf1$, and then $A([x])=\ip{x}{E_Ax}_T$ for every $x\in M$.
\end{proposition}
\begin{proof}
On $S^2$ the extreme values of $A$ are $a_0\pm\lvert a\rvert$, the eigenvalues of $E_A$.  Further
$\ip{x}{E_Ax}_T=\tr(\rho_{r_x}E_A)=a_0+\ip{a}{r_x}$ with $\rho_r=\frac12(\mathbf1+r\cdot\sigma)$.
\end{proof}
The effects of a qubit are therefore the $[0,1]$-valued affine functions of the comparison data.
Theorem~\ref{thm:instruments} shows that, under C0--C3, each of them
is realised by an ancilla, a unitary and a test, and that nothing else is.

Define $\Var_q(A):=\sum_\pm(a_0\pm\lvert a\rvert)^2\beta(q,p_\pm)-A(q)^2=\lvert a\rvert^2-\ip{a}{r}^2$,
$\Cov_q(A,B):=\ip{a}{b}-\ip{a}{r}\ip{b}{r}$ and $\{A,B\}(r):=\ip{a\times b}{r}$, where $\times$ is the
cross product of $\R^3\cong\su(2)$, half the Lie bracket for $[e_i,e_j]=2\varepsilon_{ijk}e_k$.

\begin{theorem}\label{thm:unc}
For all $A,B\in W_T$ and all rays $q$, $\Var_q(A)\Var_q(B)=\Cov_q(A,B)^2+\{A,B\}(q)^2$.
\end{theorem}
\begin{proof}
With $r$ a unit vector, $\Var_q(A)=\lvert a\times r\rvert^2$ and $\Cov_q(A,B)=\ip{a\times r}{b\times r}$;
Lagrange's identity and $(a\times r)\times(b\times r)=\ip{r}{a\times b}\,r$ give the claim.
\end{proof}
This is Schr\"odinger's strengthening \cite{Schrodinger30} of Robertson's relation \cite{Robertson}, with
equality for pure states of one system.  The bracket term is the noncommutativity of transport.

\subsection{The standard postulates}
Table~\ref{tab:postulates} compares the postulates of single-system and composite quantum kinematics
with the statements of this paper.

\begin{table}[h]
\caption{Standard postulates of quantum mechanics, the corresponding statement here, and what is assumed.}
\label{tab:postulates}
\footnotesize
\begin{tabular}{@{}p{0.27\textwidth}p{0.38\textwidth}p{0.29\textwidth}@{}}
\toprule
Standard postulate & Statement here & Assumed beyond Axiom $\Omega$\\
\midrule
States are rays of a complex Hilbert space & $P_T=M/T\cong\C P^1$ (Thm.~\ref{thm:born}); composites $\C P^{2^N-1}$ (Thm.~\ref{thm:comp}) & [L], [B] (Thm.~\ref{thm:readings}); C0--C2, [Span] for composites\\
Transition probability $\lvert\ip{\psi}{\phi}\rvert^2$ & clock-invariant second moment of $\kappa$ (Thm.~\ref{thm:born}); sums to one on complete tests; the only admissible weight (Thm.~\ref{thm:bornN}) & reading as probability: [T]; square: [W], [Comp], [N]; composites: C0--C2, [Span] and the hypotheses (a)--(c), $w\ge0$, $w=f(\beta_N)$ of Thm.~\ref{thm:bornN}(iii)\\
Observables are Hermitian operators & $W_T\cong\Herm(2)$, spectral form via $\beta$ & none\\
Effects $0\le E\le\mathbf1$, POVMs & $[0,1]$-valued elements of $W_T$ (Prop.~\ref{prop:effects}); all realised by ancilla, unitary and test (Thm.~\ref{thm:instruments}) & C0--C3, Convention T\\
Instruments, complete positivity & Kraus form from ancillas (Thm.~\ref{thm:instruments}) & C0--C3, Convention T\\
Symmetries unitary or antiunitary & Thm.~\ref{thm:wigner}, $\Theta^2=-1$ & none\\
Unitary evolution & basepoint-free recurrence (Thm.~\ref{thm:dyn}) & the generator $H$\\
Uncertainty relation & equality (Thm.~\ref{thm:unc}) & none\\
One clock; tensor product & Thm.~\ref{thm:oneclock}; or Thm.~\ref{thm:product}, Cor.~\ref{cor:ampmult}; Prop.~\ref{prop:avg} & clock-read factors with [Bil$_\cL$], [Span], [Loc], [O], [C2], [N], C1(b); or C0, C1, [Span], and C2 for $N\ge3$\\
All unit vectors are states; all unitaries are operations & Thm.~\ref{thm:comp} & C0--C2, [Span]; C3 for operations\\
Measurement update & conditional state $K_a\psi/\lVert K_a\psi\rVert$ (Prop.~\ref{prop:relative}, Thm.~\ref{thm:instruments}) & C0--C3, Convention T\\
Finite memory & Gram-window extension and weak nesting (Thm.~\ref{thm:fmschur}); reset heat $\ge k_BT\,\Delta S(L|K)$ (Prop.~\ref{prop:fmheat}); exact phase recovery iff equal inaccessible states (Thm.~\ref{thm:fmeraser}); gluing iff PSD completion, chordal clique test (Thm.~\ref{thm:fmglue}) & [Gw]; Gibbs bath, joint unitary, kept side marginal; [Ca]\\
Single outcomes & no covariant outcome rule (Prop.~\ref{prop:nooutcome}); for a record process, Born at every stage $\Leftrightarrow$ martingale weights (Thm.~\ref{thm:sorigidity}); threshold and simplex selection Born $\Leftrightarrow$ invariant auxiliary law (Thms.~\ref{thm:sosphere}, \ref{thm:sosimplex}) & [R]; [Ap] for an apparatus; the realised dynamics: open\\
\bottomrule
\end{tabular}
\end{table}

%=====================================================================================
\section{Many systems}\label{sec:many}
Let $(M_k,\kappa_k)$, $k=1,\dots,N$, be copies of $S^3$ with clocks $T_k$, spaces $(V_k,J_k)$ and
amplitudes $\ip{\cdot}{\cdot}_k$.  By Theorem~\ref{thm:obstruction} the axiom cannot apply to the
composite, and something must be added.

\subsection{Readings}\label{sec:readings}
Section~\ref{sec:one} read the states of one system as clock orbits ([R] of earlier versions); for one
system the reading of points as states is equally consistent (Remark~\ref{rem:ball}).  We now classify
all readings of the same kind and show that binarity selects the clock and that composition with
normalisation fixes the weight.

\begin{definition}[readings]\label{def:reading}
A \emph{reading} of $M=S^3$ is a closed subgroup $K\subseteq\Aut(M)$, the part of the configuration
declared unobservable, normalised by $\cL$, together with a weight on the classes $M/K$ of the form
$w([x],[y])=\mathsf F(s_K(x,y))$, $s_K(x,y):=\sup_{k\in K}c(kx,y)$, with $\mathsf F$ continuous, $\mathsf F(1)=1$ and values in
$[0,1]$.  Two classes are \emph{orthogonal} if $s_K=0$; when $\mathsf F^{-1}(0)=\{0\}$ this is the same as $w=0$.  (The
ball reading of Remark~\ref{rem:ball}, with $\mathsf F(c)=\frac12(1+c)$, is described there by its zero set $c=-1$.)  The reading is
\emph{binary} if every class has exactly one orthogonal class;
\emph{normalised} ([N]) if for every maximal family of pairwise orthogonal classes the weights of any
class against the family sum to one; and
\emph{composable} ([Comp]) if for two copies there are a real Hilbert space $\mathcal W$, a closed group $K_{12}$ of
orthogonal maps of $\mathcal W$, unit vectors $u(x,y)\in \mathcal W$ for the product configurations and a continuous
$\mathsf F_{12}$ with $\mathsf F_{12}(1)=1$, such that
$s_{12}(u(x,y),u(x',y'))=s_K(x,x')\,s_K(y,y')$ and $\mathsf F_{12}(s_{12})=w([x],[x'])\,w([y],[y'])$, where
$s_{12}(u,u'):=\sup_{k\in K_{12}}\ip{ku}{u'}$.
\end{definition}
The form $w=\mathsf F\circ s_K$ is an input, [W]: it reduces composability to Cauchy's equation.  For a clock it
costs nothing beyond continuity: $\cL$ acts on $P_T\cong S^2$ through $\SO(3)$, which is transitive on pairs of
points at a given distance, so every continuous $\cL$-invariant weight is a function of the Bloch angle,
i.e.\ of $s_T=\cos(\vartheta/2)$.  For $K=\{\pm1\}$ it is a restriction: $\beta$ is an $\cL$-invariant weight on
$\R P^3$ that is not a function of $s_{\{\pm1\}}=\lvert c\rvert$.
Normalisation of $\cL$ says that the left transports act on classes and remain symmetries.  Binarity is
the form of law $(\Omega2\mathrm d)$ at the level of states.  Composability says that the class-correlations
and the weights of independent systems multiply.

\begin{lemma}\label{lem:normaliser}
A closed subgroup $K\subseteq\Aut(M)=\cL\cR$ normalised by $\cL$ either lies in $\cR$ or contains $\cL$.
\end{lemma}
\begin{proof}
Write elements of $\Aut(M)$ as pairs $(\ell,r)$, $z\mapsto\ell zr^{-1}$, modulo $(-1,-1)$.  If $K$ contains $(\ell,r)$ with
$\ell\ne\pm1$, then for every $a\in\SU(2)$ it contains $(a\ell a^{-1},r)(\ell,r)^{-1}=(a\ell a^{-1}\ell^{-1},1)$, so
$K\cap\cL$ is a closed normal subgroup of $\cL\cong\SU(2)$ containing a non-central element, hence equal to
$\cL$.  Otherwise every element is $(\pm1,r)=(1,\pm r)$, a right transport.
\end{proof}
If $K\supseteq\cL$ there is one class.  From now on $K\subseteq\cR$, and $kx$ means the right multiplication $xk$.
Write $W_K:=\mathrm{span}_\R K\subseteq\HH$.

\begin{lemma}\label{lem:classes}
Let $\{1\}\ne K\subseteq\cR$ be closed.  Then $s_K\in[0,1]$, and $s_K(x,y)=0$ iff $y\in x\,W_K^\perp$.  The classes
orthogonal to $[x]$ in the sense $s_K=0$ are the $K$-orbits on the unit sphere of $xW_K^\perp$:
\begin{center}\footnotesize
\begin{tabular}{@{}lll@{}}
\toprule
$K$ & $\dim W_K$ & classes with $s_K=0$\\
\midrule
$\{\pm1\}$ & $1$ & $\R P^2$ \\
cyclic $C_n$, $n\ge3$ & $2$ & a circle \\
a clock $T\cong\U(1)$ & $2$ & exactly one \\
binary dihedral, tetrahedral, octahedral, icosahedral; $N(T)$; $\SU(2)$ & $4$ & none \\
\bottomrule
\end{tabular}
\end{center}
\end{lemma}
\begin{proof}
$s_K(x,y)=h_K(x^{-1}y)$ with $h_K(v)=\sup_{k\in K}\ip kv$.  The Haar mean $m=\int_Kk\,dk$ satisfies $mk_0=m$ for
$k_0\in K$, and right multiplication by $k_0\ne1$ fixes no nonzero vector, so $m=0$ and
$h_K(v)\ge\int\ip kv\,dk=0$; equality forces $\ip kv=0$ for all $k$, i.e.\ $v\perp W_K$.  $W_K$ and $W_K^\perp$ are
invariant under right multiplication by $K$, so $K$ acts freely on the unit sphere of $W_K^\perp$.  The
closed subgroups of $\SU(2)$ are the cyclic groups, the binary dihedral, tetrahedral, octahedral and
icosahedral groups, the clocks, their normalisers and $\SU(2)$.  A cyclic group of order $n\ge3$ and a clock
span a plane $\C_u=\R+\R u$, whose complement is a plane on which $e^{\theta u}$ acts as the rotation by
$\theta$: transitively for the clock, with a circle of orbits for $C_n$.  $\{\pm1\}$ spans $\R$, with $\R P^2$
of orbits on the unit sphere of $\im\HH$.  For the other groups, $W_K$ is closed under multiplication and
contains $1$, so it is a subalgebra of $\HH$: $\R$, a copy $\C_u$, or $\HH$.  If $W_K\subseteq\C_u$, $K$ lies in the unit
circle of $\C_u$ and is abelian; the other groups are non-abelian, so $W_K=\HH$.
\end{proof}

\begin{theorem}[readings]\label{thm:readings}
Let $(K,\mathsf F)$ be a reading with $K\subseteq\cR$.
(i) If it is composable and $K\ne\{1\}$, then $\mathsf F(s)=s^p$ on the range of $s_K$, for some $p\ge0$; if
$K=\{1\}$, then $\mathsf F(c)=\lvert c\rvert^p$.  For $p>0$, $w=0$ exactly where $s_K=0$.
(ii) A reading is binary if and only if $K$ is a clock.  No weight form and no composability are used.
(iii) A composable normalised reading with $K$ a clock or $K=\{\pm1\}$ has $p=2$.
(iv) The clock reading $w=s_T^2=\beta$ is binary, composable and normalised; the reading $K=\{\pm1\}$,
$w=c^2$ (the real four-level system on $\R P^3$) is composable and normalised but not binary; $K=\cR$
(the quaternionic line) has one class.
Hence, under [L] and [B] the reading is through a clock, and under [W], [Comp] and [N] in addition its
weight is $w=\beta$; without normalisation the exponent $p$ is free.
\end{theorem}
\begin{proof}
(i) Taking $y=y'$ in the composability identity gives $\mathsf F_{12}=\mathsf F$ on the range of $s_K$, and then
$\mathsf F(st)=\mathsf F(s)\mathsf F(t)$.  If $\mathsf F(s_0)=0$ for some $s_0>0$ then $\mathsf F(s_0^{1/2^n})=0$ for all $n$, contradicting
$\mathsf F(1)=1$ by continuity; so $\mathsf F>0$ on $(0,1]$, $\log \mathsf F(e^{-t})$ is additive and continuous, and $\mathsf F(s)=s^p$ with
$p\ge0$.  For $K\ne\{1\}$ the range of $s_K$ contains an interval ending at $1$ (it is $[0,1]$ when
$W_K\ne\HH$).  For $K=\{1\}$, $s=c\in[-1,1]$ and
$\mathsf F(-1)^2=\mathsf F(1)$ forces $\mathsf F(-1)=1$, so $\mathsf F(c)=\lvert c\rvert^p$.
(ii) For $K=\{1\}$ the orthogonal set of $x$ is the two-sphere $c=0$; for $K\ne\{1\}$, Lemma~\ref{lem:classes} lists the
orthogonal classes, and exactly one occurs only for a clock.
(iii) For a clock $s_T(x,y)=\lvert\ip{x}{y}_T\rvert$, since $c(xe^{\chi u},y)=\re(e^{-i\chi}\ip{x}{y}_T)$; a maximal
orthogonal family is $\{[e_1],[e_2]\}$ for an orthonormal basis, and $\lvert z_1\rvert^p+\lvert z_2\rvert^p=1$ for all
unit $(z_1,z_2)\in\C^2$ forces $p=2$ ($p=0$ gives $2$).  For $K=\{\pm1\}$, $s=\lvert c\rvert$, maximal families are orthonormal bases of
$\R^4$, and $\sum_j\lvert z_j\rvert^p=1$ on the unit sphere of $\R^4$ forces $p=2$.
(iv) For the clock take $\mathcal W=(V,J_T)\otimes_\C(V,J_T)$, $K_{12}$ its unit scalars, $u(x,y)=x\otimes y$: then
$s_{12}=\lvert\ip{x}{x'}_T\ip{y}{y'}_T\rvert$ (Theorem~\ref{thm:product}).  For $K=\{\pm1\}$ take $\mathcal W=V\otimes_\R V$,
$K_{12}=\{\pm1\}$, $u=x\otimes y$: $s_{12}=\lvert c(x,x')c(y,y')\rvert$.  Binarity is Lemma~\ref{lem:classes}.
\end{proof}
The real, complex and quaternionic structures on the same $\R^4\cong\C^2\cong\HH^1$ ($K=\{\pm1\}$, a clock,
$\cR$) thus give systems with four, two and one perfectly distinguishable states; among these, binarity
holds exactly when the division algebra has dimension two, and the square in the Born weight follows
from multiplicativity and normalisation.  With orthogonality defined instead by $w=0$, binarity alone would
not select a clock: for the binary tetrahedral group $K=2T$ the minimum of $s_K(x,\cdot)$ is $1/\sqrt2$, attained
exactly on the class $x(2O\setminus2T)$, so a weight vanishing only at $s=1/\sqrt2$ has exactly one $w$-orthogonal
class; by (i) it is not composable.  For composable readings with $p>0$ the two notions agree.  Composability in the sense of [Comp] needs class-correlations, not only weights, to
multiply: for $K=\{1\}$, $w=\kappa$ (binary and normalised on $\{y,-y\}$ in the weight-zero sense) the composite
$\mathcal W=(\R\oplus V)^{\otimes2}$, $K_{12}=\{1\}$, $u=\Phi(x)\otimes\Phi(y)$, $\mathsf F_{12}(c)=\max(c,0)$ multiplies weights
exactly, since $\ip{\Phi(x)\otimes\Phi(y)}{\Phi(x')\otimes\Phi(y')}=\kappa(x,x')\kappa(y,y')\ge0$, but not
class-correlations.  The division algebras enter here only through the readings of one $\R^4$; on the transport side,
a ternary map on a round sphere with the standard comparison that satisfies the two Mal'cev identities
and leaves $\kappa$ invariant under all left and right transports exists only on $S^0$, $S^1$, $S^3$, $S^7$, by
Hurwitz's theorem \cite[Prop.~5.13]{PaperI}, which bears on the configuration space and not on the readings.
What the theorem adds to the classifications of real, complex and quaternionic
quantum theory \cite{Dyson62,Baez12,Adler95} and to the reconstructions that derive the complex qubit from
a binary elementary system with composites \cite{DB,BruknerZeilinger09,HoehnWever17,MMAP,MMPA14,BGW20} is
the exhaustive treatment of all closed subgroups of $\SU(2)$ acting on the comparison sphere: the finite
groups and the normaliser of a clock give no binary reading.  The exponent $2$ from
multiplicativity and normalisation is close to the arguments of \cite{Aaronson04,GalleyMasanes17}; the
count of distinguishable states matches Hardy's $K=N^r$ with $r=2$ \cite{Hardy}, and real-vector-space
quantum theory, our $K=\{\pm1\}$, is discussed in \cite{Wootters12}.  The hypothesis that $\cL$ normalises $K$ keeps a
transitive group of symmetries; by transport reversal the mirror statement holds with $\cL$ and $\cR$
exchanged.  This is the counterpart, in the present setting, of \cite{MMAP,MMPA14}, where among ball
state spaces only the three-dimensional one admits the composites considered there.

Binarity can be traded for a more primitive input.  Call two classes orthogonal when $s_K=0$ (as for
every weight $\mathsf F\circ s_K$ with $\mathsf F^{-1}(0)=\{0\}$), and consider
\begin{itemize}[leftmargin=3.2em]
\item[{[Cl]}] the unobservable group $K\subseteq\cR$ contains a clock, a one-parameter group of right
transports: for a closed system the phase of its own clock is not observable (there is no external time
reference);
\item[{[NT]}] some pair of classes is orthogonal.
\end{itemize}

\begin{proposition}[binarity from an unobservable clock]\label{prop:clockB}
Under [Cl] and [NT], $K$ is a clock and the reading is binary: every class has exactly one orthogonal
class.  No composability and no weight form are used.
\end{proposition}
\begin{proof}
The identity component $K_0$ is a closed connected subgroup of $\SU(2)$ containing a circle $T$; since
$\SU(2)$ has no two-dimensional subgroups, $K_0=T$ or $K_0=\SU(2)$.  If $K_0=T$, $K$ normalises $T$, so
$K\subseteq N(T)=T\cup Tu'$ with $u'$ a unit anticommuting with the generator of $T$, and $K=T$ or $N(T)$.  For
$N(T)$ and $\SU(2)$, $W_K=\HH$ and $s_K>0$ everywhere (Lemma~\ref{lem:classes}); explicitly
$s_{N(T)}=\max(\cos\frac\vartheta2,\sin\frac\vartheta2)\ge1/\sqrt2$, with $\vartheta$ the angle between the Hopf images, and
$s_{\SU(2)}\equiv1$.  [NT] leaves $K=T$, and Lemma~\ref{lem:classes} gives exactly one orthogonal class.
\end{proof}
So [B] is a theorem under [Cl] and [NT]; conversely [B] gives a clock (Theorem~\ref{thm:readings}(ii)), hence [Cl]
and [NT].  Under [L] the two formulations are equivalent; [Cl] with [NT] states the condition through
what is unobservable.  The notion of orthogonality matters: with $w=0$ in place of $s_K=0$ the statement
fails unless $\mathsf F^{-1}(0)=\{0\}$: for
$K=N(T)$ and $\mathsf F(s)=\max\bigl(0,(\sqrt2s-1)/(\sqrt2-1)\bigr)$, $w=0$ on the whole set $s_K=1/\sqrt2$, so some classes
are $w$-orthogonal although $K$ is not a clock, and the reading is not binary; for $K=T$ and
$\mathsf F(s)=s^2\lvert2s^2-1\rvert$, $w=0$ also at $s_T=1/\sqrt2$, and the clock reading is not binary in the $w$-sense.
On either formulation the weight is then fixed as before: [W] and [Comp] give
$\mathsf F=s^p$ (Theorem~\ref{thm:readings}(i)) and [N] gives $p=2$ (Theorem~\ref{thm:readings}(iii)); [N] alone does
not suffice for one system, since in complex dimension two every $\mathsf F$ with $\mathsf F(\cos t)+\mathsf F(\sin t)=1$
is normalised, for example $\mathsf F(s)=3s^4-2s^6$.  The real reading $K=\{\pm1\}$ and every finite $K$ contain no clock and are excluded by [Cl]
directly.  The real reading, $w=c^2$ on $S^3/\{\pm1\}$, is the $\SO(3)$ solution of $(\Omega1)$, $(\Omega2\mathrm{a,b,c})$,
$(\Omega3)$ that $(\Omega2\mathrm d)$ excludes at the level of configurations \cite[Rem.~3.13]{PaperI}; at the level of
states it is excluded by binarity (Theorem~\ref{thm:readings}(ii)) and equally by [Cl].  [Cl] is the counterpart, in the comparison structure, of Stueckelberg's description of complex
quantum theory as real quantum theory with an operator $J$ commuting with all observables
\cite{Stueckelberg}: here $J$ generates the unobservable clock, and every quantity defined on classes is
invariant under it.  [Cl] is an input with this physical reading; it is not derived from the axiom.
Corollary~\ref{cor:clbridge} derives it from the invisibility of the total phase of a closed composite under the
phase bridge [Ph], which posits the complex composite at the level of configurations; the real reading
shows that closure with an invisible total phase alone does not give it (Remark~\ref{rem:clreal}).

\begin{proposition}[the comparison reading does not compose]\label{prop:nogo}
Let $w(x,y)=f(c(x,y))$ be an $\Aut(M)$-invariant weight on configurations with $f(c)=0$ iff $c=-1$ ($\kappa$ is
one).  Suppose the configurations of a composite of two systems are unit vectors of a real
inner-product space containing the product configurations and carry a weight $w_{12}=\mathsf F(c_{12})$ with
$\mathsf F(c)=0$ iff $c=-1$.  Then $w_{12}=w\otimes w$ on product configurations is impossible.  For $w=\kappa$ and
$c_{12}=2\kappa_1\kappa_2-1$ the Gram matrix on $(\pm x_1,\pm x_2)$ is $2\cdot\mathbf1-\mathbf J$, with eigenvalues
$(2,2,2,-2)$.
\end{proposition}
\begin{proof}
For $(x_1,x_2)$, $(-x_1,x_2)$, $(x_1,-x_2)$ every pair has one factor with $f(-1)=0$, so $w_{12}=0$ and
$c_{12}=-1$ for each pair: three pairwise antipodal unit vectors, which do not exist.
\end{proof}
Proposition~\ref{prop:nogo} does not assume that class-correlations multiply, but it assumes that the
composite weight vanishes only at $c_{12}=-1$; the example above, with $\mathsf F_{12}(c)=\max(c,0)$, shows that this
hypothesis is needed.  It is a second route to
the exclusion of $K=\{1\}$.

\begin{lemma}[one clock inside the real tensor product]\label{lem:oneclock}
In $V_1\otimes_\R V_2$ let $H_\pm:=\{J_1\otimes1=\pm1\otimes J_2\}$.  Let $\mathcal W'\subseteq V_1\otimes_\R V_2$ be a subspace invariant under $J_1\otimes1$
and $1\otimes J_2$, with a complex structure $J$ such that every orbit of $T_1\times T_2$ in $\mathcal W'$ lies in an
orbit of $\{e^{\theta J}\}$.  Then $\mathcal W'\subseteq H_+$ or $\mathcal W'\subseteq H_-$, and $J=\pm J_1\otimes1$ there.  Moreover
$H_+\cong(V_1,J_1)\otimes_\C(V_2,J_2)$ and $H_-\cong(V_1,J_1)\otimes_\C(V_2,-J_2)$ have real dimension $8$; a real
product vector $x\otimes y$ has components of norm $1/\sqrt2$ in both, and the normalised components have
correlations $\re(\ip{x}{x'}\ip{y}{y'})$ and $\re(\ip{x}{x'}\overline{\ip{y}{y'}})$.
\end{lemma}
\begin{proof}
Differentiating the orbit condition, $(J_1\otimes1)w$ and $(1\otimes J_2)w$ are real multiples of $Jw$, of the same
norm as $w$, so each $w\in \mathcal W'$ satisfies $(J_1\otimes1)w=\pm(1\otimes J_2)w$; a subspace contained in $H_+\cup H_-$
lies in one of them.  $P_\pm=\frac12(1\mp J_1\otimes J_2)$ projects onto $H_\pm$; $\lVert P_\pm(x\otimes y)\rVert^2=
\frac12(1\mp c(x,J_1x)c(y,J_2y))=\frac12$, and $2\ip{P_\pm(x\otimes y)}{P_\pm(x'\otimes y')}=c(x,x')c(y,y')\mp
c(x,J_1x')c(y,J_2y')$, which are the stated correlations (Proposition~\ref{prop:avg}).
\end{proof}

\subsection{One clock for a composite}\label{sec:oneclock}
Lemma~\ref{lem:oneclock} assumes that the composite has a single clock.  We now derive it.  Let
$V_1,\dots,V_N$ be elementary systems read through their clocks (Theorem~\ref{thm:readings} under [L],
[W], [Comp], [B], [N]): $V_k=\R^4$ with clock $T_k=\{e^{\theta J_k}\}$, left transports $\cL_k$ and
$G_k:=\cL_kT_k\cong\U(2)$.  The inputs are:
\begin{itemize}[leftmargin=3.6em]
\item[{[Bil]}] the product configurations are $\Phi(x_1,\dots,x_N)$ for a real-multilinear map
$\Phi\colon V_1\times\dots\times V_N\to\mathcal W$ into a real inner-product space, with unit values on unit arguments,
equivariant for an orthogonal action $\rho$ of the clock torus $T_1\times\dots\times T_N$ on $\mathcal W$:
$\Phi(t_1x_1,\dots,t_Nx_N)=\rho(t)\Phi(x_1,\dots,x_N)$;
\item[{[Bil$_\cL$]}] in addition, $\rho$ extends to an orthogonal action of $G_1\times\dots\times G_N$ for which $\Phi$ is
equivariant;
\item[{[Span]}] $\mathcal W$ is spanned by the product configurations;
\item[{[Loc]}] the unobservable group of the composite is $K:=\rho(T_1\times\dots\times T_N)$, the group generated by
the clocks of the factors: the composite hides what its parts hide and nothing more;
\item[{[O]}] the weight $w$ of the composite on its classes has values in $[0,1]$, is symmetric,
$w(u,v)=w(v,u)$, has $w(u,u)=1$, and $w=0$ exactly on orthogonal pairs, $s_K=0$,
$s_K(u,v):=\sup_{k\in K}\ip{ku}{v}$;
\item[C2] the configurations of the composite contain the products and are closed under the round-trip
reflections $s_x(w)=2\ip{x}{w}x-w$ (Postulate~\ref{post:C2} and Theorem~\ref{thm:comp}(i), stated for $\mathcal W$);
\item[{[N]}] normalisation: for every maximal family of pairwise orthogonal classes, the weights of any
class against the family sum to one.
\end{itemize}
[Loc] contains an analogue of [Cl]: the clocks of the factors stay unobservable in the composite.  This is
not a consequence of [Cl], which is justified for a closed system, while a factor in a composite is not
closed (the model $K=\rho(T_1)$ below satisfies [Cl] for each factor and violates [Loc]); [Loc] adds that
nothing else is hidden.  The weight $\mathsf F\circ s_K$ with $\mathsf F(1)=1$ and
$\mathsf F^{-1}(0)=\{0\}$, inherited from the factors, is a special case of [O].  Its symmetry is the symmetry $(\Omega2\mathrm a)$ of the comparison, carried to classes;
it is the property that does the work below.

\begin{lemma}[double counting]\label{lem:double}
Let $w$ be a symmetric weight that is normalised on two finite maximal families $A$, $B$ of pairwise
orthogonal classes.  Then $\sum_{a\in A}\sum_{b\in B}w(a,b)=\lvert A\rvert=\lvert B\rvert$.
\end{lemma}
\begin{proof}
Summing over $b$ first gives $\sum_{a\in A}1$; by symmetry, summing over $a$ first gives $\sum_{b\in B}1$.
\end{proof}

For $\varepsilon\in\{\pm1\}^N$ let $H_\varepsilon:=\{v\in V_1\otimes_\R\dots\otimes_\R V_N: J_kv=\varepsilon_1\varepsilon_kJ_1v,\ k=2,\dots,N\}$ (with
$J_k$ acting on the $k$-th factor) and $J_\varepsilon:=\varepsilon_1J_1$ on it; $H_\varepsilon=H_{-\varepsilon}$, and for $N=2$,
$H_{(1,1)}=H_+$, $H_{(1,-1)}=H_-$ of Lemma~\ref{lem:oneclock}.  Write $P_\varepsilon$ for the orthogonal projection onto
$H_\varepsilon$ and $\chi_\varepsilon(\theta):=\sum_k\varepsilon_k\theta_k$.

\begin{theorem}[one clock]\label{thm:oneclock}
Let the factors be read through their clocks.
\begin{enumerate}[label=(\roman*),leftmargin=2em]
\item $V_1\otimes_\R\dots\otimes_\R V_N=\bigoplus_\varepsilon H_\varepsilon$ over the $2^{N-1}$ classes $\pm\varepsilon$; each $(H_\varepsilon,J_\varepsilon)$ is the
complex tensor product $\bigotimes_{\C,k}(V_k,\varepsilon_kJ_k)$, of real dimension $2^{N+1}$, on which the clock torus
acts by $e^{\chi_\varepsilon(\theta)J_\varepsilon}$.  Under [Bil], $\mathcal W_0:=\operatorname{span}\Phi$ is the orthogonal sum of nonzero
components $\mathcal W_\varepsilon$, $\varepsilon\in\mathcal E$, for a set $\mathcal E$ of classes $\pm\varepsilon$, on which the torus acts by
$e^{\chi_\varepsilon(\theta)J'_\varepsilon}$ for complex structures $J'_\varepsilon$.  Under [Bil$_\cL$] and [Span] moreover there are
$a_\varepsilon>0$ with $\sum_{\mathcal E}a_\varepsilon^2=2^{N-1}$ and an isometry of $\mathcal W$ onto $\bigoplus_{\mathcal E}H_\varepsilon$ intertwining $\rho$,
under which $\Phi(x_1,\dots,x_N)=\sum_{\varepsilon\in\mathcal E}a_\varepsilon P_\varepsilon(x_1\otimes\dots\otimes x_N)$.
\item With [Loc]: $s_K(u,v)=0$ iff $\ip{u_\varepsilon}{v_\varepsilon}=0$ for every $\varepsilon\in\mathcal E$ (Hermitian products in the
components).  For $u,v\in\mathcal W_\varepsilon\oplus\mathcal W_{\varepsilon'}$, $\varepsilon'\ne\pm\varepsilon$, $s_K(u,v)=\lvert\ip{u_\varepsilon}{v_\varepsilon}\rvert+\lvert\ip{u_{\varepsilon'}}{v_{\varepsilon'}}\rvert$; for
$N\ge3$ this formula does not extend to three or more components, whose characters are dependent.
Under [Bil$_\cL$] and [Span], on products $s_K(\Phi(x),\Phi(y))=\prod_ks_{T_k}(x_k,y_k)$ for every $N$ and every $\mathcal E$.
\item With [Loc], [O], C2 and [N], the configurations of the composite are all unit vectors of $\mathcal W_0$,
and $\mathcal E$ has exactly one element.
\item Hence $K$ acts on $\mathcal W_0$ as the unit circle of one complex structure (one clock), and if
$\dim_\C\mathcal W_0\ge3$ the weight is $w(u,v)=\lvert\ip uv\rvert^2$.  Under [Bil$_\cL$] and [Span], $\mathcal W=H_\varepsilon$,
$\dim_\C\mathcal W=2^N$, $\Phi(x_1,\dots,x_N)=x_1\otimes_\C\dots\otimes_\C x_N$ in $\bigotimes_{\C,k}(V_k,\varepsilon_kJ_k)$, and on products
$w=\prod_k\beta_k$, so probabilities of independent ray tests multiply (Postulate~\ref{post:C1}(a)).  The
orientation clause C1(b) holds for some orientation of the composite clock iff all $\varepsilon_k$ are equal,
which gives the complex tensor product $\bigotimes_{\C,k}(V_k,J_k)$.
\end{enumerate}
\end{theorem}
\begin{proof}
(i) The $J_1J_k$, $k\ge2$, are commuting involutions; their joint eigenspaces are the $H_\varepsilon$, and on $H_\varepsilon$,
$J_k=\varepsilon_kJ_\varepsilon$.  Expanding $P_\varepsilon=\prod_{k\ge2}\frac12(1-\varepsilon_1\varepsilon_kJ_1J_k)$, every term other than $1$
contains some $J_k$, $k\ge2$, exactly once, and $\ip{x_k}{J_kx_k}=0$; so $\lVert P_\varepsilon(x_1\otimes\dots\otimes x_N)\rVert^2=2^{1-N}$ for
unit $x_k$.  The map $(x_k)\mapsto P_\varepsilon(x_1\otimes\dots\otimes x_N)$ is complex-multilinear from $(V_k,\varepsilon_kJ_k)$ to
$(H_\varepsilon,J_\varepsilon)$, so $\dim_\R H_\varepsilon\le2\cdot2^N$; the $2^{N-1}$ components fill $\dim=4^N$, so equality holds
and the induced map from the complex tensor product is an isomorphism.  The torus acts on $H_\varepsilon$ by the
character $\chi_\varepsilon\ne0$, and $\chi_\varepsilon\ne\pm\chi_{\varepsilon'}$ for $\varepsilon'\ne\pm\varepsilon$.  By [Bil], $\Phi=A\circ\otimes$ with $A$ linear and
torus-equivariant; $A$ maps $H_\varepsilon$ into the isotypic component of $\mathcal W$ for $\pm\chi_\varepsilon$, on which the
generator of the torus is $\chi_\varepsilon(\cdot)J'_\varepsilon$ with $J'^2_\varepsilon=-1$.  Under [Bil$_\cL$], $H_\varepsilon$ is an outer tensor
product of irreducible representations of the $\U(2)$ factors, irreducible and of complex type
(commutant $\C$), and the $H_\varepsilon$ are pairwise inequivalent; $\ker A$ is a subrepresentation, the sum of the
$H_\varepsilon$ with $\varepsilon\notin\mathcal E$, and $A$ is injective on $H_{\mathcal E}:=\bigoplus_{\mathcal E}H_\varepsilon$, onto $\mathcal W$ by [Span].  Since
$\rho$ is orthogonal, $A^*A$ commutes with the group; on each $H_\varepsilon$ it is a positive element of the
commutant $\C$, i.e.\ $a_\varepsilon^2>0$, and $A(A^*A)^{-1/2}$ is an equivariant isometry of $H_{\mathcal E}$ onto $\mathcal W$.  Unit
values give $\sum_{\mathcal E}a_\varepsilon^22^{1-N}=1$.
(ii) $s_K(u,v)=\sup_\theta\sum_{\varepsilon\in\mathcal E}\re\bigl(e^{i\chi_\varepsilon(\theta)}z_\varepsilon\bigr)$ with $z_\varepsilon=\ip{u_\varepsilon}{v_\varepsilon}$.  The functions
$e^{\pm i\chi_\varepsilon}$ are distinct characters, so the torus average of the sum is $0$ and the sum vanishes
identically iff all $z_\varepsilon=0$; hence $s_K\ge0$ with equality iff all $z_\varepsilon=0$.  For $\varepsilon'\ne\pm\varepsilon$ the characters
$\chi_\varepsilon,\chi_{\varepsilon'}$ are linearly independent, so $\theta\mapsto(\chi_\varepsilon(\theta),\chi_{\varepsilon'}(\theta))$ maps the torus onto $T^2$ and
the supremum is $\lvert z_\varepsilon\rvert+\lvert z_{\varepsilon'}\rvert$; for $N\ge3$ the characters of three or more components can be
dependent (for $N=3$, $\chi_{(1,1,1)}+\chi_{(1,-1,-1)}=\chi_{(1,1,-1)}+\chi_{(1,-1,1)}$).  On products,
$z_\varepsilon=a_\varepsilon^22^{1-N}\prod_k\ip{x_k}{y_k}^{(\varepsilon_k)}$ (conjugated where $\varepsilon_k=-\varepsilon_1$); $\theta_k=\mp\arg\ip{x_k}{y_k}$ (one sign for all $k$, fixed by the orientation conventions) makes
every term real and positive, so $s_K=\sum_{\mathcal E}a_\varepsilon^22^{1-N}\prod_k\lvert\ip{x_k}{y_k}\rvert=\prod_ks_{T_k}$.
(iii) \emph{Configurations.}  By C2 the configurations lie in $\mathcal W_0$, which $K$ preserves; replace $\mathcal W$ by
$\mathcal W_0$.  For fixed unit $x_j$, $j\ne k$, the slice $y\mapsto\Phi(x_1,\dots,y,\dots,x_N)$ is linear with unit values on unit
vectors, hence an isometry of $V_k$ onto a four-dimensional subspace $E_k(x)$ whose unit sphere consists of
products; as in the proof of Theorem~\ref{thm:comp}(ii), the reflections in its vectors generate
$\SO(E_k(x))$.  If $\mathfrak{so}(A)$ and $\mathfrak{so}(B)$ lie in the generated Lie algebra and $A\cap B\ne0$, so does
$\mathfrak{so}(A+B)$, since $[a\wedge b,b\wedge c]=a\wedge c$ for a unit $b\perp a,c$.  $E_k(x)$ and $E_l(x)$ meet in $\Phi(x)$, any two
tuples of basis vectors are joined by changes of one slot at a time, and the slices through such tuples
span $\mathcal W_0$ by multilinearity; so the Lie algebra is $\mathfrak{so}(\mathcal W_0)$, the group generated contains
$\SO(\mathcal W_0)$ \cite{Yamabe}, and every unit vector of $\mathcal W_0$ is a configuration.
\emph{One component.}  Suppose $\mathcal E$ contains $\varepsilon\ne\pm\varepsilon'$, and let $d_\eta:=\dim_\C\mathcal W_\eta$.  By (ii) and [O],
orthogonality is componentwise.  Take complex orthonormal bases $(e^\eta_i)$ of every $\mathcal W_\eta$; $K$ acts on each
$\mathcal W_\eta$ through the full circle.  Family A, all $e^\eta_i$, is pairwise orthogonal and maximal, of size
$\sum_\eta d_\eta$.  Family B replaces the bases of $\mathcal W_\varepsilon$ and $\mathcal W_{\varepsilon'}$ by $v_i=(e^\varepsilon_i+e^{\varepsilon'}_i)/\sqrt2$,
$i\le m:=\min(d_\varepsilon,d_{\varepsilon'})$, together with the remaining $e^\varepsilon_i$ or $e^{\varepsilon'}_i$, $i>m$; it is pairwise orthogonal
and maximal by componentwise orthogonality, of size $\sum_\eta d_\eta-m$.  Lemma~\ref{lem:double} with [O] and [N]
gives $\lvert A\rvert=\lvert B\rvert$, a contradiction.  Under [Bil$_\cL$], $d_\eta=2^N$ and the sizes are $\lvert\mathcal E\rvert2^N$ and
$(\lvert\mathcal E\rvert-1)2^N$.
(iv) With $\mathcal E=\{\varepsilon\}$, $K$ is the circle $e^{\theta J'_\varepsilon}$; classes are rays and the maximal orthogonal families are
the orthonormal bases of $\mathcal W_0$.  By [N], for fixed $u$ the function $v\mapsto w(u,v)$ on rays is a nonnegative frame
function of weight one; for $\dim_\C\ge3$, Gleason's theorem \cite{Gleason57} gives $w(u,v)=\ip{v}{\rho_uv}$ for a
density operator $\rho_u$, and $w(u,u)=1$ forces $\rho_u=\lvert u\rangle\langle u\rvert$.  For the special case $\mathsf F\circ s_K$ the
simplex argument suffices: $f(t):=\mathsf F(\sqrt t)$ satisfies $\sum_if(t_i)=1$ whenever $t_i\ge0$, $\sum t_i=1$, so $f(0)=0$,
$f(a)+f(b)=f(a+b)$ and $f(t)=t$ by continuity.  Under [Bil$_\cL$] and [Span], $\Phi=a_\varepsilon P_\varepsilon\circ\otimes$ with
$a_\varepsilon^2=2^{N-1}$ is the complex tensor product by (i), and $\lvert\ip{x_1\otimes\dots}{y_1\otimes\dots}\rvert^2=\prod_k\lvert\ip{x_k}{y_k}\rvert^2$.
On $H_\varepsilon$ the factor $k$ carries $J_k=\varepsilon_kJ_\varepsilon$; C1(b) says that the composite's clock, with one of its two
orientations, restricts on each factor to that factor's clock (commentary after
Postulate~\ref{post:C1}), i.e.\ $\varepsilon_k=1$ for all $k$ up to the common sign.
\end{proof}

The content of the theorem is: if every class of unit vectors reachable by round-trip reflections is a
state, with a symmetric transition probability, and nothing is hidden beyond the local clocks, then
there is no superselection, and hence one clock.  The excluded composite is a superselected one:

\begin{proposition}[the two-sector composite]\label{prop:superselection}
Let $N=2$, $\mathcal E=\{+,-\}$, $K=T_1\times T_2$ on $\mathcal W=H_+\oplus H_-$, with all unit vectors as configurations
(as C2 gives) and orthogonality $s_K=0$.
(i) The commutant of $K$ is $M_4(\C)\oplus M_4(\C)$: this is complex quantum theory with two superselection
sectors, whose pairings of the factors are the identity and the partial transpose.
(ii) Let $w\ge0$ be any weight on classes satisfying [N], not assumed symmetric.  Then for every $u$ there is a
block-diagonal density operator $\rho(u)=\rho_+(u)\oplus\rho_-(u)$ with $w(u,x)=\tr(\rho(u)\Pi_x)$, where $\Pi_x$ is the
projection onto the span of the normalised nonzero components $\hat x_\eta=x_\eta/\lVert x_\eta\rVert$: the superselection form.
(iii) If moreover $w(u,u)=1$, then $\rho(u)=\lvert u\rangle\langle u\rvert$ for $u$ in one sector, and
$\rho_\eta(u)=p_\eta(u)\lvert\hat u_\eta\rangle\langle\hat u_\eta\rvert$ with $p_++p_-=1$ for $u$ with both components nonzero; the split $p_\eta(u)$
is not fixed by [N].  The choice $p_\eta=\lVert u_\eta\rVert^2$ is the weight
$w_{\mathrm{ss}}(u,x)=\sum_\eta\lvert\ip{u_\eta}{x_\eta}\rvert^2/\lVert x_\eta\rVert^2$ (terms with $x_\eta=0$ omitted), which vanishes exactly when
$s_K=0$ and equals $\beta_1\beta_2$ on products for every $a_\pm$.
(iv) No such weight is symmetric: for $v=(e_1+f_1)/\sqrt2$ with unit $e_1\in H_+$, $f_1\in H_-$,
$w(e_1,v)=w(f_1,v)=1$, while $w(v,e_1)+w(v,f_1)=p_+(v)+p_-(v)=1$.  In particular
$(\lvert z_+\rvert+\lvert z_-\rvert)^2$ is not the weight of a normalised model.
\end{proposition}
\begin{proof}
(i) The commutant of the two independent phases is the block-diagonal complex-linear maps.  Classes in
one sector are rays, since $K$ acts on each sector through the full circle; orthogonality is
componentwise (Theorem~\ref{thm:oneclock}(ii)).
(ii) Fix $u$ and put $f(x)=w(u,x)$.  For complex-orthonormal bases $(e_i)$ of $H_+$ and $(f_j)$ of $H_-$ the eight
classes form a maximal family, so $\sum_if(e_i)=1-\sum_jf(f_j)$ does not depend on $(e_i)$: $f$ restricted to the
rays of $H_+$ is a nonnegative frame function of some weight $c_+\in[0,1]$, and by Gleason's theorem
\cite{Gleason57} ($\dim_\C H_+=4\ge3$; for $c_+=0$, $f=0$ there) $f(x)=\ip{x}{\rho_+x}$ with $\rho_+\ge0$, $\tr\rho_+=c_+$; likewise
$\rho_-$ on $H_-$, and $\tr\rho_++\tr\rho_-=1$.  For $x$ with both components nonzero, complete $\hat x_+=e_1$ and
$\hat x_-=f_1$ to bases.  The seven classes $x,e_2,e_3,e_4,f_2,f_3,f_4$ are pairwise orthogonal, and maximal: a class
orthogonal to $e_2,e_3,e_4,f_2,f_3,f_4$ has $w_+\in\C e_1$, $w_-\in\C f_1$, and orthogonality to $x$ is
$\ip{w_+}{x_+}=\ip{w_-}{x_-}=0$, so $w=0$.  Hence $f(x)=1-\sum_{i\ge2}\ip{e_i}{\rho_+e_i}-\sum_{j\ge2}\ip{f_j}{\rho_-f_j}=
\ip{\hat x_+}{\rho_+\hat x_+}+\ip{\hat x_-}{\rho_-\hat x_-}=\tr(\rho\,\Pi_x)$.
(iii) $\ip{\hat u}{\rho_\eta\hat u}\le\tr\rho_\eta$, with equality iff $\rho_\eta\propto\lvert\hat u\rangle\langle\hat u\rvert$; so $w(u,u)=\sum_\eta\ip{\hat u_\eta}{\rho_\eta\hat u_\eta}=1=\sum_\eta\tr\rho_\eta$
forces the stated form, and a single-sector $u$ has $\rho$ concentrated on $u$.  Conversely, in any maximal family
the normalised nonzero $\pm$ components form orthonormal bases of both sectors: they are pairwise orthogonal,
since orthogonality is componentwise, and if a unit $e\in H_+$ were orthogonal to all $+$ components, the class
of $(e,0)$ would be orthogonal to the whole family and extend it (likewise for $H_-$).  Hence $\sum_x\Pi_x=\mathbf1$, and
every block-diagonal density operator, in particular any $p_\eta\ge0$ with $p_++p_-=1$, gives a weight with all
properties of (ii).  For $w_{\mathrm{ss}}$, $\lvert\ip{u_\eta}{x_\eta}\rvert^2/\lVert x_\eta\rVert^2=\lVert u_\eta\rVert^2\lvert\ip{\hat u_\eta}{\hat x_\eta}\rvert^2$; for
products, $\ip{u_\eta}{x_\eta}=\tfrac{a_\eta^2}2\ip{y_1}{x_1}\ip{y_2}{x_2}^{(\eta)}$ and $\lVert x_\eta\rVert^2=\tfrac{a_\eta^2}2$, so
$w_{\mathrm{ss}}=\tfrac{a_+^2+a_-^2}2\beta_1\beta_2$.
(iv) By (iii), $\rho(e_1)=\lvert e_1\rangle\langle e_1\rvert$ and $\hat v_+=e_1$, so $w(e_1,v)=1$; likewise $w(f_1,v)=1$; and $w(v,e_1)=p_+(v)$,
$w(v,f_1)=p_-(v)$.  Symmetry would need both to be $1$.  This is Lemma~\ref{lem:double} seen directly.
\end{proof}
The excluded composite is therefore exactly complex quantum theory with two superselection sectors and
its non-symmetric Born rule, up to the split $p_\eta(u)$ of a cross-sector state between the sectors,
which normalisation leaves free; the Born choice $p_\eta=\lVert u_\eta\rVert^2$ is one of them.

\begin{proposition}[products alone]\label{prop:productsalone}
(i) $N=2$: with the configurations restricted to products, $\mathcal E=\{+,-\}$ and $w=s_K^2=\beta_1\beta_2$ satisfy
[Bil$_\cL$], [Span], [Loc], [O] and [N]: every pairwise orthogonal family of product rays of $\C^2\otimes\C^2$
extends to an orthogonal product basis (there is no unextendible product basis in $2\otimes n$
\cite{DiVincenzo03}), so the maximal families are product bases, and $\sum\beta_1\beta_2=1$ on them.
(ii) $N\ge3$: with the configurations restricted to products, even the single component $\mathcal E=\{(1,\dots,1)\}$
with $w=\prod_k\beta_k$ violates [N]: the unextendible product basis
$\{\lvert0,1,{+}\rangle,\lvert1,{+},0\rangle,\lvert{+},0,1\rangle,\lvert{-},{-},{-}\rangle\}$ of $2\otimes2\otimes2$ \cite{Bennett99} (tensored with an
orthonormal product basis of the remaining factors) is a maximal pairwise orthogonal family of product
classes, and the weights of $\lvert000\rangle$ (times a basis vector) against it sum to $\tfrac18$.
\end{proposition}
\begin{proof}
(i) Two product rays are orthogonal iff their first or their second factors are; a family of at most three
pairwise orthogonal product rays in $2\otimes2$ is completed by a case check (or by \cite{DiVincenzo03}), and on
a product basis the Born weights of a product state sum to one.  (ii) The four vectors are pairwise
orthogonal; in each party their four local rays are pairwise distinct, so a product vector orthogonal to
all four would need, by pigeonhole, one party orthogonal to two distinct local rays, which is impossible
in $\C^2$ \cite{Bennett99}.  The weights of $\lvert000\rangle$ are $0,0,0,\tfrac18$.  Tensoring with a product basis
$(f_j)$ of the other factors keeps unextendibility (a product vector orthogonal to all $b_i\otimes f_j$ must
have its first three factors orthogonal to every $b_i$) and multiplies the sums by $\sum_j\beta(f_1,f_j)=1$.
\end{proof}
So C2 is essential: for $N=2$ it is the only input that brings the non-product vectors needed by family
B into the configuration space.  For $N\ge3$, [N] alone forces configurations beyond products: for every
unextendible product basis some configuration must be orthogonal to all its members, and such a
configuration is entangled.  This gives the existence of entangled configurations, not the closure
under round-trip reflections of Postulate~\ref{post:C2} nor the full unit sphere used in (iii).

\emph{Necessity of the inputs.}  [Loc]: $K=\rho(T_1)$ on $V_1\otimes_\R V_2\cong\C^8$ (complex structure $J_1\otimes1$, $\Phi=x\otimes y$,
$w=\lvert\ip uv_{J_1}\rvert^2$) satisfies [Bil$_\cL$], [Span], [O], C2 and [N] with $\mathcal E=\{+,-\}$; on products
$s_K=s_{T_1}(x,x')\lvert c(y,y')\rvert$, so the clock phase of factor~2 is observable in the composite.  [O]: the
asymmetric weight $w_{\mathrm{ss}}$ satisfies all other inputs (Proposition~\ref{prop:superselection}).  C2: Proposition~\ref{prop:productsalone}(i).  [N]: the two-sector composite
with $w=s_K^2$ satisfies all other inputs.  [Bil$_\cL$] is not needed for (iii) and the one-clock part of
(iv): with torus equivariance only, $\Phi(x,y)=\bigl((\mathbf1-\lvert a\otimes\bar b\rangle\langle a\otimes\bar b\rvert)(x\otimes_\C y),\ \ip{a\otimes b}{x\otimes_\C\bar y}\bigr)\in\C^3\oplus\C$ for
fixed unit $a,b\in\C^2$ is unit-valued, has characters $\theta_1+\theta_2$ and $\theta_1-\theta_2$ and slices generating
$\mathfrak{so}(8)$, is not $\cL$-equivariant, and has maximal families of sizes $4$ and $3$ and is excluded by the same
argument; left-transport equivariance fixes the form $\Phi=\sum a_\varepsilon P_\varepsilon\circ\otimes$.  [Span] is needed for the form in (iv),
not for one clock: $\mathcal W=H_+\oplus H_+$ with $\Phi=x\otimes y\oplus0$ has one clock on $\mathcal W_0$.  For the non-binary real
reading $K=\{\pm1\}$, $w=c^2$, the analogous composite is $V\otimes_\R V$ with $K_{12}=\{\pm1\}$
(Theorem~\ref{thm:readings}(iv)), real quantum theory with the real tensor product; it is excluded by the
binarity of the parts (Theorem~\ref{thm:readings}(ii)) and by the network experiments of
Example~\ref{ex:real}.

\emph{Relation to the literature.}  The splitting of the real tensor product of two complex spaces
into $H_+\oplus H_-$, the identity and the conjugate pairing, is standard in Jordan-algebraic treatments of
composites, which use the universal C$^*$-algebras of spin factors \cite{HancheOlsen83} and the theorem
that JB-algebras with tensor products are C$^*$-algebras \cite{HancheOlsen85}; composites of Euclidean
Jordan algebras are studied in \cite{BGW20}.  Real-Hilbert-space descriptions in which a separate
operator carries the imaginary unit are studied in \cite{McKagueMoscaGisin09,ABW13}; Wootters
\cite{Wootters12}, Hardy and Wootters \cite{HardyWootters12} and Renou \emph{et al.}\ \cite{Renou} concern the real
case $K=\{\pm1\}$.  In operational reconstructions the complex tensor product is selected by local
tomography \cite{Hardy,Barrett07,MMAP,DB}.  On the class of composites allowed by [Bil$_\cL$] and [Span],
read as generalised probabilistic theories with $K$-invariant effects, products of local observables span
$16$ dimensions and the sector-invariant observables of $H_+\oplus H_-$ span $32$, so local tomography holds iff
$\lvert\mathcal E\rvert=1$: on this class the inputs of Theorem~\ref{thm:oneclock} select the same composites as local
tomography; outside it the two sets of inputs are incomparable.  No tomography assumption is made, but
the family sizes $8$ and $4$ are Hardy's numbers $N_{AB}$ and $N_AN_B$ of perfectly distinguishable states in
another form; Aaronson's argument fixes the exponent $2$ by normalisation \cite{Aaronson04}.
Mathematically the theorem is a reformulation of ``no superselection, with symmetric pure-state
transition probabilities'' in terms of what is unobservable; what it adds is this formulation, with
every input stated and shown to be needed.

\subsection{Independent systems}
A \emph{system} below is an elementary comparison structure with a clock, or a composite already
constructed; its pure states are the rays of a finite-dimensional complex Hilbert space ($P_T$ for an
elementary system; all rays of $V_N$ for a composite, by Theorem~\ref{thm:comp}(ii) below).  For rays
of a Hilbert space write $\tau$ for the transition probability and, for three rays, $B_{123}$ for the
Bargmann invariant of Proposition~\ref{prop:threerays}.

By Theorem~\ref{thm:readings} each factor is read through a clock.  By Theorem~\ref{thm:oneclock} the
composite is read in the same way.  We keep the name C0 for this conclusion:

\begin{corollary}[C0, uniform reading]\label{post:C0}
For clock-read factors, under [Bil$_\cL$], [Span], [Loc], [O], C2 and [N] (Section~\ref{sec:oneclock}), the composite of $N$ systems is read
through a clock: its configurations are the unit vectors of a complex Hilbert space $\mathcal H$ whose unit
circle is its unobservable group, the class of a product configuration depends only on the classes of
its factors, the product configurations are given by one map on $N$-tuples, and the two-system
statement holds for every splitting into two groups.
\end{corollary}
\begin{proof}
Theorem~\ref{thm:oneclock}(iv); for a splitting into groups $A,B$ the complex tensor product over all
factors is $(\bigotimes_{\C,A})\otimes_\C(\bigotimes_{\C,B})$ canonically, and the clock of each group is the restriction
of the composite's clock.
\end{proof}
In earlier versions C0 was a postulate.  It can still be taken as one, without [Bil], [Bil$_\cL$], [Loc], [O] and [N] for the composite:
then clause (a) of Postulate~\ref{post:C1} is an input and Theorem~\ref{thm:product} gives the tensor
product without linearity or continuity; this is the second route below.  C0 gives a map $\mathcal P\colon P(\mathcal H_1)\times P(\mathcal H_2)\to P(\mathcal H)$ on pairs of rays into
the rays of the complex Hilbert space $\mathcal H$, whose complex structure is the composite's clock: this
was the content of the postulate in earlier versions, which is therefore the uniformity of the
reading.  Multiplicativity alone does not give it: $V_1\otimes_\R V_2$ with unobservable group $T_1\times T_2$ has
$s_{12}=\lvert\ip{x}{x'}_T\rvert\lvert\ip{y}{y'}_T\rvert$, multiplying weights and with unobservable relative phases, but its
unobservable group is a two-torus; Theorem~\ref{thm:oneclock}(iii) excludes it because no symmetric weight on it is normalised
(Lemma~\ref{lem:double}, Proposition~\ref{prop:superselection}).  Real
quantum theory with the real tensor product satisfies clause (a) below and violates C0; clause (a) is
therefore not local tomography.

\setcounter{postulate}{0}

\begin{postulate}[independent systems]\label{post:C1}
The map $\mathcal P$ of C0 satisfies
\begin{enumerate}[label=(\alph*),leftmargin=2em]
\item (probabilities of independent ray tests multiply)
$\tau(\mathcal P(p,q),\mathcal P(p',q'))=\tau(p,p')\,\tau(q,q')$;
\item (each factor keeps its orientation) for three product states that differ only in one factor, the
Bargmann invariant of their images equals that of the three factor states.
\end{enumerate}
\end{postulate}
Under the hypotheses of Theorem~\ref{thm:oneclock}, clause (a) is a theorem (Theorem~\ref{thm:oneclock}(iv))
and the postulate reduces to clause (b); clause (a) is an input only on the second route, where C0 is
assumed, and there no continuity of $\mathcal P$ is assumed.  For clock readings $s=\sqrt w$, so clause (a) is the
composability of Definition~\ref{def:reading}.  Clause (b) is a statement about rays: by
Proposition~\ref{prop:threerays}, the Bargmann invariant of a factor is fixed by its oriented clock, and
(b) says that the composite's clock restricted to a factor is that factor's clock and not its
reverse.

\begin{theorem}[product rule]\label{thm:product}
Let $\mathcal P\colon P(\C^m)\times P(\C^n)\to P(\mathcal H)$, $m,n\ge2$, satisfy (a).  Then the image of $\mathcal P$ spans a
subspace $\mathcal H_0$ of dimension $mn$, and there are a unitary $U\colon\C^m\otimes\C^n\to\mathcal H_0$
and types $\sigma,\sigma'\in\{\mathrm{id},\mathrm{conj}\}$ with
$\mathcal P([\psi],[\phi])=[U(\psi^\sigma\otimes\phi^{\sigma'})]$, where $\psi^{\mathrm{conj}}$ conjugates the coordinates in a fixed
basis.  All four types satisfy (a).  Amplitudes multiply, $\ip{U(\psi\otimes\phi)}{U(\psi'\otimes\phi')}=
\ip{\psi}{\psi'}\ip{\phi}{\phi'}$, exactly for $\sigma=\sigma'=\mathrm{id}$; the type $\sigma=\sigma'=\mathrm{conj}$ conjugates all
Bargmann invariants, and the mixed types conjugate those of one factor.  Hence (a) and (b) together
force $\sigma=\sigma'=\mathrm{id}$.
\end{theorem}
\begin{proof}
\emph{Slices.}  For fixed $q$, $p\mapsto \mathcal P(p,q)$ preserves $\tau$.  Its image lies in an $m$-dimensional
subspace $\mathcal K_q$: the images of an orthonormal basis of rays are orthonormal, and every image ray has
total transition probability one to them.  Identifying $\mathcal K_q$ with $\C^m$ by a unitary, the
non-bijective Wigner theorem on one space \cite{Geher} shows that the slice is induced by a linear or
antilinear isometry $A_q\colon\C^m\to\mathcal H$; likewise $\mathcal P(p,\cdot)$ by $B_p\colon\C^n\to\mathcal H$.
\emph{Span.}  For orthonormal bases $(e_i)$, $(f_j)$ the rays $\mathcal P([e_i],[f_j])$ are pairwise orthogonal by
(a), and for every $(p,q)$, $\sum_{ij}\tau(\mathcal P(p,q),\mathcal P([e_i],[f_j]))=\sum_i\tau(p,[e_i])\sum_j\tau(q,[f_j])=1$, so
$\mathcal P(p,q)$ lies in the span $\mathcal H_0$ of the $mn$ orthonormal vectors $u_{ij}:=A_{[f_j]}e_i$.
\emph{Coefficients.}  Rephase $B_{[e_i]}$ so that $B_{[e_i]}f_1=u_{i1}$; then $B_{[e_i]}f_j=\mu_{ij}u_{ij}$ with
$\lvert\mu_{ij}\rvert=1$.  Let $\sigma_j$ be the type of $A_{[f_j]}$ and $\sigma'_i$ that of $B_{[e_i]}$.  For
$\psi=\sum\alpha_ie_i$, $\phi=\sum\gamma_jf_j$ with all coordinates nonzero, $\mathcal P([\psi],[\phi])$ is the ray of
$A_{[\phi]}\psi=\sum_i\alpha_i^{\sigma}\lambda_i\sum_j\gamma_j^{\sigma'_i}\mu_{ij}u_{ij}$ (since $A_{[\phi]}e_i$ spans
$\mathcal P([e_i],[\phi])$) and also of $B_{[\psi]}\phi=\sum_j\gamma_j^{\sigma'}\nu_j\sum_i\alpha_i^{\sigma_j}u_{ij}$, with $\sigma,\sigma'$
the types of $A_{[\phi]}$, $B_{[\psi]}$ and phases $\lambda_i,\nu_j$.  Equality of the ratios of the
$u_{ij}$- and $u_{i1}$-coefficients for two rows $i,i'$ gives
$\bigl((\gamma_j/\gamma_1)^{\sigma'_i}\mu_{ij}\bigr)/\bigl((\gamma_j/\gamma_1)^{\sigma'_{i'}}\mu_{i'j}\bigr)
=(\alpha_i^{\sigma_j}/\alpha_i^{\sigma_1})/(\alpha_{i'}^{\sigma_j}/\alpha_{i'}^{\sigma_1})$.  The left side depends only on $\phi$,
the right only on $\psi$, so both are constant; letting $\gamma_j/\gamma_1$ range over $\C\setminus0$ gives
$\sigma'_i=\sigma'_{i'}=:\sigma'_0$, letting $\alpha_i/\alpha_{i'}$ range gives $\sigma_j=\sigma_1=:\sigma_0$, and then $\mu_{ij}=\mu_j$.  The
same comparison within one row gives $\sigma'=\sigma'_0$ and $\nu_j/\nu_1=\mu_j$, and within one column $\sigma=\sigma_0$ and
$\lambda_i=\lambda_1$.  With $U(e_i\otimes f_j):=\mu_ju_{ij}$ this proves the formula for $\phi$ with nonzero
coordinates and all $\psi$ (expand $A_{[\phi]}\psi$, whose type and phases are now known), and
symmetrically for $\psi$ with nonzero coordinates and all $\phi$.  For arbitrary $\psi$ the isometry
$B_{[\psi]}$ and $\phi\mapsto U(\psi^{\sigma_0}\otimes\phi^{\sigma'_0})$ induce the same ray map on the dense set of
$\phi$ with nonzero coordinates, hence everywhere.
\emph{Types.}  $\tau$ is unchanged by conjugation, so all types satisfy (a).  The Bargmann invariant of
$U(\psi_k^\sigma\otimes\phi^{\sigma'})$, $k=1,2,3$, equals $B(\psi_1,\psi_2,\psi_3)$ or its conjugate according to $\sigma$;
since a qubit has triples with $\im B\ne0$, (b) forces $\sigma=\sigma'=\mathrm{id}$.
\end{proof}

\begin{corollary}[amplitudes multiply]\label{cor:ampmult}
Assume C0 and Postulate~\ref{post:C1}, and for $N\ge3$ Postulate~\ref{post:C2} for the
composites already constructed, and define $V_N$ as the span of the product states ([Span]).  Then
$V_N\cong(\C^2)^{\otimes N}$ with product states $[x_1\otimes\dots\otimes x_N]$, independently of the
bracketing up to a global phase; for the lift $G(x)=x_1\otimes\dots\otimes x_N$ of product configurations
the amplitude is $a(x,y)=\ip{G(x)}{G(y)}=\prod_k\ip{x_k}{y_k}_k$.  Every other lift compatible with the clocks
of the factors differs from $G$ by a function of the product ray with values in $\U(1)$, which changes no
transition probability and no Bargmann invariant.
\end{corollary}
\begin{proof}
Under the hypotheses of Theorem~\ref{thm:oneclock} and C1(b) this is Theorem~\ref{thm:oneclock}(iv),
without induction.  On the second route, induction on $N$: Theorem~\ref{thm:product} with $\mathcal H_1=V_{N-1}$, whose rays are all states by
Theorem~\ref{thm:comp}(ii), and $\mathcal H_2=V$.  Two bracketings give unitaries $U,U'$ inducing the same
map $\mathcal P_N$ on product rays; then $U'^{-1}U$ maps every product vector to a multiple of itself and is a
scalar.  A lift $G'$ with $G'(\dots,x_ke^{i\chi},\dots)=e^{i\chi}G'(x)$ and $[G']=[G]$ is $G'=c\,G$ with $c$
invariant under every factor clock.
\end{proof}
Spanning is a definition: the composite may live in a larger Hilbert space ($\mathcal P(p,q)=[\psi\otimes\phi\oplus0]$
satisfies C0 and Postulate~\ref{post:C1}), and $V_N$ is the closed span of its product states.
The literature we searched contains the non-bijective Wigner theorem \cite{Geher}, a characterisation
of the linear maps preserving separable states (local unitaries, partial transposes and swaps, for
surjective maps \cite{FLPS}), and derivations of the tensor product from conditions on product states
in lattice-theoretic and operational settings \cite{AertsDaubechies,Matolcsi75,CMA21}; the mixed types of
Theorem~\ref{thm:product} correspond to the partial transpose, and to the second of the two
inequivalent lattice products $L(\mathcal H_1\otimes\mathcal H_2)$, $L(\mathcal H_1\otimes\bar{\mathcal H}_2)$ found by
Matolcsi \cite{Matolcsi75}.  We did not find
Theorem~\ref{thm:product}, for maps on products of ray spaces without linearity or continuity, in the
scope searched.

\begin{proposition}\label{prop:avg}
Assume amplitudes multiply (Corollary~\ref{cor:ampmult}).  Let $\Delta$ be the torus of clock settings $(\chi_k)$ with $\sum_k\chi_k=0$ and normalised Haar
measure $d\chi$.  Then
$\re a(x,y)=2^{N-1}\int_\Delta\prod_kc\bigl(\rho^{(k)}_{\chi_k}x_k,y_k\bigr)\,d\chi$.
The reproducing-kernel space of $a$ is the subspace of $V_1\otimes_\R\dots\otimes_\R V_N$ on which all
$J_k$ coincide; it has real dimension $2^{N+1}$ and equals
$(V_1,J_1)\otimes_\C\dots\otimes_\C(V_N,J_N)\cong(\C^2)^{\otimes N}$.
\end{proposition}
\begin{proof}
With $\ip{x_k}{y_k}_k=\lvert a_k\rvert e^{i\alpha_k}$ the integrand is $\prod_k\lvert a_k\rvert\cos(\alpha_k-\chi_k)$;
of the $2^N$ exponentials in its expansion only the two with all signs equal are $\Delta$-invariant,
which gives $2^{1-N}\prod\lvert a_k\rvert\cos\sum\alpha_k$.  The second statement is Aronszajn's theorem on
products of kernels \cite{Aronszajn}, with the fixed space of $\Delta$ equal to $\{J_k=J_l\ \forall k,l\}$.
\end{proof}
Corollary~\ref{cor:ampmult} thus says that independent correlations multiply and that relative clock
phases between subsystems are unobservable: a composite has one clock, and its kernel space is the
complex, not the real, tensor product.

\subsection{Round trips persist}
Let $V_N$ be this space with $c_N(\psi,\phi):=\re\ip{\psi}{\phi}$, let $\Sigma_N\subset S(V_N)$ be the
set of unit product vectors, and $\kappa_N:=\tfrac12(1+c_N)$.

\begin{postulate}[round trips persist]\label{post:C2}
The set $\mathcal S_N\subset S(V_N)$ of states of the composite contains $\Sigma_N$, and for all
$x,w\in\mathcal S_N$ there is $s_x(w)\in\mathcal S_N$ with
$c_N(y,s_xw)+c_N(y,w)=2\,c_N(w,x)\,c_N(y,x)$ for every $y\in\Sigma_N$.
\end{postulate}
For $N=1$ this is Proposition~\ref{prop:linear}(iv), a consequence of the axiom.

\begin{theorem}[entanglement from reflection closure]\label{thm:comp}
Under C0--C2 and [Span]:
(i) $s_x(w)=2c_N(x,w)x-w$ and $\kappa_N(w,s_xw)=c_N(w,x)^2$;
(ii) $\mathcal S_N=S(V_N)$, so the pure states form $\C P^{2^N-1}$;
(iii) the reflections $s_x$ generate $\OO(V_N)$, and the elements commuting with the clock of the
composite form $\U(2^N)$; every unitary is a product of round-trip reflections;
(iv) for $N\ge2$ every nonnegative, normalised, additive assignment of probabilities to orthogonal
projections is $P\mapsto\tr(\rho P)$ for a unique density operator $\rho$ \cite{Gleason57}, and the
transition probability between product rays is $\prod_k\beta_k$;
(v) $S(V_N)$ carries no transport satisfying $(\Omega1)$, and the componentwise transport of the product
of the $M_k$ violates the round-trip law for $\kappa_N$.
\end{theorem}
\begin{proof}
(i) $\Sigma_N$ spans $V_N$ over $\R$, so the law determines $s_x(w)$.
(ii) The group $\Gamma$ generated by the $s_x$, $x\in\mathcal S_N$, preserves $\mathcal S_N$.  For a product
vector $x$ and an index $k$, $E_k(x):=x_1\otimes\dots\otimes V_k\otimes\dots\otimes x_N$ is a real
four-dimensional subspace whose unit sphere lies in $\Sigma_N$; for unit $a,b\in E_k(x)$, $s_as_b$ is the
rotation of $\mathrm{span}(a,b)$ by twice their angle, so $\Gamma\supset\SO(E_k(x))$.  The real
orthonormal basis $u\,e_{i_1}\otimes\dots\otimes e_{i_N}$, $u\in\{1,i\}$, lies in $\Sigma_N$; two of its
vectors that differ in one factor lie in a common $E_k(x)$, any two are joined by such a chain, and
$[a\wedge a',a'\wedge a'']=a\wedge a''$.  So the Lie algebra generated is $\mathfrak{so}(V_N)$, the group
generated is $\SO(V_N)$ \cite{Yamabe}, and $\mathcal S_N\supseteq\SO(V_N)\Sigma_N=S(V_N)$.
(iii) $s_x$ is minus the hyperplane reflection in $x^\perp$ and has determinant $-1$ since $\dim V_N$ is
even; the centraliser of the complex structure in $\OO(2n)$ is $\U(n)$.
(iv) Gleason's theorem in dimension $2^N\ge4$; $\lvert\prod\ip{x_k}{y_k}\rvert^2=\prod\beta_k$.
(v) Theorem~\ref{thm:obstruction} and Example~\ref{ex:mirror}.
\end{proof}

\begin{example}[the mirror of a product state]\label{ex:mirror}
For $N=2$, $x=\lvert00\rangle$, $w=\lvert{+}{+}\rangle$: $c_2(x,w)=\tfrac12$ and
$s_x(w)=\lvert00\rangle-\lvert{+}{+}\rangle=\tfrac12(\lvert00\rangle-\lvert01\rangle-\lvert10\rangle-\lvert11\rangle)$,
whose concurrence $2\lvert\det\rvert$ equals $1$: maximally entangled.  The componentwise reflection is
$s_{\lvert0\rangle}\lvert{+}\rangle\otimes s_{\lvert0\rangle}\lvert{+}\rangle=\lvert{-}{-}\rangle$, a product state, and
$\kappa_2(w,\lvert{-}{-}\rangle)=\tfrac12\ne\tfrac14=c_2(w,x)^2$.
\end{example}
Postulate~\ref{post:C2} is independent of Postulate~\ref{post:C1} and of the axiom for the factors: the
state set $\mathcal S_N=\Sigma_N$ satisfies both and is not closed under round trips.
Entangled vectors are the states at which the origin-free transport is lost while the round-trip law
is kept.

Postulate~\ref{post:C2} is a statement about the set of states: the vector $s_x(w)$ is a state.  To
read Theorem~\ref{thm:comp}(iii) as a statement about evolutions one more assumption is needed.  By
Corollary~\ref{cor:ampmult}, each automorphism $\ell\in\cL T$ of a factor defines the unitary
$\ell\otimes\mathbf1$ of the composite, and by Corollary~\ref{cor:pair} each pair of factors has the
partial exchange $e^{-iJtX_{kl}}$ as its origin-free evolution.

\begin{postulate}[local operations and pair exchange]\label{post:C3}
On a composite, the operators $\ell\otimes\mathbf1$ ($\ell$ a clock-preserving automorphism of one factor)
and the partial exchanges $e^{-iJtX_{kl}}$ of pairs of factors are admissible operations, and
composites of admissible operations are admissible.
\end{postulate}

\begin{proposition}\label{prop:universal}
Under C0--C3, the admissible operations of a composite of $N$ systems
include all of $\U(2^N)$.  Postulate~\ref{post:C3} is not implied by Postulates~\ref{post:C1},
\ref{post:C2} and the one-system theorems.
\end{proposition}
\begin{proof}
The Lie algebra generated by the local $\su(2)$'s, the clock $i\mathbf1$ and $iX_{kl}$ is $\mathfrak u(2^N)$
(an entangling generator together with all local ones is universal \cite{BremnerEtAl}); the group
generated by the corresponding one-parameter groups is the connected subgroup with this Lie algebra
\cite{Yamabe}, that is $\U(2^N)$.  For independence, let the admissible operations be the products of
local unitaries: C0--C2 hold, each factor's automorphisms act, and no
entangling unitary is admissible.
\end{proof}
\begin{proposition}[complex reflections]\label{prop:complexrefl}
For every unit $x\in V_N$, $s_xs_{ix}=\mathbf1-2\lvert x\rangle\langle x\rvert$, a unitary complex reflection.  These
reflections together with the clock generate $\U(2^N)$.  Hence Postulate~\ref{post:C3} is equivalent,
in the presence of composition of operations, to: for every state $x$ the product $s_xs_{ix}$ is an
operation, and so is the clock.
\end{proposition}
\begin{proof}
$\re\ip{x}{ix}=0$ gives $s_x(s_{ix}w)=-2\re\ip xw\,x-2\re\ip{ix}w\,ix+w$, and
$\re\ip xw+i\re\ip{ix}w=\ip xw$.  On $\C^2$, $\mathbf1-2\lvert x\rangle\langle x\rvert=-r_x\cdot\sigma$, so the products
of two such reflections are $(r\cdot\sigma)(r'\cdot\sigma)=\ip r{r'}+i(r\times r')\cdot\sigma$, all of $\SU(2)$; the
$\SU(2)$'s of coordinate planes generate $\SU(2^N)$; a single reflection has determinant $-1$, and the
clock supplies the scalars, so the group generated is $\U(2^N)$.  Both this group and the one of
Postulate~\ref{post:C3} equal $\U(2^N)$ (Proposition~\ref{prop:universal}).
\end{proof}
The earlier and stronger form of the postulate, that every reflection $s_x$ is an operation, also
implies the proposition (Theorem~\ref{thm:comp}(iii)) but is not needed.  For one system, each $s_x$ is
a symmetry of $(M,\kappa)$ ($s_x\in\Aut^\pm(M)$), but it does not map rays to rays.
[Comp], [Bil], [Bil$_\cL$], [Span], [Loc], [O], [N], Convention~T and Postulates~\ref{post:C1}--\ref{post:C3} (or, on the
second route, C0 and C1(a) in place of [Bil], [Bil$_\cL$], [Loc], [O] and [N] for the composite) are the only assumptions of
Sections~\ref{sec:one}--\ref{sec:velocity} that concern more than one elementary system.  The complex
Hilbert space of the composite, which local tomography provides in \cite{Hardy,CDP,MM}, is here
Theorem~\ref{thm:oneclock}.  Postulate~\ref{post:C2} is a restricted superposition
principle in $V_N$ (it is the interference law of Proposition~\ref{prop:linear}(iv), named by analogy
with the round-trip law; composites carry no transport, Theorem~\ref{thm:comp}(v)), and
Postulate~\ref{post:C3} plays the role of continuous reversibility.

\subsection{Tests, effects and instruments}
For a state space $S(V_N)$ with its clock put
$\beta_N(\psi,\phi):=\max_\chi c_N(e^{i\chi}\psi,\phi)^2$.  Two rays are \emph{orthogonal} if
$\beta_N=0$.

\begin{definition}[Convention T]\label{conv:T}
A \emph{complete test} of a system with state space $S(V_N)$ is a family of $2^N$ pairwise orthogonal
rays; its outcomes are these rays, and the \emph{weight} of the outcome $\phi$ in the state $\psi$ is
$\beta_N(\psi,\phi)$, read as its probability.  A \emph{coarse test} is a partition of a complete test
into blocks; its outcomes are the blocks, and the weight of a block is the sum of the weights of its
rays.
\end{definition}
The convention fixes what a test is and how its outcomes are read; additivity over blocks is part of
the definition of a coarse test.  The next theorem shows that these readings are consistent and that
no other weight would be.

\begin{theorem}[Born weights for composites]\label{thm:bornN}
Assume C0--C2.
(i) The weight depends only on the rays, and
\[
\beta_N(\psi,\phi)=\lvert\ip{\psi}{\phi}\rvert^2=\max_\chi\kappa_N(\phi,s_{e^{i\chi}\psi}\phi)
=\frac1\pi\int_0^{2\pi}c_N(e^{i\chi}\psi,\phi)^2d\chi .
\]
(ii) The weights of every complete test sum to one, and the weight of a block $\mathsf B$ of a coarse
test is $\ip{\psi}{P_{\mathsf B}\psi}$, with $P_{\mathsf B}$ the projection onto the span of the block; it
does not depend on the complete test that contains the block.
(iii) Let $w_N(\psi,\phi)=f_N(\beta_N(\psi,\phi))\ge0$, $N\ge1$, be nonnegative weights that are functions of
$\beta_N$, with complete tests the $\beta_N$-orthonormal bases, such that (a) the weights of every complete
test sum to one, (b) $w_N(\psi,\psi)=1$, and (c)
$w_N(\psi,\phi)=\sum_jw_{N+1}(\psi\otimes\alpha,\phi\otimes a_j)$ for every unit $\alpha$ and complete test
$(a_j)$ of an adjoined elementary system.  Then $w_N=\beta_N$ for all $N$.
\end{theorem}
\begin{proof}
(i) As in Theorem~\ref{thm:born}, with Theorem~\ref{thm:comp}(i) for the round-trip form.
(ii) Pairwise orthogonal rays are orthonormal vectors; $2^N$ of them form a basis, and Parseval's
identity gives the sum; the weight of a block is $\sum_{\phi\in\mathsf B}\lvert\ip{\phi}{\psi}\rvert^2=
\ip{\psi}{P_{\mathsf B}\psi}$.
(iii) For $N\ge2$ and fixed $\psi$, $\phi\mapsto w_N(\psi,\phi)$ is a nonnegative frame function of
weight one on a space of dimension $2^N\ge4$, so by Gleason's theorem \cite{Gleason57}
$w_N(\psi,\phi)=\ip{\phi}{\rho_\psi\phi}$ for a density operator $\rho_\psi$.  Since $w_N$ depends only on
$\lvert\ip{\psi}{\phi}\rvert$, $\rho_\psi$ commutes with every unitary fixing $\psi$, so
$\rho_\psi=\alpha P_\psi+\gamma(\mathbf1-P_\psi)$; (b) gives $\alpha=1$ and the trace gives $\gamma=0$.
For $N=1$, (c) with $w_2=\beta_2$ gives $w_1(\psi,\phi)=\sum_j\lvert\ip{\psi}{\phi}\rvert^2\lvert\ip{\alpha}{a_j}\rvert^2
=\beta_1(\psi,\phi)$.
\end{proof}
Gleason's theorem enters only in the uniqueness statement (iii), whose hypotheses (nonnegativity,
dependence on $\beta_N$ alone) are inputs.  For one system, (a) and (b) alone do not determine the weight:
$f(x)=3x^2-2x^3$ satisfies $f(x)+f(1-x)=1$; for effects instead of projections, Gleason-type uniqueness
holds in dimension two as well \cite{Busch03,CFMR04}.

Let $V_A=V_{N_A}$, $V_B=V_{N_B}$ and $V_{AB}=V_A\otimes V_B$ with its single clock.  For
$\Psi\in S(V_{AB})$ and a unit vector $a\in V_B$ define $\Psi_a\in V_A$ by
$\ip{\phi}{\Psi_a}=\ip{\phi\otimes a}{\Psi}$ for all $\phi\in V_A$.

\begin{proposition}[relative states]\label{prop:relative}
Assume C0--C2 and Convention~\ref{conv:T}.
(i) $\beta_{AB}(\Psi,\phi\otimes a)=\lvert\ip{\phi}{\Psi_a}\rvert^2$.
(ii) For complete tests $(\phi_k)$ of $A$ and $(a_j)$ of $B$, $(\phi_k\otimes a_j)$ is a complete test
of $AB$, and the weight of the block $\{\phi_k\otimes a\}_k$ is $p(a)=\lVert\Psi_a\rVert^2$, independent of
$(\phi_k)$.
(iii) If $p(a)>0$, the conditional weights $\beta_{AB}(\Psi,\phi_k\otimes a)/p(a)$ equal
$\beta_A(\hat\Psi_a,\phi_k)$, $\hat\Psi_a:=\Psi_a/\lVert\Psi_a\rVert$, for every complete test of $A$; the
ray of $\hat\Psi_a$ is the only ray with this property.  The same holds after any unitary $U_A\otimes\mathbf1$,
with $\hat\Psi_a$ replaced by $U_A\hat\Psi_a$.
\end{proposition}
\begin{proof}
(i) is the definition of $\Psi_a$.  (ii) Products of orthonormal bases are orthonormal bases, and
$\sum_k\lvert\ip{\phi_k}{\Psi_a}\rvert^2=\lVert\Psi_a\rVert^2$.  (iii) Divide (i) by $p(a)$.  If a unit
$\chi$ has $\lvert\ip{\phi}{\chi}\rvert=\lvert\ip{\phi}{\hat\Psi_a}\rvert$ for all $\phi$, take
$\phi=\hat\Psi_a$.  Finally $((U_A\otimes\mathbf1)\Psi)_a=U_A\Psi_a$.
\end{proof}

\begin{theorem}[effects, instruments and complete positivity]\label{thm:instruments}
Assume C0--C3 and Convention~\ref{conv:T}.  Join a system $V=V_N$ to
an ancilla $V_E=V_K$ in a state $\alpha_0$, apply $U\in\U(V\otimes V_E)$, and perform a complete test
$(a)$ of the ancilla.  Put $K_a\psi:=(U(\psi\otimes\alpha_0))_a$.
(i) Each $K_a$ is linear, $\sum_aK_a^\dagger K_a=\mathbf1$, and the weight of the outcome $a$ is
$\ip{\psi}{E_a\psi}$ with $E_a=K_a^\dagger K_a$.
(ii) Given $a$, the conditional state of the system is $K_a\psi/\lVert K_a\psi\rVert$ in the sense of
Proposition~\ref{prop:relative}(iii).  For a set $e$ of outcomes,
$\sum_{a\in e}\ip{K_au}{K_av}=\ip{u}{E_ev}$ with $E_e=\sum_{a\in e}E_a$, and
$\rho\mapsto\sum_{a\in e}K_a\rho K_a^\dagger$ is completely positive.
(iii) Conversely, every family $(K_a)_{a=1}^m$ of operators on $V$ with $\sum_aK_a^\dagger K_a=\mathbf1$
arises in this way with $K=\lceil\log_2m\rceil$ ancilla systems.  In particular the effects so obtained are
exactly the operators $0\le E\le\mathbf1$, and every positive operator-valued measure on $V$ with
finitely many outcomes is realised.  The number $K$ is also necessary for a test of the ancilla alone,
whose complete tests have $2^K$ outcomes.
\end{theorem}
\begin{proof}
The unit vector $\psi\otimes\alpha_0$ is a state by Theorem~\ref{thm:comp}(ii), and $U$ is an
operation by Proposition~\ref{prop:universal}.  (i) $\psi\mapsto U(\psi\otimes\alpha_0)$
is a linear isometry, so the $K_a$ are linear and $\sum_a\lVert K_a\psi\rVert^2=1$; the weight of $a$ is
the block weight $\lVert K_a\psi\rVert^2$ of Proposition~\ref{prop:relative}(ii).  (ii) Proposition~\ref{prop:relative}(iii);
the kernel identity is $\sum_{a\in e}\ip{K_au}{K_av}=\ip{u}{\sum_{a\in e}K_a^\dagger K_av}$; a map of Kraus
form is completely positive.  (iii) Pad the family with zeros to $2^K$ operators.  The map
$W\psi:=\sum_aK_a\psi\otimes a$ is an isometry $V\to V\otimes V_E$; the linear map
$\psi\otimes\alpha_0\mapsto W\psi$ is an isometry from $V\otimes\alpha_0$ onto a subspace of the same
dimension, and extends to a unitary $U$ of $V\otimes V_E$.  Then $(U(\psi\otimes\alpha_0))_a=K_a\psi$.  For an
effect $E$ take $K_1=\sqrt E$, $K_2=\sqrt{\mathbf1-E}$.
\end{proof}
What remains assumed in this subsection is C0--C3 and
Convention~\ref{conv:T}.  Proposition~\ref{prop:nooutcome} below is unaffected: the conditional state is
defined given an outcome, and nothing in the structure selects the outcome.

\subsection{Origin-free interactions}
Identify the $N$ systems with one another by clock-compatible isomorphisms, so that $\cL\cong\SU(2)$
acts on $V_N$ by $\ell\mapsto\ell^{\otimes N}$.  A generator is \emph{origin-free} if it commutes with
this action, that is, refers to no direction.

\begin{theorem}\label{thm:originfree}
Assume C0, Postulate~\ref{post:C1} and [Span], so that $V_N\cong(\C^2)^{\otimes N}$.  The origin-free
Hermitian operators on $V_N$ are: for $N=1$, the real multiples of $\mathbf1$ (the only
direction-free evolution of an elementary system is its clock); for $N=2$, the operators
$\alpha\mathbf1+JX$ with $X(x_1\otimes x_2)=x_2\otimes x_1$, $X=\tfrac12(\mathbf1+\sum_a\sigma_a\otimes\sigma_a)$,
whose eigenspaces are the symmetric vectors (dimension $3$) and the antisymmetric ray, the only ray
invariant under joint transport, with both marginals equal to the centre of the Bloch ball; for
general $N$, the Hermitian elements of the algebra generated by the exchanges $X_{kl}$.
\end{theorem}
\begin{proof}
Schur's lemma and Schur--Weyl duality for $\SU(2)$ on $(\C^2)^{\otimes N}$ \cite{GW}.
\end{proof}

\begin{corollary}[the closed pair]\label{cor:pair}
Under the assumptions of Theorem~\ref{thm:originfree}, the origin-free evolutions of two systems are $U_t=e^{-i\alpha t}\bigl(\cos(Jt)\mathbf1-i\sin(Jt)X\bigr)$.
On rays they depend on $t$ only through $Jt$: for $J\ne0$, up to the time unit and the direction of
time there is exactly one, the partial exchange.  It is periodic; at $Jt=\pi/4$ it equals
$e^{-i\pi/4}\sqrt{X}$, with $\sqrt X:=\tfrac12(\mathbf1+X)+\tfrac i2(\mathbf1-X)$, and maps the product
of two orthogonal states to the maximally entangled $(x_1\otimes x_2-i\,x_2\otimes x_1)/\sqrt2$.
\end{corollary}
\begin{proof}
$X^2=\mathbf1$ gives $e^{-iJtX}=\cos(Jt)-i\sin(Jt)X$; at $Jt=\pi/4$ the eigenvalues on the symmetric and
antisymmetric spaces are $e^{\mp i\pi/4}$, which is $e^{-i\pi/4}\sqrt X$.  The image of $x_1\otimes x_2$
has the stated form, with Schmidt coefficients $(1/\sqrt2,1/\sqrt2)$ when $x_1\perp x_2$.
\end{proof}
A generator $H=\tfrac12\omega\cdot\sigma$ with a direction refers to a ray
$\omega/\lvert\omega\rvert\in P_T$ outside the system: an external reference.  For a closed pair the
coupling $J$ is a choice of time unit.

\begin{proposition}[symmetric origin-free generators]\label{prop:symmetric}
The origin-free Hermitian operators on $V_N$ that commute with all permutations of the $N$ factors are
the real functions $f(S^2)$ of the total spin $S^2$, where $\sum_{k<l}X_{kl}=S^2+N(N-4)/4$; they form a
space of dimension $\lfloor N/2\rfloor+1$.  For $N=1$ they are scalars.  Hence, up to adding $\alpha\mathbf1$ and choosing the time unit,
there is exactly one nonzero such generator, $\sum_{k<l}X_{kl}$, if and only if $N\in\{2,3\}$; in particular a
closed triple has equal couplings $J_{12}=J_{13}=J_{23}$.  For $N\ge4$ there are
$\lfloor N/2\rfloor-1$ further free ratios (for $N=4$, the coefficient of $(S^2)^2$).
\end{proposition}
\begin{proof}
By Schur--Weyl duality \cite{GW}, $(\C^2)^{\otimes N}=\bigoplus_SV_S\otimes M_S$ with $V_S$ the spin-$S$
representation of $\SU(2)$ and $M_S$ pairwise inequivalent irreducible representations of the symmetric
group, $S\in\{N/2,N/2-1,\dots\}\cap[0,\infty)$.  An operator commuting with both actions is a scalar on each
summand, hence a function of $S^2$, whose eigenvalues $S(S+1)$ are distinct.  From
$S^2=\tfrac14(3N+2\sum_{k<l}\sigma_k\cdot\sigma_l)$ and $X_{kl}=\tfrac12(\mathbf1+\sigma_k\cdot\sigma_l)$ follows the
formula for $\sum X_{kl}$.  Functions on a set of at most two points are affine.
\end{proof}
Permutation symmetry makes the pairwise coefficients $J_{kl}$ equal for every $N$; for $N\ge4$ the invariant generators also contain further functions of $S^2$, so the generator is not fixed up to its time unit; this leaves free parameters for
$N\ge4$.  For $g'/g$ no comparable principle exists inside one structure: the $\cL T$-invariant kinetic
forms on $\su(2)_L\oplus\mathfrak u(1)$ form a two-parameter family, and the ratio is fixed only if the form
is required to be the restriction of an $\Aut^\pm(M)$-invariant form on $\su(2)_L\oplus\su(2)_R$, which
gives $g'=g$ (Remark~\ref{rem:angle}).  That requirement is the missing step.

\subsection{Systems of other dimensions}\label{sec:otherdim}
The registers constructed so far have complex dimension $2^N$.  Other dimensions occur inside them as
sets of configurations that the paper's own operations do not leave.  We first classify such sets
in a real inner-product space $W$; closedness is relative to the unit sphere $S(W)$, and $s_x(w)=2\ip xw x-w$.

\begin{proposition}[round-trip-closed sets]\label{prop:rtclosed}
Let $C\subseteq S(W)$ be closed and round-trip-closed ($s_x(w)\in C$ for $x,w\in C$).
\begin{enumerate}[label=(\roman*),leftmargin=2em]
\item For $x,y\in C$ at angle $\alpha$, $s_ys_x$ is the rotation by $2\alpha$ in the plane of $x,y$ and the identity on its
orthogonal complement; if $\alpha/\pi$ is irrational, the great circle through $x$ and $y$ lies in $C$.
\item If $S(U)\subseteq C$ for a subspace $U$ with $\dim U\ge2$, and $y\in C$ is not orthogonal to $U$, then $S(U+\R y)\subseteq C$.
\item If $C$ is connected and has more than one point, then $C=S(U)$ for a closed subspace $U$ with $\dim U\ge2$.
\item If $C$ is moreover invariant under a clock $e^{\theta J}$ ($J$ orthogonal, $J^2=-1$), then $C$ is an orthogonal union of
spheres $S(U_j)$ of closed complex subspaces; if $C$ is connected and nonempty, $C=S(U)$ with $U$ complex.
\end{enumerate}
\end{proposition}
\begin{proof}
(i) $s_x$ fixes $x$, negates the unit $v\perp x$ in the plane and negates the orthogonal complement, so $s_ys_x$ is the
product of two line reflections of the plane, a rotation by $2\alpha$, and the identity elsewhere.  The orbit of $x$
under its powers lies in $C$ and is dense in the circle for irrational $\alpha/\pi$; $C$ is closed.
(ii) Let $x\in S(U)$.  Choose a unit $v\in U\cap x^\perp$ with $(\ip yx,\ip yv)\ne(0,0)$: any $v$ if $\ip yx\ne0$, otherwise
the normalised projection of $y$ to $U$, which is then orthogonal to $x$.  Then $t\mapsto\ip y{\cos t\,x+\sin t\,v}$ is a
nonconstant trigonometric function, and its values on every interval about $0$ fill an interval; avoiding the
countable set $\{\cos\pi q:q\in\mathbb Q\}$ gives $x_n\to x$ in $S(U)$ whose angle to $y$ is an irrational multiple of $\pi$.  By
(i) the great circle through $x_n$ and $y$ lies in $C$; letting $n\to\infty$, every normalised $ax+by\ne0$ lies in $C$.  Every
unit vector of $U+\R y$ has this form.
(iii) Join $x,y\in C$ when $\ip xy\ne0$.  Points of distinct components of this graph are orthogonal, at distance
$\sqrt2$, so components are open and closed in $C$, and a connected $C$ is one component.  If some $x,y\in C$ have
$\ip xy\notin\{\cos\pi q:q\in\mathbb Q\}$, (i) gives a great circle in $C$, and (ii), applied along paths of the graph,
gives $S(\operatorname{span}F)\subseteq C$ for every finite $F\subseteq C$, hence $S(U)\subseteq C$ for the closed span $U$ of $C$;
so $C=S(U)$.  Otherwise $w\mapsto\ip xw$ maps the connected set $C$ into a countable set, so it is constant,
equal to $\ip xx=1$, and $C=\{x\}$.
(iv) The clock orbit of $x\in C$ is the great circle through $x$ and $Jx$ (note $\ip x{Jx}=0$), so every component
contains a circle and is, by the argument of (iii), the sphere $S(U_j)$ of the closed span of the component; the
quarter turn maps $x\in S(U_j)$ to $Jx$, in the same component, so $JU_j=U_j$.
\end{proof}
\begin{remark}[the residual sets]\label{rem:residual}
Without connectedness or a clock, a closed round-trip-closed set in finite dimension is an orthogonal union of
great spheres $S(U_j)$, $\dim U_j\ge2$, and a finite set $D$ orthogonal to them, with $D=-D$ if a sphere is present.
$D$ is finite: if $y_n\to x$ in $D$, $y_n\ne x$, the angles $\alpha_n\to0$ are rational multiples of $\pi$ (otherwise
Proposition~\ref{prop:rtclosed}(i) gives a circle), the orbit of $x$ under $s_{y_n}s_x$ is a set of points spaced by at most
$2\alpha_n$ on the circle through $x$ and $y_n$, and along a subsequence of converging planes these orbits fill a
great circle in the limit, which lies in the closed set $C$ and in the component of $x$.
For a nonempty finite round-trip-closed $F$ spanning $E$, the group $G$ generated by the $s_x$, $x\in F$, acts
faithfully on $F$ and is finite; the reflections $r_x=-s_x$ generate a finite reflection group $R$, the map
$g\mapsto\det(g)g$ sends $R$ onto $G$, and $F$ is a union of $G$-orbits of unit normals of the reflecting hyperplanes
of $R$; conversely every such union is round-trip-closed, since $s_x=\det(r_x)r_x$.  Examples are the vertices of a
regular $m$-gon ($s_{x_i}x_j=x_{2i-j}$) and, in $\ell^2$, the closed infinite set $\{\pm e_j\}$, which contains no circle.
\end{remark}

\begin{corollary}[sub-structures of a register]\label{cor:substructure}
Assume C0--C2 and [Span], so that the states of $N$ systems are $S(V_N)$, $V_N\cong(\C^2)^{\otimes N}$.
(i) The nonempty connected closed subsets of $S(V_N)$ that are invariant under the clock and closed under
round trips are exactly the spheres $S(U)$ of complex subspaces $U\subseteq V_N$; every dimension $1\le d\le2^N$ occurs.
(ii) On $S(U)$ the inherited weight is $\beta_N=\lvert\ip uv\rvert^2$; the maximal orthogonal families within $S(U)$ are the
orthonormal bases of $U$, and [N] holds (Parseval).  Binarity holds exactly for $d=2$.  For $d\ge3$ every nonnegative
weight with [N] and $w(u,u)=1$ is $\beta_N$, by Gleason's theorem as in Theorem~\ref{thm:oneclock}(iv); for $d=2$ it is
not unique (the weight $3s^4-2s^6$ after Proposition~\ref{prop:clockB}) unless the ancilla consistency (c) of
Theorem~\ref{thm:bornN}(iii) is imposed, which fixes it as there.
\end{corollary}
\begin{proof}
(i) is Proposition~\ref{prop:rtclosed}(iv), and $S(U)$ is closed, clock-invariant and round-trip-closed for every
complex $U$.  (ii) An orthonormal family in $U$ that is not a basis extends inside $U$; the rest is
Theorem~\ref{thm:bornN} restricted to $U$, and orthogonal rays to $[u]$ in $U$ form $\mathbb P(U\cap u^\perp)$, empty for
$d=1$, one ray for $d=2$, infinitely many for $d\ge3$.
\end{proof}

\begin{proposition}[the exchange sector]\label{prop:symsector}
Identify the $N$ copies by clock-compatible isomorphisms, as in Theorem~\ref{thm:originfree}.  The common fixed
space of the exchanges $X_{kl}$ is $\mathrm{Sym}^N(\C^2)$, of dimension $N+1$, with orthonormal basis
$D_k=\binom Nk^{-1/2}\sum_{\lvert b\rvert=k}\lvert b\rangle$; it is invariant under the clock and irreducible of spin $N/2$ under the
collective left transports $\ell^{\otimes N}$.  For $N=2$ it is the triplet of Theorem~\ref{thm:originfree}.
\end{proposition}
\begin{proof}
The $X_{kl}$ generate the permutation action, and $P=\frac1{N!}\sum_\pi X_\pi$ is the orthogonal projection onto the
fixed space; bit strings are permuted transitively within each Hamming weight, which gives the basis.  The
collective $J_z$ has the distinct eigenvalues $N/2-k$ on $D_k$, and $J_-D_k=\sqrt{(k+1)(N-k)}\,D_{k+1}$, so every
nonzero invariant subspace contains some $D_k$ and then all of them.
\end{proof}
Given the identification of the copies, which is a choice, the sector is selected by the paper's own
operations (the exchanges of C3).  A general $d$-dimensional subspace of a register is a choice.

\begin{proposition}[operations and composites of sub-structures]\label{prop:subops}
Assume C0--C3, [Span] and Convention~\ref{conv:T}, and let $U\subseteq V_N$, $U_i\subseteq V_{N_i}$ be complex subspaces.
(i) Every unitary $L$ of $U$ is realised by the operation $L\oplus\mathbf1_{U^\perp}$ (Proposition~\ref{prop:universal}).
(ii) Every finite family $(K_a)$ of operators on $U$ with $\sum_aK_a^\dagger K_a=\mathbf1_U$ is realised as in
Theorem~\ref{thm:instruments}, with an ancilla $V_K$, $2^K\ge$ the number of outcomes, and a unitary of $V_N\otimes V_K$ that
maps $\psi\otimes\alpha_0$ to $\sum_aK_a\psi\otimes a_a$ for $\psi\in U$ and is the identity on $U^\perp\otimes V_K$; all branches stay in $U$.
(iii) The smallest round-trip-closed subset of $S(V_{N_1+N_2})$ containing the products of $S(U_1)$ and $S(U_2)$ is
$S(U_1\otimes_\C U_2)$.  The factor clocks act on it through one circle, and the weights of products multiply.
\end{proposition}
\begin{proof}
(i) Proposition~\ref{prop:universal} gives all unitaries of $V_N$.  (ii) The prescribed map is an isometry from
$U\otimes\alpha_0$ into $U\otimes V_K$; complete it to a unitary of $U\otimes V_K$ and extend by the identity; then apply (i)
on $V_N\otimes V_K$ and Theorem~\ref{thm:instruments}.  (iii) Take orthonormal bases $(e_i)$, $(f_j)$; the vectors $e_i\otimes f_j$,
$ie_i\otimes f_j$ form a real orthonormal basis of $U_1\otimes_\C U_2$.  Two of them that share $i$ or $j$ span a plane whose
unit circle consists of products, so by Proposition~\ref{prop:rtclosed}(i) the round trips in its points give all
rotations of that plane; any two basis vectors are joined by such steps, and, as in the proof of
Theorem~\ref{thm:comp}(ii), these rotations generate $\SO(U_1\otimes_\C U_2)$, which is transitive on the sphere.
Conversely $S(U_1\otimes_\C U_2)$ is round-trip-closed.  The factor phases act by $e^{i(\theta_1+\theta_2)}$, and
$\lvert\ip{u_1\otimes u_2}{v_1\otimes v_2}\rvert^2=\beta_1\beta_2$.
\end{proof}

These are theorems about sub-structures of registers.  To call such a sub-structure a system one more
input is needed:
\begin{itemize}[leftmargin=3.2em]
\item[{[Sub]}] a sub-structure $S(U)$ of a register is treated as a system: only configurations in $S(U)$ are
prepared, complete tests are the maximal orthogonal families within $S(U)$, and operations are those of
the host that preserve $U$.
\end{itemize}
Under [Sub], Corollary~\ref{cor:substructure} and Proposition~\ref{prop:subops} give complex quantum theory in
every finite dimension, with its weights, operations and instruments, and composites of such systems
with one clock.  Reconstructions contain corresponding inputs.  Hardy's Axiom~3 (Subspaces) states that
``a system whose state is constrained to belong to an $M$ dimensional subspace (i.e.\ have support on only $M$
of a set of $N$ possible distinguishable states) behaves like a system of dimension $M$'' \cite{Hardy}.  Masanes
and M\"uller's Requirement~3 (Equivalence of subspaces) states that for systems $S_c$ and $S_{c-1}$ of capacities
$c$ and $c-1$ and a complete measurement $\Omega_1,\dots,\Omega_c$ on $S_c$, the states with $\Omega_c(\psi)=0$ form a state space
equivalent to $S_{c-1}$ \cite{MM}.  Daki\'c and Brukner include in their Axiom~1 that all systems of the same
information carrying capacity are equivalent, a rephrasing of Hardy's subspace axiom that they attribute to
Grinbaum, and build
higher-dimensional systems from two-dimensional ones \cite{DB}.  Those axioms assert an equivalence
between a restriction and an independently given system; [Sub] has no such system to compare with and
declares the restriction to be the system.  Nothing here produces an elementary comparison structure of
another dimension: the models of the axiom remain $S^1$ and $S^3$ (Theorem~\ref{thm:main}).

%=====================================================================================
\section{Geometry of states}\label{sec:geometry}

\subsection{Lorentz cones of states}\label{sec:cone}
Let $\Phi(x)=(1,x)/\sqrt2\in\R\oplus V$ and $\Phi_T([x])=(1,r_x)/\sqrt2\in\R\oplus\R^3$ be the feature
maps of $\kappa$ and $\beta$, and $B((t,v),(t',v'))=tt'-\ip{v}{v'}$.

\begin{theorem}\label{thm:cone}
(i) The closed convex cone $K_M$ generated by $\Phi(M)$ is the Lorentz cone of $\R^{1,4}$; the points of
$M$ are its extreme rays, which are its null rays, and $1-\kappa(x,y)=B(\Phi x,\Phi y)$.
(ii) $\Aut(K_M)=\R_{>0}\times\OO^+(1,4)$ acts on $M$ by M\"obius transformations with
$1-\kappa(fx,fy)=\mu_f(x)\mu_f(y)(1-\kappa(x,y))$ for a positive function $\mu_f$; the subgroup
preserving $\kappa$ is $\OO(4)=\Aut^\pm(M)$.
(iii) The same holds for $\Phi_T$, $P_T$, $\beta$ with $\R^{1,3}$, $\OO^+(1,3)$ and $\OO(3)$; the identity
component of $\Aut(K_T)$ acts by $\rho\mapsto\lambda g\rho g^\dagger$, $\lambda>0$, $g\in\SL(2,\C)$, on $\Herm(2)$; and the
points of $K_T$ with $t=1/\sqrt2$ form the Bloch ball of probabilistic mixtures of rays.
(iv) A quadratic form on $\R\oplus\R^3$ vanishing on all rays of $P_T$ is a multiple of $B$.
\end{theorem}
\begin{proof}
(i) The convex hull of $\Phi(M)$ is $\{(1,v)/\sqrt2:\lVert v\rVert\le1\}$ and
$B(\Phi x,\Phi y)=\tfrac12(1-\ip{x}{y})$.  (ii) $\Aut$ of the Lorentz cone is $\R_{>0}\times\OO^+(1,n)$
\cite{FK}; for $g\in\OO^+(1,4)$, $g(1,x)=\lambda(x)(1,fx)$ with $\lambda>0$, so $\mu_f=1/\lambda$, and $g$
preserves $\kappa$ iff it fixes $(1,0)$.  (iii) The same argument with $\SL(2,\C)\to\SO^+(1,3)$.  (iv)
Rigidity of the null cone.
\end{proof}

\begin{table}[h]
\caption{The three levels of one system.  The first row needs no clock; the other two need one.}
\label{tab:levels}
\small
\begin{tabular}{@{}p{0.2\textwidth}p{0.17\textwidth}p{0.25\textwidth}p{0.3\textwidth}@{}}
\toprule
Object & Space & Kernel & Cone\\
\midrule
configurations & $M\cong S^3$ & $\kappa=\cos^2(d_{S^3}/2)$, rank 5 & Lorentz cone of $\R^{1,4}$ (boundary $M$)\\
rays (pure states) & $P_T\cong S^2$ & $\beta=\cos^2(\vartheta/2)$, rank 4 & Lorentz cone of $\R^{1,3}$, $\cong$ PSD $\Herm(2)$\\
mixtures & Bloch ball $B^3$ & $\tr(\rho\rho')$ & interior of the same cone\\
\bottomrule
\end{tabular}
\end{table}

The mismatches $1-\kappa$ and $1-\beta$ are Minkowski products of null vectors, and the signature
$(1,3)$ is $1+\dim\su(2)$.  The boosts of $\OO^+(1,3)$ are not symmetries of $\beta$; they are the
invertible filters $\rho\mapsto g\rho g^\dagger$ \cite{VDD}.

\subsection{The velocity space}\label{sec:velocity}
Let $\mathbb B:=\{\rho_r:\lvert r\rvert<1\}$ be the normalised interior of $K_T$.  By
Theorem~\ref{thm:comp} every $\rho_r$ is the marginal of a pure state $\psi\in V_2$; write
$C(\psi):=2\lvert\det\psi\rvert$ for its concurrence \cite{Wootters}, $\psi$ read as a $2\times2$ matrix.

\begin{theorem}\label{thm:velocity}
Let $g\in\SL(2,\C)$ act on $\mathbb B$ by $\rho\mapsto g\rho g^\dagger/\tr(g\rho g^\dagger)$.
(i) This is a transitive action of $\SO^+(1,3)$ preserving the distance
$\cosh d(\rho_r,\rho_{r'})=(1-\ip{r}{r'})/\sqrt{(1-\lvert r\rvert^2)(1-\lvert r'\rvert^2)}$, and
$(\mathbb B,d)$ is hyperbolic three-space in the Beltrami--Klein model, whose geodesics are the chords,
that is, the mixtures.
(ii) The filter $g=g^\dagger>0$ taking the centre to $\rho_r$ corresponds under $\SL(2,\C)\to\SO^+(1,3)$
to the pure boost of $\R^{1,3}$ with velocity parameter $r$, $\lvert r\rvert<1$; two such filters compose
to the pure boost whose parameter is the Einstein sum of the two, followed by a rotation \cite{CU}; and $\cosh d(\rho_0,\rho_r)=\gamma(r):=(1-\lvert r\rvert^2)^{-1/2}
=(4\det\rho_r)^{-1/2}$.
(iii) The rays $P_T$ form the sphere at infinity of $\mathbb B$: no filter takes an interior point to a ray.
(iv) If $\rho_r$ is the marginal of $\psi$ then $\gamma(r)^{-1}=C(\psi)$, i.e.\ $\lvert r\rvert^2+C^2=1$;
for a test axis $r_p$, with $P:=\lvert\ip{r}{r_p}\rvert$ and $V:=\lvert r-\ip{r}{r_p}r_p\rvert$,
$P^2+V^2+C^2=1$ \cite{JB}.
\end{theorem}
\begin{proof}
The hyperboloid $B(\xi,\xi)=1$, $\xi_0>0$, with the metric of $-B$ is hyperbolic space with
$\cosh d=B(\xi,\xi')$, and central projection to $\xi_0=1$ is the Klein model; $\rho_r$ corresponds to
$\gamma(r)(1,r)$, and $\det\rho_r=\tfrac14(1-\lvert r\rvert^2)$.  The action is that of
Theorem~\ref{thm:cone}(iii).  For (iv), $\det\rho_r=\lvert\det\psi\rvert^2$.
\end{proof}
The state space of one system thus carries the velocity geometry of special relativity, with purity
one as the bound and the reciprocal Lorentz factor of a subsystem equal to its concurrence.  The
Minkowski space here is the kernel space of $\beta$, and the theorem concerns states; calling
$\lvert r\rvert$ a velocity is a reading, and question (Q5) is not answered by it.  Events and a metric on a manifold of events are the subject of
Section~\ref{sec:events}; Proposition~\ref{prop:solder} identifies each tangent space of that manifold
isometrically with this Minkowski space, but does not identify the filter group with a symmetry group
of events.


%=====================================================================================
\paperpart{Part II.  Observers inside}
\noindent This part treats what an observer inside the structure can know and record.  It contains a
closed universe with internal clocks, finite memory with erasure and consensus, and single outcomes; a
section on computation is in preparation.

%=====================================================================================
\section{The closed universe: a global singlet and internal clocks}\label{sec:closedreg}
Can a closed register of comparison systems have no preferred spin direction and still show change
relative to one of its own subsystems?  For $2n$ elementary systems we construct the global singlets and
determine which marginals can be maximally mixed.  An exchange constraint then gives a conditional
Schr\"odinger equation with an explicit term for the exchanges across the clock cut; eight elementary
systems suffice for an exact example in which both the clock reading and the changing observable are
measured inside local singlet sectors.  The invisibility of the total phase gives [Cl] under a named
bridge hypothesis.  Restriction to the observables across another cut gives a modular flow, whose
agreement with exchange time has an exact Gibbs criterion.

\subsection{The closed singlet sector}\label{sec:cusinglet}
Assume C0--C2 and [Span] for the composite, so that $V_{2n}=(\C^2)^{\otimes2n}$ with the complex composite reading of
Theorem~\ref{thm:comp}; the factors are identified as in Theorem~\ref{thm:originfree}, joint transport is $g^{\otimes2n}$
($g\in\SU(2)$), and $X_{kl}$ exchanges factors $k$ and $l$.  Let $Q$ be the set of the $2n$ sites ($n\ge1$); for $A\subseteq Q$
put $\bar A=Q\setminus A$, $\rho_A=\tr_{\bar A}\lvert\Psi\rangle\langle\Psi\rvert$ and $S(\rho)=-\tr\rho\log\rho$ (natural logarithm, $0\log0=0$).  Inner
products are conjugate-linear in the first argument.  A \emph{singlet} is a vector fixed by joint transport.
The closed-register model includes every measurement and reference system used in its statistics; its
origin-free kinematics is that of the heap and the round-trip law \cite[Def.~2.3, Thm.~3.1]{PaperI}.
The choice of a global singlet is a condition on the state.

\begin{theorem}[singlets and maximally mixed marginals]\label{thm:cusinglet}
The singlet space $\mathcal S_{2n}$ has dimension $d_n=\binom{2n}{n}-\binom{2n}{n-1}=\frac1{n+1}\binom{2n}{n}$.  Every perfect
matching $\mathcal M$ of $Q$ gives a unit singlet $\bigotimes_{\{k,l\}\in\mathcal M}s_{kl}$, $s_{kl}=(\lvert0_k1_l\rangle-\lvert1_k0_l\rangle)/\sqrt2$ ($k<l$).  The
singlet ray is unique exactly for $n=1$.  Every singlet ray has global entropy zero, and its one-site
marginals are $\mathbf1/2$.  For a specified cut $A\vert\bar A$, a singlet with $\rho_A$ maximally mixed exists exactly when
$\lvert A\rvert\le n$; for $2\le\lvert A\rvert\le n$ whether $\rho_A$ is maximally mixed depends on the singlet, and for an individual
singlet it is the condition $DD^*=2^{-\lvert A\rvert}\mathbf1$ on its coefficient matrix $D$ across the cut.  No singlet has maximally
mixed marginals on all pairs.
\end{theorem}
\begin{proof}
The antisymmetric two-site tensor transforms by $\det g=1$, so products over a matching are singlets.  In
the decomposition of Proposition~\ref{prop:symmetric} a spin $j$ contributes one weight vector at each $m=-j,\dots,j$, so
$\dim V[m{=}0]-\dim V[m{=}1]$ is the multiplicity of spin zero; counting bit strings with $n$ and $n-1$ ones gives
$d_n$.  For $n>1$ the matchings $(12)(34)$ and $(13)(24)$, extended by the same pairs, give distinct rays: their
$12$ marginals have rank one and four.  An invariant ray is fixed as a vector, since its character of the
connected perfect group $\SU(2)$ is trivial.  The projector has spectrum $1,0,\dots,0$.  Invariance passes to
$\rho_A$; at one site a commuting operator is scalar by irreducibility.  In product bases $\rho_A=DD^*$, so
$\rk\rho_A\le2^{2n-\lvert A\rvert}$, which excludes $\lvert A\rvert>n$; if $\lvert A\rvert\le n$, matching each site of $A$ to a distinct site of $\bar A$
gives $\rho_A=2^{-\lvert A\rvert}\mathbf1$, and a pair inside $A$ gives a non-maximal marginal.  By Proposition~\ref{prop:symmetric},
$\sum_{k<l}X_{kl}=n(n-2)\mathbf1$ on $\mathcal S_{2n}$.  If every pair marginal were $\mathbf1/4$, every swap would have expectation $\tfrac12$,
and the sum would be $n(2n-1)/2=n(n-2)+3n/2$.
\end{proof}

\begin{corollary}[equal one-site states, different correlations]\label{cor:cucorr}
One-site isotropy admits distinct global states, including $\mathbf1/2^{2n}$ of entropy $2n\log2$.  In the traceless
Hermitian operators of a bipartite system the local sectors $A_0\otimes\mathbf1$, $\mathbf1\otimes B_0$ and the correlation sector
spanned by $A_0\otimes B_0$ ($A_0$, $B_0$ traceless) are mutually orthogonal; the correlation sector has dimension
$(d_A^2-1)(d_B^2-1)$.  When both marginals are maximally mixed, the traceless part of the state lies in the
correlation sector; $\rho_A$ is maximally mixed exactly when every non-identity Pauli string supported in $A$ has
expectation zero.
\end{corollary}
\begin{proof}
Products of orthonormal Hermitian bases (normalised identity and traceless elements) are orthonormal by
factorisation of the trace, and a local coefficient of the state is the corresponding marginal
coefficient.
\end{proof}

\begin{proposition}[twirling]\label{prop:cutwirl}
For a finite-dimensional unitary representation $U$ of a compact group with normalised Haar measure
\cite{n1-haar}, $P=\int U(g)\,dg$ is the orthogonal projection onto the fixed vectors, and $\mathcal T(D)=\int U(g)DU(g)^*dg$
maps onto the commutant.  On a complex irreducible representation of dimension $d$, $\mathcal T(D)=\tr(D)\mathbf1/d$; on a
sum of pairwise inequivalent irreducibles $\mathcal T(D)=\sum_a(w_a/d_a)\mathbf1_a$ with $w_a=\tr(P_aD)$.  A tensor product of
irreducibles of independent compact groups is irreducible for the product.  For a singlet of any compact
diagonal symmetry with an irreducible single-site representation, the one-site state is maximally mixed.
\end{proposition}
\begin{proof}
Invariance of Haar measure gives $U(h)P=P$, and inversion gives $P^*=P$; the same invariance makes $\mathcal T(D)$
commute with the action and fixes every commuting operator.  Schur's lemma (kernel and image of an
intertwiner are invariant; an endomorphism of a complex irreducible minus an eigenvalue has a kernel)
gives the block form, and the trace on each block its weight.  Successive averaging over independent
factors sends every simple tensor to a multiple of $\mathbf1$, so the commutant is scalar.  With repeated
irreducible components the commutant acts on the multiplicity space and scalarity needs further data.
\end{proof}

\begin{remark}[twirling does not prepare a pure singlet; spatial symmetry is separate]\label{rem:cutransl}
An average of vectors is a projection, and a nonzero projected vector can be normalised; an average of
density operators can be mixed: for two sites the twirl of $\lvert01\rangle\langle01\rvert$ has weight $\tfrac12$ on the singlet and $\tfrac12$
on the triplet, eigenvalues $(\tfrac12,\tfrac16,\tfrac16,\tfrac16)$.  Invariance under spin rotations and under a group of site
permutations are separate requirements; the actions commute, so their joint invariants are obtained by
averaging the permutations on $\mathcal S_{2n}$.  On a ring with $2n\ge4$ and unit shift $\mathsf T$, $d=s_{12}s_{34}\cdots s_{2n-1,2n}$ has
$\mathsf T^2d=d$, and $d+\mathsf Td$ is nonzero (the $12$ marginals of $d$ and $\mathsf Td$ differ) and $\mathsf T$-invariant; for $2n=2$ the singlet
has shift eigenvalue $-1$.  For $2n\ge4$ no ray is invariant under all permutations: all transpositions are
conjugate, so they would act by a common sign, putting the ray in the symmetric tensors (spin $n$) or the
alternating ones (zero on more than two factors of $\C^2$).  \emph{Isotropy does not give spatial homogeneity}:
the two-site marginals of a dimer product are pure on a matched pair and $\mathbf1/4$ on an unmatched pair.
\end{remark}

\smallskip\noindent\textbf{Hypothesis U (closed-universe state).}  A cosmological use of the register specifies
the physical sites and their local representations, the spatial translations and the joint invariant
state; identifying a closed universe with such a state, and choosing a common comparison frequency or an
observer window, are physical inputs.  Total momentum and gauge charge require representations of their
own groups.  Reading the compact comparison manifold of each system as a closed space, or identifying
field, wave and particle descriptions of one excitation, requires a dynamical map from that geometry to
spacetime and matter (cf.\ Remark~\ref{rem:skintrinsic}).

\subsection{Conditioning on a clock}\label{sec:cuclock}
Membership in $\mathcal S_{2n}$ gives rotational invariance.  For stationarity choose real $J_{kl}=J_{lk}$ (energy units) and
$H_{\mathrm{tot}}=\sum_{k<l}J_{kl}X_{kl}$; Theorem~\ref{thm:originfree} allows the whole Hermitian exchange algebra, and the linear choice
specifies dynamics.  A Hamiltonian commuting with joint transport preserves $\mathcal S_{2n}$.

\begin{proposition}[compatibility with the constraint]\label{prop:cukernel}
Unit constrained singlets ($H_{\mathrm{tot}}\Psi=0$) exist exactly when $0$ is in the spectrum of $H_{\mathrm{tot}}|_{\mathcal S_{2n}}$, and form one
ray exactly when this eigenvalue is simple.  The couplings admitting one form the zero set of a nonzero
polynomial.  For uniform coupling $J$ the restriction is $Jn(n-2)\mathbf1$.  For $n=2$ the kernel on $\mathcal S_4$ has dimension
$0$ or $2$.  For every $n\ge3$ some unshifted couplings select a unique singlet ray.  With a scalar offset, any
singlet eigenvector of eigenvalue $E$ solves $(H_{\mathrm{tot}}-E)\Psi=0$, and the unconstrained evolution changes by the phase
$e^{iEt/\hbar}$.
\end{proposition}
\begin{proof}
The determinant of the restriction is a polynomial in the $J_{kl}$, nonzero at $H_{\mathrm{tot}}=X_{12}$, whose restriction
squares to $\mathbf1$; a nonzero polynomial does not vanish on an open set.  The uniform value is Proposition~\ref{prop:symmetric}.
On $\mathcal S_4$ every swap is conjugate to $X_{12}=\mathrm{diag}(-1,1)$ (Example~\ref{ex:cueight}), so the restriction is a traceless
Hermitian $2\times2$ matrix, invertible or zero.  For $n\ge3$ take the sum of the swaps of one perfect matching
plus $(n-2)^{-1}\sum_{k<l}X_{kl}$: the second term is $n$ on singlets, the first is at least $-n$ with equality exactly
on the product singlet of that matching.
\end{proof}

For a nonempty proper clock register $C\subset Q$ and $S=Q\setminus C$, the edges give
\begin{equation}\label{eq:cusplit}
H_{\mathrm{tot}}=H_C\otimes\mathbf1+\mathbf1\otimes H_S+V,\qquad H_D=\sum_{k<l\in D}J_{kl}X_{kl}\ (D=C,S),\qquad V=\sum_{k\in C,\,l\in S}J_{kl}X_{kl}.
\end{equation}
A clock reading uses a unit seed $c_0\in V_C$ and the covariant vectors $c_t=e^{-itH_C/\hbar}c_0$; put $L_t=\langle c_t\rvert\otimes\mathbf1$,
$\phi_t=L_t\Psi$, $p_t=\lVert\phi_t\rVert^2$ and, for $p_t>0$, $\chi_t=\phi_t/\sqrt{p_t}$.  This is Proposition~\ref{prop:relative} for the rank-one tests
parametrised by $t$; a scalar weight of a rank-one effect density cancels in $\chi_t$.

\begin{theorem}[exchange constraint and conditional dynamics]\label{thm:cucond}
Let $H_{\mathrm{tot}}\Psi=0$, $\lVert\Psi\rVert=1$, and $r_t=L_tV\Psi$.  Then
\begin{align}
i\hbar\dot\phi_t&=H_S\phi_t+r_t,\qquad \dot p_t=\tfrac2\hbar\im\ip{\phi_t}{r_t},\label{eq:cusource}\\
i\hbar\dot\chi_t&=H_S\chi_t+\frac{r_t}{\sqrt{p_t}}-\frac{i\hbar\dot p_t}{2p_t}\chi_t\qquad(p_t>0).\label{eq:cunormalized}
\end{align}
For $V=0$, $p_t=p_0$ and $\phi_t=e^{-itH_S/\hbar}\phi_0$; the vector equation with $H_S$ holds whenever $r_t=0$, and at positive
weight the conditional rays follow $H_S$ on an interval exactly when $r_t$ is parallel to $\phi_t$ there.  For any
unit $\Psi$, with $d=H_{\mathrm{tot}}\Psi$ and $\eta(t)=\lvert t\rvert(\lVert V\rVert+\lVert d\rVert)/\hbar$,
\begin{equation}\label{eq:cuerr}
\lVert\phi_t-e^{-itH_S/\hbar}\phi_0\rVert\le\eta(t),\qquad
\lVert\chi_t-e^{-itH_S/\hbar}\chi_0\rVert\le\frac{2\eta(t)}{\sqrt{p_0}}\quad(\eta(t)<\sqrt{p_0}).
\end{equation}
For a differentiable imperfect readout $\widetilde L_t=\langle\tilde c_t\rvert\otimes\mathbf1$ with unit $\tilde c_t$, the additional source is $B_t\Psi$,
$B_t=i\hbar\dot{\widetilde L}_t+\widetilde L_t(H_C\otimes\mathbf1)$, and $\hbar^{-1}\int_{\min(0,t)}^{\max(0,t)}\lVert B_u\rVert\,du$ is added to $\eta(t)$.
\end{theorem}
\begin{proof}
$i\hbar\dot\phi_t=-L_t(H_C\otimes\mathbf1)\Psi$, and \eqref{eq:cusplit} gives $-L_tH_C\Psi=H_S\phi_t+r_t-L_td$.  Differentiating the norm and
dividing by $\sqrt{p_t}$ give the other identities.  In the frame rotating with $H_S$ the derivative of the vector is
$-(i/\hbar)e^{itH_S/\hbar}(r_t-L_td)$; the ray is constant exactly when this is a multiple of the vector, and integration
with $\lVert L_t\rVert=1$ gives the first bound.  For nonzero $u,v$, $\lVert u/\lVert u\rVert-v/\lVert v\rVert\rVert\le2\lVert u-v\rVert/\lVert v\rVert$ (insert $u/\lVert v\rVert$), with
$\lVert v\rVert=\sqrt{p_0}$; $\eta<\sqrt{p_0}$ ensures $u\ne0$.  For $\widetilde L_t$, $i\hbar\,\tfrac{d}{dt}(\widetilde L_t\Psi)=B_t\Psi-\widetilde L_t(H_C\otimes\mathbf1)\Psi$, and the same
substitution applies.  Each swap has norm one, so $\lVert V\rVert\le\sum_{k\in C,l\in S}\lvert J_{kl}\rvert$.
\end{proof}

With sitewise Pauli matrices, $X_{kl}=\tfrac12(\mathbf1+\sum_a\sigma^a_k\sigma^a_l)$ (Theorem~\ref{thm:originfree}) gives
\begin{equation}\label{eq:cuexsource}
r_t=\tfrac12\sum_{k\in C,\,l\in S}J_{kl}\Bigl(\phi_t+\sum_{a=1}^3\sigma^a_l\,(\langle c_t\rvert\sigma^a_k\otimes\mathbf1)\Psi\Bigr),
\end{equation}
which involves clock contractions other than $\phi_t$.  If $V\Psi=(\mathbf1\otimes W)\Psi$ with $W$ Hermitian, the evolution is
autonomous with $H_S+W$.  On a subspace $\mathcal P$ of constrained vectors an operator $W_t$ with $L_tV|_{\mathcal P}=W_tL_t|_{\mathcal P}$ exists
exactly when $\ker(L_t|_{\mathcal P})\subseteq\ker(L_tV|_{\mathcal P})$ (the map is then well defined on the range and extends), and a
Hermitian one exactly when the form $(\Phi,\Psi)\mapsto\ip{L_t\Phi}{L_tV\Psi}$ is Hermitian on $\mathcal P$ (block extension on the range and
its complement; necessity by taking adjoints).  The internal $S$ edges give the free Hamiltonian; the cross
edges give the corrections through their clock matrix elements.

\begin{proposition}[energy matching and finite readouts]\label{prop:cuspectrum}
Let $V=0$ and $\Psi=\sum_E\Psi_E$, $\Psi_E\in V_C(E)\otimes V_S(-E)$ ($V_D(e)$ the eigenspace of $H_D$).  On a branch with $p_t>0$ the
conditional rays are nonconstant exactly when $(\langle c_0\rvert\otimes\mathbf1)\Psi_E\ne0$ for at least two distinct $E$.  If the occupied
clock energies are $E_0+k\hbar\omega$ ($k\in\Z$), a seed with amplitudes $m^{-1/2}$ on one unit eigenvector for each of $m$
selected energies gives a normalised covariant measurement on their span, with outcome measure $dt/T$
modulo $T=2\pi/\omega$ and effect density $m\lvert c_t\rangle\langle c_t\rvert$; for $m$ consecutive levels the ticks $t_j=jT/m$ form an orthonormal
basis.  For incommensurate gaps the uniform average of $m\lvert c_t\rangle\langle c_t\rvert$ over $[0,\ell]$ has off-diagonal entries of modulus at
most $2\hbar/(\lvert E'-E\rvert\ell)$.  No normalised translation-covariant real-time measurement and no pair with
$[T_C,H_C]=i\hbar\mathbf1$ exists in finite dimension.
\end{proposition}
\begin{proof}
The product eigenbasis diagonalises the additive constraint; the contractions for different $E$ are
orthogonal on $S$ and carry phases $e^{iEt/\hbar}$, so the ray changes exactly when two survive.  With amplitudes
$m^{-1/2}$ the diagonal orbit integrals are $1/m$ and off-diagonal ones vanish by integer frequency
differences; the geometric sum gives orthogonality of the ticks; $\lvert(e^{ix}-1)/(ix)\rvert\le2/\lvert x\rvert$.  A covariant
normalised measurement on $\R$ would give, after the normalised trace, a translation-invariant probability
measure on $\R$, and the integer translates of $[0,1)$ would have equal mass.  The trace of a commutator
is zero.
\end{proof}
\noindent A finite resolution of the reading (an independent offset $u$ with law $\mu$, $V=0$) multiplies the matrix
element of $\chi_t$ between system energies $e,e'$ by $\int e^{-iu(e-e')/\hbar}d\mu(u)$, of modulus at most one; a uniform offset over
a period removes all coherences between distinct matched energies.  This classical smearing differs
from the dynamical cross-edge source.

\begin{example}[a periodic clock made of exchanges]\label{ex:cueight}
On four sites use the unit singlets $\lvert-\rangle=s_{12}s_{34}$ and $\lvert+\rangle=(t_+t_--t_0t_0+t_-t_+)/\sqrt3$, with $t_+=\lvert00\rangle$,
$t_0=(\lvert01\rangle+\lvert10\rangle)/\sqrt2$, $t_-=\lvert11\rangle$ on the pairs $(12)$, $(34)$; the total raising operator sends $t_-\mapsto\sqrt2t_0$ and
$t_0\mapsto\sqrt2t_+$ on each pair, and the raised terms cancel.  On the code basis $(\lvert-\rangle,\lvert+\rangle)$, $X_{12}=X_{34}=\mathrm{diag}(-1,1)$ and $X_{23}$ has
diagonal $(\tfrac12,-\tfrac12)$ and off-diagonal entries $\tfrac{\sqrt3}2$, so $(2X_{23}+X_{12})/\sqrt3$ restricts to $\sigma_x$.  Take $C=\{1,2,3,4\}$,
$S=\{5,6,7,8\}$, $H_{\mathrm{tot}}=J(X_{12}+X_{56})$, $J>0$, and for $0<p<1$
\[
\Psi=\sqrt p\,\lvert+\rangle_C\lvert-\rangle_S+\sqrt{1-p}\,\lvert-\rangle_C\lvert+\rangle_S ,
\]
a singlet with $H_{\mathrm{tot}}\Psi=0$.  With $c_0=(\lvert+\rangle+\lvert-\rangle)/\sqrt2$, $p_t=\tfrac12$ and $\chi_t=\sqrt p\,e^{iJt/\hbar}\lvert-\rangle_S+\sqrt{1-p}\,e^{-iJt/\hbar}\lvert+\rangle_S$, and the system
observable $\lvert-\rangle\langle+\rvert+\lvert+\rangle\langle-\rvert$ has conditional mean $2\sqrt{p(1-p)}\cos(2Jt/\hbar)$.  The seed is the positive eigenvector of
$(2X_{23}+X_{12})/\sqrt3$ on the code, and conjugation by $e^{-itH_C/\hbar}$ gives the readout projectors; the clock
measurement on $C$ is the covariant family of Proposition~\ref{prop:cuspectrum} on the code ($m=2$, period $\pi\hbar/J$, two
orthogonal ticks $\pi\hbar/(2J)$ apart), completed by the projector onto the complement of the code.  All these
effects act inside spin-zero sectors and commute with joint rotations; they are realised by the
instruments of Theorem~\ref{thm:instruments}.  At $p\in\{0,1\}$ the conditional ray is constant; for a one-site clock $H_C=0$
and the orbit is constant even for this constrained entangled state.
\end{example}

\begin{example}[a cross-edge term that changes conditional rays]\label{ex:cucross}
On four sites let $C=\{1,2\}$, $S=\{3,4\}$, $H_C=JX_{12}$, $H_S=0$, $V=\tfrac J2\sum_{k\in C,l\in S}X_{kl}$.  On $\mathcal S_4$, $\sum_{k\in C,l\in S}X_{kl}=-(X_{12}+X_{34})$ and
$X_{12}=X_{34}$, so $V=-H_C$ there, and $(\lvert+\rangle+\lvert-\rangle)/\sqrt2$ solves the constraint.  The seed $c_0=(t_0+s_{12})/\sqrt2$ gives
$\phi_t=\tfrac12e^{-iJt/\hbar}s_{34}-\tfrac1{2\sqrt3}e^{iJt/\hbar}t_0$, $p_t=\tfrac13$: the relative phase changes although $H_S=0$.  This seed mixes inequivalent
spin sectors of $C$; if admissible clock effects must commute with rotations of $C$, it needs an internal
spin reference, which the spin-zero clock of Example~\ref{ex:cueight} avoids.  On two sites $H_{\mathrm{tot}}=JX_{12}$, $J\ne0$, has no
zero-energy singlet; with $H_{\mathrm{tot}}=0$ the singlet solves the constraint and both proper clocks ($H_C=0$) have constant
readouts.  Entanglement, stationarity, covariance and a nontrivial clock are distinct conditions.
\end{example}

\begin{remark}[relational time]\label{rem:cupw}
Page and Wootters introduced conditional evolution from a stationary joint state \cite{n1-pw}, Wootters
the correlation interpretation \cite{n1-wootters}, Giovannetti, Lloyd and Maccone a history-state
formulation \cite{n1-glm}; H\"ohn, Smith and Lock relate the Page--Wootters reduction to relational Dirac
observables \cite{n1-hoehn}, and Smith and Ahmadi treat interacting clocks \cite{n1-smith}.  Here
\eqref{eq:cusource} and \eqref{eq:cuexsource} apply conditioning to the origin-free exchange generators, specify which couplings
enter the conditional equation, and Example~\ref{ex:cueight} gives a finite clock with rotation-invariant readout.  A
change of clock changes the split, the readout and the conditional probabilities; in the additive case a
shift of the seed shifts the time origin, since $e^{-i\tau H_C/\hbar}\otimes e^{-i\tau H_S/\hbar}$ fixes a zero-energy vector.  These are
conditional probabilities at prescribed readings (Proposition~\ref{prop:relative}); single outcomes remain the
question of Section~\ref{sec:outcomes}.
\end{remark}

\subsection{Total phase and the elementary clock}\label{sec:cuphase}

\begin{proposition}[total phase is invisible]\label{prop:cuphase}
In the complex composite reading with the instruments of Theorem~\ref{thm:instruments} (C0--C3, Convention~\ref{conv:T}), $\Psi$ and
$e^{i\gamma}\Psi$ have the same statistics for every internal observer, including conditional and adaptive statistics
on events of positive weight.  Relative phases between coherent components are measurable.
\end{proposition}
\begin{proof}
The projectors coincide; an instrument assigns its outputs from the projector, composition along each
branch gives the same weights and updated states, and partial traces and division by a common positive
weight preserve equality.  Observers whose ancillas are part of the register are included.  The mixing
observable of Example~\ref{ex:cueight} measures a relative phase.
\end{proof}

\smallskip\noindent\textbf{Hypothesis Ph (phase bridge).}  Fix a clock $T=\{\rho_\gamma\}\subseteq\cR$ of an elementary system and its
complex structure, in which $\rho_\gamma$ acts as the scalar $e^{i\gamma}$ (Proposition~\ref{prop:clock}(ii)).  (a) The statistics of
elementary configurations inside every admissible composite are those of the complex composite reading
built on these complex structures, so that $\rho_\gamma$ on one factor acts as the scalar $e^{i\gamma}$ on the composite.
(b) A right transport under which all internal statistics in every admissible surrounding preparation
are unchanged belongs to the unobservable group $K$.

\begin{corollary}[{[Cl]} from the phase bridge]\label{cor:clbridge}
Under Hypothesis~Ph, $T\subseteq K$, which is [Cl].  With [NT], Proposition~\ref{prop:clockB} gives $K=T$ and binarity.
\end{corollary}
\begin{proof}
By multilinearity, a scalar on one factor multiplies every simple tensor, hence every vector, by that
scalar.  Proposition~\ref{prop:cuphase} and Ph(b) place each $\rho_\gamma$ in $K$.
\end{proof}

\begin{remark}[what the corollary does and does not give]\label{rem:clreal}
\emph{Closure and the invisibility of the total phase do not give [Cl]}: the real reading $K=\{\pm1\}$ of
Section~\ref{sec:readings} has closed real composites whose statistics are insensitive to the total sign, and $K$
contains no right-transport circle.  The content of Ph is (a): that composite statistics are computed in
a complex realisation in which the elementary clock is the scalar phase.  In this paper the complex
composite reading is derived for factors that are already read through clocks
(Theorem~\ref{thm:oneclock}); Ph(a) posits it at the level of configurations, so Corollary~\ref{cor:clbridge} replaces the
input [Cl] by the input Ph and does not derive [Cl] from the axiom.  Three circles have different roles:
the scalar circle acts trivially on all rays; the clock of Example~\ref{ex:cueight} has two energies and its reading
varies through a relative phase; the common time shift generated by $H_{\mathrm{tot}}$ is the identity on its kernel.
\end{remark}

\subsection{Restriction to an observer's horizon}\label{sec:cuhorizon}

\smallskip\noindent\textbf{Hypothesis A (accessible algebra).}  An observer has the full matrix algebra on a site set $A$,
embedded by $O\mapsto O\otimes\mathbf1_{\bar A}$.  This models a horizon or a restriction of available records; a geometric
identification with a spacetime horizon requires a causal model.

\begin{proposition}[restriction and conditional ensembles]\label{prop:curestrict}
The restriction of the closed state to the accessible algebra is the unique density operator $\rho_A$ with
$\tr(\rho_AO)=\ip{\Psi}{(O\otimes\mathbf1)\Psi}$.  With Schmidt coefficients $\sqrt{\lambda_a}$ across $A\vert\bar A$ the nonzero spectra of $\rho_A$ and
$\rho_{\bar A}$ are $(\lambda_a)$, so $S(\rho_A)=S(\rho_{\bar A})$ and $\tr\rho_A^2=\sum_a\lambda_a^2$.  After a complete rank-one test of $\bar A$ with
orthonormal $b_i$, the unread marginal is $\sum_ip_i\lvert\chi_i\rangle\langle\chi_i\rvert$, $\phi_i=(\langle b_i\rvert\otimes\mathbf1)\Psi$, $p_i=\lVert\phi_i\rVert^2$, $\chi_i=\phi_i/\sqrt{p_i}$.  A
conditional vector is pure on the whole unmeasured register; its further restrictions can be mixed.
\end{proposition}
\begin{proof}
The partial trace satisfies the identity and matrix units determine $\rho_A$; a singular value decomposition
gives the Schmidt form; expanding in the test basis gives the ensemble, null outcomes included.
\end{proof}
\noindent For a singlet pair, conditioning on $\lvert0\rangle$ has weight $\tfrac12$ and leaves $\lvert1\rangle$, and a further test of
$\cos\alpha\lvert0\rangle+\sin\alpha\lvert1\rangle$ has conditional weight $\sin^2\alpha$.  In a product of two singlets a one-site test leaves the other pair
entangled with maximally mixed one-site marginals: a nonzero branch neither purifies every unmeasured
site nor breaks the homogeneity of the untouched pair.  A finite Gram window or a precision threshold
can enter as a specification of admissible effects.

\begin{example}[entropy of a cut through independent pairs]\label{ex:cupage}
For a product of pure pairs an uncut pair contributes $0$ to $S(\rho_A)$ and a cut pair with Schmidt weights $p,1-p$
contributes $h(p)=-p\log p-(1-p)\log(1-p)$.  For $12$ sites in the nested matching $(1,12),(2,11),\dots,(6,7)$ the prefix
of $m$ sites cuts $\min(m,12-m)$ pairs, with entropy $\min(m,12-m)h(p)$, peak $6\log2$ for singlets.  Pairing every site
of a half-register across the cut gives $n\log2$ irrespective of any geometric boundary assigned to the cut.
Page's average concerns random pure states \cite{n1-page}, and the holographic area law a geometric
setting \cite{n1-rt}; \emph{restriction alone gives no area law}: a spatial area law or a Page curve
requires a relation between dynamics, state and geometry.
\end{example}

\begin{theorem}[finite-dimensional modular flow and KMS]\label{thm:cumodular}
Let $\rho=\rho_A>0$ and $K_A=-\log\rho$.  In the Hilbert--Schmidt representation by left multiplication with vector
$\Omega=\rho^{1/2}$, the Tomita operator $\mathcal S(O\Omega)=O^*\Omega$ has polar decomposition $\mathcal S=\mathcal J\Delta^{1/2}$, $\mathcal J(Y)=Y^*$, $\Delta(Y)=\rho Y\rho^{-1}$,
and the modular automorphisms are $\sigma_s(O)=\rho^{is}O\rho^{-is}=e^{-isK_A}Oe^{isK_A}$.  For $F(z)=\tr(\rho B\sigma_z(D))$, $F$ is entire, bounded
on $-1\le\im z\le0$, and $F(s-i)=\tr(\rho\,\sigma_s(D)B)$.  For singular $\rho$ the statements hold on the support algebra.
\end{theorem}
\begin{proof}
$\mathcal S(Y)=\rho^{-1/2}Y^*\rho^{1/2}$; on matrix units in an eigenbasis of $\rho$, $\Delta$ has eigenvalues $\lambda_i/\lambda_j$ and $\mathcal J\Delta^{1/2}=\mathcal S$.
Conjugation by $\Delta^{is}$ on left multipliers gives $\sigma_s$.  $F$ is a finite sum $\sum\lambda_c^{1-iz}\lambda_b^{iz}(\cdot)$, bounded for
$-1\le\im z\le0$; at $z=s-i$, $\sigma_{s-i}(D)=\rho\sigma_s(D)\rho^{-1}$ and cyclicity gives the identity.  For singular $\rho$ the kernel
projection annihilates $\Omega$, which is then not separating on the full algebra \cite{n1-takesaki,n1-hhw}.
\end{proof}

\begin{corollary}[when modular flow is exchange time]\label{cor:cugibbs}
For Hermitian $H_A$ on the full accessible factor and $0<\beta<\infty$ the following are equivalent for faithful $\rho$:
$\rho=e^{-\beta H_A}/Z$; $K_A=\beta H_A+(\log Z)\mathbf1$; $\sigma_s=\alpha_{-\hbar\beta s}$ with $\alpha_t(O)=e^{itH_A/\hbar}Oe^{-itH_A/\hbar}$; the KMS condition for $\alpha$ at $i\hbar\beta$.
\end{corollary}
\begin{proof}
Logarithms and substitution give the first equivalences; equality of generators makes $K_A-\beta H_A$ central,
hence scalar, fixed by normalisation.  KMS at $t=0$ and nondegeneracy of the trace pairing give
$e^{-\beta H_A}De^{\beta H_A}\rho=\rho D$ for all $D$, so $e^{\beta H_A}\rho$ is scalar.  For singular $\rho$ the criterion applies on a support
preserved by the dynamics; at $\beta=0$ the Gibbs state is $\mathbf1/d$ and its modular flow is trivial.
\end{proof}
\noindent The unread register $S$ of Example~\ref{ex:cueight} has energies $\mp J$ with populations $p,1-p$ on its code, which $H_S$
preserves; for $p>\tfrac12$ the criterion holds at $\beta=(2J)^{-1}\log(p/(1-p))$.  At $p=\tfrac12$ the marginal has trivial modular flow
while its conditioned observable is time dependent.  For $\cos\theta\lvert00\rangle+\sin\theta\lvert11\rangle$, $0<\theta<\pi/4$, and $H_A=\mathrm{diag}(0,\Delta E)$,
$\beta\Delta E=2\log\cot\theta$: with $\beta=1/k_BT$ the matched temperature tends to $0$ as $\theta\downarrow0$ and to $\infty$ as $\theta\uparrow\pi/4$.  The
energy scale is an input; an $H_A$ not commuting with $\rho_A$ has no matched $\beta$.

\begin{remark}[thermal time and continuum horizons]\label{rem:cuthermal}
Connes and Rovelli propose the modular group as physical time \cite{n1-connes}; the finite construction
supplies an accessible algebra and a restricted state, and calibration requires the matching of
Corollary~\ref{cor:cugibbs}.  Bisognano and Wichmann identify vacuum wedge modular flow with boosts \cite{n1-bw};
the Unruh and Hawking temperatures \cite{n1-unruh,n1-hawking} use their spacetime dynamics and time
normalisations.  A comparison-register analogue requires a continuum algebra, a state and a geometric
modular action with its scale; the event geometry of Section~\ref{sec:eventspace} is a candidate kinematic setting
(programme).
\end{remark}

\subsection{Entropy balances and tests}\label{sec:cuentropy}

\begin{proposition}[finite entropy balances]\label{prop:cubalances}
For faithful $\rho$, any density $\tau$ and $D(\tau\Vert\rho)=\tr\tau(\log\tau-\log\rho)$,
\begin{equation}\label{eq:cufirstlaw}
S(\tau)-S(\rho)=\tr\bigl((\tau-\rho)K_A\bigr)-D(\tau\Vert\rho),\qquad K_A=-\log\rho,
\end{equation}
so $\delta S=\tr(\delta\rho\,K_A)$ to first order \cite{n1-blanco}.  For a pinching $\Pi$ onto orthogonal blocks,
$\lVert D_0\rVert_{\mathrm{HS}}^2=\lVert\Pi D_0\rVert_{\mathrm{HS}}^2+\lVert D_0-\Pi D_0\rVert_{\mathrm{HS}}^2$ and $S(\Pi\tau)-S(\tau)=D(\tau\Vert\Pi\tau)\ge0$, with equality exactly when $\tau=\Pi\tau$; along
unitary steps each followed by pinching the entropy increase telescopes into these relative entropies.
\end{proposition}
\begin{proof}
Expanding the traces gives \eqref{eq:cufirstlaw}; $\frac{d}{d\epsilon}\tr f(\rho+\epsilon\delta)=\tr f'(\rho)\delta$ for $f(x)=-x\log x$, and $\tr\delta=0$.  Block-diagonal and
off-block matrices are orthogonal; $\log\Pi\tau$ is block diagonal on its support (which contains that of $\tau$), so
$\tr\tau\log\Pi\tau=\tr\Pi\tau\log\Pi\tau$.  Nonnegativity and equality: with eigenvalues $a_i$, $b_j$ and squared overlaps $w_{ij}$,
$D$ is the classical divergence of $(a_iw_{ij})$ and $(b_jw_{ij})$, and $-\log x\ge1-x$ with equality iff the arrays agree,
which forces equal states.  Unitary steps preserve entropy.
\end{proof}
\noindent For qubit pinching the discarded mass is $\lvert(D_0)_{01}\rvert^2+\lvert(D_0)_{10}\rvert^2$; a state diagonal in a basis mutually unbiased
to the pinching basis goes to $\mathbf1/d$, gaining $\log d-S(\tau)$.

\begin{proposition}[finite-bath erasure]\label{prop:culandauer}
Let memory $M$ and reservoir $R$ start in $\rho_M\otimes\omega_R$, $\omega_R=e^{-\beta H_R}/Z$, $0<\beta<\infty$, and let a joint unitary produce $\rho'_{MR}$.
With $Q_R=\tr H_R(\rho'_R-\omega_R)$ and $\Delta S_D=S(\rho'_D)-S(\rho_D)$,
\[
\beta Q_R=-\Delta S_M+I(M{:}R)'+D(\rho'_R\Vert\omega_R)\ge-\Delta S_M ,
\]
with equality exactly when the output is uncorrelated and $\rho'_R=\omega_R$.  A mixed memory cannot be made pure by
such a protocol \cite{n1-landauer,n1-reeb}.
\end{proposition}
\begin{proof}
Unitary invariance and the product input give $I(M{:}R)'=\Delta S_M+\Delta S_R$; \eqref{eq:cufirstlaw} with $\rho=\omega_R$ gives
$\beta Q_R=\Delta S_R+D(\rho'_R\Vert\omega_R)$.  A pure output memory forces rank at most $\dim V_R$, while the input rank
$\rk(\rho_M)\dim V_R$ is preserved.  For a correlated input with reservoir marginal $\omega_R$, $I'-I$ replaces $I'$.
\end{proof}
\noindent With a qubit reservoir $\mathrm{diag}(0,\Delta E)$ of excited weight $0<q<\tfrac12$, swapping with a maximally mixed memory
gives $Q_R=\Delta E(\tfrac12-q)$ and an excess of $\beta Q_R$ over $\log2-h(q)$ equal to $D(\mathbf1/2\Vert\omega_R)>0$.  A controlled bit flip
copies $(\lvert0\rangle\pm\lvert1\rangle)/\sqrt2$ to $(\lvert00\rangle\pm\lvert11\rangle)/\sqrt2$; without the record both signs give the same local state, with it the
inverse operation restores them, and a surviving copy outside the recovery algebra removes the cross
terms.  Entropy rises under repeated nonselective pinching and falls when a copying operation is
reversed.  In the stationary model of Example~\ref{ex:cueight}, for $A=\{5,7\}$, $X_{57}=X_{56}X_{67}X_{56}$ has diagonal $(\tfrac12,-\tfrac12)$ and
off-diagonal $-\tfrac{\sqrt3}2$ on the code, the singlet weight $w_t=(1-\ip{\chi_t}{X_{57}\chi_t})/2$ is periodic and nonconstant for $0<p<1$,
and the conditional entropy $h(w_t)+(1-w_t)\log3$, strictly concave in $w_t$, oscillates.  A record arrow needs an
initial record condition and a mechanism for continued loss of access.

\begin{remark}[equilibrium, cooling and a fixed point]\label{rem:cuarrow}
Two diagonal qubits with the same gap and different populations each satisfy KMS for their own
modular flow, yet a swap transfers energy (gap times population difference): \emph{modular KMS is no
zeroth law}, and a physical temperature needs shared calibrated dynamics and a class of contacts.  The
rank obstruction concerns closed unitary erasure with a full-rank bath: a swap with an additional pure
register prepares an exactly pure state, and entanglement can persist elsewhere while a selected site is
pure \cite{n1-mo}.  A gravitational use of the entanglement first law needs local equilibrium, a
geometric entropy law and a continuum modular generator \cite{n1-jacobson}.  A finite benchmark: $x=\cos x$
has a unique solution in $[0,1]$ ($\cos$ maps $[0,1]$ into itself with Lipschitz constant $\sin1<1$), lying in
$(0,\pi/4)$, so it is a Schmidt angle of the two-level example: $x\simeq0.73908513$, $h(\sin^2x)\simeq0.688864$,
$k_BT/\Delta E\simeq5.390324$, $\log2-h(\sin^2x)\simeq0.004284$.  A dissipative model selecting this angle would have to derive its update
law, its records and its free-energy supply (programme).
\end{remark}

\begin{remark}[interference and spatial observables]\label{rem:cuinterf}
Equal amplitudes with relative phase $q=\Delta k\cdot x-\omega t$ give $\lvert1+e^{iq}\rvert^2/4=(1+\cos q)/2$; for $\Delta k\ne0$ a constant-phase
plane moves with signed speed $\omega/\lvert\Delta k\rvert$; two planar reciprocal vectors of length $2\pi/a$ at angle $\vartheta$ give the
period $a/(2\lvert\sin(\vartheta/2)\rvert)$; a centred aperture of length $L$ along $\Delta k$ multiplies the cosine by $\sin(\lvert\Delta k\rvert L/2)/(\lvert\Delta k\rvert L/2)$.  An
irrational rotation of the circle averages every continuous function to its Haar mean.  Band structure,
signal velocity or a stress tensor require dynamics.  The trace-free part $T_{\mu\nu}-g_{\mu\nu}T^\alpha{}_\alpha/4$ is unchanged by $T_{\mu\nu}\mapsto T_{\mu\nu}+cg_{\mu\nu}$;
a cosmology using this cancellation must derive its field equations, their integration constant and a
positive observed value.  A sector annihilated by a readout is invisible to it; a dark-matter reading
would need a conserved stress tensor, clustering, lensing and abundances.  A spatially invariant global
state could supply homogeneous initial data; explaining the horizon problem requires a mechanism
selecting it and a derivation of the primordial spectrum.  These are a programme.
\end{remark}

\begin{conjecture}[local equilibration in the singlet sector]\label{conj:cueth}
For a periodic chain of $2n$ qubits with independent nearest- and next-nearest-neighbour exchange couplings
uniform in $[0.9,1.1]$ and $[0.45,0.55]$, consider the singlet eigenvectors within $n^{1/4}$ of the normalised singlet
trace mean.  The probability that this window is nonempty tends to one, and, conditional on that, for a
uniformly chosen eigenvector of an orthonormal eigenbasis in the window the marginal of two specified
neighbouring sites converges to $\mathbf1/4$ in trace norm in probability.
\end{conjecture}
\noindent Evidence.  The normalised singlet projector commutes with permutations and rotations; its swap
expectation is $(n-2)/(2n-1)$, its pair singlet weight $(n+1)/(2(2n-1))$, and its trace-norm distance to $\mathbf1/4$ is
$3/(2(2n-1))$.  For $n=3,4,5,6$ and $20$ coupling samples each (generator seed $441$), the window contains $2$, $4$--$5$,
$13$--$15$ and $43$--$47$ levels, and the mean pair distances are $0.432733$, $0.251560$, $0.209589$, $0.165086$.  A proof
needs control of the level count and of eigenvector fluctuations in the window \cite{n1-eth}.

\begin{remark}[what does not hold]\label{rem:cunogo}
(i) The two-site marginals of a global singlet are not all maximally mixed: a matched pair of a
dimer product is pure, and the swap sum excludes jointly flat pairs (Theorem~\ref{thm:cusinglet}).  (ii) Isotropy does
not give spatial homogeneity (Remark~\ref{rem:cutransl}).  (iii) Twirling does not prepare a pure singlet: the
twirl of $\lvert01\rangle\langle01\rvert$ has spectrum $(\tfrac12,\tfrac16,\tfrac16,\tfrac16)$.  (iv) A singlet does not solve the unshifted exchange
constraint in general: the two-site singlet has energy $-J$ for $H=JX_{12}$ (Proposition~\ref{prop:cukernel}).  (v) Not every
subsystem evolves by its complement's Hamiltonian: a one-site clock has $H_C=0$ and a constant orbit, and
in Example~\ref{ex:cucross} the conditional ray changes with $H_S=0$.  (vi) A nonzero branch need not purify every
unmeasured site (after Proposition~\ref{prop:curestrict}).  (vii) Closure with invisible total phase does not give [Cl]
(Remark~\ref{rem:clreal}).  (viii) Restriction alone gives neither an area law (Example~\ref{ex:cupage}) nor a physical
temperature (Remark~\ref{rem:cuarrow}).  (ix) A stationary pure register need not have frozen observations
(Example~\ref{ex:cueight}), and local entropy obeys no one-way law under reversible protocols.  (x) Remaining
entanglement elsewhere does not preclude a pure local state, and a pure ancilla can prepare one.
\end{remark}

\begin{remark}[programme]\label{rem:cuprogramme}
Derive a record boundary condition and loss mechanism for an arrow of time; a continuum observable
algebra with a geometric modular action and a physical scale on the event geometry; a gravitational
first-law limit; a primordial spectrum with acoustic phase correlations; field equations with the
residual vacuum energy; a conserved dark sector with clustering and lensing; and, for a spacetime area
law or a black-hole Page curve, a controlled local model with a continuum limit preserving entropy or
modular information.
\end{remark}


%=====================================================================================
\section{Finite memory: Gram windows, the arrow of forgetting, erasure and consensus}\label{sec:finitemem}
How does a finite record constrain the next comparison, and when can another observer recover or agree
with what was recorded?  A Gram window determines an exact extension set, and retaining more comparisons
weakly decreases the squared norm available outside their span, with an explicit equality criterion and a
projection limit.  Physical forgetting costs heat only through a memory reset; loss of access to coherent
records has a different criterion, set by the states left in the inaccessible register.  Local
comparison windows of several observers have a common realisation exactly when the partial Gram matrix
has a positive semidefinite completion; a common Gram matrix gives neither a joint state nor joint
outcomes.  These results refine the restriction model of Proposition~\ref{prop:curestrict}.

\subsection{Gram windows and loss of constraints}\label{sec:fmwindows}

\smallskip\noindent\textbf{Hypothesis Gw (finite Gram window).}  An observer stores, at time $t$, the Gram matrix
$G_M=(\ip{v_i}{v_j})_{i,j}$ of the unit representatives $v_j$ of its last $M$ records in a real or complex Hilbert space $V$
(phase conventions are part of the stored data); records are ordered newest to oldest, so depth $M$ uses the
first $M$ of a fixed list.  Write $B_Ma=\sum_{j\le M}v_ja_j$, $S_M=\mathrm{ran}\,B_M$, $P_M$ the orthogonal projection onto $S_M$, and for
proposed overlaps $c\in\mathrm{ran}\,G_M$ the \emph{Schur slack} $\Delta(M;c)=1-c^*G_M^+c$ ($+$: Moore--Penrose inverse; $\Delta(0;\cdot)=1$).  Finite
depth is a storage or sampling assumption.

\begin{theorem}[extension, weak nesting and saturation]\label{thm:fmschur}
For $G=G_M$: $\bigl(\begin{smallmatrix}G&c\\ c^*&1\end{smallmatrix}\bigr)\succeq0$ iff $c\in\mathrm{ran}\,G$ and $\Delta(M;c)\ge0$; the bordered rank is $\rk G+\mathbf1_{\{\Delta>0\}}$, the
bordered determinant is $\det G\cdot\Delta$ for invertible $G$, and the unit vectors with overlaps $c$ are $BG^+c+w$, $w\in S_M^\perp$,
$\lVert w\rVert^2=\Delta$ (so positive slack in a fixed ambient space also needs $S_M^\perp\ne0$).  For a fixed unit $x$ and $c_M=B_M^*x$,
\[
\Delta(M;c_M)=\lVert(\mathbf1-P_M)x\rVert^2,\qquad \Delta(M;c_M)-\Delta(M{+}1;c_{M+1})=\lVert(P_{M+1}-P_M)x\rVert^2\ge0 ,
\]
so the corridors $\{x:\lVert(\mathbf1-P_M)x\rVert^2\ge\delta\}$ are \emph{weakly} nested (deleting anchors reverses the inclusion), and the slack
decreases to $\lVert(\mathbf1-P_\infty)x\rVert^2$, $P_\infty$ the projection onto the closed union of the spans (eventually constant in finite
dimension).
\end{theorem}
\begin{proof}
For $a\in\ker G$, positivity on $(sa,1)$ for all scalars $s$ forces $a^*c=0$.  With $b=G^+c$ the form on $(a,z)$ is
$(a+bz)^*G(a+bz)+\lvert z\rvert^2\Delta$, and the shear has determinant one.  $B^*B=G$ gives $\ker B=\ker G$, $B^*BG^+c=c$, $\lVert BG^+c\rVert^2=c^*G^+c$, and
other solutions differ by $w\in\ker B^*=S_M^\perp$.  For actual overlaps the minimum-norm solution is $P_Mx$, and
$S_{M+1}=S_M\oplus(S_{M+1}\cap S_M^\perp)$ gives the increment; approximation of $P_\infty x$ from some $S_{M_0}$ gives the limit
\cite{n2-schur,n2-albert}.
\end{proof}

\begin{corollary}[overlap and cap corridors]\label{cor:fmcaps}
With one extra dimension available the feasible overlaps form the ellipsoid $G^{1/2}\{z\in\mathrm{ran}\,G:\lVert z\rVert\le1\}$ (volume
$\pi^{M/2}\sqrt{\det G}/\Gamma(1+M/2)$ for real invertible $G$), and the deeper ellipsoid projects onto the old one; if the
anchors span $V$, only $\lVert z\rVert=1$ is allowed.  For real directions on $S^2$ and $0<\alpha<\pi/2$, the caps
$\Omega_M=\{x:v_j\cdot x\ge\cos\alpha,\ j\le M\}$ satisfy $\Omega_{M+1}\subseteq\Omega_M$ and $\lvert\Omega_M\rvert\downarrow\lvert\bigcap\Omega_M\rvert$; two anchors more than $2\alpha$ apart give $\Omega=\varnothing$,
and a nested chain with empty intersection is empty at a finite stage; a nonempty $\Omega_M$ ($M\ge1$) is
geodesically convex and can have zero area.  Removing an anchor enlarges $\Omega$, strictly iff a newly
admitted direction violated its inequality.
\end{corollary}
\begin{proof}
Spectral coordinates and Theorem~\ref{thm:fmschur}; continuity of measure from above; compactness; the angular
triangle inequality; the normalised chord between two feasible points lies on their shorter arc and
normalisation only increases positive lower bounds.
\end{proof}

\begin{remark}[growth does not inevitably lock; sliding need not tighten]\label{rem:fmnolock}
Nesting is weak and can be stationary: identical anchors give a constant cap of area $2\pi(1-\cos\alpha)$, and in $\ell^2$
the orthogonal candidate $x=e_1$ with anchors $v_j=e_{j+1}$ has slack one at every depth; the limit is a squared projection
residual, and a physical Zeno locking requires dynamics (Example~\ref{ex:fmzeno}).  A sliding window need not
tighten: two consecutive antipodal single-anchor caps are disjoint, hence incomparable; deleting a
duplicate anchor changes nothing.
\end{remark}

\begin{example}[saturation, release and fixed dimension]\label{ex:fmgram}
For real $a=\ip{v_1}{v_2}$, $b=\ip{v_1}x$, $c=\ip{v_2}x$ the bordered determinant is $1+2abc-a^2-b^2-c^2=(1-a^2)(1-b^2)-(c-ab)^2$.  For $a=b=\tfrac45$
the feasible range is $\tfrac7{25}\le c\le1$, with determinant $\tfrac{13}{125}$ at $c=\tfrac45$ and $-\tfrac7{25}$ at $c=0$; deleting the second anchor gives
$1-b^2=\tfrac9{25}$ and releases $c$.  For $a=b=\tfrac12$ the endpoints are $-\tfrac12$ and $1$.  For $c_M=\tfrac12\mathbf1$ and $G_M=\tfrac23I+\tfrac13\mathbf1\mathbf1^*$:
$\det G_M=(\tfrac23)^{M-1}(M+2)/3$ and $\Delta=(M+8)/(4(M+2))\downarrow\tfrac14$; for orthogonal anchors $\Delta=1-M/4$, with the unique solution $\tfrac12\sum v_j$
at $M=4$ and none for $M\ge5$.  Three independent real anchors span $\R^3$ and their Gram matrix fixes their
geometry up to $\OO(3)$, reflections included \cite[Prop.~0.25]{PaperI}: three anchors fix no chirality, and
their caps can contain an open patch while every actual direction has zero slack.  Two caps whose
centres are exactly $2\alpha$ apart meet in one point.
\end{example}

\begin{proposition}[memory capacity]\label{prop:fmcapacity}
A $d$-dimensional register holds at most $d$ perfectly distinguishable density matrices; a Gram window of
pure descriptions can index more than $d$ records and has rank at most $d$; $\ip uv=1$ iff $u=v$ for unit vectors.
\end{proposition}
\begin{proof}
A perfectly discriminating POVM forces pairwise orthogonal supports; $\lVert u-v\rVert^2=2-2\re\ip uv$.
\end{proof}

\subsection{Forgetting and heat}\label{sec:fmheat}

\begin{proposition}[restriction and thermodynamic reset]\label{prop:fmheat}
(i) Deleting comparison data implements the readout restriction of a partial trace when the records lie
in distinct tensor factors and the retained comparisons act on the retained factors; then every retained
expectation is preserved (Proposition~\ref{prop:curestrict}).
(ii) Let $L$ be an old memory, $K$ retained side information, $R$ a finite reservoir, initially $\rho_{LK}\otimes\omega_R$ with
$\omega_R=e^{-\beta H_R}/Z$ (Gibbs initialisation), $\beta=1/k_BT>0$; let a joint unitary give $\rho'_{LKR}$ with $\rho'_K=\rho_K$ (side marginal
kept).  With $Q_R=\tr H_R(\rho'_R-\omega_R)$ and $S(L|K)=S(LK)-S(K)$,
\[
\beta Q_R+\Delta S(L|K)=D(\rho'_{LKR}\Vert\rho'_{LK}\otimes\omega_R)\ge D(\rho'_R\Vert\omega_R)\ge0,\qquad\text{so}\qquad Q_R\ge k_BT\,[S(L|K)-S'(L|K)],
\]
with equality iff $\rho'_{LKR}=\rho'_{LK}\otimes\omega_R$.  For a correlated input with Gibbs reservoir marginal, $-I(LK{:}R)_{\mathrm{in}}$ is added on the
right of $\beta Q_R+\Delta S(L|K)=I(LK{:}R)'+D(\rho'_R\Vert\omega_R)$.
\end{proposition}
\begin{proof}
(i) is the defining property of the partial trace.  (ii) On the support, $D=-S(\rho'_{LKR})+S(\rho'_{LK})+\beta\tr H_R\rho'_R+\log Z$; unitary invariance,
$S(\omega_R)=\beta\tr H_R\omega_R+\log Z$ and $\rho'_K=\rho_K$ give the identity; monotonicity of relative entropy under $\tr_{LK}$ and under the
trace gives the inequalities \cite{n2-lindblad,n2-uhlmann-dpi}, and the equality case is that of
Proposition~\ref{prop:cubalances}.  For trivial $K$ this is Proposition~\ref{prop:culandauer} and Reeb--Wolf
\cite[Thm.~3]{n1-reeb}; the one-bit case is Landauer's bound $k_BT\log2$ \cite{n1-landauer,n2-bennett}.
\end{proof}

\begin{remark}[index restriction costs no heat]\label{rem:fmzeroheat}
\emph{Pure index restriction or discarding a record can cost zero heat}: with degenerate Hamiltonians a
controlled copy into a pure record realises the pinching $\Pi(\rho)=\sum P_j\rho P_j$ by a joint unitary, and restriction then
raises the local entropy by $h(\sin^2\theta)$ for $(\cos\theta,\sin\theta)$ (and removes squared Hilbert--Schmidt mass $\tfrac12\sin^22\theta$; maximal
$\log2$ and $\tfrac12$ at $\theta=\pi/4$; in dimension $d$ a state diagonal in a basis unbiased to the pinching goes to $\mathbf1/d$) with zero
reservoir energy change; the copy can be uncomputed while accessible.  The bound of
Proposition~\ref{prop:fmheat} concerns a physical reset: reusing a maximally mixed bit costs at least $k_BT\log2$, and
exact purification by a finite faithful bath is excluded by the rank argument of
Proposition~\ref{prop:culandauer}.  This is the arrow of forgetting: a record arrow under repeated discard and
reuse, with these resources.
\end{remark}

\smallskip\noindent\textbf{Hypothesis Ue (uniform phase-space encoding).}  If deletion enlarges a cap corridor from area $A>0$ to
$A'$, the old memory was encoded and reset with $S(L|K)-S'(L|K)=q:=\log(A'/A)$ (e.g.\ uniform classical memories on $m$ and $m'$
cells with $m/m'=A'/A$).  Under Gibbs initialisation, joint unitarity and a kept side marginal the reset then costs
at least $k_BTq$.  The identification of entropy drop and geometry is this hypothesis; it is unavailable at $A=0$.

\subsection{Coherent erasure and the accessible algebra}\label{sec:fmerasure}

\smallskip\noindent\textbf{Hypothesis Ca (causal record access).}  In the tensor-factor form of Hypothesis~A
(Section~\ref{sec:cuhorizon}), the observer acts jointly on the system $S$ and an accessible record $A$, and the
inaccessible record $E$ is traced out.  In the static-patch model the two-way region of an observer is
$W_x=I^+(\gamma_x)\cap I^-(\gamma_x)$ (Proposition~\ref{prop:static}), and records outside $J^-(\gamma_x)$ cannot signal to it; assignment to $A$ also
needs availability within the experiment and coherent control.  The cap geometry specifies neither a
register factorisation nor dynamics.

\begin{theorem}[exact recovery of an encoded phase]\label{thm:fmeraser}
Let a controlled isometry encode a qubit by $\lvert i\rangle\mapsto\lvert i\rangle\lvert r_i\rangle_{AE}$, $X=\tr_E\lvert r_0\rangle\langle r_1\rvert$, $\sigma_i=\tr_A\lvert r_i\rangle\langle r_i\rvert$.  For the
equal-amplitude inputs of opposite sign,
$\tfrac12\lVert\rho^+_{SA}-\rho^-_{SA}\rVert_1=\lVert X\rVert_1=\lVert\sigma_0^{1/2}\sigma_1^{1/2}\rVert_1$.  One channel on $SA$ recovers every input qubit exactly iff $\sigma_0=\sigma_1$ (equal
inaccessible conditional states); orthogonal supports make all phases indistinguishable on $SA$.  For product
records $r_i=a_i\otimes e_i$, $\gamma=\ip{e_1}{e_0}$, input $\sqrt p\lvert0\rangle+e^{i\chi}\sqrt{1-p}\lvert1\rangle$ and a record effect $0\le F\le\mathbf1$ on $A$ with
$a=\ip{a_0}{Fa_0}$, $b=\ip{a_1}{Fa_1}$, $d=\ip{a_1}{Fa_0}$, the outcome has probability $pa+(1-p)b$ and conditional fringe visibility
$2\sqrt{p(1-p)}\lvert d\gamma\rvert/(pa+(1-p)b)\le\lvert\gamma\rvert$; orthogonal accessible markers read in the diagonal basis give two outcomes of probability $\tfrac12$
with opposite fringes of visibility $2\sqrt{p(1-p)}\lvert\gamma\rvert$.  The unread system marginal is invariant under every
record-only channel.
\end{theorem}
\begin{proof}
The difference of the two accessible states is $\bigl(\begin{smallmatrix}0&X\\ X^*&0\end{smallmatrix}\bigr)$.  With coefficient matrices $C_i$, $X=C_0C_1^*$ and $\sigma_i=(C_i^*C_i)^{\mathsf T}$,
whose polar factors preserve the singular values of $\lvert C_0\rvert\lvert C_1\rvert$.  $\lVert\sqrt{\sigma_0}\sqrt{\sigma_1}\rVert_1=\max_U\lvert\tr U\sqrt{\sigma_0}\sqrt{\sigma_1}\rvert\le1$ with equality iff
$\sigma_0=\sigma_1$; a channel contracts trace distance, so exact recovery needs distance one, and for $\sigma_0=\sigma_1$ a
unitary on $A$ conditioned on label $1$ maps $r_1$ to $r_0$.  For $F$ the unnormalised state has diagonal $(pa,(1-p)b)$ and
off-diagonal $e^{-i\chi}\sqrt{p(1-p)}\,d\gamma$; $\lvert d\rvert\le\sqrt{ab}$ and the arithmetic--geometric mean inequality.  Record-only channels fix $\mathbf1$ in
the dual.
\end{proof}

\begin{remark}[forgetting or replay does not universally restore interference]\label{rem:fmcopies}
A discarded orthogonal marker makes opposite phases identical on the accessible algebra, and no later
operation on $SA$ restores them.  With product copies the coherence is multiplied by $\prod_j\gamma_j$; controlled
inverse couplings on accessible copies remove their factors; a conjugate replay on a subset multiplies by
$\prod\overline{\gamma_j}$, so the all-copy replay gives $\prod\lvert\gamma_j\rvert^2$, which restores interference only if every $\lvert\gamma_j\rvert=1$.  For two copies with
records $((1,0),(i,0))$ the coherence $\tfrac12$ becomes $-\tfrac12$, then $-\tfrac i2$ after one replay and $\tfrac12$ after both; a zero overlap
makes replay ineffective.  Orthogonal accessible markers instead support the conditional eraser
protocol of Theorem~\ref{thm:fmeraser} \cite{n2-scully,n2-herzog,n2-kim}; conditional fringes need the outcome
label, and summing outcomes returns the original marginal.
\end{remark}

\subsection{Agreement between observers}\label{sec:fmagreement}

\begin{theorem}[gluing local windows]\label{thm:fmglue}
Let observers supply Hermitian Gram matrices on windows of a finite labelled record set, agreeing on
common entries, with all diagonals specified, and let the graph join labels whose comparison is
specified.  A simultaneous realisation by unit vectors exists iff the partial matrix has a positive
semidefinite completion with unit diagonal (of rank at most $d$ in dimension $d$).  If the graph is chordal, a
completion exists iff every fully specified clique submatrix is positive semidefinite, and the minimum
rank is the largest clique rank \cite{n2-dirac,n2-grone}; on every nonchordal graph there are locally
positive definite data with no completion.
\end{theorem}
\begin{proof}
A completion's square root gives a realisation.  For chordal graphs eliminate a simplicial vertex,
realise the rest in dimension $d$ (the largest clique rank) by induction, realise the vertex's clique, and
match the two realisations of the common neighbours by an isometry of their spans extended to a unitary
(\cite[Thm.~7]{n2-grone} and closure).  Every clique of the union graph must be tested: separate windows can
supply the three edges of a triangle that no window contains.  For an induced cycle of length $k\ge4$, put
$r=1-1/(2k^2)$ on all edges but one and $-r$ on the closing edge (other off-diagonals zero): cliques are positive
definite, consecutive distances are $\sqrt{2-2r}=1/k$, so the endpoints are within $(k-1)/k<1$, whereas $-r$ requires distance
$\sqrt{2+2r}>1$.
\end{proof}
\noindent This is the completion structure of the companion's Part~0: a triangle-flat comparison on a chordal
graph has a potential \cite[Thm.~0.12, Prop.~0.18]{PaperI}, and cycles carry the obstruction
\cite[Sec.~0.4]{PaperI}; Gram realisation and its uniqueness up to a unitary are
\cite[Props.~0.24, 0.25]{PaperI}.

\begin{example}[comparison, marginal and outcome consistency]\label{ex:fmoutcomes}
A four-cycle with three edges $+1$ and one edge $-1$ has positive semidefinite edge windows and no realisation
(the first three force equal vectors); $I_4$ has a completion but none in dimension three.  Two realised
directions $v,-v$ already give $\Omega_A\cap\Omega_B=\varnothing$ for every $\alpha<\pi/2$.  \emph{A common Gram does not yield a joint state}: if $AB$
and $AC$ are both assigned the pure Bell state, a positive extension with pure $AB$ marginal is a product
across $AB|C$, so its $AC$ marginal is a product, a contradiction.  \emph{Nor joint outcomes}: for $\lvert\Phi^+\rangle$ with
$A_1=Z$, $A_2=X$, $B_{1,2}=(Z\pm X)/\sqrt2$ the CHSH sum is $2\sqrt2$, exceeding the bound $2$ of any joint $\pm1$ assignment
\cite{n2-chsh}, although the four correlations have a real Gram realisation by the corresponding
axes in a plane.  Extended Wigner-friend protocols locate the outcome question with further assumptions
\cite{n2-brukner,n3:Bong}.
\end{example}

\begin{example}[contextuality]\label{ex:fmcontext}
In the Peres--Mermin square with rows $(X\otimes\mathbf1,\mathbf1\otimes X,X\otimes X)$, $(\mathbf1\otimes Z,Z\otimes\mathbf1,Z\otimes Z)$, $(X\otimes Z,Z\otimes X,Y\otimes Y)$ every row and
column consists of commuting observables, with products $+\mathbf1$ except the last column ($-\mathbf1$); a sign assignment
respecting the context products would give $+1$ and $-1$ for the product of all six context products
\cite{n2-peres,n2-mermin}.  The obstruction concerns signs on overlapping contexts, not inner products.
\end{example}

\begin{theorem}[perfect records select an operational basis]\label{thm:fmpointer}
If the reduced channel to $S$ and fragments $F_1,\dots,F_N$ is $\rho\mapsto\sum_i\rho_{ii}\lvert i\rangle\langle i\rvert\otimes\bigotimes_j\sigma_i^j$ with $\sigma_i^j$ of pairwise
orthogonal supports for each $j$, every observer reads the same label by a local measurement without
disturbing the output, and among rank-one bases preserved by the system marginal the pointer projectors
are unique up to permutation.  The same holds for a coherent controlled copy when at least one disjoint
perfect copy is discarded.
\end{theorem}
\begin{proof}
Project on the conditional supports (completing the test); the marginal on $S$ is pinching, whose rank-one
fixed projectors are diagonal; a discarded orthogonal copy kills the cross terms.
\end{proof}

\begin{remark}[generic broadcast; cited]\label{rem:fmbph}
Environment-induced selection and quantum Darwinism \cite{n2-zurek81,n2-zurek09,n2-blumekohout}
motivate the distinction between stability and redundancy.  Brand\~ao, Piani and Horodecki's
generic-channel theorem \cite[Thm.~1]{n2-bph} gives, for fixed input dimension and $0<\delta<1$, one POVM and at
least $(1-\delta)N$ of $N$ outputs whose marginal channels are $o(1)$-close in diamond norm to measure-and-prepare channels
$\rho\mapsto\sum_i\tr(M_i\rho)\sigma_i^j$; a constant-output channel satisfies this without revealing a label.  Perfect distinguishability
and nondemolition (Theorem~\ref{thm:fmpointer}) give the operational claim, in the spectrum-broadcast setting
\cite{n2-horodecki}.
\end{remark}

\begin{proposition}[redundancy in a two-branch model]\label{prop:fmredundancy}
For the product-branch encoding of Theorem~\ref{thm:fmeraser} with copies $J$, weights $p,1-p$, overlaps $\gamma_j$, $g_F=\prod_{j\in F}\gamma_j$ and
$H(z)=h\bigl((1-\sqrt{1-4p(1-p)(1-\lvert z\rvert^2)})/2\bigr)$: $I(S{:}F)=H(g_J)+H(g_F)-H(g_{J\setminus F})$, which is $0$ at $F=\varnothing$, $2S(\rho_S)$ at $F=J$, and tends to $h(p)$ when both
$g_F$ and $g_{J\setminus F}$ tend to zero.  For $p=\tfrac12$, $N=16$, $\lvert\gamma_j\rvert=0.3$ and $3\le\lvert F\rvert\le13$, $\lvert I(S{:}F)-S(\rho_S)\rvert<5.26\times10^{-4}$ bits.
\end{proposition}
\begin{proof}
The reduced fragment state is $CC^*$ with $C^*C$ of trace one and determinant $p(1-p)(1-\lvert g_F\rvert^2)$; purity gives $S(SF)=S(J\setminus F)$.  At
$p=\tfrac12$, $H(0)-H(z)=\int_0^{\lvert z\rvert}\mathrm{artanh}\,t\,dt\le\lvert z\rvert^2/2+\lvert z\rvert^4/[12(1-\lvert z\rvert^2)]$, and both moduli are at most $0.3^3=0.027$.
\end{proof}

\begin{theorem}[Bures--Wasserstein barycentre]\label{thm:fmbary}
For positive definite $A_j$ and weights $w_j>0$ summing to one, $f(X)=\sum_jw_j[\tr X+\tr A_j-2\tr(A_j^{1/2}XA_j^{1/2})^{1/2}]$ has a unique positive
semidefinite minimiser $W$, which is positive definite and characterised by $T_jWT_j=A_j$, $T_j>0$, $\sum_jw_jT_j=\mathbf1$
\cite[Thm.~8, Eq.~(57)]{n2-bjl}; on the trace-one slice the minimiser is unique and positive definite, with
$\sum_jw_jT_j=\lambda\mathbf1$.
\end{theorem}
\begin{proof}
$2\tr(A^{1/2}XA^{1/2})^{1/2}\le\tr(XT)+\tr(AT^{-1})$ for $T>0$, with equality iff $TXT=A$ (minimise $\lVert T^{1/2}X^{1/2}-T^{-1/2}A^{1/2}U\rVert_{\mathrm{HS}}^2$ over unitaries), so twice the
trace square root is an infimum of affine majorants and strictly concave on segments of positive definite
matrices.  $f$ is coercive; at singular $X$, $X+t\mathbf1$ gains at least $\sqrt{t\lambda_{\min}(A_j)}$ per term against a linear cost, also after
renormalisation to trace one.  The derivative is $\tr(T\,\cdot)$, giving the balance conditions, and a positive solution
gives a supporting affine bound, hence the minimiser.
\end{proof}

\begin{remark}[a unique barycentre does not single out a pointer basis]\label{rem:fmaxisbary}
For commuting inputs with eigenvalues $a_{ji}$, $s_i=\sum_jw_j\sqrt{a_{ji}}$, the barycentre has eigenvalues $s_i^2$ (normalised on the trace-one
slice); a simple spectrum fixes the common eigenbasis.  But with all inputs $\mathbf1/d$ the unique barycentre is $\mathbf1/d$
and every basis is an eigenbasis: uniqueness of the barycentre does not select a basis, and objective
records need a recording channel (Theorem~\ref{thm:fmpointer}).  Scalar inputs $1,9$ with equal weights give
$f=(\sqrt x-2)^2+1$, minimum $1$ at $x=4$, whereas the arithmetic mean $5$ gives $10-4\sqrt5>1$.  For independently chosen qubit
test axes $n_j$ with weights $w_j$, mixed with the identity at transverse multiplier $\lambda$, the mean Bloch vector is
$(\lambda\mathbf1+(1-\lambda)M)^kr_0$, $M=\sum w_jn_jn_j^{\mathsf T}$; for $\lambda<1$ a simple top eigenvalue gives a slowest-decaying line, an exactly
preserved direction exists iff all axes are parallel, isotropic axes ($M=\mathbf1/3$) give none, and ordered steps
differ from mean steps by $(n\cdot m)(nm^{\mathsf T}-mn^{\mathsf T})$.
\end{remark}

\subsection{Sequential control and geometric phase}\label{sec:fmcontrol}

\begin{proposition}[instruments, readout and intervention]\label{prop:fminstrument}
A history-dependent instrument $\{\mathcal I_i\}$ followed by feedback $U_i$ updates $\rho\mapsto U_i\mathcal I_i(\rho)U_i^*$; traces are branch weights,
the sum is the unread channel, and composition gives the joint law and the conditional states of
Proposition~\ref{prop:relative} and Theorem~\ref{thm:instruments}.  First-step effects do not determine the joint
law (Kraus operators $P_i$ and $XP_i$ have equal effects, and a following test repeats or reverses the
outcome).  Kraus operators with $\sum K_i^*K_i=P$ a projection complete to the channel
$X\mapsto\sum K_iXK_i^*+\tr((\mathbf1-P)X)\lvert v\rangle\langle v\rvert$, agreeing on inputs supported in $P$ (for $K_1=P=\lvert0\rangle\langle0\rvert$, $v=\lvert0\rangle$: the reset).
\end{proposition}
\begin{proof}
Trace preservation and cyclicity; the added Kraus operators $\lvert v\rangle\langle u_j\rvert(\mathbf1-P)$.
\end{proof}

\begin{example}[control with a fixed reference]\label{ex:fmaction}
For $U_\theta=e^{-i\theta Y/2}$, $P=(\mathbf1+Z)/2$, $\rho=(\mathbf1+X)/2$: rotating the state gives readout $(1-\sin\theta)/2$, rotating the effect $(1+\sin\theta)/2$;
pinching before the rotation gives $\tfrac12$; the reset $\sigma\mapsto\tr(\sigma)P$ gives $1$ and entropy change $-\log2$, at the cost of
Proposition~\ref{prop:fmheat}.  For $\rho=\mathrm{diag}(\tfrac34,\tfrac14)$ and probes at polar angles $\pi/3$, $2\pi/3$ (both at Fubini--Study distance
$\pi/12$ from the equatorial probe), the outcome probabilities are $\tfrac58,\tfrac38$, the cross-entropies against $(\tfrac34,\tfrac14)$ are
$\log4-\tfrac58\log3$ and $\log4-\tfrac38\log3$, and the differences of cross- and relative entropy are both $\tfrac14\log3$; a controller
minimising either chooses the first probe, given a specified reference and objective, as in a
generative model \cite{n2-friston}; in a causal model an intervention also specifies which mechanism it
changes \cite{n2-pearl}.
\end{example}

\begin{proposition}[reachability and speed]\label{prop:fmspeed}
Rotations about two nonparallel axes generate $\SO(3)$ by finite words, and invariance under two moving-axis
actions implies invariance under all rotations.  \emph{Positive Schur slack does not give continuous
reachability}: on $S^2$ the exact constraints $x\cdot e_1=x\cdot e_2=0$ leave only $\pm e_3$, with no feasible path between them
(deleting one restores a great circle), and $x\cdot e_i=0$ for $i=1,2,3$ admits no unit solution although the bordered
Gram matrix is $I_4$.  The Fubini--Study speed of $i\hbar\dot\psi=H\psi$ is $\Delta_\psi H/\hbar$, and two protocols from one ray reaching
two orthogonal record rays in time $\tau$ with time-averaged uncertainty at most $E$ need $\tau\ge\pi\hbar/(4E)$, attained by two qubit
meridians from their midpoint.
\end{proposition}
\begin{proof}
Conjugation produces a third independent axis and the inverse function theorem a neighbourhood of
the identity; in a moving frame left rotations about $Oa$ act as right rotations about $a$.  The squared ray speed is
$\lVert\dot\psi\rVert^2-\lvert\ip\psi{\dot\psi}\rvert^2$; the triangle inequality gives $\pi/2\le2E\tau/\hbar$.
\end{proof}

\begin{example}[a dynamical Zeno limit]\label{ex:fmzeno}
For $H=\hbar\omega X$, initial $\lvert0\rangle$ and $n$ tests of $\lvert0\rangle\langle0\rvert$ in $[0,t]$, the all-success probability $\cos^{2n}(\omega t/n)$ has failure at most
$(\omega t)^2/n$ \cite{n3:Zeno}; this uses repeated instruments and a time budget, whereas the nesting of
Theorem~\ref{thm:fmschur} holds without them.
\end{example}

\begin{proposition}[anchored scalar transport]\label{prop:fmgraphphase}
On a connected multigraph with $n$ vertices and $m$ edges carrying phases $u_e\in\U(1)$, with $k\ge1$ anchored vertices
(including a root) where gauge transformations are trivial, the gauge quotient is a torus of dimension
$m+k-n$, with coordinates the tree transports to anchored vertices and the loop phases of non-tree edges.  For
unitary matrices the loop products transform by conjugation ($XZ=-ZX$); compare the group-valued
analysis of \cite[Sec.~0.4]{PaperI}.
\end{proposition}
\begin{proof}
Tree normalisation; every coordinate list is realised, and equal coordinates give the allowed gauge;
$m-(n-1)$ loops and $k-1$ open coordinates.
\end{proof}

\begin{theorem}[perception--action holonomy]\label{thm:fmberry}
For a closed piecewise smooth loop of pure qubit rays with parallel transport for $\ip\psi{d\psi}$ and a spanning
surface of signed area $\mathcal A$ on the Bloch sphere, the holonomy is $e^{-i\mathcal A/2}$, so the principal phase is at most half the
unsigned area; for a geodesic polygon with spinors $\psi_1,\dots,\psi_m$ it is $e^{-i\arg B}$, $B=\prod\ip{\psi_i}{\psi_{i+1}}$; for Bloch length $\ell<2\pi$,
$\lvert\gamma\rvert\le(\ell/2)\tan(\ell/4)$; and $\tr(P_aP_bP_c)=\tfrac14[1+a\cdot b+b\cdot c+c\cdot a+i\,a\cdot(b\times c)]$, which is $(1+i)/4$ for the positive octant,
with phase $-\pi/4$.  \emph{The holonomy bounds a spanning area; it neither forces nor fixes work}: $H=i\hbar[\dot P,P]$
transports the ray exactly with $\tr(PH)=0$ and zero isolated driving work $\int\tr(P\dot H)dt$ on a loop with
nonzero phase, and a separate reset actuator adds work $a>0$ (exciting and releasing an auxiliary qubit of
gap $a$) without changing the loop.
\end{theorem}
\begin{proof}
Horizontal lifts on arcs; the connection $i\sin^2(\theta/2)d\phi$ has curvature $\tfrac i2\sin\theta\,d\theta\wedge d\phi$, and Stokes; for the length bound
$\lvert\sin^2(\theta/2)d\phi\rvert\le\tfrac12\tan(\ell/4)ds$ in polar coordinates about the base point; $(a\cdot\sigma)(b\cdot\sigma)=(a\cdot b)\mathbf1+i(a\times b)\cdot\sigma$.  This is the
spinor form of the excess of \cite[Thm.~9.9, Cor.~9.11]{PaperI}
\cite{n2-pancharatnam,n2-berry,n2-simon,n2-sb}.  Differentiating $P^2=P$ gives $P\dot PP=0$, $[H,P]=i\hbar\dot P$, $\tr(PH)=0$ and
$\tr(P\dot H)=0$.
\end{proof}

\begin{remark}[three free energies; coverage]\label{rem:fmenergies}
For a reference model $p(s,o|u)$,
\[
\textstyle\sum_sq(s)\log[q(s)/p(s,o|u)]=D(q\Vert p(\cdot|o,u))-\log p(o|u)
\]
is dimensionless; $\tr\rho H-k_BTS(\rho)$ has energy units; an energy model with passive increment $\Delta E_{\mathrm p}>0$ and entropy tax $\Delta s$ gives $\Delta E_{\mathrm p}-\Delta s/\beta$, positive iff $\beta>\Delta s/\Delta E_{\mathrm p}$; a
measured deposit needs dynamics.  For a fixed sequence of full axis tests $P_j=n_jn_j^{\mathsf T}$ (or partial, single-step loss
$(1-\lambda_j^2)\lVert(\mathbf1-P_j)r_{j-1}\rVert^2$), the loss $D=\lVert r_0\rVert^2-\lVert r_L\rVert^2$ and the initial coverage $C=\sum_{j\le L}\lVert(\mathbf1-P_j)r_0\rVert^2$ satisfy $D\ge2C/[L(L+1)]$.
\end{remark}

\begin{conjecture}[small-phase control in a cap corridor]\label{conj:fmcontrol}
For an interior direction $p$ and control axes $a,b$ with $J=\lvert p\cdot(a\times b)\rvert>0$, using piecewise rotations of unit speed
about either axis, the least total rotation $L(\varphi)$ of closed loops at $p$ in the corridor with spinor holonomy
$e^{i\varphi}$ satisfies $L(\varphi)/\sqrt{\lvert\varphi\rvert}\to4\sqrt{2/J}$ as $\varphi\to0$.
\end{conjecture}
\noindent Evidence: in the tangent model with coordinates along $a\times p$, $b\times p$ the cost is $V_1+V_2$ (total variations),
the signed area is at most $V_1V_2/4$, attained by a square, physical area is $J$ times this and phase half the
area (Theorem~\ref{thm:fmberry}).  A proof needs uniform control of arbitrary switching and closure corrections.

\subsection{Reproducible examples}\label{sec:fmexamples}

\begin{proposition}[area of a cap intersection]\label{prop:fmarea}
For finitely many caps $v_i\cdot x\ge h$ ($0<h<1$) with nonempty interior and transverse two-cap corners,
$A=2\pi-h\sum_{\mathrm{arcs}}\delta_i-\sum_{\mathrm{corners}}\arccos\frac{v_i\cdot v_j-h^2}{1-h^2}$, $\delta_i$ the angular sweeps.
\end{proposition}
\begin{proof}
Gauss--Bonnet for a convex disk \cite[Ch.~4, \S5]{n2-docarmo}: the small circle has geodesic curvature $h/\sqrt{1-h^2}$, and adjacent
inward normals have inner product $(v_i\cdot v_j-h^2)/(1-h^2)$.
\end{proof}

\begin{example}[an expanding spiral of anchors]\label{ex:fmspiral}
For $h=\tfrac12$ and $v_j$ at azimuth $(j-1)\pi(3-\sqrt5)$ and polar angle $(12+6(j-1))^\circ$, the solid angles for $M=1,\dots,9$ are
$3.141593$, $2.297861$, $1.802272$, $1.412802$, $1.049640$, $0.720881$, $0.436215$, $0.208320$, $0.054604$ (boundary arcs $1,2,3,3,\dots,3$); the
first three anchors have Gram determinant $0.041133$; $\log(A_8/A_9)=1.338960$ nats ($1.931710$ bits), $\log(4\pi/A_9)=5.438666$, $A_3/A_9=33.0060$.
For $M\ge10$ the corridor is empty: $0.2388v_8+0.4261v_9+0.3351v_{10}$ has norm $<0.496$, whereas a feasible unit vector has inner
product at least $\tfrac12$ with it.  For a second observer with $w_j=R_y(10^\circ)v_j$ the common areas are $1.570105$ ($M=3$) and
$0.010419$ ($M=9$), ratios to the first observer's $0.871181$ and $0.190813$; the nonzero eigenvalues of the combined Gram
matrices are $(0.167199,0.456613,5.376189)$ and $(2.744723,3.881656,11.373621)$.  At $25^\circ$ and $30^\circ$ the $M=8$ common areas are $0.043804$
and $0.022659$, and at $M=9$ the means of $(v_8,w_9)$ and $(v_5,w_9)$ have norms $0.498642$ and $0.468178$, certifying emptiness.  The
anchor count alone supplies neither a feasible corridor nor an optimal memory policy.
\end{example}

\begin{example}[replacing versus deleting anchors]\label{ex:fmcentroid}
Replacing the three coordinate anchors by $(1,1,1)/\sqrt3$, keeping the last, or keeping the last two gives slacks
$1-(q_1+q_2+q_3)^2/3$, $1-q_3^2$, $q_1^2$: $(\tfrac2{27},\tfrac89,\tfrac49)$ at $q=(2,2,1)/3$ and $(\tfrac{26}{27},\tfrac59,\tfrac19)$ at $q=(1,2,-2)/3$.  Which compression is
better depends on the future queries.
\end{example}

\begin{example}[elimination and fill]\label{ex:fmfill}
For the divisibility graph on $\{2,\dots,6\}$ (edges $24,26,36$) and $L+\mathbf1/4$ with terminals $5,6$, eliminating $2$ first creates the
fill entry $-\tfrac49$ between $4$ and $6$; the order $(4,3,2)$ creates none; both give the terminal diagonal $(\tfrac14,\tfrac{441}{580})$
($\tfrac{29}{20}-\tfrac{20}{29}$).  A truncated Schur block $\tilde S$ leaves the residual $(S-\tilde S)\tilde x_T$ on the terminals.  For maximum degree $d$,
$P=\mathbf1-L/(1+d)$ is stochastic with spectrum in $(-1,1]$, and $P^k$ is positive exactly on pairs joined by walks of length $\le k$.
\end{example}

\begin{proposition}[a separability witness]\label{prop:fmseparable}
For Hermitian $R$ and its Hilbert--Schmidt projection $P$ onto the closed cone of separable matrices, $W=P-R$
pairs nonnegatively with every separable matrix and $\tr(WR)=-\lVert R-P\rVert_{\mathrm{HS}}^2$, a strict witness when $R$ is outside the cone.
\end{proposition}
\begin{proof}
The cone over the compact convex hull of product projectors is closed; the variational inequality at
the nearest point, tested at $S=0$ and $S=2P$.
\end{proof}

\begin{remark}[programme]\label{rem:fmprogramme}
Assign operational algebras to regions of the cap geometry with a causal composition rule and identify
the access restriction of Theorem~\ref{thm:fmeraser}; give a record-writing mechanism, an initial record condition
and a heat or work account with retained side information; quantify recovery by restricted local
instruments with matching bounds; and give a prediction target that distinguishes useful compression
from overlap deletion.  These connect the Gram model to Conjecture~\ref{conj:cueth} without assuming
equilibration.
\end{remark}

\begin{remark}[what does not hold]\label{rem:fmnogo}
(i) Growth does not inevitably lock and sliding does not always tighten; deletion need not reopen
(Remark~\ref{rem:fmnolock}).  (ii) Three anchors fix no chirality and give no nonempty optimal window
(Examples~\ref{ex:fmgram}, \ref{ex:fmspiral}).  (iii) Index restriction or discarding can cost zero heat
(Remark~\ref{rem:fmzeroheat}).  (iv) Forgetting or conjugate replay does not universally restore interference
(Remark~\ref{rem:fmcopies}).  (v) Local consistency does not give a global Gram, and a common Gram gives neither a joint
state nor joint outcomes (Theorem~\ref{thm:fmglue}, Examples~\ref{ex:fmoutcomes}, \ref{ex:fmcontext}).  (vi) A unique barycentre
or a generic shared POVM does not single out a pointer basis (Remarks~\ref{rem:fmaxisbary}, \ref{rem:fmbph}).  (vii) Positive
Schur slack does not give continuous reachability (Proposition~\ref{prop:fmspeed}).  (viii) Holonomy neither forces
nor fixes work (Theorem~\ref{thm:fmberry}).
\end{remark}

%=====================================================================================
\section{Single outcomes: bounded martingales and finite-tick halting}\label{sec:outcomes}
The weights $\beta$ of a complete test are the complete covariant description of the test; an outcome needs
further data.  This section determines what those data must satisfy.  If the trajectories of a
stochastic measurement tend almost surely to single orthogonal records, the Born rule at every stage is
equivalent to the current weights being martingales in the observer's filtration; the Born rule at
preparation alone is weaker.  On the Bloch sphere and on the simplex, a hidden threshold selects outcomes
with Born probabilities exactly when its law is the invariant one.  Finite certification is a separate
question: an inverse-tick resolution bound holds when a transcript distinguishes at most linearly many
alternatives, and adaptive binary comparisons beat it exponentially.

\subsection{What the structure does not supply}

\begin{proposition}[outcome rules]\label{prop:nooutcome}
(i) For one system there is no map $F$ assigning to a ray $q$ and a test $\{p,p^\perp\}$ with
$\beta(q,p)=\tfrac12$ one of $p,p^\perp$ such that $F(gq,\{gp,gp^\perp\})=gF(q,\{p,p^\perp\})$ for all
$g\in\SO(3)$, the image of $\cL$.  Away from ties the rule ``take the outcome with $\beta>\frac12$'' is
covariant.
(ii) For a composite ($\dim V_N\ge4$) there is no assignment of one outcome to every
complete test that is noncontextual, i.e.\ chooses a ray independently of the test containing it
\cite{KS67,Bell66}.
\end{proposition}
\begin{proof}
(i) Choose $q$ with $r_q\perp r_p$ and let $g$ be the rotation by $\pi$ about $r_q$: $gq=q$ and $g$ exchanges
$p$ and $p^\perp$, so $F(q,\{p,p^\perp\})$ would be fixed by $g$, which fixes neither.  (ii) Such an assignment
is a $\{0,1\}$-valued function on rays summing to one on every orthonormal basis, which the
Kochen--Specker theorem excludes in dimension $\ge3$.
\end{proof}
For one system the obstruction is only at ties; for composites it is the Kochen--Specker theorem.  An
outcome requires data that $(M,\hp{\cdot}{\cdot}{\cdot},\kappa,T)$ does not contain.  Given an outcome, the conditional
state is determined (Proposition~\ref{prop:relative}, Theorem~\ref{thm:instruments}).  The models below use additional
random data and the context of the implemented test, so Proposition~\ref{prop:nooutcome}, which concerns
deterministic assignments on the original data, does not apply to them.

\subsection{Weights, terminal records and martingales}\label{sec:sorecords}
Let $P_1,\dots,P_d$ be nonzero orthogonal record projections with $\sum_iP_i=\mathbf1$ on the space of Theorems~\ref{thm:born} and
\ref{thm:bornN}, $w_i(\rho)=\tr(P_i\rho)$, and $e_i$ the vertices of the simplex $\Delta_{d-1}$; the $w_i$ are nonnegative and sum to
one, they are the weights $\beta$, $\beta_N$ (summed over blocks for coarse tests), and Convention~\ref{conv:T} reads
them as probabilities.  A law for conditioned trajectories lets us prove this reading for the
frequencies of a fixed test.

\begin{definition}[record process]\label{def:soprocess}
Let $(\mathcal F_t)$ be the observer's filtration, in discrete time $t\in\mathbb N_0$ or in continuous time $t\ge0$ (then complete and
right-continuous, with right-continuous weight paths with left limits).  An adapted density process $\rho_t$ has
weights $W_i(t)=w_i(\rho_t)$.  It is a \emph{record process} ([R]) if $W(t)\to e_J$ almost surely for a record index $J$
measurable in $\mathcal F_\infty$ (one terminal vertex), and a vertex once attained is absorbing (absorbing record
faces).  The limit may be asymptotic.
\end{definition}

\begin{theorem}[Born weights and conditional probabilities]\label{thm:sorigidity}
For a record process the following are equivalent:
(i) $W_i(t)=\mathbb P(J=i\mid\mathcal F_t)$ almost surely for every deterministic $t$ and every $i$ (conditional Born
calibration at every time); (ii) every $W_i$ is a bounded $(\mathcal F_t)$-martingale.
Either gives the same formula at every almost surely finite stopping time $\tau$ (optional stopping); at
a deterministic preparation the winning probabilities are $w_i(\rho_0)$, and for independent repetitions the
frequencies converge to them almost surely.  In discrete time (ii) is zero conditional one-step drift; in
continuous time, for a special semimartingale $W_i=W_i(0)+M_i+A_i$ with predictable finite-variation $A_i$, it is $A_i=0$.
\end{theorem}
\begin{proof}
(i)$\Rightarrow$(ii) is the tower property.  Conversely, $\mathbb E[W_i(s)\mid\mathcal F_t]=W_i(t)$, $0\le W_i\le1$ and $W_i(s)\to\mathbf1_{\{J=i\}}$ give (i)
by conditional dominated convergence as $s\to\infty$.  For a bounded finite-valued $\tau$ test on each slice $\{\tau=t\}$;
in continuous time round a bounded $\tau$ up on dyadic grids, use right continuity and dominated
convergence; for finite $\tau$ test on $B\cap\{\tau\le m\}$ and let $m\to\infty$.  The frequencies follow from the strong
law.  The drift statement is uniqueness of the canonical decomposition, a bounded local martingale
being a martingale \cite{n3:Williams}.
\end{proof}

\begin{corollary}[fresh-start calibration and restart]\label{cor:sorestart}
Suppose the conditional continuation given $\mathcal F_t$ is the same law restarted at $\rho_t$ (and at time $t$ for a
time-inhomogeneous family); a Markov family with the density as state and the natural filtration has this
property at deterministic times.  Then Born terminal probabilities at every admissible start (fresh-start
calibration) imply the martingale property, and conversely.  A declaration $\hat J$ at a finite stopping time
with $W_{\hat J}(\tau)\ge1-\alpha$ has conditional error at most $\alpha$, also on the declaration event of a rule with an
undecided output.
\end{corollary}
\begin{proof}
Restart gives $\mathbb P(J=i\mid\mathcal F_t)=w_i(\rho_t)$, which is (i); the converse is Theorem~\ref{thm:sorigidity} at each start.  On
$\{\hat J=i\}$ the error is $1-W_i(\tau)$.
\end{proof}

\begin{remark}[Born at preparation does not force zero drift]\label{rem:somemory}
For two rank-one outcomes, $p=w_1(\rho_0)\in(0,1)$ and $a=p(1-p)$, put $W_1(1)=p+Sa$ with an equiprobable sign $S$, and from
tick two on $W_1=\mathbf1_{\{J=1\}}$ with $J$ drawn with probability $p$ independently of $S$ (diagonal densities).
$W_1(1)\in(0,1)$ since $p-a=p^2$ and $1-p-a=(1-p)^2$.  Both unread maps from preparation are the pinching channel, which
is CPTP and affine, and every preparation has Born winning probabilities; yet $\mathbb E[W_1(2)\mid\mathcal F_1]=p\ne W_1(1)$.
\emph{Born calibration at preparation plus affine unread evolution does not force zero conditional drift};
the correct statement is Theorem~\ref{thm:sorigidity} (calibration at every time) or Corollary~\ref{cor:sorestart}
(fresh-start calibration plus restart).  The conditional increments cancel in expectation from preparation;
a drift of fixed sign (a strict sub- or supermartingale) does change the terminal probabilities.
\end{remark}

\begin{proposition}[unread channels and record weights]\label{prop:sochannels}
If $\mathbb E[\rho_t\mid\mathcal F_s]=\Phi_{s,t}(\rho_s)$ ($s\le t$) with CPTP maps preserving every record face, the weights are martingales.
Face-preserving CPTP maps from deterministic preparations to the unread averages give Born terminal
probabilities at preparation only.  With remote steering and local tomography, absence of signalling by
a remote choice of ensemble requires an affine unread map on density matrices; local CPTP maps satisfy
it.
\end{proposition}
\begin{proof}
$E_i=\Phi^*_{s,t}(P_i)$ satisfies $0\le E_i\le\mathbf1$, and face preservation gives $\ip v{E_iv}=1$ on face $i$ and $0$ on the others;
a vector with zero expectation of a positive operator is in its kernel, so $E_i=P_i$, and tracing gives the
martingale identity.  From preparations the same argument gives $\mathbb EW_i(t)=w_i(\rho_0)$, and bounded convergence
gives the terminal law.  For an ensemble $\sum_aq_a\rho_a=\rho$, the purification with coefficients $\sqrt{\lambda_j}$ and the
effects $(\rho^{-1/2}q_a\rho_a\rho^{-1/2})^{\mathsf T}$ on the purifying side steer to $q_a\rho_a$; local tomography makes every such
ensemble give the same unread density, which is the affine law.  The dual of a local CPTP map fixes $\mathbf1$, so
remote effects $\mathbf1\otimes E$ keep their expectations.
\end{proof}

\begin{proposition}[unread and selected states]\label{prop:sopinching}
$\mathcal D(\rho)=\sum_iP_i\rho P_i=\sum_{i:w_i>0}w_i\rho_i$ with $\rho_i=P_i\rho P_i/w_i$; $\mathcal D$ is idempotent, fixes exactly the block-diagonal
matrices, and is the reduced channel of the isometry $v\mapsto\sum_iP_iv\otimes\lvert i\rangle$.  If a rank-one coupling sends $\lvert i\rangle$
to $\lvert i\rangle\otimes E_i$ and $m$ fresh registers are used, the $(i,j)$ entry becomes $\rho_{ij}\ip{E_j}{E_i}^m$, and if $\lvert\ip{E_j}{E_i}\rvert\le q<1$ for
$i\ne j$, $\lVert\rho_m-\mathcal D(\rho)\rVert_{\mathrm{HS}}\le q^m\lVert\rho-\mathcal D(\rho)\rVert_{\mathrm{HS}}$.
\end{proposition}
\begin{proof}
Blocks are positive with trace $w_i$; orthogonality gives $P_i\mathcal D(\rho)P_j=\delta_{ij}P_i\rho P_i$; the partial trace removes the terms
$i\ne j$.  Inner products multiply on tensor products.  (A reused environment needs its own analysis.)
\end{proof}

\begin{remark}[one surviving record needs more than linearity, pinching or stabilisation]\label{rem:sononselection}
At ray angle $30^\circ$ the unread density $\mathrm{diag}(\tfrac34,\tfrac14)$ is fixed by pinching with both weights in $(0,1)$: linearity
and pinching keep both candidates.  A map $U$ on subsets with $U(A)\subseteq A$ stabilises within $\lvert S_0\rvert$ updates, but
$U(A)=A$ with $\lvert S_0\rvert=2$ keeps both candidates; only if $U$ preserves nonemptiness and strictly shrinks every
nonsingleton does a nonempty $S_0$ reach a singleton (within $\lvert S_0\rvert-1$ updates).
\end{remark}

\begin{example}[ideal continuous nondemolition measurement]\label{ex:sodiffusion}
For $A=\sum_ia_iP_i$ with distinct real $a_i$ and no Hamiltonian mixing of record faces, the conditional diffusion
$d\rho_t=-\gamma[A,[A,\rho_t]]dt+\sqrt{2\gamma}(A\rho_t+\rho_tA-2\bar a_t\rho_t)dB_t$, $\bar a_t=\tr(A\rho_t)$, is a record process with martingale weights.
Construction: with $\eta=\sqrt{2\gamma}$ let $Y_t$ be Brownian motion with drift $2\eta a_i$ chosen with probability $w_i(\rho_0)$;
$K_t=\exp(\eta AY_t-\eta^2A^2t)$, $X_t=K_t\rho_0K_t$, $\rho_t=X_t/\tr X_t$ gives the posterior weights (the likelihood of drift $2\eta a_i$
is $\exp(2\eta a_iY_t-2\eta^2a_i^2t)$), and $B_t=Y_t-\int_0^t2\eta\bar a_sds$ is Brownian in the observation filtration.  It\^o's rule gives
$dX=\eta(AX+XA)dY-\gamma[A,[A,X]]dt$ and the quotient rule the stated equation.  Then
$dW_i=2\sqrt{2\gamma}(a_i-\bar a_t)W_idB_t$; the It\^o isometry bounds $\mathbb E\int_0^\infty8\gamma(a_i-\bar a_t)^2W_i^2dt$ by one, so $(a_i-\bar a_\infty)W_i(\infty)=0$ and exactly
one limit is positive; the stochastic-exponential form keeps positive weights positive at finite times.
The mean is the dephasing channel, block $(i,j)$ multiplied by $e^{-\gamma t(a_i-a_j)^2}$, the average of $e^{-iAX}\cdot e^{iAX}$ over real
Gaussian $X$ of variance $2\gamma t$.  The construction uses the projective weights in the observation law.
\end{example}

\begin{remark}[relation to stochastic measurement theory]\label{rem:sohistory}
Gisin's ensemble argument concerns dependence of the dynamics on the density, and Simon, Bu\v{z}ek and
Gisin extend it to composites \cite{n3:Gisin,n3:Simon}; Gisin and Percival represent unread evolution by
quantum-state diffusion, and Di\'osi models continuous measurement by stochastic evolution
\cite{n3:GisinPercival,n3:Diosi}; Jacobs and Steck review continuous monitoring \cite{n3:Jacobs}; Adler,
Brody, Brun and Hughston use martingales for energy-driven reduction and its terminal probabilities
\cite{n3:Adler}; Bassi and Ghirardi, and Bassi et al., discuss dynamical reduction and its empirical
distinctions from unitary dynamics \cite{n3:BassiGhirardi,n3:Bassi}.  Theorem~\ref{thm:sorigidity} states, independently of a
choice of diffusion, which filtration and restart conditions give the Born rule across stages, and
Remark~\ref{rem:somemory} shows that the weaker condition at preparation does not.
\end{remark}

\subsection{Attraction basins on the sphere and the simplex}\label{sec:sobasins}
Orient the Bloch axis $u$ (Theorem~\ref{thm:born}) towards record one; for ray angle $\theta$ from it, $z=r\cdot u=2w_1-1=\cos2\theta$.

\begin{theorem}[threshold selection on the Bloch sphere]\label{thm:sosphere}
Let $k>0$ and draw $\lambda\in(-1,1)$ from one Borel law for all input rays.  For fixed $\lambda$ the flow $\dot r=k(z-\lambda)(u-zr)$
on $\lvert r\rvert=1$, with $\dot z=k(z-\lambda)(1-z^2)$, sends $z>\lambda$ to $+u$ and $z<\lambda$ to $-u$ (asymptotically) and leaves the latitude
$z=\lambda$ fixed.  The north probability is $(1+z)/2$ for every input iff $\lambda$ is uniform, and the uniform law is the
push-forward of the rotation-invariant area on $S^2$ under $r\mapsto r\cdot u$ (Archimedes' hat-box theorem).
\end{theorem}
\begin{proof}
$r\cdot(u-zr)=0$, so the field is tangent; the equilibria $1,-1,\lambda$ of the scalar equation have derivatives $-2k(1-\lambda)$,
$-2k(1+\lambda)$, $k(1-\lambda^2)$, and $z(t)$ is strictly monotone towards its pole in each basin; $\lvert r\mp u\rvert^2=2(1\mp z)$.  The north
probability is $\mathbb P(\lambda<z)$; equality with $(1+z)/2$ for all $z$ fixes the law on half-open intervals and excludes atoms.  In
coordinates $(\sqrt{1-z^2}\cos\varphi,\sqrt{1-z^2}\sin\varphi,z)$ the area is $dz\,d\varphi$ \cite{n3:Archimedes}; averaging over $\SO(3)$ shows uniqueness of the
invariant probability.
\end{proof}

\begin{remark}[symmetry alone is not enough; no affine flow]\label{rem:sothresholds}
The deterministic $\lambda=0$ leaves the equator undecided.  The north--south symmetric density $3\lambda^2/2$ gives north
probability $(1+z^3)/2$, i.e.\ $9/16$ at $z=\tfrac12$ (Born: $3/4$) and $7/16$ at $z=-\tfrac12$; an even atomless law only fixes $\tfrac12$ at the
equator.  Even the uniform law gives no affine collapse flow: the mean initial velocity is $kz(1-z^2)u$ for a mixture
of two rays at height $z$ with opposite transverse components and $0$ for the mixture of poles with weights
$(1\pm z)/2$, although both have the same density; the terminal law is the pinched density.
\end{remark}

\begin{theorem}[simplex selection and its sampling law]\label{thm:sosimplex}
For $d\ge2$ let $U$ be a random point of the open simplex $\Delta_{d-1}$ with law independent of the input weights $p$, and let
the winner be the strict minimiser of $U_i/p_i$ over $p_i>0$ (ties undecided).  Every input has winning
probabilities $p_i$ iff $U$ is uniform ($\mathrm{Dirichlet}(1,\dots,1)$), which is the push-forward of the invariant probability on $\C P^{d-1}$
by the squared moduli in the record basis.
\end{theorem}
\begin{proof}
For a standard complex Gaussian vector $G$ the ray is unitary invariant and $X_i=\lvert G_i\rvert^2$ are independent rate-one
exponentials; $U_i=X_i/s$, $s=\sum X_i$, has Jacobian $s^{d-1}$ and $\int_0^\infty e^{-s}s^{d-1}ds=(d-1)!$, the uniform density.  $X_i/p_i$ is
exponential of rate $p_i$, and $i$ wins with probability $\int_0^\infty p_ie^{-t}dt=p_i$.  Conversely, for interior $p$ with $p_i=t_ip_d$,
the winning event of $d$ is $\{U_i/U_d>t_i\ \forall i<d\}$, so the required probabilities fix the joint survival function of
the ratios, hence the law of $U$.
\end{proof}

\begin{remark}[scope of the simplex rule]\label{rem:sobasinscope}
For $d=2$, $\lambda=2U_1-1$ recovers Theorem~\ref{thm:sosphere}.  With $b_i=p_i/U_i$, $x_i(t)=p_ie^{b_it}/\sum_jp_je^{b_jt}$ solves $\dot x_i=x_i(b_i-\sum_jx_jb_j)$
and tends to the winner, a continuous realisation of the basins with frozen scores.  The rule is the
exponential-race form of the vertex-opposite partition of the simplex; Aerts and Sassoli de Bianchi use
uniform simplex volumes \cite{n3:Aerts}, and complex Gaussian coordinates for invariant rays are used
in \cite{n3:Zyczkowski}.  For coarse tests sample in a fine basis and group labels; raw block masses are
$\mathrm{Dirichlet}$ with the block ranks and do not implement the coarse race: for ranks $(2,1)$ and weights $(\tfrac12,\tfrac12)$
the first block wins with probability $\int_0^{1/2}2u\,du=\tfrac14$.  For an actual apparatus the invariant auxiliary law is
an input (Hypothesis~Ap below).
\end{remark}

\begin{example}[five-state comparison]\label{ex:sochain}
On $p_j=j/4$ ($0\le j\le4$) with absorbing ends and interior up-probability $q$, the absorption time has tail at most
$(1-q^4)^m$ after $4m$ ticks but no deterministic bound from $j=2$.  The upper hitting probability is $h_j=j/4$ for $q=\tfrac12$
(mean time $j(4-j)$) and $h_j=(1-2^{-j})/(1-2^{-4})$ for $q=\tfrac23$, i.e.\ $8/15,4/5,14/15$ against Born $1/4,1/2,3/4$, with drift
$(2q-1)/4=1/12$: a drift of fixed sign changes the probabilities.
\end{example}

\subsection{Finite comparison ticks}\label{sec:soticks}

\begin{definition}[certified resolution]\label{def:sotranscript}
An unknown parameter ranges over an interval of length $L>0$; a protocol of at most $N_{\mathrm{ticks}}\ge1$ ticks reports a
terminal transcript, and certifies resolution $\varepsilon_{\mathrm{self}}$ if for every transcript the compatible inputs lie in
an interval of length (diameter) at most $\varepsilon_{\mathrm{self}}$; certification of absolute error by one estimate uses
the radius.
\end{definition}

\smallskip\noindent\textbf{Hypothesis Tr (linear transcript capacity).}  For the measurement and parameter in question
at most $AN_{\mathrm{ticks}}$ distinguishable terminal transcripts are available for uniform certification, $A>0$ independent
of $N_{\mathrm{ticks}}$; the transcript contains all information used, including any recorded stopping time.

\begin{proposition}[inverse-tick resolution]\label{prop:soresolution}
With at most $M$ transcripts, $\varepsilon_{\mathrm{self}}\ge L/M$ (diameter); under Hypothesis~Tr, $\varepsilon_{\mathrm{self}}\ge c_{\mathrm{res}}/N_{\mathrm{ticks}}$ with
$c_{\mathrm{res}}=L/A$; for absolute error the constants are halved.  Equal cells attain the bound (e.g.\ diameter $0.05$ for
$L=1$, $M=N_{\mathrm{ticks}}=20$).
\end{proposition}
\begin{proof}
The certified intervals cover the input interval.  Equal cells with $M-1$ fixed cuts attain the bound, their
midpoints the radius bound.  For randomised protocols with the seed known to the certifier, apply the
argument at each seed.
\end{proof}

\begin{remark}[the unrestricted inverse-tick bound is false]\label{rem:sobisection}
Adaptive bisection halves the interval at each binary comparison and reaches diameter $2^{-N_{\mathrm{ticks}}}$ on $[0,1]$,
optimal for depth-$N_{\mathrm{ticks}}$ binary trees, and $2^{-N}<c_{\mathrm{res}}/N$ eventually for every $c_{\mathrm{res}}>0$.  A discrete label can be
read exactly in one tick.  A finite tick count or a speed limit therefore does not by itself give an
inverse-tick floor; Hypothesis~Tr does.
\end{remark}

\begin{remark}[clock and causal budgets]\label{rem:soclock}
With one comparison per period $1/\nu_0$ the serial budget in a window $\Delta\tau$ is $\lfloor\nu_0\Delta\tau\rfloor$; a round trip over distance $R$
at speed at most $c_{\mathrm{light}}$ needs $2R\le c_{\mathrm{light}}\Delta\tau$; a lattice with hop length $\ell$ per period has speed $\ell\nu_0$.  With $L=1$, $A=1$,
Proposition~\ref{prop:soresolution} gives $\varepsilon_{\mathrm{self}}\ge1/\lfloor\nu_0\Delta\tau\rfloor\ge\ell/(c_{\mathrm{light}}\Delta\tau)$ if $c_{\mathrm{light}}=\ell\nu_0$; a universal clock and the
capacity are physical inputs (a nearest-neighbour cubic lattice has a different reachable set from a light
cone).  Separately, under an assumed Fubini--Study angular speed $E/\hbar$, a stage within angle $E\tau/\hbar$ of two
orthogonal rays needs $\pi/2\le2E\tau/\hbar$, so $n\le4E\Delta\tau/(\pi\hbar)$ such stages fit in $\Delta\tau$; using this as a precision bound
needs a link to the transcript capacity.  The geodesic recurrence of \cite[Thm.~6.2]{PaperI} advances along
an arc; a measurement comparison also needs the observable and the information returned.
\end{remark}

\begin{proposition}[residual certificate]\label{prop:soresidual}
For a contraction $F$ with constant $q<1$ of a complete metric space, $x_n=F(x_{n-1})$ and $r_n=d(x_n,x_{n-1})$: $r_n\le q^{n-1}r_1$
and $d(x_n,x_*)\le\frac q{1-q}r_n$, so a residual test with positive tolerance halts in finite time with this error bound.
\end{proposition}
\begin{proof}
$r_{n+1}\le qr_n$ and the geometric tail; uniqueness by contraction.
\end{proof}
\noindent For $\cos$ on $[0,1]$ ($q=\sin1$) from $x_0=0$: $1$, $0.540302$, $0.857553$, $0.654290$, $0.793480$; the fifth residual $0.139191$ is the
first below $0.150$, with actual error $0.054395$; the set $\{\lvert\cos x-x\rvert\le0.150\}$ is $[0.6476,0.8270]$.  At tolerance zero there is
no finite halt from $0$ (injectivity), whereas a start at the fixed point halts at once, so zero residual does
not universally prevent termination.  The shift $x\mapsto\{10x\}$ reveals the digits $0.739$, $0.390$, $0.908$, $0.085$, $0.851$ of
$x_*=0.7390851332\ldots$, since $\{10\{10^tx\}\}=\{10^{t+1}x\}$.

\subsection{Spectral truncation and repeated readout}\label{sec:sospectra}

\begin{proposition}[tails and finite observability]\label{prop:sotails}
For an orthogonal projection $P$ and a skew-adjoint unitary $R$, $R^2=-\mathbf1$ and $\lVert Pv\rVert^2+\lVert R(\mathbf1-P)v\rVert^2=\lVert v\rVert^2$; the rotated tail
stays in the discarded subspace for all $v$ iff that subspace is $R$-invariant (the quarter turn $R(x,y)=(-y,x)$ with $P$
the first coordinate is a counterexample), and $\re\ip v{Rv}=0$.  For $T\colon V\to V$, $C\colon V\to Y$ finite-dimensional, the tower
$K_m=\mathrm{span}\{(T^\dagger)^kC^\dagger y:k\le m\}$ has at most $\dim V-\rk C$ strict inclusions, and the first equality is permanent.
\end{proposition}
\begin{proof}
$RR^\dagger=\mathbf1$ with $R^\dagger=-R$; orthogonality and norm preservation.  Each strict inclusion raises the dimension, and
$T^\dagger K_m\subseteq K_{m+1}=K_m$ propagates.  (The identity update saturates at once, so saturation is no ranking.)
\end{proof}

\begin{proposition}[power filtering and truncation]\label{prop:sopower}
For $G\ge0$ with eigenpairs $(\sigma_i,u_i)$, $\sigma_1\ge\dots\ge\sigma_d$, and $v=\sum c_iu_i$: $\sigma_1^{-n}G^nv=\sum(\sigma_i/\sigma_1)^nc_iu_i$; if $\sigma_1>\sigma_2$ and $c_1\ne0$ the
normalised iterates converge to $u_1$ with orthogonal part at most $(\sigma_2/\sigma_1)^n(\sum_{i>1}\lvert c_i\rvert^2)^{1/2}/\lvert c_1\rvert$; with a repeated
top eigenvalue the limit is the normalised leading projection, and with zero leading projection the
largest populated positive eigenspace controls the limit \cite{n3:vonMises}.  The truncated singular-value
decomposition minimises spectral and Frobenius errors among rank-$\le r$ matrices, with errors $s_{r+1}$ and
$(\sum_{i>r}s_i^2)^{1/2}$ \cite{n3:EckartYoung,n3:Mirsky}.
\end{proposition}
\begin{proof}
Diagonalise.  For rank $\le r$, $\ker B$ meets the span of the first $r+1$ right singular vectors; for Frobenius, a projection
$Q$ onto $d-r$ kernel vectors gives $\sum_is_i^2\ip{v_i}{Qv_i}$ with coefficients in $[0,1]$ summing to $d-r$.
\end{proof}
\noindent For $\sigma=(1,0.76,0.54,0.35,0.20,0.10)$ the fourth powers are $(1,0.333622,0.085031,0.015006,0.001600,0.000100)$; the shares of the first
two modes for $k=2,8,10$ in $\sum_{i\le2}\sigma_i^k/\sum\sigma_i^k$ are $0.772689$, $0.993334$, $0.997997$ (operator, fourth power, fourth-power output of
an input $\propto\sigma_i$).  The masses $\sigma_i^2/\sum\sigma_j^2$ above $0.100$ retain $0.9155$; releasing one discarded mode per tick gives
$0.9951$ after two ticks (a specified schedule).

\begin{remark}[basis-dependent weights]\label{rem:sobasis}
Weights are Rayleigh quotients $\ip{b_i}{\rho b_i}$: $\psi=(1,2,5)/\sqrt{30}$ has weights $(1,4,25)/30$ in the coordinate basis and $(1,0,0)$ in a
basis starting with $\psi$.  Interpolating from the coordinate basis to the real orthonormal basis $F$ with
columns $(1,1,1)/\sqrt3$, $(1,0,-1)/\sqrt2$, $(-1,2,-1)/\sqrt6$ ($\det F=1$) by $F(t)=\mathbf1+\sin(t\omega)K+(1-\cos(t\omega))K^2$, $\omega=\arccos((\tr F-1)/2)$,
$K=(F-F^{\mathsf T})/(2\sin\omega)$, gives at $t=65/90$ the weights $0.6273$, $0.2525$, $0.1202$.  Inner-product scoring also occurs in
attention, $\mathrm{softmax}(QK^{\mathsf T}/\sqrt{d_k})V$ \cite[Sec.~3.2.1]{n3:Attention}; quantum weights use squared moduli and
orthogonal completeness, and the similarity is an analogy.
\end{remark}

\begin{proposition}[repeated readout and Zeno survival]\label{prop:sozeno}
For a qubit in record zero rotated through angle $a/N$ between tests, the probability of record zero at all $N$
tests is $S_N=\cos^{2N}(a/N)=1-a^2/N+O(N^{-2})$, increasing in $N$ for $0<a<\pi/2$; with unread intermediate tests the final
probability is $[1+\cos^N(2a/N)]/2$.  For $a=\pi/4$: $N=1$ both $\tfrac12$; $N=16$ (step loss $0.002408$) $0.962166$ and $0.962838$;
$S_N=0.728553,0.856232,0.925676,0.990408$ for $N=2,4,8,64$ \cite{n3:Zeno}.
\end{proposition}
\begin{proof}
Each success resets the state; the unread transition matrix has eigenvalues $1$, $\cos(2a/N)$; $2\log\cos x=-x^2+O(x^4)$; the
derivative of $2N\log\cos(a/N)$ in $N$ is $2(\log\cos x+x\tan x)>0$ at $x=a/N$.
\end{proof}
\noindent The all-success history and the unread final population are different observables.  For first
returns, $h_n=((\mathbf1-P)U)^nv$, $s_n=\lVert h_n\rVert^2$ and $p_{n+1}=\lvert\ip v{Uh_n}\rvert^2$ satisfy $p_{n+1}=s_n-s_{n+1}$ and $\sum_{n<N}p_{n+1}+s_N=1$.

\subsection{Physical origin of the record law}\label{sec:soscope}

\begin{example}[interference with a retained record]\label{ex:sorecoherence}
For $\Psi=\sqrt p\,v_0+\sqrt{1-p}\,v_1$ with orthonormal joint record vectors, $\Psi$ and $p\lvert v_0\rangle\langle v_0\rvert+(1-p)\lvert v_1\rangle\langle v_1\rvert$ have the same
system marginal, but the joint test of $\lvert\Psi\rangle\langle\Psi\rvert$ (undo the coupling and test the initial ray) gives $1$ and $p^2+(1-p)^2$,
differing by $2p(1-p)$.
\end{example}
\noindent Decoherence suppresses interference; the selection of a realised record is a separate
question \cite{n3:Schlosshauer}.  Unread pinching and the selected process agree on the reduced ensemble
(Proposition~\ref{prop:sopinching}) and differ on a conditioned run.  Theorem~\ref{thm:sorigidity} is a process criterion
for the frequencies of realised records, given that a terminal record exists; the absorption law, the
observer's information and the coupling to the apparatus are physical ingredients.

\smallskip\noindent\textbf{Hypothesis Ap (preparation-independent apparatus sampling).}  An implemented test has
apparatus degrees of freedom inducing an auxiliary ray independent of the system preparation, with a
unitary-invariant law before the test-dependent selection rule; in a binary test the height is then
uniform, and a fine test can be modelled by the simplex rule of Theorem~\ref{thm:sosimplex}.

\begin{conjecture}[invariant sampling from apparatus dynamics]\label{conj:somixing}
There is a regime of apparatus dynamics driven by local comparisons in which the conditional law of the
auxiliary ray, before the fixed test is applied, converges to the invariant projective probability
uniformly in the preparations.
\end{conjecture}
\noindent Evidence: independent isotropic complex Gaussian amplitudes give this limit law exactly
(Theorem~\ref{thm:sosimplex}).  A derivation must specify the microscopic measure and dynamics independently of the
target frequencies, an apparatus sequence, a topology of convergence and control uniform in the
preparation.

\begin{remark}[programme]\label{rem:soprogramme}
An effective account must specify the observed record, the omitted bath degrees of freedom and when
conditional instruments approximate the stochastic dynamics; an objective-collapse account must specify a
modified joint evolution, verify the unread maps on entangled composites and predict recoherence tests; a
nested-observer account must assign filtrations to actual observations and interventions and compare
joint predictions, under the constraints of extended Wigner-friend protocols \cite{n3:FR,n3:Bong}.
Reading a rotated tail as a source of physical time, a tolerance as causal control, or a coarse-graining
gain in effective information \cite{n3:Hoel} as a property of an apparatus each needs a dynamics and an
intervention model.
\end{remark}

\begin{remark}[what this settles]\label{rem:sosettles}
Settled, as theorems: the statistics of any single-outcome dynamics of the kind considered (a record
process with one terminal vertex and absorbing record faces in the observer's filtration) agree with the
weights at every stage iff the weights are martingales (Theorem~\ref{thm:sorigidity}); unread evolution compatible
with no-signalling is affine (Proposition~\ref{prop:sochannels}), and a face-preserving CPTP conditional law gives the
martingale property; a hidden-threshold or simplex selection is Born iff its auxiliary law is the invariant
one (Theorems~\ref{thm:sosphere}, \ref{thm:sosimplex}).  Not settled: which dynamics, if any, is realised; why records
are absorbing; and whether an apparatus produces the invariant law (Hypothesis~Ap, Conjecture~\ref{conj:somixing}).
Refuted: Born calibration at preparation with affine unread evolution forces zero drift
(Remark~\ref{rem:somemory}); north--south symmetry gives the Born law, or the uniform threshold an affine flow
(Remark~\ref{rem:sothresholds}); linearity, pinching or stabilisation give one record (Remark~\ref{rem:sononselection}); a tick
budget alone gives an inverse-tick floor (Remark~\ref{rem:sobisection}).
\end{remark}


%=====================================================================================
\paperpart{Part III.  Bounds and budgets}
\noindent This part is reserved for speed limits, the Tsirelson bound and an entanglement budget; the
sections are in preparation.  The velocity geometry of states (Section~\ref{sec:velocity}) is its present
starting point.


%=====================================================================================
\paperpart{Part IV.  Fields, particles and gravity}

%=====================================================================================
\section{The space of events}\label{sec:eventspace}

\subsection{Events as dichotomies}\label{sec:events}
This subsection uses no postulate; only Propositions~\ref{prop:solder}(iv),(vi) use a clock.  The
elements $\xi=(t,v)$ of the kernel space $\R\oplus V$ of $\kappa$ define the \emph{affine comparison
functions} $f_\xi(x)=t+c(v,x)$ on $M$.  The set $C_\xi:=\{x\in M:f_\xi(x)>0\}$ is neither empty nor all of $M$
iff $\lvert v\rvert>\lvert t\rvert$, or $t=\lvert v\rvert\ne0$, in which case $C_\xi=M\setminus\{-v/\lvert v\rvert\}$; we call
$C_\xi$ a \emph{dichotomy} when $\lvert v\rvert>\lvert t\rvert$ (the punctured spheres belong to the null cone and
the conformal boundary) and normalise $\lvert v\rvert^2-t^2=1$.  A dichotomy is a threshold comparison of
configurations with a reference point, $\{x:\kappa(n,x)>\lambda\}$: a sharp yes/no observable on $M$.  It is
not a test in the sense of Convention~\ref{conv:T}, since it is not a union of clock orbits.  Let
$\cD:=\{\xi:\lvert v\rvert^2-t^2=1\}$ with the metric $g$ induced by $\lvert dv\rvert^2-dt^2$, and
$P(\xi,\eta):=\ip{v}{w}-ts$ for $\eta=(s,w)$.

\begin{theorem}\label{thm:desitter}
(i) $\xi\mapsto C_\xi$ is a bijection from $\cD$ onto the open caps of $M$: $C_\xi$ is the geodesic
ball of radius $a\in(0,\pi)$ about $n$, where $v=n/\sin a$, $t=-\cot a$; equivalently
$C_\xi=\{x:\kappa(n,x)>\cos^2\tfrac a2\}$.
(ii) $(\cD,g)$ is de~Sitter space $dS_4$: geodesically complete, Lorentzian, of constant curvature $1$;
its Ricci tensor is $3g$, so it solves $G+\Lambda g=0$ with $\Lambda=3$ in units of the radius of $M$.
(iii) The M\"obius group $\OO^+(1,4)$ acts on $\cD$ transitively by isometries, compatibly with its
action on caps; $\Aut^\pm(M)=\OO(4)$ preserves $t$.
(iv) For $\xi\ne\eta$:
(a) $C_\xi\subseteq C_\eta$ iff $P(\xi,\eta)\ge1$ and $t_\xi\le t_\eta$, iff $\eta$ lies in the causal
future of $\xi$;
(b) the boundary spheres of $C_\xi$ and $C_\eta$ cross iff $\lvert P(\xi,\eta)\rvert<1$; for $\eta\ne-\xi$ this holds iff $\xi$ and
$\eta$ are joined by a spacelike geodesic, while $\eta=-\xi$ has $P=-1$, coincident boundary spheres, and is
joined by spacelike semicircles;
(c) $P(\xi,\eta)\le-1$ iff $C_\xi\cap C_\eta=\emptyset$ or $C_{-\xi}\cap C_{-\eta}=\emptyset$, where
$C_{-\xi}$ is the interior of the complement of $C_\xi$; such pairs are not causally related.
For distinct, non-coincident boundary spheres, tangency corresponds to $\lvert P\rvert=1$.
(v) As $a\to0$ the caps shrink to points and as $a\to\pi$ their complements do: $M$ is the past and
the future conformal boundary of $\cD$.
\end{theorem}
\begin{proof}
(i) $t+\ip{v}{x}>0$ iff $\ip{n}{x}>\cos a$, iff $\kappa(n,x)>\cos^2\tfrac a2$.  (ii) $\cD$ is the
hyperboloid of unit spacelike vectors in $\R^{1,4}$ \cite{ONeill,HawkingEllis}.  (iii) $\OO^+(1,4)$
preserves $B$, and $f_\xi(x)=B(\epsilon\xi,(1,x))$ with $\epsilon(t,v)=(t,-v)$; if $g(1,x)=\lambda(1,y)$
with $\lambda>0$, then $x\in C_\xi$ iff $y\in C_{\epsilon g\epsilon\xi}$, so $g$ acts on $\cD$ by
$\xi\mapsto\epsilon g\epsilon\,\xi$, again an element of $\OO^+(1,4)$.  (iv) With $\theta$ the angle between the centres,
$P=(\cos\theta-\cos a\cos b)/(\sin a\sin b)$.  (a) Inclusion is $\theta\le b-a$, i.e.\
$\cos\theta\ge\cos(b-a)$ and $a\le b$, i.e.\ $P\ge1$ and $t_\xi\le t_\eta$; in $dS_4$ the causal future of
$\xi$ is $\{\eta:P\ge1,\ t_\eta\ge t_\xi\}$ \cite{ONeill}.  (b) The spheres cross iff
$\lvert a-b\rvert<\theta<\min(a+b,2\pi-a-b)$ iff $\lvert P\rvert<1$, and points of $dS_4$ with
$\lvert P\rvert<1$ are joined by a spacelike geodesic.  (c) $f_{-\xi}=-f_\xi$ and $P(-\xi,\eta)=-P(\xi,\eta)$;
by (a), $P(-\xi,\eta)\ge1$ iff $C_\eta\subseteq C_{-\xi}$ or $C_{-\xi}\subseteq C_\eta$, and an open cap
lies in $C_{-\xi}$ iff it misses $C_\xi$, while $C_{-\xi}\subseteq C_\eta$ iff $C_{-\xi}$ misses $C_{-\eta}$.
Causally related points have $P\ge1$.  (v) follows from (i).
\end{proof}

Pairs with $P\le-1$ are causally unrelated as well, so crossing of boundary spheres is equivalent to
joinability by a spacelike geodesic and not to causal unrelatedness; (c) records the remaining case.

\begin{proposition}[cosmological form]\label{prop:cosmo}
In the coordinates $(a,n)\in(0,\pi)\times M$ of Theorem~\ref{thm:desitter}(i),
\[
g=\frac{-da^2+g_M}{\sin^2a},
\]
with $g_M$ the round metric of $M$.  The cap radius $a$ is conformal time; the curves $n=\mathrm{const}$
are timelike geodesics with proper time $s=\ln\tan\frac a2$; and $g=-ds^2+\cosh^2s\,g_M$, the global
slicing of de~Sitter space, whose spatial sections are copies of $M$ with scale factor
$R=1/\sin a=\cosh s$.
\end{proposition}
\begin{proof}
With $\xi=(-\cot a,n/\sin a)$, $dt=da/\sin^2a$ and $dv=dn/\sin a-n\cos a\,da/\sin^2a$ with $\ip{n}{dn}=0$, so
$\lvert dv\rvert^2-dt^2=\lvert dn\rvert^2/\sin^2a+(\cos^2a-1)da^2/\sin^4a$.  The curve $n=\mathrm{const}$ is the
fixed set of the isometries in $\OO(4)$ fixing $n$, hence a geodesic, and $ds=da/\sin a$ integrates to
$\ln\tan\frac a2$, with $\cosh s=1/\sin a$, $\sinh s=-\cot a$.
\end{proof}
A dichotomy $C_\xi=\{x:\kappa(n,x)>\cos^2\frac a2\}$ is a yes/no question about $x$: is its comparison
with the test point $n$ above the threshold $\cos^2\frac a2$?  By Theorem~\ref{thm:desitter}(iv)(a), $\eta$
lies in the causal future of $\xi$ iff every $x$ that answers yes to $\xi$ answers yes to $\eta$, that is,
iff the first question implies the second.  Along a comoving curve the threshold decreases as $a$
grows, so later questions are coarser.

\begin{proposition}[the static patch of a binary alternative]\label{prop:static}
Fix $x\in M$ and let $\gamma_x$ be the curve of caps centred at $x$, a timelike geodesic with proper time
$s=\ln\tan\frac a2$ (Proposition~\ref{prop:cosmo}).
(i) $\cD$ is the union, up to the hypersurfaces $\{\ip vx=\pm t\}$, of the four open sets
$W_x=\{\ip vx>\lvert t\rvert\}$ (caps containing $x$ whose closure misses $x^\perp$),
$W_{x^\perp}=\{-\ip vx>\lvert t\rvert\}$, $F=\{t>\lvert\ip vx\rvert\}$ (caps containing both) and
$P=\{-t>\lvert\ip vx\rvert\}$ (closures containing neither).
(ii) $I^+(\gamma_x)=\{\xi:x\in C_\xi\}$ and $I^-(\gamma_x)=\{\xi:x^\perp\notin\overline{C_\xi}\}$; hence
$W_x=I^+(\gamma_x)\cap I^-(\gamma_x)$, the set of events that can both receive a signal from $\gamma_x$ and
send one to it.  The horizons $\partial I^\pm(\gamma_x)$ are the dichotomies whose boundary sphere passes
through $x$, respectively $x^\perp$.
(iii) The boosts of $\R^{1,4}$ in the plane of $(1,0)$ and $(0,x)$ fix the null rays of $x$ and $x^\perp$,
preserve $W_x$ and act on $\gamma_x$ by $s\mapsto s+\mathrm{const}$; they are M\"obius maps that do not
preserve $\kappa$.  The stabiliser of $x$ in $\Aut(M)$, a copy of $\SO(3)$, acts on $W_x$ by rotations.
(iv) On $W_x$ define $\varrho\in[0,1)$, $\tau\in\R$, $\omega\in S^2\subset x^\perp$ by
$1-\varrho^2=-f_\xi(x)f_\xi(x^\perp)$, $e^{2\tau}=-f_\xi(x)/f_\xi(x^\perp)$ and
$\omega=(v-c(v,x)x)/\lvert v-c(v,x)x\rvert$ (undefined on $\gamma_x$).  For a cap with centre $n$ and radius $a$, with
$\cos\theta_0=c(n,x)$: $\varrho=\sin\theta_0/\sin a$ and $e^{2\tau}=(\cos\theta_0-\cos a)/(\cos\theta_0+\cos a)$.
The metric is $-(1-\varrho^2)d\tau^2+(1-\varrho^2)^{-1}d\varrho^2+\varrho^2d\omega^2$, the static patch of
$dS_4$ \cite{GibbonsHawking}, and on $\gamma_x$ ($\varrho=0$), $\tau=s=\tfrac12\ln\frac{1-\lambda}\lambda$ for
$C=\{y:\kappa(x,y)>\lambda\}$.
\end{proposition}
\begin{proof}
(i) $f_\xi(x)=t+\ip vx$ and $f_\xi(x^\perp)=t-\ip vx$ since $x^\perp=-x$; the four sign patterns give the
four sets.  (ii) The chronological relation is $\xi\ll\eta$ iff $P(\xi,\eta)>1$ and $t_\xi<t_\eta$, i.e.\ iff the
closed cap of $\xi$ lies in the open cap of $\eta$ (proof of Theorem~\ref{thm:desitter}(iv)).  If $x\in C_\eta$, a
small closed cap about $x$ lies in $C_\eta$; conversely such a cap contains $x$.  If
$x^\perp\notin\overline{C_\eta}$, compactness gives $\varepsilon>0$ with $\overline{C_\eta}$ outside the ball of
radius $\varepsilon$ about $x^\perp$, hence inside the open cap of radius $\pi-\varepsilon/2$ about $x$;
conversely a cap of radius $<\pi$ about $x$ misses $x^\perp$.  The horizons are $f_\xi(x)=0$ and
$f_\xi(x^\perp)=0$.  (iii) The boost fixes $(1,\pm x)$ up to the factors $e^{\pm r}$ and preserves
$\ip vx^2-t^2$ and the sign of $\ip vx$; on $\gamma_x=\{(\sinh s,\cosh s\,x)\}$ it adds $r$ to $s$.  It
changes $\kappa$ by the conformal factor of Theorem~\ref{thm:cone}(ii).  Rotations fixing $x$ fix $t$ and
$\ip vx$.  (iv) With $T=t$, $X_4=c(v,x)$, $X_\perp=v-X_4x$: $W_x=\{X_4>\lvert T\rvert\}$, and
$X_4=\sqrt{1-\varrho^2}\cosh\tau$, $T=\sqrt{1-\varrho^2}\sinh\tau$, $X_\perp=\varrho\,\omega$ give the metric by
substitution in $-dT^2+dX_4^2+\lvert dX_\perp\rvert^2$.  The formulas in $(n,a)$ follow from
$f_\xi(\pm x)=(\pm\cos\theta_0-\cos a)/\sin a$.
\end{proof}

\noindent\emph{Reading under Hypothesis E.}  An observer is an elementary binary alternative
$\{x,x^\perp\}$.  The events it can exchange signals with are exactly the dichotomies that separate its two
alternatives ($x$ lies in the open cap and $x^\perp$ outside the closed cap); its horizon is reached where a test stops separating them;
its time translation is the one-parameter group of M\"obius maps fixing $x$ and $x^\perp$, which acts on
the state cone $K_M$ as the filters act on $K_T$ (Section~\ref{sec:velocity}).  With a clock the same
construction applied to $\beta$ on $P_T\cong S^2$ gives $dS_3$, and the kernel space
$\R\oplus\R^3\cong\Herm(V,J_T)$, ordered by $K_T$, is Minkowski space, whose order automorphisms are
generated by the orthochronous Lorentz group, translations and dilations \cite{AO53,Zeeman}.

\smallskip\noindent\textbf{Hypothesis E.}  Physical events are the dichotomies of an elementary
comparison structure, and their physical causal order is inclusion.

\begin{proposition}[the metric is forced by symmetry]\label{prop:invmetric}
(i) Every symmetric $2$-tensor field on $\cD$ invariant under the M\"obius group $\SO^+(1,4)$ is a constant
multiple of $g$.  Hence every invariant Lorentzian metric of signature $(-,+,+,+)$ is $g'=\ell^2g$ with $\ell>0$, and it satisfies
$G'+\Lambda'g'=0$ with $\Lambda'\ell^2=3$.
(ii) Every bijection of $\cD$ that preserves the inclusion order of caps in both directions is a smooth
conformal isometry of $(\cD,g)$.
\end{proposition}
\begin{proof}
(i) $\SO^+(1,4)$ acts transitively with isotropy $\SO^+(1,3)$ acting on $T_\xi\cD=\xi^\perp\cong\R^{1,3}$ by its
defining representation, which is absolutely irreducible; an invariant field is fixed by its value at
one point, an isotropy-invariant symmetric form, and by Schur's lemma such forms are multiples of the
Minkowski form.  Then $\mathrm{Ric}(\ell^2g)=\mathrm{Ric}(g)=3g$.  (ii) By Theorem~\ref{thm:desitter}(iv) the
inclusion order is the causal order $\le$ of the globally hyperbolic spacetime $(\cD,g)$.  It determines the
chronological order: for $\xi\le\eta$, $\xi\ne\eta$, one has $\xi\ll\eta$ iff the order interval
$\{\zeta:\xi\le\zeta\le\eta\}$ is not a chain.  Indeed, if $\xi\ll\eta$ the interval contains the open set
$I^+(\xi)\cap I^-(\eta)$, which contains causally unrelated points; if $\xi\le\eta$ but not $\xi\ll\eta$, then $P=1$,
$\eta-\xi$ is a null vector of $\R^{1,4}$, the null geodesic from $\xi$ to $\eta$ is unique, every causal curve
from $\xi$ to $\eta$ is that geodesic \cite{ONeill}, and the interval is its segment, a chain.  So an order
bijection preserves $\ll$ in both directions, and a chronological isomorphism between distinguishing
spacetimes is a smooth conformal isometry \cite{HKM76,Malament77}; see also \cite{Levichev87} for the causal
relation directly.
\end{proof}
Under Hypothesis~E the dimension four, the Lorentzian signature and the causal order follow from the
axiom, and every metric that respects the symmetry of the structure is de~Sitter, hence solves the
vacuum equations with $\Lambda\ell^2=3$; this is a property of the invariant metric and does not select the
field equations, which also have non-invariant solutions.  The dimensionless number $\Lambda\ell^2=3$ is fixed; $\Lambda$ itself only fixes the
unit of length, which the structure does not supply.  By (ii), the order of dichotomies alone determines the
conformal class of the metric.  Hypothesis~E itself is not derived.

\subsection{A soldering form}\label{sec:solder}
For $\xi=(t,v)\in\cD$ let $n:=v/\lvert v\rvert$ be the centre of the cap $C_\xi$.  A tangent vector at $\xi$
is $(\tau,w)$ with $t\tau=\ip{v}{w}$; put $w_\perp:=w-\ip{w}{n}n$.  On $\R\oplus\im\HH$ use the form
$\eta(a_0,q):=-a_0^2+\lvert q\rvert^2$.

\begin{proposition}\label{prop:solder}
Define $\sigma^L_\xi(\tau,w):=(\tau/\lvert v\rvert,\ \bar n\,w_\perp)$ and
$\sigma^R_\xi(\tau,w):=(\tau/\lvert v\rvert,\ w_\perp\bar n)$.
(i) $\bar nw_\perp$ and $w_\perp\bar n$ are imaginary, and $\sigma^L_\xi$, $\sigma^R_\xi$ are linear
isometries from $(T_\xi\cD,g)$ onto $(\R\oplus\im\HH,\eta)$.
(ii) $\sigma^L$ is invariant under left transports and equivariant under right transports,
$\sigma^L_{\xi b}(\tau,wb)=(\tau/\lvert v\rvert,\bar b\,(\bar nw_\perp)\,b)$; $\sigma^R$ is invariant under right
and equivariant under left transports; the transport reversal $z\mapsto\bar z$ carries $\sigma^L$ to
$\sigma^R$ up to the sign of the spatial part.
(iii) $E_0:=(\lvert v\rvert,t\,n)$ and $E_q:=(0,nq)$, $q\in\im\HH$ a unit, satisfy
$\sigma^L(E_0)=(1,0)$, $\sigma^L(E_q)=(0,q)$: a global orthonormal frame (tetrad) of $\cD\cong\R\times S^3$.
With $t=\sinh s$, $v=\cosh s\cdot n$, the metric is $-ds^2+\cosh^2s\,g_{S^3}$, $E_0=\partial_s$, and
$E_q=(\cosh s)^{-1}X_q$ with $X_q(n)=nq$ the left-invariant fields, which generate the right transports
and are parallel for the left parallelism of the heap \cite[Prop.~5.3]{PaperI}.
(iv) The clock field, the tangent $nu$ of the orbit $n\,e^{\chi u}$, satisfies $\sigma^L(0,nu)=(0,u)$, a
constant spacelike unit vector, and $\sigma^R(0,nu)=(0,nu\bar n)$.
(v) The flat metric connection for which the tetrad is parallel has torsion
\[
T(E_0,E_q)=\frac{t}{\lvert v\rvert}\,E_q=\tanh s\;E_q,\qquad
T(E_p,E_q)=-\frac{2}{\lvert v\rvert}\,E_{p\times q}=-\frac1{\cosh s}\,E_{pq-qp}.
\]
(vi) For a clock $T$ with generator $u$, the map
$\Sigma_\xi(\tau,w):=(\tau/\lvert v\rvert)\mathbf1+J_T\ell_{w_\perp\bar n}$, with $\ell_q$ left multiplication,
is a linear isometry from $(T_\xi\cD,-g)$ onto $(\Herm(V,J_T),\det)$, equivariant under $\cL T$.
\end{proposition}
\begin{proof}
(i) $\re(\bar nw_\perp)=\ip{n}{w_\perp}=0$.  Since $\ip{w}{n}=t\tau/\lvert v\rvert$ and
$\lvert v\rvert^2-t^2=1$, $\lvert w\rvert^2-\tau^2=\lvert w_\perp\rvert^2+t^2\tau^2/\lvert v\rvert^2-\tau^2
=\lvert w_\perp\rvert^2-\tau^2/\lvert v\rvert^2$; an isometry between nondegenerate spaces of equal dimension
is bijective.  (ii) Under $z\mapsto az$, $n\mapsto an$ and $w_\perp\mapsto aw_\perp$, so $\bar nw_\perp$ is
unchanged; under $z\mapsto zb$ it becomes $\bar b\bar nw_\perp b$; conjugation gives
$\overline{w_\perp\bar n}=n\bar w_\perp=-w_\perp\bar n$ for the reversed point.  (iii) $(\lvert v\rvert,tn)$ is
tangent since $t\lvert v\rvert=\ip{v}{tn}$, and has $w_\perp=0$; $(0,nq)$ is tangent with $w_\perp=nq$.
In the coordinates $(s,n)$, $\partial_s=(\cosh s,\sinh s\,n)=E_0$ and a displacement $\delta n$ moves $v$ by
$\cosh s\,\delta n$.  (iv) $\bar n\,nu=u$.  (v) For a parallel frame the torsion is $T(E_a,E_b)=-[E_a,E_b]$.
$[\partial_s,(\cosh s)^{-1}X_q]=-\tanh s\,E_q$, and $[X_p,X_q]=X_{pq-qp}=2X_{p\times q}$ gives
$[E_p,E_q]=(\cosh s)^{-1}E_{pq-qp}$.  (vi) $J_T$ is right multiplication by $u$ and commutes with
$\ell_q$; $\ell_q$ is skew-Hermitian on $(V,J_T)$ with $\ell_q^2=-\lvert q\rvert^2$, so $J_T\ell_q$ is Hermitian,
traceless, with eigenvalues $\pm\lvert q\rvert$, and $\det(a_0\mathbf1+J_T\ell_q)=a_0^2-\lvert q\rvert^2=-\eta(a_0,q)$.
For $a\in\cL$, $\ell_aJ_T\ell_q\ell_a^{-1}=J_T\ell_{aq\bar a}$, which matches (ii); $T$ commutes with
$J_T\ell_q$ and fixes $\sigma^R$.
\end{proof}
\begin{proposition}[coframe and uniqueness]\label{prop:solderunique}
(i) In the coordinates of Proposition~\ref{prop:cosmo}, $\sigma^L=(e^0,e)$ with $e^0=da/\sin a=ds$ and
$e=\theta/\sin a=\cosh s\,\theta$, where $\theta=n^{-1}dn$ is the $\im\HH$-valued Maurer--Cartan form of the
left parallelism of the heap.  Cartan's structure equation $d\theta=-\theta\wedge\theta$ gives
$de=\tanh s\;e^0\wedge e-(\cosh s)^{-1}\,e\wedge e$, which is the torsion of
Proposition~\ref{prop:solder}(v).
(ii) Let $\sigma'$ be a continuous field of linear isometries $(T_\xi\cD,g)\to(\R\oplus\im\HH,\eta)$ that
is invariant under left transports and equivariant under right transports, through $\mathrm{Ad}$ on
$\im\HH$ and trivially on $\R$.  Then $\sigma'=D\sigma^L$ with $D=\mathrm{diag}(\pm1,\pm\mathbf1_3)$.  The
same holds for $\sigma^R$ with left and right exchanged.
\end{proposition}
\begin{proof}
(i) $\sigma^L$ of $\partial_a$ is $(\lvert v\rvert^{-1}\sin^{-2}a,0)=(1/\sin a,0)$, and of a spatial displacement
$dn$ it is $(0,\bar n\,dn/\sin a)$.  Then $de=\sinh s\,ds\wedge\theta+\cosh s\,d\theta$, and for a parallel frame
the torsion is $de$; $(e\wedge e)(E_p,E_q)=pq-qp=2p\times q$.
(ii) $\Lambda:=\sigma'\circ(\sigma^L)^{-1}$ takes values in $\OO(1,3)$.  Invariance of both under left
transports, which act transitively on $n$ at fixed $a$, makes $\Lambda$ a function of $a$ alone; equivariance
of both under right transports makes $\Lambda(a)$ commute with $\mathrm{Ad}(\SU(2))=\SO(3)$ on the spatial part.
The commutant of $\SO(3)$, irreducible on $\R^3$ and trivial on $\R$, is $\{\mathrm{diag}(\lambda,\mu\mathbf1_3)\}$,
and $\Lambda(a)\in\OO(1,3)$ forces $\lambda,\mu\in\{\pm1\}$; by continuity $\Lambda$ is constant.
\end{proof}
The spatial part of the torsion is the torsion $-[X_p,X_q]$ of the left parallelism of the unit
three-sphere \cite[Prop.~5.3(i)]{PaperI}, divided by the scale factor $\cosh s$; the temporal part is the expansion rate
$\tanh s$.  The soldering is therefore unique up to the signs $D$ and the choice of side, which
transport reversal exchanges.

%=====================================================================================
\section{Fields of comparison structures}\label{sec:fields}

\smallskip\noindent\textbf{Hypothesis F.}  (i) A manifold $X$ is given with a locally trivial bundle
$E\to X$ whose fibres are comparison structures isomorphic to $S^3$; a \emph{configuration} is a section
$\sigma$.  (ii) A \emph{clock field}, a smooth choice of clock $T_p$ in each fibre, is given.  (iii) The
connection preserves the clock field.  (iv) The kinetic form of configurations is
$\Aut(M)$-invariant, i.e.\ built from the Hessian form $h$ of $\kappa$.

\begin{theorem}\label{thm:gauge}
Assume Hypothesis~F(i); (ii) uses F(ii), (iii); (iv) uses F(iv).
(i) $E$ is associated to a principal $\Aut(M)=\SO(4)$ bundle; no isomorphism $E_p\to E_{p'}$ is
distinguished, so $\kappa(\sigma(p),\sigma(p'))$ is undefined for $p\ne p'$, and comparison along a
path requires a connection, locally an $\su(2)_L\oplus\su(2)_R$-valued one-form.
(ii) A clock field reduces the group to $\cL T\cong\U(2)$, acting on $(V,J_T)\cong\C^2$ by the defining
representation; clock-preserving connections are pairs $(W,B)$ with values in $\su(2)\oplus\mathfrak u(1)$.
(iii) The stabiliser of $o\in M$ in $\cL T$ is a circle embedded diagonally, generated by $Q=T_3+Y$,
with $Y$ the generator of the clock and $T_3$ the generator of $\cL$ that $o$ carries to it; the
stabiliser of $o$ in $\Aut(M)$ is $\SO(3)$ and acts irreducibly on $T_oM$.
(iv) Up to a factor $v^2$ there is exactly one density quadratic in the covariant derivative of a
configuration whose fibre form is $\Aut(M)$-invariant, namely $v^2h(D\sigma,D\sigma)$ with $h$ the
Hessian form of $\kappa$.  In the gauge $\sigma\equiv o$, with connection components $gW^a$, $g'B$ in an
$h$-orthonormal basis of $\su(2)$ whose third vector generates the clock,
$h(D\sigma,D\sigma)=g^2((W^1)^2+(W^2)^2)+(gW^3-g'B)^2$; its null direction $g'W^3+gB$ is the stabiliser
of (iii), and its nonzero eigenvalues are in the ratio $g^2:(g^2+g'^2)$.
\end{theorem}
\begin{proof}
(i), (ii) restate Propositions~\ref{prop:aut} and \ref{prop:clock} fibrewise.  (iii) In the model
$(a,t)\cdot o=aot^{-1}$, so the stabiliser of $o=1$ is the diagonal $\{(t,t)\}$.  (iv) An $\Aut(M)$-invariant form on
$T_oM$ is $\SO(3)$-invariant, hence a multiple of $h$; the infinitesimal motion of $o=1$ under
$(\xi,y)\in\su(2)\oplus\R$ is $\xi-ye_3$.
\end{proof}

\begin{proposition}[invariance under the clock-preserving group]\label{prop:berger}
A quadratic form on $T_oM$ invariant only under the stabiliser $\U(1)_Q$ of $o$ in $\cL T$ is
$q=\lambda_1(x_1^2+x_2^2)+\lambda_3x_3^2$ in the basis of Theorem~\ref{thm:gauge}(iv) (the Berger family), and
the corresponding mass form gives $m_W^2/m_Z^2=(\lambda_1/\lambda_3)\cos^2\theta_W$ with $\tan\theta_W=g'/g$.
\end{proposition}
\begin{proof}
$\U(1)_Q$ acts on $T_1M=\im\HH$ by conjugation with $e^{i\phi}$, the rotations about the clock axis
$e_3$; the invariant symmetric forms of such rotations are $\lambda_1(x_1^2+x_2^2)+\lambda_3x_3^2$.  The mass
form is then $\lambda_1g^2((W^1)^2+(W^2)^2)+\lambda_3(gW^3-g'B)^2$, with eigenvalues $\lambda_1g^2$ (twice),
$\lambda_3(g^2+g'^2)$ and $0$.
\end{proof}

\begin{remark}[electroweak structure]\label{rem:ew}
Theorem~\ref{thm:gauge} is the symmetry structure of the electroweak theory with a fixed-length scalar
doublet \cite{Weinberg67,AB80}: gauge group $\U(2)$, hypercharge acting as the clock, residual $\U(1)$
with $Q=T_3+Y$, three massive directions and one massless, and the ratio $g^2:(g^2+g'^2)$ of the mass-form
eigenvalues, which with F$_{\mathrm g}$(a) is $m_W=m_Z\cos\theta_W$ at tree level.  By
Proposition~\ref{prop:berger}, this ratio holds because the fibre form is invariant under both
transports $(\Omega2\mathrm c)$, i.e.\ $\lambda_1=\lambda_3$; invariance under the reduced group $\U(2)$ alone
would leave the ratio $\lambda_1/\lambda_3$, the $\rho$ parameter, free.  The stabiliser $\SO(3)$ of (iii) is
the custodial symmetry \cite{Veltman77,SSVZ}.  Not derived: $g$, $g'$, $v$, the radial mode (it lives in
$V\supset M$), matter fields, and the dynamics of the connection; reading the eigenvalues of the mass form as
masses requires a gauge kinetic term (Hypothesis~F$_{\mathrm g}$).
\end{remark}

\begin{remark}[the mixing angle for one structure]\label{rem:angle}
Transport reversal lies in $\Aut^\pm(M)$ and exchanges $\cL$ with $\cR$.  If the kinetic form of the
connection is $\Aut^\pm(M)$-invariant on $\su(2)_L\oplus\su(2)_R$, the two couplings are equal; since the
clock lies in $\cR$, $g'=g$ and $\sin^2\theta_W=\tfrac12$ at the scale at which this form applies.  The
low-energy value is $\hat s_Z^2=0.23122(6)$ \cite{ex-PDG26}.  A boundary value at a high scale does not contradict it;
Section~\ref{sec:sm} gives the value $\tfrac38$ for three structures and the scales at which the
Standard Model couplings reach either value (Remark~\ref{rem:scales}).
\end{remark}

\begin{remark}[gravitation]\label{rem:gravity}
Under Hypothesis~F with $\dim X=4$ the cones $K_{T_p}$ form a bundle of Lorentz cones of rank four, and
a metric on $X$ requires a soldering form.  For $X=\cD$ with a constant clock field,
Proposition~\ref{prop:solder}(vi) supplies one, unique up to signs and side under full left and
right equivariance (Proposition~\ref{prop:solderunique}); under the clock-preserving equivariance alone,
boosts in the plane of time and the clock axis commute with the isotropy and give further solderings, and the metric it induces is the de~Sitter metric.
\end{remark}

\subsection{The centre field on the space of events}\label{sec:centre}
Assume Hypotheses~E and F with $X=\cD$ and the trivial bundle $\cD\times M$.  The \emph{centre field}
$\sigma_0(\xi):=n$, the centre of the cap $C_\xi$, is a configuration determined by the structure; it is
equivariant under $\OO(4)$ and its restriction to each slice $a=\mathrm{const}$ is the identity map of $M$,
of degree one.  Take the action
$S_v[\sigma]=-\tfrac{v^2}2\int g^{\mu\nu}\ip{\partial_\mu\sigma}{\partial_\nu\sigma}\,d\mathrm{vol}_g$, with
$\ip{\cdot}{\cdot}$ the round metric $2h$ on $M\subset V$: by Theorem~\ref{thm:gauge}(iv) the fibre form is
unique up to the factor $v^2$.  Its Euler--Lagrange equation for maps into the unit sphere of $V$ is
$\Box_g\sigma+g^{\mu\nu}\ip{\partial_\mu\sigma}{\partial_\nu\sigma}\,\sigma=0$ (wave maps).

\begin{proposition}\label{prop:centre}
(i) $\sigma_0$ is a wave map.
(ii) With the flat connection $A:=dn\,n^{-1}$, the right-invariant Maurer--Cartan form of the heap
pulled back by $\sigma_0$ and acting by left transports, $D\sigma_0=d\sigma_0-A\sigma_0=0$; the gauge
transformation $n^{-1}$ maps $(\sigma_0,A)$ to $(1,0)$.  In the gauged theory with this flat connection
($D\sigma_0=0$, $F_A=0$), $\sigma_0$ carries no energy and the de~Sitter metric is undisturbed; gauged
configurations in general carry stress.
(iii) With the trivial connection, the stress tensor
$T_{\mu\nu}=v^2\bigl(\ip{\partial_\mu\sigma_0}{\partial_\nu\sigma_0}-\tfrac12g_{\mu\nu}g^{\rho\lambda}
\ip{\partial_\rho\sigma_0}{\partial_\lambda\sigma_0}\bigr)$ is that of a perfect fluid at rest in the
coordinates of Proposition~\ref{prop:cosmo}, with energy density $\rho=\tfrac32v^2\sin^2a=\tfrac32v^2/R^2$
and pressure $p=-\rho/3$.
\end{proposition}
\begin{proof}
(i) $\sqrt{\lvert g\rvert}g^{\mu\nu}=\sin^{-2}a\,\sqrt{\lvert\hat g\rvert}\hat g^{\mu\nu}$ for
$\hat g=-da^2+g_M$, and $\sigma_0$ does not depend on $a$, so $\Box_g\sigma_0=\sin^2a\,\Delta_M n=-3\sin^2a\,n$,
the coordinate functions of the unit $S^3$ being eigenfunctions of $\Delta_M$ with eigenvalue $-3$; and
$g^{\mu\nu}\ip{\partial_\mu n}{\partial_\nu n}=3\sin^2a$.
(ii) $dn-dn\,n^{-1}n=0$, and $dA=A\wedge A$.
(iii) In the orthonormal frame $(\sin a\,\partial_a,\sin a\,e_i)$ the matrix $\ip{\partial\sigma_0}{\partial\sigma_0}$ is
$\mathrm{diag}(0,\sin^2a\,\mathbf1_3)$, with trace $3\sin^2a$.
\end{proof}

\smallskip\noindent\textbf{Hypothesis GR.}  The metric of the space of events is dynamical and obeys
$G_{\mu\nu}+\Lambda g_{\mu\nu}=8\pi G_NT_{\mu\nu}$.

\smallskip\noindent Without sources every invariant metric $\ell^2g$ solves these equations with $\Lambda=3/\ell^2$
(Proposition~\ref{prop:invmetric}); this is a condition on $\Lambda$ for that metric.  The vacuum equations also
have non-invariant solutions, and a source proportional to the metric only shifts $\Lambda$.  Lovelock's
theorem \cite{Lovelock} states that in four dimensions the symmetric divergence-free tensors built locally
from the metric and its first two derivatives are combinations of $G_{\mu\nu}$ and $g_{\mu\nu}$; it does not
supply a matter action.  Hypothesis~GR is not a consequence of the axiom.

\begin{proposition}[pointwise obstruction]\label{prop:nullobstruction}
Assume Hypotheses~E, F with the trivial connection, and GR with $G_N>0$, the invariant metric $\ell^2g$,
and the centre field $\sigma_0$ with $v\ne0$ as the only source besides $\Lambda$.  Then the field equations fail
at every event, for every $\Lambda$: for every null vector $k$, $T(k,k)=v^2\lvert d\sigma_0(k)\rvert^2>0$, while
$(G+\Lambda g)(k,k)=0$.  Additional sources satisfying the null energy condition do not remove the
failure.  For any configuration $\sigma$ with $v\ne0$ on this metric, the equations on null vectors force
$d\sigma=0$; a constant configuration solves them for every $v$ when $\Lambda\ell^2=3$, and its slices have degree zero, while the
centre field has degree one.
\end{proposition}
\begin{proof}
$G=-3g$ for the invariant metric, so $(G+\Lambda g)(k,k)=0$ for null $k$.  For null $k$ the trace term of
$T$ drops out and $T(k,k)=v^2\ip{d\sigma(k)}{d\sigma(k)}$.  For $\sigma_0$ and $k=e_0\pm e_i$ in an orthonormal
frame, $d\sigma_0(k)=d\sigma_0(e_i)\ne0$ (Proposition~\ref{prop:centre}(iii)).  Null vectors span each tangent space.
\end{proof}

\begin{proposition}[the centre field as a source]\label{prop:friedmann}
Assume Hypotheses~F (trivial connection) and GR, the centre field $\sigma_0$ with $\gamma:=4\pi G_Nv^2>0$, and
put $L=\Lambda/3$.
(i) If the metric $g$ on the cap manifold $(0,\pi)\times M$ is smooth and its causal order is inclusion of caps
(Hypothesis~E), then $g=\Omega(a)^2(-da^2+g_M)$ with $\Omega$ depending on $a$ alone: the metric is of
Friedmann--Lema\^itre--Robertson--Walker form $-ds^2+R(s)^2g_M$, $ds=\Omega\,da$, $R=\Omega$.
(ii) For this form $\sigma_0$ is a wave map for every $R$, $\rho=\tfrac32v^2/R^2$, $p=-\rho/3$, and the field
equations are $\dot R^2=LR^2+\gamma-1$, $\ddot R=LR$; the scalar curvature is $4\Lambda+6\gamma/R^2$.  The maximal
solutions, up to time translation and reversal, and their conformal durations $\int ds/R$ are:
\begin{center}\footnotesize
\begin{tabular}{@{}lll@{}}
\toprule
& solution & conformal duration\\
\midrule
$L>0$, $\gamma<1$ & $R=R_0\cosh(s/\ell)$, $\ell=L^{-1/2}$, $R_0=\ell\sqrt{1-\gamma}$, $s\in\R$ & $\pi/\sqrt{1-\gamma}>\pi$\\
$L>0$, $\gamma=1$ & $R=R_1e^{s/\ell}$, $s\in\R$ & infinite\\
$L>0$, $\gamma>1$ & $R=\ell\sqrt{\gamma-1}\,\sinh(s/\ell)$, $s>0$ & infinite (logarithmic at $s=0$)\\
$L=0$, $\gamma=1$ & $R$ constant, $s\in\R$ & infinite\\
$L=0$, $\gamma>1$ & $R=\sqrt{\gamma-1}\,s$, $s>0$ & infinite\\
$L<0$, $\gamma>1$ & $R=\ell\sqrt{\gamma-1}\,\sin(s/\ell)$, $\ell=(-L)^{-1/2}$, $0<s<\pi\ell$ & infinite\\
\bottomrule
\end{tabular}
\end{center}
For $L=0$, $\gamma<1$ and for $L<0$, $\gamma\le1$ there is no solution.
(iii) No maximal solution of (ii) is order-isomorphic, by any bijection, to the caps ordered by
inclusion.  By (i), every solution of Hypotheses~E, F (ungauged) and GR with $v\ne0$ is the restriction of a
branch of (ii) to an open conformal-time interval of length $\pi$; it carries the full cap order through
$(\eta,n)\mapsto C(n,\eta-\eta_0)$, and it is not maximal.
\end{proposition}
\begin{proof}
(i) Futures and pasts of the cap order are closed and separate points (Theorem~\ref{thm:desitter}(iv)), so the
spacetime is distinguishing, and the identity map is a causal isomorphism onto $(\cD,g_{dS})$; by Levichev's
theorem \cite{Levichev87} $g$ is conformal to $g_{dS}$, hence $g=\Omega^2\bar g$ with $\bar g=-da^2+g_M$ and $\Omega>0$ a
function of $(a,n)$.  In dimension four $\Box_g\sigma+\lvert d\sigma\rvert_g^2\sigma=\Omega^{-2}\bigl(\Box_{\bar g}\sigma+\lvert d\sigma\rvert_{\bar g}^2\sigma+2\,d\sigma(\mathrm{grad}_{\bar g}\log\Omega)\bigr)$.  For
$\sigma_0(a,n)=n$ the first two terms cancel ($\Delta_Mn=-3n$), and the field equation of $\sigma_0$ leaves
$d\sigma_0(\mathrm{grad}_{g_M}\log\Omega)=0$; $d\sigma_0$ is injective on spatial vectors, so $\Omega=\Omega(a)$.
(ii) The computation of Proposition~\ref{prop:centre}(i), (iii) with $\sin a$ replaced by $1/R$ gives the
wave-map equation and $\rho$, $p$.  For spatial curvature $+1$ and a perfect fluid, Einstein's equations
are the Friedmann equations $\dot R^2/R^2=\tfrac{8\pi G_N}3\rho+L-R^{-2}$ and
$\ddot R/R=-\tfrac{4\pi G_N}3(\rho+3p)+L$ \cite{HawkingEllis}, with $\rho+3p=0$.  The second equation gives
$R=Ae^{s/\ell}+Be^{-s/\ell}$, $R=A+Bs$ or $R=A\sin(s/\ell)+B\cos(s/\ell)$, and the first fixes the constants.  The
conformal times are $(2\ell/R_0)\arctan e^{s/\ell}$, $-(\ell/R_1)e^{-s/\ell}$, $(\gamma-1)^{-1/2}\log\tanh(s/2\ell)$,
$s/R$, $(\gamma-1)^{-1/2}\log s$ and $(\gamma-1)^{-1/2}\log\tan(s/2\ell)$, up to constants.
(iii) In conformal time each solution is $(I\times M,R^2(-d\eta^2+g_M))$ with $\lvert I\rvert>\pi$, and $(\eta,n)\le(\eta',m)$ iff
$\eta'-\eta\ge d(n,m)$.  If $\lvert I\rvert>\pi$, some point $p$ has more than $\pi$ of conformal time to its future, and
every point $y$ has a common upper bound with $p$ (a point over the position of $y$, later than $y$ and than
$\eta_p+\pi$).  In the cap order this fails: for $p=(a,n)$ and $y=(b,-n)$ with $a+b>\pi$, an upper bound $(c,m)$
needs $2c-a-b\ge d(n,m)+d(-n,m)=\pi$, while $c<\pi$ gives $2c-a-b<\pi$.  An order isomorphism preserves this
property.  The restriction to an open interval of length $\pi$ is the cap cylinder with the conformal factor
$R^2$, which does not change the causal order; on $[\eta_0,\eta_0+\pi)$ a point of the initial
slice has a common upper bound with every point, and on $(\eta_0,\eta_0+\pi]$ a point of the final slice has a
common lower bound with every point; the caps have neither property, so neither half-open restriction is
order-isomorphic to them.
\end{proof}

\begin{proposition}[linear stability of the centre field]\label{prop:stability}
Assume Hypotheses~E and F (trivial connection).  A perturbation $V(s,n)\in T_nM$ of $\sigma_0$ on
$(\cD,-ds^2+\cosh^2s\,g_M)$ solves, at linear order, $\ddot V+3\tanh s\,\dot V+\cosh^{-2}s\,\mathcal JV=0$,
where $\mathcal J=\nabla^*\nabla-2$ is the Jacobi operator of the identity map of $M$.  The spectrum of $\mathcal J$ is:
on gradients of scalar harmonics of degree $l\ge1$, $l(l+2)-4$ with multiplicity $(l+1)^2$; on divergence-free
fields, $(l+1)^2-4$ with multiplicity $2l(l+2)$, $l\ge1$.  So the index is four ($\lambda=-1$, gradients of first
harmonics), the kernel is the six-dimensional space of Killing fields, and all other eigenvalues are
positive.  For every eigenvalue $\lambda$ and data at $s_0\ge0$, the mode $V(s)Y$ stays bounded for $s\ge s_0$; for
$\lambda=-1$, $V=\tanh s$ is a solution.  In the contracting phase some data are amplified: for $\lambda=0$,
$\dot V(s_0)=\delta$ at $s_0<0$ gives $V(0)-V(s_0)=\delta\cosh^3s_0\int_{s_0}^0\cosh^{-3}$, asymptotic to
$\delta\,\tfrac{\pi}{32}e^{3\lvert s_0\rvert}$.  On the static cylinder of radius $R_*$ the four modes with $\lambda=-1$
grow like $e^{t/R_*}$, the zero modes are affine in $t$, and the others oscillate.
\end{proposition}
\begin{proof}
The linearised wave-map operator at $\sigma_0$ is $-\cosh^{-3}s\,\partial_s(\cosh^3s\,\partial_sV)+\cosh^{-2}s
(\Delta V+\mathrm{Ric}_M V)$ with $\mathrm{Ric}_M=2$; this is the stated equation.  Spectrum: in the left
frame $X_i(n)=nI_i$, $\nabla_{X_i}=2L_i+S_i$ on frame components, with $L_i=X_i/2$ and $S_i$ the spin-one
matrices; both satisfy $[A_i,A_j]=\epsilon_{ijk}A_k$ and commute, so
$\nabla^*\nabla=2C_L+2C_{L+S}-C_S$ with Casimirs $C=j(j+1)$.  By the Peter--Weyl theorem the spin $j$ of $L$ runs
over $0,\tfrac12,1,\dots$ with multiplicity $2j+1$, and $L+S$ has spin $k$ with $\lvert j-1\rvert\le k\le j+1$ in unit steps (the triangle rule); hence
$\lambda=2j(j+1)+2k(k+1)-4$ with multiplicity $(2j+1)(2k+1)$.  The block $k=j$, $l=2j$, consists of the gradients
(they are equivariant and fill its dimension $(l+1)^2$); the blocks $k=j\pm1$ are divergence-free and give
$(l+1)^2-4$ with $l=2j\pm1$.  Killing fields have $\nabla^*\nabla K=2K$.  Write the time equation as
$(\cosh^3s\,\dot V)^\cdot=-\lambda\cosh s\,V$.  For $\lambda\ge0$ the energy
$\tfrac12\dot V^2+\tfrac12\lambda\cosh^{-2}s\,V^2$ is nonincreasing for $s\ge0$, so $\dot V$ is bounded and $V$ grows
at most linearly; then $\lvert\dot V(s)\rvert\le C(1+s)e^{-2s}$, which is integrable, so $V$ is bounded.  For
$\lambda=-1$, with $m(s)=\max_{[s_0,s]}\lvert V\rvert$, $\lvert\dot V\rvert\le(\cosh^3s_0\lvert\dot V(s_0)\rvert+m\sinh s)/\cosh^3s$,
and integrating gives $m\le A+m/2$.  For $\lambda=0$, $\cosh^3s\,\dot V$ is constant, and
$\int_{-\infty}^0\cosh^{-3}=\pi/4$.  On the cylinder the equation is $\ddot V+R_*^{-2}\mathcal JV=0$.
\end{proof}
The identity map of $S^3$ is unstable as a harmonic map \cite{SmithHarmonic,Xin80}, and on a static
closed universe the one-texture configuration has a collapse instability \cite{CHKL97}; this is our
eigenvalue $\lambda=-1$, whose eigenspace is four-dimensional (one unstable mode up to $\OO(4)$).  Chen
\emph{et al.}\ also find that expansion lengthens the instability time; on the expanding half of $\cD$ the
modes stay bounded and the $\lambda=-1$ mode stops growing, for data in that half.  These are statements on
the fixed background; the coupled metric--texture perturbations and nonlinear stability are open.  A
four-derivative (Skyrme) term changes this picture on compact slices: Proposition~\ref{prop:skcentre}.

The centre field restricted to a slice is the identity map of the three-sphere, the configuration of
a global texture wound once around a closed universe; its energy density scales as $R^{-2}$, like
spatial curvature, as found by Davis \cite{Davis87,Davis87b} and used by Kamionkowski and Toumbas
\cite{KT96}.  It renormalises the curvature term of the Friedmann equation.  By
Propositions~\ref{prop:nullobstruction} and \ref{prop:friedmann}, for the ungauged centre field with $v\ne0$ the
invariant metric is not a solution, and every solution of Hypotheses~E, F and GR is the restriction of a
longer FLRW branch to an open conformal-time interval of length $\pi$; this restriction carries the full cap
order on the whole set of caps, and it is extendible.  So no maximal solution carries the full cap order;
maximality is the additional premise of this statement.  Dropping it means accepting a spacetime that can
be extended beyond the events supplied by the comparison structure, the added events not being
dichotomies.

\smallskip\noindent\emph{Flat gauging.}  With a dynamical connection, take the action
$S=-\tfrac{v^2}2\int\lvert D\sigma\rvert^2+S_{\mathrm{YM}}$ with a positive Yang--Mills term.  The pair $(\sigma_0,A)$ with $A=dn\,n^{-1}$ of
Proposition~\ref{prop:centre}(ii) is gauge-equivalent to $(1,0)$, by the global gauge transformation $n^{-1}$;
both stress tensors vanish, and $\ell^2g$ solves the field equations with $\Lambda=3/\ell^2$.  The transformation $n^{-1}$ is
large: it has degree $-1$ on each slice, so the degree of the centre field moves into the Chern--Simons
number of $A$, and only their difference is gauge invariant; $(\sigma_0,A)$ and $(1,0)$ are the same vacuum only if
large gauge transformations are identified, that is, if configurations related by them are declared
physically equal.  Under F$_{\mathrm g}$(b) below the centre field is pure gauge.  Conversely, on the
invariant metric the field equations require the stress to be proportional to $g$, so it vanishes on null
vectors; both stresses are nonnegative there, the sigma term vanishes on all null $k$ only if $D\sigma=0$, and
the Yang--Mills term $\sum\lvert F(k,\cdot)\rvert^2$ over the screen of $k=e_0\pm e_i$ vanishes for all $i$ only if
$F=0$.  In unitary gauge $\sigma\equiv1$, with canonically normalised fields and generators $I_a/2$ of the two
transport factors (the $h$-orthonormal basis of Theorem~\ref{thm:gauge}(iv) is $e_a=\sqrt2I_a$, so $g$ here is $2\sqrt2$ times
the $g$ there, and likewise $g'$), the quadratic part is $-\tfrac{v^2}8\bigl[g^2((W^1)^2+(W^2)^2)+(gW^3-g'B)^2\bigr]$, so
$m_W=\lvert v\rvert g/2$, $m_Z=\tfrac{\lvert v\rvert}2\sqrt{g^2+g'^2}$ and $m_\gamma=0$.  We state this as an option:

\smallskip\noindent\textbf{Hypothesis F$_{\mathrm g}$ (gauged reading).}  (a) The connection of Hypothesis~F is
dynamical, with a positive Yang--Mills term and the covariant action $-\tfrac{v^2}2\int\lvert D\sigma\rvert^2$.  (b) The cap
geometry with the centre field is the flat, covariantly constant vacuum of (a).

\smallskip\noindent Under F the eigenvalues of the mass form are in the ratio $g^2:(g^2+g'^2)$; under F and
F$_{\mathrm g}$(a) they are tree-level masses, and $m_W=m_Z\cos\theta_W$ with a massless photon.  Under E,
F$_{\mathrm g}$(a), (b) and GR the full cap order and the de~Sitter metric with $\Lambda\ell^2=3$ hold together.  Open
under F$_{\mathrm g}$: the coupled stability of the vacuum, a calibration of the length $\ell$, and the
matching scale and spectrum that would fix $g$, $g'$ and $v$.  Not derived: Hypothesis~GR itself, the
value of $G_Nv^2$, and Hypothesis~F$_{\mathrm g}$.

\smallskip\noindent\emph{No scalar.}  Under Hypotheses~F and F$_{\mathrm g}$(a) the configuration is $S^3$-valued, and the
electroweak sector is the gauged $\OO(4)/\OO(3)$ nonlinear sigma model, that is, the Standard Model with the
radial (Higgs) mode removed \cite{AB80}.  It contains no scalar particle; the tree-level amplitudes for
scattering of longitudinally polarised $W$ and $Z$ bosons grow with energy and violate partial-wave
unitarity at the TeV scale \cite{LQT77}; and its loop corrections, those to $\rho$ included, depend on the
cutoff \cite{AB80}.  A scalar of mass $125.13\pm0.11$\,GeV is observed, with coupling modifiers
$\kappa_V=1.023\pm0.026$ to vector-boson pairs and $\kappa_{2V}=1.02\pm0.23$ to the quartic vector-boson coupling
\cite{ex-PDG26}.  The electroweak sector under F and F$_{\mathrm g}$(a), as it stands, contains no scalar and is
contradicted by the observed one; a further input is required, for instance a radial mode: a varying
radius of the comparison sphere, that is, a configuration with values in the cone
$\R_{>0}\cdot S^3=\R^4\setminus\{0\}\subset V$ instead of in $S^3$, with a potential for the radius that nothing in the axiom
supplies.  What the data require is a scalar that restores unitarity with couplings close to those of the
Standard Model; the measured modifiers do not single out the radial mode among such scalars.  Section~\ref{sec:radial}
adds the radius as Hypothesis~F$_{\mathrm r}$; Section~\ref{sec:hopfbase} gives a curved alternative, Hypothesis~F$_{\mathrm H}$.

\smallskip\noindent\emph{F and G on this point.}  Hypothesis~G(ii) lists a bidoublet with values in $V_3$, that is,
a field with a radial mode; it is the cone-valued field that F lacks, for the third structure.  Adopting
G therefore replaces the $S^3$-bundle of F(i) by a bundle with fibre $V_3$, whose angular part is the
configuration of F, and requires a potential for the radius; G(iv) fixes only the vacuum value, a point of
$M_3$, and specifies no potential.  So the statement of the previous paragraph concerns F and F$_{\mathrm g}$(a)
alone and does not apply to G as it stands; whether the radial mode of G's bidoublet matches the
observed scalar depends on that potential.  The two hypotheses are consistent once F(i) is read, under
G, with the fibre $V_3$.


%=====================================================================================
\section{A variable radius: the cone-valued field}\label{sec:radial}
The electroweak sector under F and F$_{\mathrm g}$(a) has no scalar (Section~\ref{sec:centre}).  This section
asks what the comparison structure says about a variable radius of the comparison sphere, that is,
about the punctured space $V\setminus\{0\}=\HH^\times\cong\R_{>0}\times S^3$, and adds the radius as an optional hypothesis.

\subsection{The round-trip law on the cone}
On $\HH^\times$ the heap $\hp xyz=xy^{-1}z$ satisfies $(\Omega1)$: in polar form $q=rn$ the radial factor is $r_xr_z/r_y$
and the direction is $n_xn_y^{-1}n_z$.  A comparison invariant under all left and right multiplications and
symmetric is $\kappa(x,y)=\varphi(x^{-1}y)$ with $\varphi$ invariant under conjugation and inversion, hence a function of
\[
s=\log\bigl(\lvert y\rvert/\lvert x\rvert\bigr),\qquad \theta=\arccos\ip{n_x}{n_y}\in[0,\pi],
\]
even in $s$; we extend it evenly and $2\pi$-periodically in $\theta$.  Since $w^{-1}\hp xwx=(w^{-1}x)^2$, the round-trip
law $(\Omega3)$ reads, for $F=2\kappa-1$,
\begin{equation}\label{eq:conelaw}
F(2s,2\theta)=2F(s,\theta)^2-1,\qquad -1\le F\le1 .
\end{equation}

\begin{proposition}[comparisons on the cone]\label{prop:cone}
Let $F$ be a continuous solution of \eqref{eq:conelaw}.
(i) If $F$ does not depend on $s$, then $F\equiv1$, $F\equiv-\tfrac12$, or $F=\cos(m\theta)$ with $m\in\{1,2,\dots\}$.  The comparison
of directions $\kappa_{\mathrm{dir}}(x,y)=\tfrac12(1+\ip{n_x}{n_y})$ is $m=1$, the only one that restricts to the comparison of $M$ on
the unit sphere.
(ii) Comparisons that detect radii exist, for example $F=\cos(c\,s)$, $c>0$.  No continuous solution is
faithful with a unique zero partner: if $F(0,0)=1$, the radial profile $F(s,0)$ equals $1$ at some $s\ne0$, and a
nonconstant solution takes the value $-1$ at two values $\pm s$ of the ratio of radii.
(iii) If $F$ has a second-order expansion $F=1-\tfrac12(As^2+B\theta^2)+o(s^2+\theta^2)$ at the identity, then
$F=\cos\sqrt{As^2+B\theta^2}$, $\sqrt B$ is an integer, and $B>0$ forces $A=0$: the Hessian is degenerate, and the
round-trip law supplies an angular quadratic form and nothing radial.
\end{proposition}
\begin{proof}
(i) The three families solve \eqref{eq:conelaw}.  For a nonconstant continuous even $2\pi$-periodic solution put
$\alpha=\arccos F$; on an interval of length $2\pi/2^n$, every level $j\pi/2^n$ crossed by $\alpha$ gives alternate
values $\pm1$ of $F$ at $2^n\theta$, and uniform continuity bounds the number of such levels independently of
$n$, so $\alpha$ is Lipschitz.  At a point $\theta_0$ where $\alpha'(\theta_0)=b\ne0$, a subsequence of $(2^n\theta_0,2^n\alpha(\theta_0))$
converges modulo $2\pi$ to $(u,v_0)$, and $F(u+z)=\lim\cos(2^n\alpha(\theta_0+2^{-n}z))=\cos(v_0+bz)$; periodicity
makes $b$ an integer and evenness makes $v_0\in\{0,\pi\}$, and $v_0=\pi$ gives $F(0)=-1$, against \eqref{eq:conelaw}.
A constant solution is $1$ or $-\tfrac12$.
(ii) $\cos(cs)$ solves \eqref{eq:conelaw}.  If $F(0,0)=1$, near $s=0$ the radial profile is $\cos a(s)$ with $a(2s)=2a(s)$; so
$a$ is either identically zero or unbounded, and in both cases reaches a multiple of $2\pi$ at some $s\ne0$.  For
a nonconstant solution the values of $\arccos F$ on the half-cylinder $s>0$ fill an interval, whose dyadic
multiples contain $\pi$; by \eqref{eq:conelaw} and evenness $F=-1$ occurs at some $s\ne0$ and at $-s$.
(iii) Iterating \eqref{eq:conelaw}, $F(s,\theta)=\lim_n\cos\bigl(2^n\arccos F(2^{-n}s,2^{-n}\theta)\bigr)=\cos\sqrt{As^2+B\theta^2}$.
Periodicity at $s=0$ makes $\sqrt B$ an integer; if $B>0$, the values at $\theta=2\pi$ and $\theta=0$ agree for all
small $s$ only if $A=0$.
\end{proof}
Part (ii) agrees with the classification of Paper~I: faithful solutions live on compact groups
\cite[Thm.~3.11]{PaperI}, and the cone is not compact.

\begin{proposition}[linear structure on the cone]\label{prop:conelinear}
The centred comparison $c=\ip xy$ of $M$ has the unique bilinear extension $C(x,y)=\ip xy$ to $V$, and for
$w,x\ne0$, $m:=\hp xwx=x\bar wx/\lvert w\rvert^2$,
\[
m=\frac{2C(w,x)x-\lvert x\rvert^2w}{\lvert w\rvert^2},\qquad \lvert m\rvert=\frac{\lvert x\rvert^2}{\lvert w\rvert},\qquad
C(w,m)+\lvert x\rvert^2=\frac{2C(w,x)^2}{\lvert w\rvert^2};
\]
dividing by $\lvert x\rvert^2=\lvert w\rvert\lvert m\rvert$ gives \eqref{eq:conelaw} for $\kappa_{\mathrm{dir}}$.  No homogeneous extension of
nonzero degree satisfies the dimensionless round-trip law, and the heap has no continuous extension
to a middle argument $0$ ($\hp1t1=1/t$).
\end{proposition}
\begin{proof}
$x\bar w+w\bar x=2C(w,x)$ gives $x\bar wx=2C(w,x)x-w\lvert x\rvert^2$; the rest is substitution.  A simultaneous
extension of degree $D$ agreeing with $\kappa$ on the sphere has diagonal value $r^D$ at radius $r$, while the
round trip on the diagonal forces values in $\{1,\tfrac14\}$, so $D=0$.
\end{proof}

\subsection{Hypothesis F$_{\mathrm r}$ and the scalar doublet}
\smallskip\noindent\textbf{Hypothesis F$_{\mathrm r}$ (variable radius).}  In Hypothesis~F the configuration takes values
in $V\setminus\{0\}$ (equivalently, in $V$), with the unit transports acting by $q\mapsto\ell qr^{-1}$, $\ell,r\in S^3$; the clock
of F(ii) and the connection of F(iii) are as before, and the action is the canonical renormalisable one,
\begin{equation}\label{eq:radialaction}
\mathcal L=-\tfrac14W^a_{\mu\nu}W^{a\mu\nu}-\tfrac14B_{\mu\nu}B^{\mu\nu}+\tfrac12\lvert D_\mu q\rvert^2-U(\lvert q\rvert),\qquad
U(r)=U_0+\mu^2r^2+\lambda r^4 .
\end{equation}

\smallskip\noindent\emph{Dictionary.}  Write $q=z_0+jz_1$ with $z_0,z_1\in\C$ and represent $i\mapsto\mathrm{diag}(i,-i)$,
$j\mapsto\bigl(\begin{smallmatrix}0&-1\\1&0\end{smallmatrix}\bigr)$.  Then
\[
\Sigma(q)=\begin{pmatrix}z_0&-\bar z_1\\ z_1&\bar z_0\end{pmatrix}=\bigl(\sqrt2\,\widetilde H,\ \sqrt2\,H\bigr),\qquad
H=\tfrac1{\sqrt2}\binom{-\bar z_1}{\bar z_0},\quad \widetilde H=i\sigma_2\bar H,
\]
with $\det\Sigma=\lvert q\rvert^2$, $\Sigma^\dagger\Sigma=\lvert q\rvert^2\mathbf1$, $H^\dagger H=\lvert q\rvert^2/2$, and $\tfrac12\lvert Dq\rvert^2=(DH)^\dagger DH$.  The action
$q\mapsto\ell qr^{-1}$ is $\Sigma\mapsto\ell\Sigma r^\dagger$, which has image $\SO(4)$ and kernel $\{\pm(1,1)\}$; the vacuum $\Sigma=v\mathbf1$ has
the diagonal stabiliser, the custodial $\SO(3)$.  The clock acts on the two columns with opposite phases,
so $H$ is a doublet of $\SU(2)$ with hypercharge $\tfrac12$ (up to the orientation of the clock), and
$Q=T_{3L}+Y$ annihilates the vacuum $H=(0,v/\sqrt2)$.  This is the Standard Model scalar doublet in its
bidoublet form with custodial $\SO(4)$ \cite{SSVZ}.

\begin{proposition}[what F$_{\mathrm r}$ fixes and what it leaves free]\label{prop:radial}
Assume Hypotheses~F, F$_{\mathrm g}$(a) and F$_{\mathrm r}$.
(i) A minimum on a sphere $\lvert q\rvert=v>0$ exists iff $\lambda>0$ and $\mu^2=-2\lambda v^2<0$; then
$U(r)-U(v)=\lambda(r^2-v^2)^2$ and the radial mode $q=(v+h)n$ has mass $m_h^2=8\lambda v^2$.
(ii) In unitary gauge the quadratic and Higgs-vector terms are
$\bigl(1+\tfrac{2h}v+\tfrac{h^2}{v^2}\bigr)\bigl(m_W^2W^+W^-+\tfrac12m_Z^2Z^2\bigr)$ with $m_W=gv/2$, $m_Z=\sqrt{g^2+g'^2}\,v/2$, a massless
photon, $\rho=1$ at tree level, and coupling modifiers $\kappa_V=\kappa_{2V}=1$ at tree level.
(iii) In the gaugeless limit the longitudinal amplitude is $\mathcal M=\delta_{ab}\delta_{cd}\mathcal A(s)+\delta_{ac}\delta_{bd}\mathcal A(t)+\delta_{ad}\delta_{bc}\mathcal A(u)$ with
$\mathcal A(x)=x/v^2-x^2/(v^2(x-m_h^2))=-m_h^2x/(v^2(x-m_h^2))$: the radial exchange cancels the growth, while without
it the $J=0$ singlet wave $a_0^0=s/(16\pi v^2)$ exceeds $\tfrac12$ above $\sqrt s=\sqrt{8\pi}\,v\approx1.23$\,TeV.
(iv) Not fixed by anything in the comparison structure: $\mu^2$ and $\lambda$, hence $m_h$; the Yukawa couplings;
$g'/g$.  The requirement that the vacuum manifold be exactly a comparison sphere gives only $\lambda>0$ and
$\mu^2<0$; the comparison, its identities, the clock and the caps are the same for every $(\mu^2,\lambda)$ in this
range, and a potential read as a diagonal comparison obeying the round-trip law would be constant.
\end{proposition}
\begin{proof}
(i) Minimise $U_0+\mu^2z+\lambda z^2$ over $z=r^2\ge0$; the second derivative of $\lambda(r^2-v^2)^2$ at $r=v$ is $8\lambda v^2$.
(ii) $q=(v+h)n$ with $\lvert n\rvert=1$ gives $\tfrac12\lvert Dq\rvert^2=\tfrac12(\partial h)^2+\tfrac12(v+h)^2\lvert Dn\rvert^2$, and in unitary gauge $\tfrac12\lvert Dn\rvert^2$
is the quadratic form of the flat-gauging paragraph of Section~\ref{sec:centre}; $\rho=1$ is the equality of the
two eigenvalue ratios there.  (iii) With $n=(\sqrt{1-\pi^2/v^2},\pi^a/v)$ the relevant terms are
$\tfrac12(\partial\pi)^2+\tfrac1{8v^2}(\partial(\pi^2))^2+\tfrac hv(\partial\pi)^2$; the contact term gives $x/v^2$ and the $h$ exchange
$-x^2/(v^2(x-m_h^2))$; the equivalence theorem \cite{CLT74,LQT77} transfers this to longitudinal $W$, $Z$.  The
singlet amplitude without exchange is $(3s+t+u)/v^2=2s/v^2$, with $J=0$ projection $s/(16\pi v^2)$.
(iv) The comparison of directions and the caps depend on $q/\lvert q\rvert$ only (Proposition~\ref{prop:cone}(i) and the
next paragraph).
\end{proof}
\noindent\emph{Comparison with data.}  The quoted values $\kappa_V=1.023\pm0.026$ and $\kappa_{2V}=1.02\pm0.23$
\cite{ex-PDG26} contain the tree-level value $1$, and $m_h=125.13$\,GeV is fitted by
$\lambda=m_h^2/(8v^2)\approx0.032$ for $v=246.2$\,GeV; the electroweak comparisons of Example~\ref{ex:ew} are those of the
Standard Model.

\smallskip\noindent\emph{Relation to F and to G.}  The limit $\lambda\to\infty$ at fixed $v$, $g$, $g'$ freezes the radius:
eliminating $h=v\lvert Dn\rvert^2/m_h^2+\dots$ gives F's angular action $\tfrac12v^2\lvert Dn\rvert^2$ plus $v^2(\lvert Dn\rvert^2)^2/(2m_h^2)$, so F
with F$_{\mathrm g}$(a) is the leading term of the frozen-radius limit; at finite $m_h$ the scalar pole and its
couplings to $W$ and $Z$ distinguish the two.  The bidoublet of G(ii) is the same real four-component
field for the third structure; the remaining G sector (the field $\Xi$, mixing, thresholds) is not
specified, so matching G to a single light doublet with the couplings of (ii) needs further input.

\smallskip\noindent\emph{Events.}  Hypothesis~E factors through directions: the pullback of a cap under $q\mapsto q/\lvert q\rvert$
is $\{q:t\lvert q\rvert+\ip uq>0\}$, homogeneous of degree one, so every positive ray answers each dichotomy in the
same way, and the space of events and its order are unchanged.  The origin has no direction; it
answers no dichotomy, and a configuration through $q=0$ has no event reading there.  Radial
comparisons such as $\cos(cs)$ give threshold sets that are not pullbacks of caps; reading them as events
would be a new hypothesis.

\smallskip\noindent\emph{Status.}  With F$_{\mathrm r}$ the electroweak sector is the custodial Standard Model scalar sector:
compatible at tree level with the quoted measurements, with the scalar mass fitted, and with no
statement that distinguishes it from the Standard Model.  The price is one hypothesis and the free
parameters $\mu^2$, $\lambda$ (and the Yukawa couplings and $g'/g$, as before).  What remains open is a principle
that selects the potential.  An alternative in which the scalar arises from a larger comparison
geometry, with a coupling modifier below one, is Section~\ref{sec:hopfbase}.

%=====================================================================================
\section{A radial mode from two structures: the four-sphere $\HH P^1$, composite-Higgs geometry and its tests}\label{sec:hopfbase}
Can the state sphere of a pair supply the scalar missing from the fixed-length electroweak model of
Section~\ref{sec:centre}?  Choosing a quaternionic structure on the pair and taking its Hopf base $\HH P^1\cong S^4$ as the
scalar target gives the minimal composite-Higgs geometry $\SO(5)/\SO(4)$ of Agashe, Contino and Pomarol
\cite{n7:ACP}.  Its round metric gives one physical scalar in a broken vacuum and the tree relation
$\kappa_{2V}=2\kappa_V^2-1$ with $\kappa_V<1$, a one-parameter deviation from the Standard Model couplings that the data can
test.  The flat completion F$_{\mathrm r}$ of Section~\ref{sec:radial} is its limit $f\to\infty$.  The curved completion is a field
hypothesis; neither the pair nor Hypothesis~G forces it, and its scalar mass requires a matched potential.

\subsection{Scalar target and gauge couplings}

\smallskip\noindent\textbf{Hypothesis F$_{\mathrm H}$ (round four-sphere target).}  Take the pair space $V_2\cong\C^4$ of
Proposition~\ref{prop:pairsu3} and choose an antiunitary $J_{\mathrm H}$ with $J_{\mathrm H}^2=-1$.  Use the right quaternionic lines of this
structure as the fibres of a scalar bundle over the base $X$ of Hypothesis~F, so that locally the scalar target
is $\HH P^1\cong S^4=\SO(5)/\SO(4)$ with its round metric scaled to radius $f>0$ (the kinetic scale).  The electroweak group
gauges $\SU(2)_L$ and the right clock $\U(1)_{T_{3R}}$ inside the $\SO(4)$ fixing a pole of $S^4$, with canonical gauge kinetic
terms as in F$_{\mathrm g}$(a) and minimal coupling, truncated at two derivatives.  A gauge-invariant potential
selects a broken vacuum; choose the pole and the orientation of the polar angle $\chi$ so that $0<\chi_0\le\pi/2$ and the
scalar points towards increasing $\chi$.

\noindent F$_{\mathrm H}$ changes the fibre and kinetic form of F(i), (iv), keeps the clock reduction of F(ii)--(iii) and the
gauge dynamics of F$_{\mathrm g}$(a), and is a curved alternative to F$_{\mathrm r}$ for the scalar kinetic term.

\begin{theorem}[vector couplings]\label{thm:hbvector}
Under F$_{\mathrm H}$ there are three broken generators and one canonically normalised scalar $h=f(\chi-\chi_0)$.  With
$v=f\sin\chi_0$, $\xi=v^2/f^2$ and $g_Z=\sqrt{g^2+g'^2}$, the two-derivative tree theory gives
\begin{equation}\label{eq:hbvector}
m_W=\tfrac12gv,\quad m_Z=\tfrac12g_Zv,\quad m_\gamma=0,\quad \rho=1,\quad \kappa_V=\sqrt{1-\xi},\quad \kappa_{2V}=1-2\xi,
\end{equation}
where the modifiers multiply $2h/v$ and $h^2/v^2$ in the scalar-dependent vector mass operator normalised to its
vacuum value.  At the equator $(\kappa_V,\kappa_{2V})=(0,-1)$; at a pole the gauge symmetry is unbroken.
\end{theorem}
\begin{proof}
Write $\Phi=(\sin\chi\,U,\cos\chi)\in S^4\subset\R^4\oplus\R$, with $U$ a unit quaternion read as an $\SU(2)$ matrix by the dictionary
of Section~\ref{sec:radial}.  Differentiating $\lvert U\rvert^2=1$ removes the cross term, so
\begin{equation}\label{eq:hbmetric}
ds^2=f^2\bigl(d\chi^2+\sin^2\chi\,\lvert DU\rvert^2\bigr),\qquad \lvert DU\rvert^2=\tfrac12\tr\bigl((DU)^\dagger DU\bigr),\qquad
D_\mu U=\partial_\mu U-igW^a_\mu T_aU+ig'B_\mu UT_3,
\end{equation}
$T_a=\sigma_a/2$.  Left multiplication is transitive on each latitude, so $U=\mathbf1$ is a local unitary gauge in the
broken phase; the stabiliser is the diagonal clock, generating $Q=T_{3L}+Y$, $Y=T_{3R}$.  In this gauge the
scalar Lagrangian is $\tfrac12(\partial h)^2+\tfrac18f^2\sin^2(\chi_0+h/f)\bigl[g^2(W^1W^1+W^2W^2)+(gW^3-g'B)^2\bigr]$; with $W^\pm=(W^1\mp iW^2)/\sqrt2$,
$Z=(gW^3-g'B)/g_Z$, $A=(g'W^3+gB)/g_Z$ the mass term is $(m_W^2W^+W^-+\tfrac12m_Z^2Z^2)\sin^2(\chi_0+h/f)/\sin^2\chi_0$.  The first two
Taylor coefficients of the ratio are $2c_0/v$ and $(c_0^2-s_0^2)/v^2$ ($s_0=\sin\chi_0$, $c_0=\cos\chi_0\ge0$), which are \eqref{eq:hbvector}.
On the other hemisphere reversing the orientation of $h$ gives the same pair.
\end{proof}

\begin{corollary}[longitudinal scattering and the flat limit]\label{cor:hbscatter}
At leading tree order, fixed angle and $m_h^2,m_W^2,m_Z^2\ll s$ below the matching scale, the energy-growing
Goldstone amplitudes have coefficient $1/f^2$; in particular $\mathcal M(\pi^a\pi^b\to hh)|_{E^2}=\delta_{ab}\,s/f^2$, and
$1-\kappa_V^2=\kappa_V^2-\kappa_{2V}=\xi$.  At fixed $v$, $g$, $g'$ and finite $h$, $f\to\infty$ gives the flat kinetic theory of
Section~\ref{sec:radial} (full matching also needs a limit of the potential); at finite $f$, freezing $h=0$ gives the
angular theory of Section~\ref{sec:centre}.
\end{corollary}
\begin{proof}
In the gaugeless theory put $\Phi=\sqrt{1-\lvert\varphi\rvert^2/f^2}\,\Phi_0+\sum_{I\le4}\varphi^Ie_I/f$ ($e_4$ along increasing $\chi$).  The kinetic
term is $\tfrac12\lvert\partial\varphi\rvert^2+(\partial\lvert\varphi\rvert^2)^2/(8f^2)+O(f^{-4})$, with no cubic vertex; the three pairings give the
contact amplitude $(s\delta_{IJ}\delta_{KL}+t\delta_{IK}\delta_{JL}+u\delta_{IL}\delta_{JK})/f^2$ (contraction as in Proposition~\ref{prop:radial}(iii)), and
$I,J=a,b$, $K=L=4$ give the mixed amplitude; potential vertices carry no momenta.  The equivalence theorem
used in Proposition~\ref{prop:radial}(iii) transfers the leading power to longitudinal channels \cite{CLT74,n7:SILH}.  For
the limit, $f\sin(\chi_0+h/f)=v\cos(h/f)+fc_0\sin(h/f)\to v+h$, and in $r=f\sin\chi$ the metric is $dr^2/(1-r^2/f^2)+r^2\lvert DU\rvert^2$.  If
$m_h>0$ and the background is slowly varying, eliminating $h$ at tree level gives $h=vc_0\lvert DU\rvert^2/m_h^2$ and adds
$v^2c_0^2(\lvert DU\rvert^2)^2/(2m_h^2)$ to the angular Lagrangian.
\end{proof}

\subsection{Fermion embeddings and the potential}

\smallskip\noindent\textbf{Hypothesis F$_{\mathrm{Hf}}$ (fermion embeddings).}  For a massive Dirac fermion use the partial-compositeness
spurions of $\mathrm{Spin}(5)\times\U(1)_X$ with $Y=T_{3R}+X$ \cite{n7:ACP,n7:CPW,n7:CGG,n7:PW}; the chirality-changing mass comes from
the indicated leading invariant (for the $\mathbf{14}$ a linear combination of two), and the kinetic normalisations
do not depend on $h$ at this order.  ($X$ here is a fermion charge.)

\begin{proposition}[tree Yukawa modifiers]\label{prop:hbyuk}
Under F$_{\mathrm H}$ and F$_{\mathrm{Hf}}$, at a vacuum with $m_f\ne0$ and $\kappa_f:=v\,\partial_hm_f(0)/m_f(0)$: the spinor $\mathbf4$ gives $\kappa_f=c_0=\sqrt{1-\xi}$; the
vector $\mathbf5$ and antisymmetric $\mathbf{10}$ with one invariant give $\kappa_f=(1-2\xi)/\sqrt{1-\xi}$; the $\mathbf{14}$ with a left bidoublet and a
right singlet gives $m_f=\sin\chi\cos\chi\,(A_f+B_f\sin^2\chi)$ and
\begin{equation}\label{eq:hbyuk14}
\kappa_f=\frac{1-2\xi}{\sqrt{1-\xi}}+\frac{2B_f\xi\sqrt{1-\xi}}{A_f+B_f\xi}\qquad(A_f,B_f\ \text{real}).
\end{equation}
\end{proposition}
\begin{proof}
In unitary gauge $\Phi=s_\chi e_4+c_\chi e_5$.  (i) Clifford multiplication by $\Phi$ on the spinor $(2,1)\oplus(1,2)$ is the
quaternionic matrix with diagonal $c_\chi,-c_\chi$ and off-diagonal $s_\chi U,s_\chi U^\dagger$; the left-to-right contraction takes
its off-diagonal block, $m_f\propto s_\chi$.  (ii) For the vector, $(u^\dagger\Phi)(\Phi^Te_5)\propto s_\chi c_\chi$.  (iii) With $L=ue_5^T-e_5u^T$ and $R$
in the right triplet of the four-plane, $(L\Phi)^\dagger R\Phi=s_\chi c_\chi u^\dagger Re_4$.  (iv) With $L=ue_5^T+e_5u^T$ and $R=\mathrm{diag}(1,1,1,1,-4)$,
$(L\Phi)^\dagger R\Phi=-3s_\chi c_\chi\bar u_4$ and $(\Phi^TL\Phi)^*(\Phi^TR\Phi)=2s_\chi c_\chi(1-5c_\chi^2)\bar u_4$, whose combinations are $s_\chi c_\chi(A+Bs_\chi^2)$.
Then $\kappa_f=s_0\,d\log\lvert m_f\rvert/d\chi$.  Embeddings with $m_f\propto s_\chi c_\chi$ have $m_f=0$ at the equator.  Higgs-dependent
kinetic factors would add the derivative of $-\tfrac12\log(Z_LZ_R)$.
\end{proof}

\begin{remark}[one-loop matching]\label{rem:hbloops}
The Coleman--Weinberg prescription \cite{n7:CW,n7:ACP,n7:PW} applied to a transverse gauge kernel $K_T(p,\chi)$ and top
form factors $\Pi_L,\Pi_R,M$ gives $V_g(\chi)-V_g(0)=\tfrac32\int\frac{d^4p}{(2\pi)^4}\log\frac{\det K_T(p,\chi)}{\det K_T(p,0)}$ (a real Gaussian per transverse
component) and $V_t(\chi)-V_t(0)=-2N_c\int\frac{d^4p}{(2\pi)^4}\log\frac{p^2\Pi_L\Pi_R+\lvert M\rvert^2}{(\cdot)|_{\chi=0}}$ (a Grassmann determinant, two spin components,
each with determinant $-p^2\Pi_L\Pi_R-\lvert M\rvert^2$).  For canonical local kernels the gauge determinant is
$(p^2+m_W^2(\chi))^2(p^2+m_Z^2(\chi))p^2$ and the top factor $p^2+\lvert m_t(\chi)\rvert^2$; $\log(1+m^2/p^2)=m^2/p^2-m^4/(2p^4)+O(p^{-6})$ shows where
renormalisation enters.  The spherical metric fixes the background gauge masses; the momentum
dependence, the spurions and the finite matching terms are dynamical inputs.
\end{remark}

\begin{proposition}[vacuum and scalar mass]\label{prop:hbmass}
For a smooth invariant $V(\chi)$, a nondegenerate interior broken minimum has $V'(\chi_0)=0$ and $m_{\mathrm{curv}}^2=V''(\chi_0)/f^2>0$.  For
$V=V_0+\alpha x+\beta x^2$, $x=\sin^2\chi$, such a minimum with $0<\xi<1$ exists exactly when $\beta>0$ and $-2\beta<\alpha<0$, and then
\begin{equation}\label{eq:hbvev}
\xi=-\frac\alpha{2\beta},\qquad m_{\mathrm{curv}}^2=\frac{8\beta}{f^2}\xi(1-\xi),\qquad \kappa_\lambda:=\frac{v\,\partial_h^3V(0)}{3m_{\mathrm{curv}}^2}=\frac{1-2\xi}{\sqrt{1-\xi}} .
\end{equation}
\end{proposition}
\begin{proof}
$d/dh=f^{-1}d/d\chi$.  A quadratic in $x\in[0,1]$ has an interior nondegenerate minimiser exactly in the stated range.  At it,
$d^2V/d\chi^2=2\beta(x')^2$ and $d^3V/d\chi^3=6\beta x'x''$, with $x'=2s_0c_0$, $x''=2(c_0^2-s_0^2)$.  $\kappa_\lambda$ compares with the tree
Standard Model value $3m^2/v$ for the same mass and $v$.
\end{proof}

\begin{example}[the geometry does not determine the scalar mass]\label{ex:hbmasscounter}
For fixed $f$ and $0<\xi<1$, $V_C=C(\sin^2\chi-\xi)^2$ has the same broken latitudes and the same modifiers \eqref{eq:hbvector} for every
$C>0$, while $C$ and $4C$ give masses in the ratio $1:2$ by \eqref{eq:hbvev}.  Adding $D(x-\xi)^3$ with $\lvert D\rvert<C$ keeps the vacuum
and the mass ($C+D(x-\xi)>0$ on $[0,1]$, zero second derivative at $x=\xi$) and adds $6D(2s_0c_0/f)^3$ to $\partial_h^3V(0)$.  What fixes
the mass is an input: a matched curvature, $\beta=f^2m_{\mathrm{curv}}^2/[8\xi(1-\xi)]$, $\alpha=-2\beta\xi$, or a specified ultraviolet action
that determines the potential (Remark~\ref{rem:hbloops}); predicting the observed pole mass also requires the
momentum-dependent two-point function at its pole, with tadpoles and wave functions treated
consistently.  A mass near $125$\,GeV is therefore not a consequence of the geometry.
\end{example}

\subsection{Coupling and precision tests}

\begin{example}[bounds from the quoted Higgs couplings]\label{ex:hbdata}
Section~\ref{sec:centre} quotes $\kappa_V=1.023\pm0.026$, $\kappa_{2V}=1.02\pm0.23$, $m_h=125.13\pm0.11$\,GeV \cite{ex-PDG26}.  In the tree input
scheme without new contributions to muon decay, $G_F/\sqrt2=1/(2v^2)$, and $G_F=1.1663785(6)\times10^{-5}$\,GeV$^{-2}$ gives
$v=246.22$\,GeV \cite[Sec.~10.3.2 of the electroweak review]{ex-PDG26}.  \emph{Statistical treatment}: each quoted
modifier is treated separately as a Gaussian estimator with fixed standard deviation $\sigma$, its expectation
approximated by \eqref{eq:hbvector}; no covariance and no joint likelihood is used.  The one-sided lower limit
$L=\hat\kappa-z\sigma$, $z=1.64485$ (95\%), lies below the true expectation with probability $\Phi(z)$.  Both predictions
decrease in $\xi$; here $0<L_V<1$ and $-1<L_{2V}<1$, so $\xi\le1-L_V^2$ and $\xi\le(1-L_{2V})/2$, and $f=v/\sqrt\xi$:
\begin{center}\footnotesize
\begin{tabular}{@{}lcccc@{}}
\toprule
input & $\xi$ (one-sided, $z=1.64485$) & $f$ & $\xi$ (central, $z=1.95996$) & $f$\\
\midrule
$\kappa_V$ & $\le0.0391$ & $\ge1.245$\,TeV & $\le0.0551$ & $\ge1.049$\,TeV\\
$\kappa_{2V}$ & $\le0.179$ & $\ge0.582$\,TeV & $\le0.215$ & $\ge0.531$\,TeV\\
\bottomrule
\end{tabular}
\end{center}
These are approximate individual inversions of reported summaries; a joint constraint needs the
experimental likelihood with its correlations and profiling, and a fit with floating Yukawa and
self-coupling modifiers may differ.  The limit $\xi=0$ lies in both one-standard-deviation intervals.
\end{example}

\begin{remark}[electroweak precision]\label{rem:hbprecision}
The oblique parameters $T$ and $S$ \cite{n7:PT} measure custodial breaking and neutral vacuum polarisation
\cite[Sec.~10.6, Eq.~(10.84), Table~10.8 of the electroweak review]{ex-PDG26}.  The two-derivative action has $\rho=1$;
hypercharge and unequal fermion couplings are custodial-breaking spurions.  In the oblique
approximation the reduced Higgs coupling leaves a positive infrared logarithm in $S$ and a negative one in $T$
\cite{n7:BBRV,n7:PW}, and heavy vector resonances contribute to $S$.  Local operators are independent of the
metric: with $H_\Phi=(f\sin\chi/\sqrt2)\,U(0,1)^T$, $(H_\Phi^\dagger T_aH_\Phi)W^a_{\mu\nu}B^{\mu\nu}$ contains $-\tfrac14v^2W^3_{\mu\nu}B^{\mu\nu}$ and $\lvert H_\Phi^\dagger D_\mu H_\Phi\rvert^2$
contains $\tfrac1{16}v^4(gW^3-g'B)^2$ at the vacuum, shifting neutral mixing and the neutral mass independently.  A
precision fit needs heavy-state matching.
\end{remark}

\begin{table}[h]
\caption{The two completions of the missing scalar.}\label{tab:twoscalars}
\footnotesize
\begin{tabular}{@{}p{0.2\textwidth}p{0.36\textwidth}p{0.36\textwidth}@{}}
\toprule
& F$_{\mathrm r}$: flat radial mode (Sec.~\ref{sec:radial}) & F$_{\mathrm H}$: round $S^4$ target (this section)\\
\midrule
Target, kinetic metric & cone $\R^4\setminus\{0\}$, flat & $\HH P^1\cong\SO(5)/\SO(4)$, round of radius $f$\\
Vector couplings & $\kappa_V=\kappa_{2V}=1$ & $\kappa_V=\sqrt{1-\xi}$, $\kappa_{2V}=1-2\xi$, $\xi=v^2/f^2$\\
Longitudinal growth & cancels at $E^2$ & coefficient $1/f^2$ remains\\
Scalar mass & $m_h^2=8\lambda v^2$, $\lambda$ fitted & not fixed by the geometry (Ex.~\ref{ex:hbmasscounter}); matched input\\
Distinguishing statement & none (Standard Model couplings) & $\kappa_{2V}=2\kappa_V^2-1$, $\kappa_V\le1$; present bound $\xi\le0.039$, $f\ge1.24$\,TeV (Ex.~\ref{ex:hbdata})\\
Limit & --- & $f\to\infty$ at fixed $v$ gives F$_{\mathrm r}$\\
\bottomrule
\end{tabular}
\end{table}

\subsection{What the pair supplies}

\begin{proposition}[quaternionic quotient and metric]\label{prop:hbhopf}
A compatible quaternionic structure on $V_2$ makes $S(V_2)\cong S^7$ a principal $\mathrm{Sp}(1)$-bundle over $\HH P^1\cong S^4\cong\SO(5)/\SO(4)$.  The
quotient metric of the unit $S^7$ is the round metric of radius $\tfrac12$; F$_{\mathrm H}$ fixes the physical radius $f$.
Block-diagonal unit quaternions act on the base by $y\mapsto LyR^{-1}$, $z\mapsto z$, the custodial action of \eqref{eq:hbmetric}.
\end{proposition}
\begin{proof}
With right multiplication by $i$ the complex structure, $J_{\mathrm H}$ is right multiplication by $j$; for a unit $a$,
$\ip a{J_{\mathrm H}a}=0$, and the complement of the complex span of $a,J_{\mathrm H}a$ is $J_{\mathrm H}$-invariant, which gives a quaternionic
orthonormal basis of $\HH^2$.  For $q=(q_1,q_2)$, $\lvert q\rvert=1$, put $y=2q_1\bar q_2$, $z=\lvert q_1\rvert^2-\lvert q_2\rvert^2$; then $\lvert y\rvert^2+z^2=(\lvert q_1\rvert^2+\lvert q_2\rvert^2)^2=1$,
right multiplication by unit quaternions fixes $(y,z)$, and the sections $(a,\bar y/(2a))$, $a=\sqrt{(1+z)/2}$ ($z>-1$), and
$(y/(2b),b)$, $b=\sqrt{(1-z)/2}$ ($z<1$), trivialise the bundle.  $qq^\dagger=\tfrac12\mathbf1+\tfrac12\bigl(\begin{smallmatrix}z&y\\ \bar y&-z\end{smallmatrix}\bigr)$, and $\mathrm{Sp}(2)$ acts by
conjugation, preserving $z^2+\lvert y\rvert^2$, with kernel $\{\pm1\}$ and image $\SO(5)$ (dimension ten, connected).  At $q=(1,0)$ the
horizontal vectors $(0,w)$ map to $(2\bar w,0)$, so the submersion metric is $\tfrac14$ of the unit-base metric.
\end{proof}

\begin{proposition}[the quotient requires extra structure]\label{prop:hbobstruction}
The transitive action of $\SU(V_2)$ (Proposition~\ref{prop:pairsu3}) admits no equivariant smooth submersion of $S(V_2)$
onto a four-manifold; equivariance under the $\SU(3)$ fixing the closed-pair vector $\omega$ already excludes it
at $\omega$.  The local tensor action of $\SU(2)\times\SU(2)$ admits no invariant compatible quaternionic structure.
\end{proposition}
\begin{proof}
At $\omega$ the tangent space is $i\R\omega\oplus\omega^\perp$, $\omega^\perp\cong\C^3$, and the kernel of an equivariant submersion is a real
three-plane invariant under the stabiliser $\SU(3)$.  Its projection to $\C^3$ is invariant; $\SU(3)$ is transitive on the
unit sphere of $\C^3$, so the projection is $0$ or of dimension six, and neither fits a three-plane.  For the
tensor action, $C=(\epsilon\otimes\epsilon)K$ ($\epsilon=i\sigma_2$, $K$ complex conjugation) commutes with it ($\epsilon\bar L=L\epsilon$) and $C^2=\mathbf1$; an
invariant antiunitary $J_{\mathrm H}$ makes $J_{\mathrm H}C^{-1}$ a complex-linear element of the commutant, which is scalar, so $J_{\mathrm H}^2=\mathbf1$.
\end{proof}

\begin{example}[the quaternionic line is not a marginal readout]\label{ex:hbmixed}
The complex rays form $\C P^3=S^7/\U(1)$, and the quaternionic choice gives a further map to $\HH P^1$ with fibre
$\mathrm{Sp}(1)/\U(1)\cong S^2$; the two-qubit Hopf description is used for entanglement in \cite{MD}.  In the convention
$q_1=z_{00}+jz_{01}$, $q_2=z_{10}+jz_{11}$, the vector $q=(1,i)/\sqrt2$ and its right multiple by $(1+j)/\sqrt2$ have coefficient matrices
$\frac1{\sqrt2}\bigl(\begin{smallmatrix}1&0\\ i&0\end{smallmatrix}\bigr)$ and $\frac12\bigl(\begin{smallmatrix}1&1\\ i&-i\end{smallmatrix}\bigr)$; the first marginals $ZZ^\dagger$ are pure and $\mathbf1/2$, on one quaternionic line.  So the
line determines neither the marginal nor the purity: one fibre contains a product state and a
maximally entangled one.  The marginal is the tensor partial trace: for $(\lvert0\rangle\lvert e_0\rangle+\lvert1\rangle\lvert e_1\rangle)/\sqrt2$ with real
$\ip{e_0}{e_1}=\cos\theta$ it is $(\mathbf1+\cos\theta\,\sigma_1)/2$, eigenvalues $(1\pm\cos\theta)/2$, and for an orthogonal recording $\sum_aP_a\otimes\lvert a\rangle$ the
nonselective marginal is $\sum_aP_a\rho P_a$.  The orbit radius $f\sin\chi$ of the scalar is a kinetic quantity of a separate
field, distinct from the state cone of Theorem~\ref{thm:cone} and the cone-valued field of Section~\ref{sec:radial}.
\end{example}

\begin{remark}[the pair, Hypothesis G and the four-sphere]\label{rem:hborigin}
Paper~I's $S^7$ family concerns ternary operations after relaxing the single-system axioms
\cite[Thm.~5.18, Props.~5.16, 5.17]{PaperI}; the single-system solutions have dimensions one and three.  In
Hypothesis~G(ii) the real bidoublet is $V_3$ of a third structure, on which $\SU(V_2)$ acts trivially.  The tangent
space at a pole of $S^4$ is the same real bidoublet under $\SO(4)$ (the action on $y$ in Proposition~\ref{prop:hbhopf}), and
the extra coordinate $z$ is a custodial singlet.  Agreement of representations leaves the curved metric and
the potential undetermined; \emph{neither the pair nor G forces the curved scalar sector}: the stabiliser
obstruction (Proposition~\ref{prop:hbobstruction}) excludes a descent equivariant under the pair's colour
symmetry, and F$_{\mathrm r}$, or a real five-component field with $\lvert\Phi\rvert^2=1$ and Euclidean kinetic term, gives a scalar
without the pair.  A match inside G requires an electroweak identification compatible with the pair
colour action, the breaking field and the scalar mixing and thresholds of G.  A linear unitary $R$ with
$R^\dagger=-R$ has $R^2=-1$, but supplies no antiunitary quaternionic structure.
\end{remark}

\begin{lemma}[transport of a chirality-changing field]\label{lem:hbtransport}
For unitary edge transports $U_L,U_R$ on left and right fermion fibres and chirality-changing maps $F_i$ at
vertices, $U_LF_i-F_jU_R=(U_LF_iU_R^{-1}-F_j)U_R$; for $U_{L,R}=\mathbf1-\delta A_{L,R}+O(\delta^2)$ and $F_j=F_i+\delta\partial F+O(\delta^2)$ the coefficient of $-\delta$
is the covariant derivative $\partial F+A_LF-FA_R$.
\end{lemma}
\begin{proof}
Distribute and cancel $U_R^{-1}U_R$; expand to first order.  The commutator $[F,F]$ of one endomorphism vanishes and
cannot be a scalar potential; an internal curvature needs further operator and action data, as in the
finite Dirac operator and spectral action \cite{n7:Connes,n7:CC,n7:CC12}.
\end{proof}

\begin{conjecture}[microscopic matching of the pair scalar]\label{conj:hbmicro}
An edge-transport model with a dynamical selection of a compatible quaternionic structure has a
continuum phase with a gap to all other scalar modes and a round two-derivative $\HH P^1$ target in the limit
of vanishing gauge and Yukawa spurions.
\end{conjecture}
\noindent The evidence is compatibility only: the Hopf quotient, its custodial tangent representation and
Lemma~\ref{lem:hbtransport}.  No microscopic action or continuum construction is given.

\begin{remark}[tests and programme]\label{rem:hbtests}
\emph{Test.}  At leading order the round model predicts $0\le\kappa_V\le1$, $\kappa_{2V}=2\kappa_V^2-1$ and a common $\kappa_V$ for $W$ and $Z$; a
significant deviation after radiative, matching and derivative corrections refutes F$_{\mathrm H}$.  Present bound
(Example~\ref{ex:hbdata}): $\xi\le0.039$, $f\ge1.24$\,TeV.  Longitudinal scattering tests the coefficient $1/f^2$ against the
deficit in single- and double-Higgs couplings; for a specified embedding a Yukawa fit tests
Proposition~\ref{prop:hbyuk}; with the two-term potential the trilinear tests \eqref{eq:hbvev}.  These tests do not test the
pair origin: the constrained five-component field gives the same results without the pair, and a test of
the origin needs an observable fixed by microscopic matching.  \emph{Programme}: a microscopic action selecting
a reduction of the pair symmetry; the continuum spectrum and kinetic term; the potential, the chiral
Yukawa matrices and the thresholds.  A Dirac mass $y_fv/\sqrt2$ follows from a vacuum matrix element of the
chirality-changing field.  Mass is not a universal internal oscillation: free Dirac evolution at rest,
$e^{-i\gamma^0mt}$, leaves the chirality probabilities of a positive-energy eigenvector constant, and a beat at $2m$
needs coherence between opposite energies.
\end{remark}

\begin{remark}[what does not hold]\label{rem:hbnogo}
(i) The geometry does not determine the scalar mass (Example~\ref{ex:hbmasscounter}); a matched curvature or an
ultraviolet action fixes it.  (ii) The pair or Hypothesis~G does not force the curved scalar sector
(Proposition~\ref{prop:hbobstruction}, Remark~\ref{rem:hborigin}).  (iii) The quaternionic-line quotient is not a universal
marginal or purity readout (Example~\ref{ex:hbmixed}).  (iv) A self-commutator is no potential, and rest-frame
chirality of a positive-energy state does not oscillate (Lemma~\ref{lem:hbtransport}, Remark~\ref{rem:hbtests}).
\end{remark}

%=====================================================================================
\section{Self-contained particles: the Skyrme term, topological charge, Hopf solitons and spin one half}\label{sec:skyrme}
Can a comparison field carry a localised particle of finite size?  On a flat spatial slice the
two-derivative energy of Hypothesis~F decreases under contraction (Derrick); a square of the current
commutator supplies an energy that grows under contraction.  With it the degree bounds the energy
from below, a minimiser exists in degree one in a weak class, the clock quotient carries an integer
Hopf charge with its own bound, and the configuration space admits spinorial quantisation without
forcing it.  A dynamical connection, or the radial mode of Section~\ref{sec:radial}, removes the
topological protection.  This section extends Theorem~\ref{thm:gauge} and Section~\ref{sec:centre} under the following
optional input.

\subsection{The current and its area term}

\smallskip\noindent\textbf{Hypothesis F$_{\mathrm S}$ (Skyrme term).}  Take the quaternionic solution (Theorem~\ref{thm:main}) in
the Euclidean realisation of Proposition~\ref{prop:linear}.  On flat spacetime with signature $(-,+,+,+)$, in an
ungauged trivialisation with constant clock, represent the configuration by $U\colon X\to\SU(2)$, put
$L=U^{-1}dU$ and take
\begin{equation}\label{eq:skyrmeaction}
\mathcal L_{\mathrm S}=\frac{v^2}{4}\tr(L_\mu L^\mu)+\frac1{32e_{\mathrm S}^2}\tr\bigl([L_\mu,L_\nu][L^\mu,L^\nu]\bigr),
\end{equation}
with $v$ the scale of the action of F$_{\mathrm g}$(a) (the tangent norm is the unit-round one, so the first term is
$-\tfrac{v^2}2\lvert dU\rvert^2$), $e_{\mathrm S}>0$ an independent coupling, the fundamental trace, and indices summed over ordered
pairs.  Localised flat-space fields are required to tend to $U=1$ uniformly at spatial infinity; homotopy
is taken in the compact-open topology on based continuous maps, the integral formulas for maps smooth
on the compactification.

\begin{proposition}[the area square]\label{prop:areasq}
In the unit frame $e_a=-i\sigma_a$ write $L_\mu=\ell_\mu^ae_a$.  Then $e_ae_b=-\delta_{ab}\mathbf1+\epsilon_{abc}e_c$,
$[L_\mu,L_\nu]=2(\ell_\mu\times\ell_\nu)^ae_a$ and $-\tr([L_\mu,L_\nu]^2)=8\lvert\ell_\mu\wedge\ell_\nu\rvert^2$.  The Maurer--Cartan connection is flat
and $dL=-L\wedge L$ has components $-[L_\mu,L_\nu]$; up to scale, the second term of \eqref{eq:skyrmeaction} is the
parity-even metric square of this two-form with an $\SO(3)$-invariant fibre pairing.
\end{proposition}
\begin{proof}
The Pauli product gives the scalar and vector parts; $\tr(e_ae_b)=-2\delta_{ab}$.  Differentiating $U^{-1}dU$ gives the
structure equation.  An $\SO(3)$-invariant fibre pairing is a multiple of the Euclidean one, and the parity-even
contraction of two antisymmetric index pairs is unique up to scale.  This is the area square of Skyrme's
model \cite{n6:skyrme}.
\end{proof}
\begin{remark}[the area square is not the unique four-derivative invariant]\label{rem:notunique}
The invariant $(\tr L_\mu L^\mu)^2$ is a different four-derivative density: on a nonconstant static rank-one
current, for which the commutator vanishes, it is positive.  Uniqueness holds only for the metric square
of the commutator two-form.  The area of brackets also enters the Levi-Civita curvature of the round
sphere \cite[Prop.~5.4, Thm.~9.12]{PaperI}, which, like the curvature of a dynamical connection, is a
different object from $dL$.  On a symmetric space $G/H$ with lift $\omega+P$ ($\omega\in\mathfrak h$, $P\in\mathfrak m$) the isotropy
part of flatness reads $d\omega+\omega\wedge\omega=-(P\wedge P)_{\mathfrak h}$, a related area curvature.
\end{remark}

\subsection{Degree, scaling and minimisation}

\begin{lemma}[vacuum limit and degree]\label{lem:skdegree}
The uniform vacuum limit extends $U$ to a based map $S^3_{\mathrm{space}}=\R^3\cup\{\infty\}\to\SU(2)\cong S^3$, whose based homotopy
class is its degree $\mathcal Q\in\Z$.  With the target oriented by the left frame and $\eta=\theta^1\wedge\theta^2\wedge\theta^3$ the dual
coframe volume, for a smooth extension
\begin{equation}\label{eq:skdegree}
\mathcal Q=\frac1{2\pi^2}\int U^*\eta=-\frac1{24\pi^2}\int\tr(L\wedge L\wedge L),
\end{equation}
and $\mathcal Q$ is conserved under continuous based histories.
\end{lemma}
\begin{proof}
Stereographic projection compactifies; the degree theorem identifies $\pi_3(S^3)$ with $\Z$ \cite[\S\S2.2, 4.2]{n6:hatcher}.
Replace the normalised volume form by a form of integral one supported near a regular value; the
difference is exact and integrates to zero after pullback.  Here $\tr(e_ae_be_c)=-2\epsilon_{abc}$ and the volume of the
unit $S^3$ is $2\pi^2$; the oriented identity has degree $+1$ (integral $\tfrac12$ on either hemisphere).  The
pullback of $\eta$ to spacetime is closed, which gives current conservation without the field equation.
\end{proof}
\begin{remark}[finite energy does not force a vacuum limit]\label{rem:novacuumlimit}
Take a smooth bump $\chi$ supported in the unit ball with $\chi(0)=1$, balls at $x_k=(k,0,0)$ of radii $r_k=2^{-k-2}$, and
$U=\exp\bigl(e_3\sum_k\chi((x-x_k)/r_k)\bigr)$.  All currents commute, so the quartic energy is zero, and the quadratic
integral is $\lVert\nabla\chi\rVert_2^2\sum_kr_k<\infty$; yet $U=\exp e_3$ at the centres and $U=1$ between the balls.  The uniform
vacuum limit is therefore an input.
\end{remark}

\begin{theorem}[static bound and scaling]\label{thm:skbound}
Let $\mu_1,\mu_2,\mu_3\ge0$ be the singular values of $dU$ in orthonormal frames.  The static energy of
\eqref{eq:skyrmeaction} is
\[
E_{\mathrm S}=E_2+E_4,\qquad E_2=\frac{v^2}2\int\sum_i\mu_i^2\,d^3x,\qquad E_4=\frac1{2e_{\mathrm S}^2}\int\sum_{i<j}\mu_i^2\mu_j^2\,d^3x,
\]
and
\begin{equation}\label{eq:skbound}
E_{\mathrm S}\ge6\pi^2\frac v{e_{\mathrm S}}\lvert\mathcal Q\rvert ,
\end{equation}
with equality on flat space, for a smooth field, only for the vacuum.  For $U_\tau(x)=U(x/\tau)$, $\mathcal Q[U_\tau]=\mathcal Q[U]$ and
$E_{\mathrm S}[U_\tau]=\tau E_2+\tau^{-1}E_4$; a field of nonzero degree has a strict scale minimum at $\tau=\sqrt{E_4/E_2}$, every
stationary field of nonzero degree has $E_2=E_4$, and in the two-derivative model the infimum in each degree is
zero (every nonconstant field can lower its energy by contraction) \cite{n6:derrick}.
\end{theorem}
\begin{proof}
By Proposition~\ref{prop:areasq}, counting both orders of each spatial pair, the static energy has the stated form,
and its density equals $\tfrac12\sum_{\mathrm{cyclic}}(v\mu_i-e_{\mathrm S}^{-1}\mu_j\mu_k)^2+3(v/e_{\mathrm S})\mu_1\mu_2\mu_3$.  The product is the absolute
Jacobian, whose integral is at least $2\pi^2\lvert\mathcal Q\rvert$; this is Faddeev's bound \cite{n6:faddeev}.  If the squares vanish at a
point with a positive singular value, all singular values equal $ve_{\mathrm S}$ there, and on a neighbourhood $U$ is a local
isometry from a flat metric onto the positively curved sphere, which is impossible; otherwise $dU=0$ and the
boundary condition gives $U=1$.  Dilation multiplies derivatives by $\tau^{-1}$ and volume by $\tau^3$; nonzero degree
makes both energies positive (smooth based representatives exist in every degree), and the second
derivative $2E_4/\tau^3$ is positive.
\end{proof}

\begin{remark}[two normalisations of the same bound]\label{rem:sknorm}
In the convention with coefficients $f^2$ and $e^{-2}$ on the two sums of the static energy,
$E=\int(f^2\sum\mu_i^2+e^{-2}\sum_{i<j}\mu_i^2\mu_j^2)$, the same completed square gives $E\ge12\pi^2(f/e)\lvert\mathcal Q\rvert$.  The two are the same
bound: $f=v/\sqrt2$ and $e=\sqrt2\,e_{\mathrm S}$, so $12\pi^2f/e=6\pi^2v/e_{\mathrm S}$.  With coefficients $F_\pi^2/16$ and $1/(32e_{\mathrm S}^2)$ on the two
traces of \eqref{eq:skyrmeaction} ($F_\pi=2v$) the bound reads $3\pi^2F_\pi/e_{\mathrm S}$.  The coefficient $12\pi^2v/e_{\mathrm S}$ in the convention of
this paper would exceed the energy of the trial field below and is wrong.  The energy unit is $v/e_{\mathrm S}$ and
the length unit $(ve_{\mathrm S})^{-1}$.
\end{remark}

\begin{example}[a stereographic trial field]\label{ex:sktrial}
For $r=\lvert x\rvert$ and $a>0$ put $U_a(x)=\frac{r^2-a^2}{r^2+a^2}\mathbf1-\frac{2a}{r^2+a^2}\sum_ix_ie_i$.  It has $\mathcal Q=+1$ and all singular values equal to
$2a/(r^2+a^2)$; with $r=a\tan t$, $4\pi\int_0^\infty r^2\mu^p\,dr=4\pi^2a,\ 2\pi^2,\ 2\pi^2/a$ for $p=2,3,4$.  So $E_2=6\pi^2v^2a$,
$E_4=3\pi^2/(e_{\mathrm S}^2a)$, and the optimal trial energy is $6\sqrt2\,\pi^2v/e_{\mathrm S}\approx1.414$ times the bound, at $a=1/(\sqrt2\,ve_{\mathrm S})$.
\end{example}

\begin{remark}[existence of minimisers; cited]\label{rem:esteban}
For the Skyrme functional on $\R^3$ use maps in $W^{1,2}_{\mathrm{loc}}(\R^3,S^3)$ with $\nabla U$ and $\bigwedge^2\nabla U$ in $L^2$.  For such maps the
weak degree is an integer \cite{n6:esteban-muller}.  Esteban's variational treatment
\cite{n6:esteban1986}, completed in \cite{n6:esteban2004}, gives a minimiser in degree $q$ when the infimum $I_q$ is
strictly below $\sum_jI_{q_j}$ for every split of $q$ into at least two nonzero integers; as cited in the later
literature, this yields minimisers in degrees $\pm1$.  The binding inequality in degree $\pm1$ is checked here:
a nontrivial split has total absolute charge at least $3$, hence energy at least $18\pi^2v/e_{\mathrm S}$ by \eqref{eq:skbound}
(which holds for almost-everywhere derivatives), while Example~\ref{ex:sktrial} gives $6\sqrt2\pi^2v/e_{\mathrm S}$.  Smoothness,
uniqueness, the perturbation spectrum and the time evolution of the minimiser are further questions.
\end{remark}

\subsection{The clock quotient and Hopf charge}

\begin{proposition}[quotient and energy]\label{prop:skhopf}
With the clock of Proposition~\ref{prop:clock} along $e_3$, $h(U)=Ue_3U^{-1}$ realises $\SU(2)/\U(1)\cong S^2$ and induces
isomorphisms $\pi_3(S^3)\to\pi_3(S^2)$ and $\pi_4(S^3)\to\pi_4(S^2)$.  For $n=h\circ U$ and $\mathcal F_n=n^*\omega$ ($\omega$ the area form of the unit
$S^2$), the transverse part of the static energy is
\[
E_{\mathrm H}[n]=\int\bigl(a_2\lvert dn\rvert^2+a_4\lvert\mathcal F_n\rvert^2\bigr)d^3x,\qquad a_2=\frac{v^2}8,\quad a_4=\frac1{32e_{\mathrm S}^2},
\]
with the two-form norm summed over unordered pairs, and $E_{\mathrm H}[h\circ U]\le E_{\mathrm S}[U]$.  Taking $n$ as an autonomous field
with energy $E_{\mathrm H}$ is a separate dynamical choice, of Faddeev--Skyrme type \cite{n6:faddeev-niemi}.
\end{proposition}
\begin{proof}
The stabiliser of $e_3$ is the clock circle (the Hopf bundle \cite{n6:hopf-original}); in the long exact sequence
$\pi_k(S^1)=0$ for $k\ge2$.  $dn=\Ad_U[L,e_3]$ gives $\lvert dn\rvert^2=4\sum_i((\ell^1_i)^2+(\ell^2_i)^2)$ and
$(\mathcal F_n)_{ij}=4(\ell^1_i\ell^2_j-\ell^2_i\ell^1_j)$; substituting into $E_{\mathrm S}$ keeps the transverse components of $L$ and the $e_3$
component of each cross product, and the omitted squares are nonnegative.  For a nonconstant compactly
supported $\phi$, the lift $U=\exp(e_3\phi)$ has $dn=0$ and $E_2>0$, while the constant lift of the same $n$ has zero
energy: the full energy depends on the lift.
\end{proof}

\begin{definition}[Hopf charge]\label{def:hopfcharge}
With $d b=\mathcal F_n$ on the compactification, $\mathcal Q_{\mathrm H}=-\frac1{(4\pi)^2}\int b\wedge db$.
\end{definition}
\noindent Every smooth based $n$ has a smooth based lift (the pulled-back circle bundle over $S^3$ is trivial
\cite[\S1.2]{n6:bundles}).  In the coframe of Lemma~\ref{lem:skdegree}, $h^*\omega=4\theta^1\wedge\theta^2=-2d\theta^3$, so $b=-2U^*\theta^3$,
$b\wedge db=-8U^*\eta$ and $\mathcal Q_{\mathrm H}=\mathcal Q[U]$: Whitehead's integral \cite{n6:whitehead}, normalised so that the Hopf map has charge
$+1$, and the linking number of two preimages \cite{n6:hopf-original}.  A degree-one lift projects to Hopf
charge one, so topology alone supplies no second, independent integer.

\begin{theorem}[Vakulenko--Kapitanski\u{\i} bound]\label{thm:vk}
Let $C_6$ be a constant with $\lVert b\rVert_6\le C_6\lVert\nabla b\rVert_2$ on $\R^3$.  For smooth based fields,
\[
E_{\mathrm H}[n]\ge2\sqrt{2a_2a_4}\,(16\pi^2/C_6)^{3/4}\lvert\mathcal Q_{\mathrm H}\rvert^{3/4},
\]
and for $\mathcal Q_{\mathrm H}\ne0$ there is a strict minimum along dilations.
\end{theorem}
\begin{proof}
In the Coulomb gauge $d^*b=0$, $\lVert\nabla b\rVert_2=\lVert\mathcal F_n\rVert_2$, and the integral of $b\wedge\mathcal F_n$ is unchanged (the removed gradient is
orthogonal to $*\mathcal F_n$).  With $X=\int\lvert dn\rvert^2$, $Y=\int\lvert\mathcal F_n\rvert^2$ and $\lvert\mathcal F_n\rvert\le\lvert dn\rvert^2/2$, H\"older and interpolation give
$16\pi^2\lvert\mathcal Q_{\mathrm H}\rvert\le\lVert b\rVert_6\lVert\mathcal F_n\rVert_{6/5}\le C_6Y^{1/2}\lVert\mathcal F_n\rVert_1^{2/3}\lVert\mathcal F_n\rVert_2^{1/3}\le C_6(XY/2)^{2/3}$, and $a_2X+a_4Y\ge2\sqrt{a_2a_4XY}$
\cite{n6:vk-original}.  No optimality of the constant is claimed.
\end{proof}
\begin{remark}[Hopf minimisers; cited]\label{rem:linyang}
Lin and Yang proved the existence of minimisers in an infinite set of Hopf sectors \cite{n6:lin-yang} and
attainment at unit Hopf charge \cite{n6:lin-yang2}, in their finite-energy variational setting.
\end{remark}

\subsection{Rotation loops and spinorial quantisation}

\begin{proposition}[based configuration spaces]\label{prop:skconfig}
Let $\mathcal C_{\mathcal Q}$ be the degree-$\mathcal Q$ component of $\mathrm{Map}_*(S^3,\SU(2))$.  Then $\pi_1(\mathcal C_{\mathcal Q})\cong\pi_4(S^3)\cong\Z/2$, and a
$2\pi$ spatial rotation represents $\mathcal Q\bmod2$.  The clock projection induces an isomorphism of the fundamental
groups of the corresponding based mapping components and sends rotation class to rotation class.
\end{proposition}
\begin{proof}
Multiplication by the inverse of a fixed degree-$\mathcal Q$ map identifies $\mathcal C_{\mathcal Q}$ with the identity component;
adjunction identifies its loops with maps $S^4\to S^3$, and $\pi_4(S^3)=\Z/2$ \cite[Ch.~I]{n6:toda}.  The rotation
statement is Giulini's \cite{n6:giulini}.  The map of based mapping spaces induced by the Hopf bundle is a
fibration with fibre $\mathrm{Map}_*(S^3,S^1)$, which is contractible (each map lifts to a based real function, which
can be scaled to zero); the long exact sequence gives the isomorphism, and postcomposition commutes with
rotations.
\end{proof}
\begin{remark}[Finkelstein--Rubinstein quantisation]\label{rem:fr}
A wave function on the universal cover may transform under either character of the deck group $\Z/2$.
For the nontrivial character the $2\pi$ rotation in odd degree changes the sign; the same loop represents
the exchange of two separated unit textures \cite{n6:fr-original,n6:krusch}.  In the rotational
collective-coordinate model of the equivariant field of Example~\ref{ex:sktrial} the nontrivial character
requires $\psi(-g)=-\psi(g)$ on $\SU(2)$, which permits spin $\tfrac12$ and the other half-integer spins.  The
trivial character is equally compatible with the same classical $\SU(2)$ action: \emph{the bosonic action does
not force fermionic quantisation}.  What would force it is an additional term, as in the chiral flavour
$\SU(3)$ model, where the Wess--Zumino--Witten term at level $N_c$ gives the phase $(-1)^{N_c}$ to the rotation of the
embedded unit texture \cite{n6:witten-current,n6:witten-baryons}; the colour symmetry of
Proposition~\ref{prop:pairsu3} and Theorem~\ref{thm:sm} supplies no such flavour term.  The spinor sign of
Example~\ref{ex:neutron} concerns the state space of one system; its double cover and the configuration
cover here have different domains.
\end{remark}

\subsection{The centre field with the Skyrme term}

\begin{proposition}[the compact centre field]\label{prop:skcentre}
The identity from a round spatial sphere of radius $R$ to the unit target has $E_2=3\pi^2v^2R$ and
$E_4=3\pi^2/(e_{\mathrm S}^2R)$.  At $R=(ve_{\mathrm S})^{-1}$ it attains the bound \eqref{eq:skbound} and is the energy minimum in degree one on
that sphere.  Along the conformal modes (the eigenvalue $-1$ of Proposition~\ref{prop:stability}) its second variation
is positive iff $R<\sqrt2\,(ve_{\mathrm S})^{-1}$.  On round slices of a time-dependent metric, $\rho_2=3v^2/(2R^2)$, $p_2=-\rho_2/3$,
$\rho_4=3/(2e_{\mathrm S}^2R^4)$, $p_4=\rho_4/3$, and the null contraction of the stress is positive: the obstruction of
Proposition~\ref{prop:nullobstruction} persists.
\end{proposition}
\begin{proof}
The singular values are $1/R$ and the volume is $2\pi^2R^3$; at the given radius each completed square vanishes,
with positive Jacobian.  For a conformal dilation of the unit sphere with factor
$\Omega_t=(\cosh t-\sinh t\cos\chi)^{-1}$ the energy is $\tfrac32\bigl(\tilde R\int\Omega_t^2+\tilde R^{-1}\int\Omega_t^4\bigr)$ in the units of
Remark~\ref{rem:sknorm} ($\tilde R=ve_{\mathrm S}R$), and its second derivative at $t=0$ changes sign at $\tilde R=\sqrt2$ (computed in
\texttt{check\_physics.py}).  The pressures follow from $p=-dE/dV$ at fixed field and $\rho=E/V$; the null contraction
is $(\rho_2+p_2+\rho_4+p_4)(k^0)^2>0$.
\end{proof}
\noindent This changes the conclusion of Proposition~\ref{prop:stability} as follows.  Without the Skyrme term the
identity has index four on every static round slice, and its conformal modes grow; with it, the conformal
modes are stabilised on slices with $R<\sqrt2/(ve_{\mathrm S})$, and the identity attains the topological bound at
$R=1/(ve_{\mathrm S})$.  Manton showed that the identity is stable exactly up to the radius $\sqrt2/(ve_{\mathrm S})$ \cite{Manton87};
we have checked the conformal modes only.  On the expanding de~Sitter slices $R$ grows without bound, so
the conformal instability returns at late times on the fixed background; the coupled metric--field
problem with the quartic stress, which scales like radiation, is open.

\subsection{The dynamical connection and the radial mode}

\smallskip\noindent\textbf{Hypothesis F$_{\mathrm{Sg}}$ (covariant area square).}  In addition to F$_{\mathrm g}$(a), put $DU=dU+AU-UC$, with $A$
the left $\su(2)$ connection and $C$ the right clock connection, and use \eqref{eq:skyrmeaction} with $L_D=U^{-1}DU$ in place of $L$,
together with the positive Yang--Mills energy (the left potential in the convention $d+A$ is the negative of
that of Section~\ref{sec:centre}).  Localised configurations use the trivial based spatial bundle, fields and
connections equal to the vacuum and to zero near infinity.

\begin{proposition}[covariance and screening]\label{prop:skscreen}
Under $U\mapsto aUt^{-1}$, $L_D\mapsto tL_Dt^{-1}$, so both densities are gauge invariant; the flat vacuum of F$_{\mathrm g}$(b) keeps zero
stress and the mass form of Theorem~\ref{thm:gauge} still applies.  In every degree there is a configuration of
zero energy, $A=-dU\,U^{-1}$, $C=0$; under large based gauge transformations the degree and the projected Hopf
charge change.
\end{proposition}
\begin{proof}
$A\mapsto aAa^{-1}-da\,a^{-1}$ and $C\mapsto tCt^{-1}-dt\,t^{-1}$ give the covariance.  The curvature identity
$dL_D+C\wedge L_D+L_D\wedge C+L_D\wedge L_D=U^{-1}F_AU-F_C$ shows that the commutator square is of fourth order at $DU=0$.  For
$A=-dU\,U^{-1}$, $DU=0$ and $F_A=0$, so all three energies vanish; $a=U^{-1}$ sends the pair to $(1,0)$, and
$\mathcal Q[aUt^{-1}]=\mathcal Q[U]+\mathcal Q[a]$.
\end{proof}
\begin{remark}[degree and Chern--Simons number]\label{rem:skcs}
With $\mathrm{CS}(A)=-\frac1{8\pi^2}\int\tr(A\wedge dA+\tfrac23A\wedge A\wedge A)$, $\mathrm{CS}(A^a)=\mathrm{CS}(A)+\mathcal Q[a]$, so $\mathcal Q[U]-\mathrm{CS}(A)$ is gauge invariant
(this is the remark on large gauge transformations in Section~\ref{sec:centre}, now with the Skyrme term).  For
$A=tP$, $P=-dU\,U^{-1}$, $\mathrm{CS}(tP)=(3t^2-2t^3)\mathcal Q[U]$, so the difference takes non-integer values along finite-energy
paths.
\end{remark}
\begin{example}[a barrier in a trial family]\label{ex:skbarrier}
With $C=0$, the Yang--Mills energy $-\int\tr(F_{ij}F_{ij})/(2g^2)$ (ordered indices), $U_\tau(x)=U(x/\tau)$ of degree one and
$A=-t\,dU_\tau U_\tau^{-1}$, $0\le t\le1$, the energy is $\tau(1-t)^2E_2+\tau^{-1}E_4[(1-t)^4+\kappa t^2(1-t)^2]$, $\kappa=16e_{\mathrm S}^2/g^2$.  At the
optimal scale it is $2\sqrt{E_2E_4}\,P_0(t)^{1/2}$ with $P_0=(1-t)^4[(1-t)^2+\kappa t^2]$ and
$P_0'=2(1-t)^3[-3+(6+\kappa)t-3(1+\kappa)t^2]$.  For $\kappa>24$ the bracket has two roots in $(0,1)$ (discriminant
$\kappa(\kappa-24)$): a local minimum and a barrier before the screened vacuum at $t=1$.
\end{example}
\begin{conjecture}[weak-gauging metastability]\label{conj:metastable}
At fixed $v,e_{\mathrm S}>0$, with the clock connection decoupled, the flat-space gauged model has a smooth non-vacuum
local energy minimiser for sufficiently small $g>0$ (modulo based gauge and Euclidean symmetries).
\end{conjecture}
\noindent The evidence is Example~\ref{ex:skbarrier}; a proof must control deformations outside the trial family.

\smallskip\noindent\textbf{Hypothesis F$_{\mathrm{Sr}}$ (area square with the radial mode).}  With F$_{\mathrm r}$ and the potential $\lambda(r^2-v^2)^2$, extend
the area square to $q\in V$ by the polynomial density $-\frac1{4e_{\mathrm S}^2v^4}\ip{D_\mu q\wedge D_\nu q}{D^\mu q\wedge D^\nu q}$ (Euclidean bivector pairing).

\begin{proposition}[matching and radial unwinding]\label{prop:skradial}
On $\lvert q\rvert=v$ the polynomial area term equals the covariant Skyrme density; around the vacuum $m_h^2=8\lambda v^2$ and
$\kappa_V=\kappa_{2V}=1$ for the nonderivative vertices, with additional momentum-dependent $(\partial h)^2V^2$ and $V^4$ terms of order
$1/(e_{\mathrm S}^2v^4)$.  The homotopy $q_s=v[(1-s)U+s\mathbf1]$ has finite energy, passes through $q=0$ at $s=\tfrac12$ when $\mathcal Q[U]\ne0$, and
changes the winding of $q/\lvert q\rvert$ and its Hopf projection.
\end{proposition}
\begin{proof}
At $r=v$ left translation identifies $D_\mu U$ with $(L_D)_\mu$ isometrically, and Proposition~\ref{prop:areasq} gives the
coefficient $-\frac1{4e_{\mathrm S}^2}\lvert\ell_\mu\times\ell_\nu\rvert^2$.  Proposition~\ref{prop:radial} gives the rest of the vacuum expansion.  A map of
nonzero degree is onto, so $U$ takes the value $-\mathbf1$, where $q_{1/2}=0$.
\end{proof}
\noindent In the gaugeless Goldstone expansion the area operator adds four-Goldstone contacts with four powers
of momenta and coefficient of order $1/(e_{\mathrm S}^2v^4)$ (for two Goldstone fields, the square of the antisymmetrised
product of their gradients); the two-derivative cancellation of Proposition~\ref{prop:radial}(iii) is unchanged,
and the four-derivative contacts remain in the effective theory.  Matching under G requires the
potential, mixing, the higher-operator coefficients and the thresholds (programme).  The fixed-length
model has no radial particle; the statement of Section~\ref{sec:centre} on the observed scalar is unchanged.

\begin{remark}[what does not hold]\label{rem:sknogo}
(i) Finite energy does not force a vacuum limit (Remark~\ref{rem:novacuumlimit}).  (ii) The area square is not the
unique four-derivative invariant (Remark~\ref{rem:notunique}).  (iii) Winding does not ensure stability once the
connection or the radius is dynamical: the screened pairs of Proposition~\ref{prop:skscreen} and the homotopy of
Proposition~\ref{prop:skradial} are explicit counterexamples; nor does every homogeneous target carry an integer
charge, e.g.\ $\pi_3(S^1)=0$.  (iv) The bosonic action does not force spin $\tfrac12$ (Remark~\ref{rem:fr}).  (v) A compact
slice with a parallel frame does not make every configuration a unit-charge particle: a constant field has
charge zero, and a hemisphere carries charge density $\tfrac12$ although the total charge is an integer.
\end{remark}

\smallskip\noindent\textbf{Hypothesis P (particle reading).}  A localised, quantised degree texture is a candidate for a
baryon-like excitation; a configuration of the autonomous clock-quotient dynamics is a candidate for a
lepton-like excitation if an independent physical rule distinguishes it from a projected degree texture
(the projection equates the two labels).

\begin{remark}[programme]\label{rem:skprogramme}
A particle identification needs an isolated solution and its fluctuation spectrum, a physical choice of
quantisation character, couplings to the multiplets of Theorem~\ref{thm:sm}, and a mechanism for the observed
charges, chirality, masses and lifetimes; replacing the fermion content of Hypothesis~G would require
those results.  Exact stability in the gauged or radial model requires an obstruction to the paths above;
metastability requires a barrier in the full problem and a decay rate.  The gravitational programme is to
solve the coupled equations with the quartic stress and check the event and maximality conditions of
Section~\ref{sec:centre}.
\end{remark}

\subsection{Intrinsic sphere geometry}

\begin{proposition}[frame, wavefronts and first modes]\label{prop:sksphere}
The round sphere $S^3(R)$ has a divergence-free left frame and $\dim\bigwedge^kT^*_xS^3=1,3,3,1$.  About a centre $p$, the
wavefront at distance $r\in[0,\pi R]$ has area $4\pi R^2\sin^2(r/R)$ and encloses the volume $\pi R^3(2r/R-\sin(2r/R))$; at
$r=\pi R/2$ the area is maximal and the volume is half the total, and at $r=\pi R$ all radial geodesics meet at $-p$.
The two closed hemispheres are balls glued along their boundary $S^2$.  The first nonconstant eigenspace of
$\Delta_{S^3(R)}$ has dimension four and eigenvalue $-3/R^2$.
\end{proposition}
\begin{proof}
The frame vectors generate right translations, which preserve the metric and volume.  The wavefront
$\cos(r/R)p+R\sin(r/R)\xi$ has angular metric factor $R^2\sin^2(r/R)$, and integration gives the volume.  Homogeneous
harmonic polynomials of degree $k$ restrict to eigenfunctions with eigenvalue $-k(k+2)/R^2$, and their
restrictions are dense \cite[Ch.~2]{n6:axler}.
\end{proof}
\begin{remark}[field, wave and observation]\label{rem:skintrinsic}
A parallel frame, a Laplace eigenmode and a map of nonzero degree are a frame, a function and a
configuration; the first two exist even for a constant configuration of charge zero.  Identifying a wave, a
field and a particle as descriptions of one system requires a model linking their variables, and a rule
replacing a finite volume cell by a point at another scale requires a coarse-graining map with control of
the effective dynamics.  The comparison vacuum of F$_{\mathrm S}$ is a boundary condition, movable by the homogeneous
target action (Proposition~\ref{prop:aut}; its point stabiliser is the $\SO(3)$ of Theorem~\ref{thm:gauge}(iii)).  For unit
points $A$, $B$ at angle $\chi$ the affine midpoint has squared norm $\tfrac12(1+\cos\chi)$ and $2B-A$ has squared norm
$5-4\cos\chi$, so affine continuation leaves the sphere, and anchored addition $A+B-C$ changes by $-\delta$ when the
anchor moves by $\delta$; the heap continuation $BA^{-1}B$ stays on the sphere, and at antipodes returns $A$.  A boundaryless round sphere need not exhaust a larger space or state set (the round
$S^3\subset\R^4$ has an ambient complement); admitting all non-constant comparison profiles and excluding constant
ones would be additional closure rules.  These items are a programme.
\end{remark}

%=====================================================================================
\section{Three structures and the Standard Model group}\label{sec:sm}
This section uses C0--C3 for a closed pair and, from
Hypothesis~G on, a labelled hypothesis.

\begin{proposition}[parity and chirality]\label{prop:parity}
(i) Transport reversal $\iota$ ($z\mapsto\bar z$ in the model) lies in $\Aut^\pm(M)$ and satisfies
$\iota\cL\iota=\cR$.
(ii) $V$ with $z\mapsto\ell z r^{-1}$ is the real bidoublet $(2,2)$ of $\cL\times\cR$; its unit sphere is
$M$, the stabiliser of a point is the diagonal $\SU(2)$ (acting on $V$ through $\SO(3)$), and, with a
clock $T\subset\cR$, the stabiliser of a point in $\cL T$ is the circle generated by $Q=T_{3L}+T_{3R}$.
(iii) $V$ as a module for $\cL$ alone, $z\mapsto\ell z$, with the complex structure $J_T$ of a clock, is
the doublet $(2,1)$; $V$ as a module for $\cR$ alone, $z\mapsto zr^{-1}$, with the complex structure of a
one-parameter subgroup of $\cL$, is $(1,2)$; $\iota$ exchanges them, and
$V\otimes_\R\C\cong(2,1)\otimes_\C(1,2)$ as representations of $\cL\times\cR$.
(iv) The automorphisms commuting with $J_T$ form $\cL T\cong\U(2)$, and $\iota$ maps them to
$\cR\,\iota T\iota$.
\end{proposition}
\begin{proof}
(i) $\overline{x\bar yz}=\bar zy\bar x$, and $\overline{a\bar z}=z\bar a$.  (ii) Proposition~\ref{prop:aut} and
Theorem~\ref{thm:gauge}(iii).  (iii) Proposition~\ref{prop:clock}(ii) and its mirror; the complexification
of the vector representation of $\SO(4)=(\SU(2)\times\SU(2))/\Z_2$ is the tensor product of the two
half-spin representations.  (iv) Theorem~\ref{thm:dyn}.
\end{proof}
In this structure, which transports act unitarily and on which side the clock lies are one choice, and
chirality is the statement of which transport acts on a $V$-valued field.

Let $V_2=(V,J_T)\otimes_\C(V,J_T)$ be the kernel space of a closed pair (one clock, by
C0, Corollary~\ref{post:C0}; C1(b) fixes the relative orientation) and $\omega$ the antisymmetric unit vector.

\begin{proposition}[colour from the closed pair]\label{prop:pairsu3}
(i) $[\omega]$ is the only ray invariant under joint transports, and $\omega^\perp=\mathrm{Sym}^2V\cong\C^3$.
(ii) The stabiliser of $\omega$ in $\SU(V_2)$ is $\SU(\omega^\perp)\cong\SU(3)$, acting on $\omega^\perp$ by
its defining representation; the stabiliser of the ray is $\mathrm S(\U(1)\times\U(3))\cong\U(3)$, with
Lie algebra $\su(3)\oplus\mathfrak u(1)_{B-L}$, where $B-L:=-P_\omega+\tfrac13P_{\omega^\perp}$; and
$S(V_2)=\SU(4)/\SU(3)\cong S^7$.  The stabiliser of $\omega$ in $\U(V_2)$ is $\U(\omega^\perp)\cong\U(3)$.
(iii) The joint transports act on $\omega^\perp$ through $\SO(3)\subset\SU(3)$.
(iv) Under Postulates~\ref{post:C2}, \ref{post:C3} every element of $\SU(V_2)$ is an operation.
\end{proposition}
\begin{proof}
(i) Theorem~\ref{thm:originfree}.  (ii) An element fixing $\omega$ preserves $\omega^\perp$ and has
determinant one there; the elements of $\su(4)$ commuting with $\su(\omega^\perp)$ are
$i\,\mathrm{diag}(\lambda,\mu,\mu,\mu)$ with $\lambda+3\mu=0$; $\SU(4)$ is transitive on $S^7$.  (iii) The
spin-one representation is of real type.  (iv) Theorem~\ref{thm:comp}(iii).
\end{proof}
With the pair's clock generator normalised to $\tfrac12$ on $V_2$, the charge $\tfrac12(B-L)+\tfrac12$ is
$\mathrm{diag}(0,\tfrac23,\tfrac23,\tfrac23)$ on $(\omega,\omega^\perp)$ and generates the centre of
$\U(\omega^\perp)$; with the clock reversed one obtains $\mathrm{diag}(-1,-\tfrac13,-\tfrac13,-\tfrac13)$.

\smallskip\noindent\textbf{Hypothesis G.}  (i) A field of composites has gauge group
$G:=\SU(V_2)\times\cL_3\times\cR_3$, where $\cL_3,\cR_3$ are the transports of a third elementary
structure $M_3$ with kernel space $V_3$.  (ii) Its scalar multiplets are the $(4,1,2)$ field of (iv) and the bidoublet $V_3$ of
Proposition~\ref{prop:parity}(ii), on which $\SU(V_2)$ acts trivially.  (iii) Its matter multiplets are
$S_L:=V_2\otimes_\C V_3^{(2,1)}$ and $S_R:=V_2\otimes_\C V_3^{(1,2)}$, of types $(4,2,1)$ and $(4,1,2)$,
with the one-sided modules of Proposition~\ref{prop:parity}(iii) for a clock $T\subset\cR_3$ and its
reversal $\iota T\iota\subset\cL_3$.  (iv) The breaking configuration is $\Xi:=\omega\otimes v\in S_R$ for a
point $v\in M_3$, the value of a scalar field of type $(4,1,2)$; the bidoublet takes the value of a
point of $M_3$ (used for $\U(1)_Q$ in Theorem~\ref{thm:sm}(i)).  (v) The kinetic form of the gauge
fields is the single trace form $-\tr_{S_L\oplus S_R}$ with equal couplings $g_4=g_L=g_R$ at a matching
scale.

\smallskip
That the pair carries $\SU(V_2)$ while the third structure carries its transports
$\cL_3\times\cR_3\cong\mathrm{Spin}(4)$ is part of clause (i); we have no derivation of this asymmetry.
Clause (iii) uses the one-sided modules: on $(V_3,J_T)$ with both transports acting, $T$ acts by
scalars, every vector carries the same nonzero $T_{3R}$, and the hypercharges below would be shifted.
Clause (iv) is the standard $(4,1,2)$ breaking of Pati--Salam models and is a choice beyond the
structural data $[\omega]$, $T$ and a point of $M_3$; the field carrying it is a scalar of type
$(4,1,2)$, in addition to the bidoublet of clause (ii).

\begin{theorem}\label{thm:sm}
Assume C0--C3 and Hypothesis~G.  The $r\in\cR_3$ with
$v\,r^{-1}\in\C v$ in $V_3^{(1,2)}$ form a clock $T_v$ (equal to $\bar vTv$ in the model, and to $T$ for
$v=1$); let $T_{3R}\in\mathrm{Lie}(T_v)$ be its generator, normalised with eigenvalues $\pm\frac12$ on
$V_3^{(1,2)}$ and $T_{3R}v=\frac12v$.  Put $Y:=T_{3R}+\tfrac12(B-L)$ and $Q:=T_{3L}+Y$.
(i) (Stabilisers of the structural data.)  The ray $[\omega]$ reduces $\SU(V_2)$ to $\U(3)$, locally
$\SU(3)\times\U(1)_{B-L}$; the elements of $\cR_3$ commuting with the clock $T_v$ form $T_v=\U(1)_{T_{3R}}$; a point of $M_3$ in the bidoublet
reduces $\cL_3\times\U(1)_{T_{3R}}$ to $\U(1)_Q$ with $Q=T_{3L}+T_{3R}$.
(ii) The stabiliser of $\Xi$ in $G$ is $H\cong\U(3)\times\SU(2)_L$, with Lie algebra
$\su(3)\oplus\su(2)_L\oplus\mathfrak u(1)_Y$.
(iii) The kernel of $G$ on $S_L\oplus S_R\oplus V_3$ is $\Z_2=\langle(-1,-1,-1)\rangle\subset H$; the group
acting faithfully is $G/\Z_2$, the Pati--Salam group, and $H/\Z_2\cong\mathrm S(\U(3)\times\U(2))\cong
(\SU(3)\times\SU(2)\times\U(1))/\Z_6$.
(iv) $Y$ takes the values $-\tfrac12$ on $\omega\otimes V_3^{(2,1)}$ and $\tfrac16$ on $\omega^\perp\otimes V_3^{(2,1)}$;
$0,-1$ on $\omega\otimes V_3^{(1,2)}$ and $\tfrac23,-\tfrac13$ on $\omega^\perp\otimes V_3^{(1,2)}$; $Q$ takes the
electric charges of one generation of quarks and leptons with a right-handed neutrino, leptons along
$\omega$ and quarks along $\mathrm{Sym}^2V$.  On the bidoublet $Y=T_{3R}$ takes the values $\pm\frac12$.
\end{theorem}
\begin{proof}
(i) Proposition~\ref{prop:pairsu3}(ii), Theorem~\ref{thm:wigner}(i) ($T$ is its own centraliser in
$\cR\cong\SU(2)$) and Theorem~\ref{thm:gauge}(iii).
(ii) $\cL_3$ acts trivially on $S_R$.  For $(g,r)\in\SU(V_2)\times\cR_3$, $g\omega\otimes vr^{-1}=\omega\otimes v$ in
the complex tensor product iff $g\omega=\lambda\omega$ and $vr^{-1}=\lambda^{-1}v$ for some $\lambda\in\U(1)$.
Such $r$ form the circle generated by $T_{3R}$ and are determined by $\lambda$; $g$ is
$\mathrm{diag}(\lambda,U)$ with $U\in\U(\omega^\perp)$, $\det U=\lambda^{-1}$.  So $(g,r)\mapsto U$ is an
isomorphism onto $\U(3)$.  The circle $\lambda=e^{i\theta}$, $U=e^{-i\theta/3}\mathbf1$, has generator
$-i(B-L)-2iT_{3R}=-2iY$.
(iii) If $(g,\ell,r)$ acts trivially on $S_L$ then $g=\mu\mathbf1$, $\ell=\mu^{-1}\mathbf1$ with
$\ell\in\SU(2)$, so $\mu=\pm1$; $S_R$ gives $r=\mu^{-1}$ and the bidoublet $\ell=r$.  Identify $H$ with
$\U(3)\times\SU(2)$ by (ii).  The map $f(U,\ell):=(U/\det U,\ \det U\cdot\ell)$ is a homomorphism into
$\mathrm S(\U(3)\times\U(2))$; it is onto (given $(A,B)$, take $\delta$ with $\delta^2=\det B$, $U=\delta A$,
$\ell=B/\delta$), and its kernel is $\{(\pm\mathbf1,\pm\mathbf1)\}=\Z_2$.  The map
$(C,\ell,z)\mapsto(z^2C,\,z^{-3}\ell)$ from $\SU(3)\times\SU(2)\times\U(1)$ onto
$\mathrm S(\U(3)\times\U(2))$ has kernel $\{(z^{-2}\mathbf1,z^3\mathbf1,z):z^6=1\}\cong\Z_6$ \cite{BaezHuerta}.  Under
$f$ the hypercharge circle ($\lambda=e^{i\theta}$, $U=e^{-i\theta/3}\mathbf1$) goes to
$(e^{2i\theta/3}\mathbf1,e^{-i\theta}\mathbf1)$, the image of $z=e^{i\theta/3}$.
(iv) $T_{3R}=0$ on $S_L$ and $\pm\frac12$ on $S_R$; $\frac12(B-L)=-\frac12$ on $\omega$, $\frac16$ on
$\omega^\perp$; $T_{3L}=\pm\frac12$ on $S_L$ and $0$ on $S_R$.
\end{proof}

\begin{proposition}[the mixing angle for three structures]\label{prop:angle38}
Assume Hypothesis~G including clause (v): $\mathrm{Lie}(G)$ is normalised by the single trace form
$-\tr_{S_L\oplus S_R}(XY)$ and the three factors have equal couplings $g_4=g_L=g_R$.  Then $\tfrac12(B-L)=\sqrt{2/3}\,T_{15}$ with $T_{15}$
normalised like $T_{3R}$, $1/g_Y^2=1/g_R^2+\tfrac23/g_4^2$, and
$\sin^2\theta_W=g_Y^2/(g_L^2+g_Y^2)=\tr T_{3L}^2/\tr Q^2=\tfrac38$.  If the hypercharge is the bare clock
generator $T_{3R}$ (one structure, Remark~\ref{rem:angle}), the same normalisation gives $\tfrac12$.
\end{proposition}
\begin{proof}
With $T_{15}=\mathrm{diag}(1,1,1,-3)/(2\sqrt6)$ in a basis $(\omega^\perp,\omega)$, $\tr_{\C^4}T_{15}^2=\tfrac12$
and $\tfrac12(B-L)=\mathrm{diag}(\tfrac16,\tfrac16,\tfrac16,-\tfrac12)=\sqrt{2/3}\,T_{15}$; both $T_{15}$ and
$T_{3R}$ have trace square $2$ on $S_L\oplus S_R$.  For $Y=\sum_ic_iT_i$ with equally normalised $T_i$,
$1/g_Y^2=\sum_ic_i^2/g_i^2$, which gives $g_Y^2=\tfrac35g^2$ and $\sin^2\theta_W=\tfrac38$; directly,
$\tr T_{3L}^2=2$ and $\tr Q^2=\tfrac{16}3$ over $S_L\oplus S_R$.  For $Y=T_{3R}$, $g_Y=g$.
\end{proof}

\begin{remark}[what fixes $g'/g$]\label{rem:couplings}
The structure fixes the form of the mass matrix and the tree-level relation $\rho=1$
(Theorem~\ref{thm:gauge}(iv)), not $g'/g$.  The gauge kinetic term is a separate invariant form on the Lie
algebra; on a product of simple factors and an abelian factor an invariant pairing between different
factors vanishes, so each factor has its own coefficient.  Each value of the angle comes from a choice
of that form:
\begin{center}\footnotesize
\begin{tabular}{@{}lll@{}}
\toprule
kinetic form and hypercharge & $g_Y/g_L$ & $\sin^2\theta_W$\\
\midrule
invariant under the transports, or under the clock-reduced group; $Y=T_{3R}$ & free & free in $(0,1)$\\
one coefficient on both transport factors (reversal-symmetric); $Y=T_{3R}$ & $1$ & $\tfrac12$\\
one trace on $S_L\oplus S_R$ (G(v)); $Y=T_{3R}+\tfrac12(B-L)$ & $\sqrt{3/5}$ & $\tfrac38$\\
one coefficient times the Killing form of $\su(4)\oplus\su(2)\oplus\su(2)$; same $Y$ & $\sqrt{3/7}$ & $\tfrac3{10}$\\
\bottomrule
\end{tabular}
\end{center}
The last row uses that the Killing norm of $\su(N)$ is $2N$ times the fundamental trace, so $g_4^2=g_L^2/2$.
With the inputs of Remark~\ref{rem:scales}, a boundary value $S$ at scale $\mu$ requires threshold shifts
$\Delta_i$ of $\alpha_i^{-1}$ with $\Delta_Y-r\Delta_2=r\,\alpha_2^{-1}-\alpha_Y^{-1}$, $r=(1-S)/S$; at the meeting $\alpha_2=\alpha_3$ this is
$17.64$ for $S=\tfrac38$ and $-13.71$ for $S=\tfrac12$.  The model supplies no such shifts.
\end{remark}

\begin{remark}[scales]\label{rem:scales}
With one-loop running $\alpha_i^{-1}(\mu)=\alpha_i^{-1}(M_Z)-\frac{b_i}{2\pi}\ln\frac\mu{M_Z}$ of the Standard
Model couplings ($b_Y=\tfrac{41}6$, $b_2=-\tfrac{19}6$, $b_3=-7$, no thresholds),
$\alpha_Y^{-1}(M_Z)=\cos^2\theta_W\,\alpha^{-1}$, $\alpha_2^{-1}(M_Z)=\sin^2\theta_W\,\alpha^{-1}$,
$\sin^2\theta_W(\mu)=\alpha_2^{-1}/(\alpha_Y^{-1}+\alpha_2^{-1})$, $M_Z=91.1879$\,GeV and $\hat s_Z^2=0.23122$ from PDG 2026 \cite{ex-PDG26}, and $\hat\alpha^{-1}(M_Z)=127.95$ ($\overline{\mathrm{MS}}$)
and $\alpha_s(M_Z)=0.1180$ from PDG 2024 \cite{PDG} (PDG 2026 gives $\alpha_s(M_Z)=0.1177(9)$ and quotes the
on-shell $\alpha(M_Z^2)$), $\sin^2\theta_W$ reaches $\tfrac38$ at
$\mu\approx1.0\times10^{13}$\,GeV and $\tfrac12$ at $\mu\approx5.4\times10^{20}$\,GeV, above the Planck
mass; at the Planck mass $1.22\times10^{19}$\,GeV it is $0.471$.  At the crossing, $1.031\times10^{13}$\,GeV, $\alpha_2^{-1}=\tfrac35\alpha_Y^{-1}=42.4$ but
$\alpha_3^{-1}=36.8$: the equal-coupling boundary condition of Proposition~\ref{prop:angle38}, which
requires $g_4=g_3$, fails for the strong coupling by $5.6$ in $\alpha^{-1}$, as is known for the
non-supersymmetric Standard Model.  Two-loop terms, thresholds and intermediate states are ignored.
The strong and weak couplings meet at $9.7\times10^{16}$\,GeV, where $\sin^2\theta_W=0.436$ (with
$\alpha_s=0.1177$: $\alpha_3^{-1}=36.85$ at the $\tfrac38$ crossing, $1.031\times10^{13}$\,GeV, and the meeting at $9.4\times10^{16}$\,GeV); so with
Standard Model content no scale realises the equal couplings of Hypothesis~G(v) at one loop.  Neither
boundary value is a prediction of the low-energy value $0.231$.
\end{remark}

\smallskip
Under Hypothesis~G(i)--(iv), the stabiliser of $\Xi$, its global form and the hypercharge table are
theorems; $\sin^2\theta_W=\tfrac38$ follows from clause (v), which one-loop running does not realise.  Not obtained: why there are
three structures (a pair and one more) and not another number; three generations; the breaking
scale; masses and Yukawa couplings; unification of the strong coupling; any dynamics.
Theorem~\ref{thm:sm} and Proposition~\ref{prop:angle38} reproduce the Pati--Salam embedding with
lepton number as the fourth colour \cite{PatiSalam}, the left--right symmetric structure
\cite{MohapatraPati,SenjanovicMohapatra} and the value of \cite{GQW}; see \cite{BaezHuerta} for the
global forms.  What is new is only the identification of the ingredients with objects of the
comparison structure: $\SU(2)_L\times\SU(2)_R$ with the left and right transports, parity with
transport reversal, $\SU(4)$ with the unitaries of a closed pair, the lepton direction with its
origin-free singlet, and the colour triplet with $\mathrm{Sym}^2V$.

\begin{remark}[octonions]\label{rem:octonions}
Without the third identity of $(\Omega1)$ the geometry is already a round $S^1$, $S^3$ or $S^7$
\cite[Thm.~5.18, Cor.~5.19]{PaperI}; in this paper $S^7$ is excluded by that identity, equivalently by the group
structure of the base loops, and nothing is claimed for it.  Weakening the third identity to subsystems
generated by two elements admits the unit octonions $S^7$ with $\hp{x}{y}{z}=(x\bar y)z$ and $\kappa=\lVert\tfrac12(x+y)\rVert^2$, which satisfy
$(\Omega2)$, $(\Omega3)$ and violate the third identity (alternativity; $(x\bar w)x=2\ip{x}{w}x-w$); by
Adams \cite{Adams} a sphere with the first two identities and a jointly continuous operation has dimension $1$, $3$ or $7$.  On a round sphere
with the standard comparison, Paper~I proves more by Hurwitz's theorem: a ternary operation with the two
Mal'cev identities and $\kappa$ invariant under left and right transports exists only in dimensions $1,2,4,8$
of the ambient space (spheres $S^0,S^1,S^3,S^7$) \cite[Prop.~5.13]{PaperI}; it is unique on $S^1$ \cite[Rem.~5.14]{PaperI}, it is $a\bar bx$ or its
opposite on $S^3$ under continuity \cite[Thm.~5.15]{PaperI}, and on $S^7$ the continuous solutions form an
infinite-dimensional family, from which one origin-free identity, the identity that recovers a ternary
Moufang operation from its base loop \cite{BK12}, selects $(a\bar b)x$ up to isometry
\cite[Props.~5.16, 5.17]{PaperI}; the weakened third identity does not give di-associativity
\cite[Prop.~5.10]{PaperI}.  These statements concern the configuration space of one system; they imply
nothing for physics beyond kinematics, and in particular do not make $S^7$ the state space of a system.  The sphere
$S(V_2)\cong S^7$ is homogeneous under $\SU(4)$ with stabiliser $\SU(3)$
(Proposition~\ref{prop:pairsu3}) and under $\mathrm{Spin}(7)$ with stabiliser $G_2$, inside which the
stabiliser of a unit imaginary octonion is again $\SU(3)$ \cite{GG,Baez}.  Here colour arises from the
pair without weakening the axiom; a canonical identification of $S(V_2)$ with the unit octonions is
not available \cite{Dixon,Furey}.
\end{remark}

\begin{remark}[a pair and one more]\label{rem:hopf}
The state sphere of $N$ systems is $S(V_N)=S^{2^{N+1}-1}$.  For $N=1,2,3$ these are the total spaces of
the Hopf fibrations $S^3\to S^2$, $S^7\to S^4$, $S^{15}\to S^8$ over the projective lines of $\C$, $\HH$ and
$\mathbb O$, which have been related to one, two and three qubits \cite{MD,BC}; by Adams' theorem
\cite{Adams} there is no fibration $S^{2k-1}\to S^k$ with fibre $S^{k-1}$ for $k\notin\{1,2,4,8\}$, so the
sequence stops at $N=3$.  The matter space $V_2\otimes_\C V_3^{(2,1)}\cong\C^8$ of Hypothesis~G has the
dimension of three qubits.  This is an observation.  To turn it into a derivation of Hypothesis~G one
would need a principle selecting, as carriers of fields, the composites whose state sphere fibres over
the projective line of a normed division algebra, and a reason why the third structure enters through
a one-sided module rather than as a third tensor factor.

Stabilisers do not single out $N=3$: in $\U(2^N)$ every unit vector has stabiliser $\U(2^N-1)$, and in
a construction with more structures, e.g.\ two closed pairs with group $\SU(V_2)\times\SU(V_2)$, the
stabiliser of $\omega\otimes\omega$ is locally $\SU(3)\times\SU(3)\times\U(1)$, of dimension $17$, and no
stabiliser property distinguishes it from the case $N=3$.  Two properties single out small $N$: the Hopf fibrations above ($N\le3$), and
Proposition~\ref{prop:symmetric}, by which a permutation-symmetric origin-free dynamics is fixed up to
its time unit exactly for $N\in\{2,3\}$.  Neither selects $3$ over $2$; a principle that does is missing.
\end{remark}


%=====================================================================================
\paperpart{Part V.  Status, tests and open problems}

%=====================================================================================
\section{Physical examples}\label{sec:examples}

\subsection{What is closed, and what the examples test}\label{sec:closed}
\emph{Theorems from the axiom alone} ($(\Omega)$ and [NC]): the carrier $S^3$ with $\kappa(x,y)=\lVert(x+y)/2\rVert^2$
and its linear space $\R^4$ (Paper~I); the equivalence of the round-trip law with two-readout
determinacy; the complex structures commuting with $\cL$ are the quarter turns of clocks, all conjugate;
for a fixed clock, the Born kernel $\beta$ as a function on $M$, the unitary and antiunitary form of the
clock-preserving and clock-reversing symmetries, $\Theta^2=-1$ for one system, the basepoint-free form of
Schr\"odinger evolution and the clock holonomy; the Lorentz cones of states; the de~Sitter space of
dichotomies with inclusion as causal order, the vacuum equations for every invariant metric and the
static patch of a binary alternative.
\emph{Theorems under named inputs}: states are clock orbits ([L], [B]; equivalently [L], [Cl], [NT],
Proposition~\ref{prop:clockB}) and the weight is
$\beta$ ([W], [Comp], [N]), read as a probability ([T]) (Theorem~\ref{thm:readings}); for clock-read factors, one clock for
a composite ([Bil], [Loc], [O], C2, [N]), and with [Bil$_\cL$], [Span] the multiplication of independent
probabilities (Theorem~\ref{thm:oneclock}), in which C2 is essential (Proposition~\ref{prop:productsalone}) and the
symmetry of the weight excludes superselection (Proposition~\ref{prop:superselection}); the complex tensor
product (in addition C1(b)); entangled states (C2); Born weights of composites (C0--C2, [Span]) and their
uniqueness among nonnegative functions of $\beta_N$ with (a)--(c) of Theorem~\ref{thm:bornN}(iii); effects, Kraus
instruments and conditional states (C0--C3, [T]); sub-structures of every finite dimension (C0--C2, [Span]),
and their operations and composites (C0--C3, [T]), as systems under [Sub]; the singlet sector of a closed
register, the conditional dynamics relative to an internal clock with its cross-cut source, the finite
modular flow and its Gibbs criterion (C0--C3, [Span], [T]), and [Cl] (under [Ph], Corollary~\ref{cor:clbridge}); for single
outcomes, the statistics of any record process: Born at every stage iff martingale weights, with the
restart converse, affine unread maps from no-signalling, and the Born threshold and simplex laws exactly
for the invariant auxiliary law ([R], Theorems~\ref{thm:sorigidity}, \ref{thm:sosphere}, \ref{thm:sosimplex}), and the inverse-tick
bound ([Tr]); for finite memories, Gram extension and weak nesting ([Gw]), the reset bound (Gibbs bath, joint
unitarity, kept side marginal), exact phase recovery and gluing on chordal patterns; $\U(2)$ and the mass-form eigenvalue ratio (F), and $m_W=m_Z\cos\theta_W$ as tree-level masses (F, F$_{\mathrm g}$(a)); the Pati--Salam
group, its breaking and the hypercharges (C0--C3, [Span], G(i)--(iv)); $\sin^2\theta_W=\tfrac38$ (G(v));
the round-$S^4$ couplings $\kappa_V=\sqrt{1-\xi}$, $\kappa_{2V}=1-2\xi$ (F, F$_{\mathrm g}$(a), F$_{\mathrm H}$); the obstruction for the ungauged back-reacting centre field and the zero-stress flat vacuum (E, F, GR; F$_{\mathrm g}$).
\emph{Inputs not derived}: [Bil], [Bil$_\cL$], [Span], [Loc], [O], [N], Postulates C1(b), C2, C3, [Sub], the reading
inputs [L], [W], [Comp], [T], and [B] or, equivalently, [Cl] with [NT], the hypotheses $w\ge0$, $w=f(\beta_N)$ and (a)--(c) of
Theorem~\ref{thm:bornN}(iii), [Ph] (which gives [Cl]), [R], [Tr], [Gw], [Ue], [Ca], Hypotheses Ap, U, A, E, F, F$_{\mathrm g}$, F$_{\mathrm r}$, F$_{\mathrm H}$ (with F$_{\mathrm{Hf}}$), F$_{\mathrm S}$ (with F$_{\mathrm{Sg}}$, F$_{\mathrm{Sr}}$), P, G, GR; on the second route C0 and C1(a) in place of [Bil],
[Bil$_\cL$], [Loc], [O] and [N] for the composite.
\emph{Open}: a unit of length; $G_Nv^2$; Hypothesis~GR with sources and its consistency with
Hypothesis~E; the ratio $g'/g$; the number of structures; three generations; scales, masses and
Yukawa couplings; which single-outcome dynamics, if any, is realised (Hypothesis~Ap, Conjecture~\ref{conj:somixing}); small-phase control
in a cap corridor (Conjecture~\ref{conj:fmcontrol}); a derivation of [Loc], [O], C2, [Cl] (other than from [Ph]) or [B]; local equilibration in the
singlet sector (Conjecture~\ref{conj:cueth}); a principle selecting the radial potential; a microscopic origin of the round $S^4$ target
(Conjecture~\ref{conj:hbmicro}).
\emph{Testability} (Table~\ref{tab:falsify} lists it statement by statement).  The kinematic part (one system under the clock reading, composites under the inputs of
Theorem~\ref{thm:oneclock} including C2, C1(b) and C3) reproduces standard quantum mechanics exactly, so no example below distinguishes the inputs
of this paper from standard quantum mechanics.  The field part has no confirmed numerical consequence
beyond the tree-level relation $m_W=m_Z\cos\theta_W$ and the group and hypercharge assignments, which the
Standard Model shares.  The electroweak sector under F and F$_{\mathrm g}$(a), as it stands, contains no scalar and is contradicted by
the observed one (Section~\ref{sec:centre}).  With F$_{\mathrm r}$ it is the custodial Standard Model scalar sector,
compatible at tree level with the quoted couplings, with the scalar mass fitted and nothing that
distinguishes it from the Standard Model; the potential is not selected (Section~\ref{sec:radial}).  With F$_{\mathrm H}$
instead it has a testable one-parameter deviation, $\kappa_{2V}=2\kappa_V^2-1$ with $\kappa_V=\sqrt{1-\xi}\le1$; the quoted $\kappa_V$ gives
$\xi\le0.039$ ($f\ge1.24$\,TeV) under separate Gaussian errors, and the same deviation follows from any round
$S^4$ target, so it does not test the pair origin (Section~\ref{sec:hopfbase}, Table~\ref{tab:twoscalars}).  Otherwise a
distinguishing consequence would have to come from Hypothesis~E (events as
dichotomies, which needs a length relating the radius of $M$ to a physical scale, Example~\ref{ex:desitter})
or from the angle $\tfrac38$ versus $\tfrac12$ together with the scales at which they would hold
(Remark~\ref{rem:scales}).  The texture supplies none: ungauged it admits no maximal solution under E and
GR, and under F$_{\mathrm g}$ it is pure gauge.

\subsection{Conventions and tags}
Throughout, a point of $M=S^3$ is a normalised spinor $\psi\in\mathbb C^2$ (global phase included),
$c(\psi,\phi)=\operatorname{Re}\langle\psi|\phi\rangle$ is the Euclidean inner product of $\mathbb R^4$,
$\kappa=\tfrac12(1+c)=\cos^2(d/2)$ with $d$ the distance on $S^3$, and $\beta(\psi,\phi)=\lvert\langle\psi|\phi\rangle\rvert^2$
is the Born weight (Theorem~\ref{thm:born}).  A spin rotation by the angle $\alpha$ about the unit axis $n$ is
$U=e^{-i\alpha\,n\cdot\sigma/2}$; it acts on $M$ as a left transport.  For every $\psi$,
\begin{equation}\label{eq:ex-key}
c(\psi,U\psi)=\operatorname{Re}\langle\psi|U|\psi\rangle=\cos\tfrac\alpha2-\sin\tfrac\alpha2\,\operatorname{Re}\bigl(i\langle\psi|n\cdot\sigma|\psi\rangle\bigr)=\cos\tfrac\alpha2 ,
\end{equation}
because $\langle\psi|n\cdot\sigma|\psi\rangle$ is real.  So $\kappa(\psi,U\psi)=\tfrac12(1+\cos\frac\alpha2)=\cos^2(\alpha/4)$ for every $\alpha$, independently of the state, and
$d(\psi,U\psi)=\alpha/2$ for $\alpha\in[0,2\pi]$ ($d=2\pi-\alpha/2$ for $\alpha\in[2\pi,4\pi]$).  This identity underlies Examples~\ref{ex:neutron}, \ref{ex:larmor} and
\ref{ex:diffusion}.  Each example is tagged (a) exact instance of a theorem: the quantity is given by
the formula of the theorem; (b) instance under a named postulate or hypothesis of
Section~\ref{sec:assumptions}; or (c) analogy.  When an identity is (a) and its physical reading (b),
both tags are given.  Reading a measured relative frequency as a weight is [T] throughout.  Symbols
introduced inside an example are local to it.

\subsection{The comparison as an interference statistic}

\begin{example}[Neutron interferometry and the $4\pi$ periodicity of spinors; (a) for the identity, (b) for the reading]\label{ex:neutron}
\emph{System.}  A perfect-crystal neutron interferometer; in one path a static field $B$ over a length
$L$ rotates the spin by $\alpha=\gamma_{\mathrm n}BL/v_{\mathrm n}$ ($\gamma_{\mathrm n}$ the gyromagnetic ratio, $v_{\mathrm n}$ the
neutron speed), the other path is field-free \cite{ex-AS67,ex-Bernstein67}.  The experiments are those of
Rauch \emph{et al.}\ \cite{ex-Rauch75} and Werner \emph{et al.}\ \cite{ex-WCOE75} (1975).
\emph{Measured quantity.}  The intensity of the forward exit beam, $I_0\propto\lVert\tfrac12(\psi+U\psi)\rVert^2$ for
zero residual phase.
\emph{Formula.}  By \eqref{eq:ex-key},
\[
\lVert\tfrac12(\psi+U\psi)\rVert^2=\kappa(\psi,U\psi)=\tfrac12\bigl(1+\cos\tfrac\alpha2\bigr)=\cos^2\tfrac{d}{2},\qquad d=d(\psi,U\psi)=\begin{cases}\alpha/2,&\alpha\in[0,2\pi],\\2\pi-\alpha/2,&\alpha\in[2\pi,4\pi],\end{cases}
\]
the comparison of Proposition~\ref{prop:linear}(ii) read by a Hadamard test \cite{ex-CEMM98,AJL}.  Because the
right side does not depend on $\psi$, an unpolarised beam gives the same fringe; both 1975 experiments
used unpolarised neutrons \cite{ex-WCOE75,ex-Rauch88}, and with residual phase $\delta$ the intensity has
the form $I_0\propto1+\cos\delta\cos(\alpha/2)$ \cite{ex-WCOE75}.  The period in $\alpha$ is $4\pi$: at $\alpha=2\pi$,
$U=-\mathbf1$ and $\kappa(\psi,-\psi)=0$, while the Born weight is $\beta(\psi,-\psi)=1$.  The fringe therefore
follows $\kappa$ on $S^3$, which separates $\psi$ from $-\psi$, and not $\beta$ on $S^2$, which identifies them
(the distinction drawn after Theorem~\ref{thm:born}).  The zero is not specific to spin: a scalar
phase shifter of $\pi$ in one arm maps $\psi$ to $e^{i\pi}\psi=-\psi$ and gives the same $\kappa(\psi,-\psi)=0$.  What
the experiment shows is that a rotation by $2\pi$ produces this sign.  The round-trip law $(\Omega3)$ reads
$\kappa(\psi,U^2\psi)=(2\kappa(\psi,U\psi)-1)^2$: doubling the field integral squares the fringe correlation
$\cos(\alpha/2)$.
\emph{Numbers.}  Werner \emph{et al.}: wavelength $1.445$\,\AA, field $0$--$500$\,G over $L\approx2$\,cm
\cite{ex-WCOE75}.  With CODATA 2022 values \cite{ex-CODATA22} ($\gamma_{\mathrm n}=1.83247174\times10^{8}\,\mathrm{s^{-1}T^{-1}}$,
$m_{\mathrm n}$, $h$) a $4\pi$ precession needs $BL=4\pi h/(\gamma_{\mathrm n}m_{\mathrm n}\lambda_{\mathrm n})=271.3\,\mathrm{G\,cm}\,/\lambda_{\mathrm n}[\text{\AA}]$,
i.e.\ $187.7$\,G\,cm at $\lambda_{\mathrm n}=1.445$\,\AA.  Rauch \emph{et al.}\ measured the period as
$716.8\pm3.8^\circ$ \cite{ex-Rauch78}, against $720^\circ$.
\emph{Status.}  (a) for the identity $\kappa(\psi,U\psi)=\cos^2(\alpha/4)$ (\eqref{eq:ex-key} and
Proposition~\ref{prop:linear}(ii)).  (b) for reading the measured exit probability as $\kappa$: the spin
alone is read through a clock (Section~\ref{sec:readings}) and does not distinguish $\psi$ from $-\psi$; the
interferometer is a composite of path and spin, and the exit probability is the Born weight of the
composite, which equals $\kappa$ of the two spin configurations.  This uses the composite reading
(C0, C1 and [Span]) and [T].  [B] and [W] concern the reading of a closed system and are
not affected.  Real interferometers have contrast below one, $I_0=a_0+a_1\cos(\alpha/2)$ with $a_1<a_0$; $\kappa$
is the ideal-contrast case.
\end{example}

\begin{example}[Ramsey fringes; (a)]\label{ex:ramsey}
\emph{System.}  Ramsey's separated-oscillatory-field method \cite{ex-Ramsey50}: two $\pi/2$ pulses separated
by free evolution with phase $\varphi=\Delta\,\tau_{\mathrm R}$ (detuning $\Delta$, free time $\tau_{\mathrm R}$).
\emph{Measured quantity.}  The transition probability $q(\varphi)=\tfrac12(1+\cos\varphi)$.
\emph{Formula.}  $q(\varphi)$ is the Born weight between two states whose Bloch vectors lie on one great circle
of the Bloch sphere at angle $\varphi$.  The rays of a great circle form a real two-level system, and by
Proposition~\ref{prop:rebit} its transition probability is the comparison of the commutative solution
$S^1$: $q(\varphi)=\kappa_{S^1}(e^{i\varphi},1)=\cos^2(\varphi/2)$.  The round-trip law of $S^1$ is the
fringe-duplication identity $q(2\varphi)=(2q(\varphi)-1)^2$, and two-readout determinacy on $S^1$
(Theorem~\ref{thm:trd}, Remark~\ref{rem:trd}(ii)) gives, for samples at $\varphi_n=\varphi_0+n\,\delta$,
\[
q_{n+1}+q_{n-1}=2\cos\delta\;q_n+(1-\cos\delta);
\]
the samples $e^{i\varphi_n}$ are equally spaced on the great circle $S^1$, the case (i) of the geodesic
recurrence \cite[Thm.~6.2(i)]{PaperI}.  This inhomogeneous three-term recurrence of uniformly sampled fringes is the identity
on which Prony's method \cite{ex-Prony1795} and linear-prediction analysis of NMR time-domain signals
\cite{ex-BdBBvO85} rest.
\emph{Status.}  (a): Proposition~\ref{prop:rebit} with Theorem~\ref{thm:trd} for the commutative solution;
\eqref{eq:ex-key} is not used.  On a single great circle these statements reduce to
$\cos2\theta=2\cos^2\theta-1$; their content on $S^3$ is that the same law holds for all pairs
(Example~\ref{ex:larmor}).
\end{example}

\subsection{Uniform transport and two-readout determinacy}

\begin{example}[Stroboscopic Larmor precession; (a) for the spinor identities]\label{ex:larmor}
\emph{System.}  A spin-$\tfrac12$ in a static field $B$ along $z$, sampled at times $n\tau$: $\psi_n=U^n\psi_0$ with
$U=e^{-i\omega\tau\sigma_z/2}$, $\omega=\gamma_{\mathrm p}B$ for a proton.
\emph{Formula.}  The samples satisfy the heap recurrence $\psi_{n+1}=\hp{\psi_n}{\psi_{n-1}}{\psi_n}$ (\cite[Thm.~6.2]{PaperI};
Theorem~\ref{thm:dyn}(i)), which in spinor form is
\[
\psi_{n+1}=2\,c\,\psi_n-\psi_{n-1},\qquad c=c(\psi_n,\psi_{n-1})=\cos\tfrac{\omega\tau}2\quad\text{(for every }\psi_0),
\]
the Cayley--Hamilton identity $U^2=(\operatorname{tr}U)\,U-\mathbf 1$ applied to $\psi_{n-1}$.  Hence
$\psi_n=U_{n-1}(c)\psi_1-U_{n-2}(c)\psi_0$ with Chebyshev polynomials of the second kind \cite[Thm.~6.2]{PaperI}.  For every test
spinor $y$ the comparisons $r_n=\kappa(y,\psi_n)$ obey $r_{n+1}=2c\,r_n-r_{n-1}+1-c$ with coefficients
independent of $y$ (Remark~\ref{rem:trd}(ii)).  The Born weights $b_n=\beta(y,\psi_n)=\tfrac12(1+r_y\cdot r_{\psi_n})$
obey a different recurrence, $b_{n+1}=2\cos(\omega\tau)\,b_n-b_{n-1}+(1-\cos\omega\tau)(1+r_y^z r_{\psi_0}^z)$: the
step angle is the full Bloch angle $\omega\tau$ and the offset depends on $y$.  Two-readout determinacy
with test-independent coefficients holds for $\kappa$ and fails for $\beta$; the half angle $\omega\tau/2$ in the
first recurrence is the $4\pi$ periodicity of Example~\ref{ex:neutron}.
\emph{Numbers.}  Proton, $B=1$\,T: $\gamma_{\mathrm p}/2\pi=42.577478461$\,MHz/T \cite{ex-CODATA22}; Bloch period
$23.4866$\,ns, spinor period $46.9732$\,ns.  For $\tau=1$\,ns, $\omega\tau=0.267522$\,rad and $c=0.99106732$.
\emph{Status.}  (a) for the spinor recurrence and the recurrence of $\kappa$, which are identities.  Conventional
single-spin NMR and ESR record Bloch-vector components; the transverse components obey the homogeneous
full-angle recurrence with $\cos(\omega\tau)$, and the longitudinal one is constant, so it violates that
recurrence by $(2-2\cos\omega\tau)\lvert n_z\rvert$.  The half-angle (spinor) recurrence appears in phase-referenced
data: neutron interferometry and NMR interferometry with a coupled spin \cite{ex-SVV77}; its reading is
(b) as in Example~\ref{ex:neutron}.
\end{example}

\begin{example}[Grover search as a reflection recurrence; (a) for the identities, (b) under C3 for the realisation]\label{ex:grover}
\emph{System.}  Grover's algorithm \cite{ex-Grover97} on $N$ items, realised in NMR with $N=4$ \cite{ex-CGK98} and
with trapped ions with $N=8$ \cite{ex-Figgatt17}.
\emph{Formula.}  The Grover iterate is $G=-(\mathbf1-2|s\rangle\langle s|)(\mathbf1-2|w\rangle\langle w|)$.  By
Proposition~\ref{prop:complexrefl} each factor is a product of two round-trip reflections,
$\mathbf1-2|x\rangle\langle x|=s_xs_{ix}$, so $G=-s_ss_{is}s_ws_{iw}$; the sign $-\mathbf1=e^{\pi J_T}$ is an element of the
clock, so $G$ and the product of four reflections act in the same way on rays.  The orbit $x_k=G^ks$ stays
in the real plane spanned by $s$ and $w$, and satisfies $x_{k+1}=s_{x_k}(x_{k-1})=2c\,x_k-x_{k-1}$ with
$c=\operatorname{Re}\langle x_k,x_{k-1}\rangle=1-2/N$, the reflection form (d) of \cite[Thm.~6.2]{PaperI}; hence
$\lvert\langle w|x_k\rangle\rvert^2=\sin^2((2k+1)\theta)$, $\sin\theta=N^{-1/2}$.
\emph{Numbers.}  $N=4$: one iteration gives success probability $1$; $N=8$: $0.7812$ after one and $0.9453$
after two iterations (ideal values; the measured ones are in \cite{ex-Figgatt17}).
\emph{Status.}  (a) for the identities, which hold in any complex Hilbert space; the composite carries no
transport (Theorem~\ref{thm:obstruction}), so only the reflection form of the recurrence applies.  (b): that
these reflections are operations of a register is Postulate~C3 (Proposition~\ref{prop:complexrefl}).
\end{example}

\subsection{Brownian motion on $S^3$: rotational diffusion}

\begin{example}[Isotropic rotational diffusion of a spin-$\tfrac12$; (a) for Brownian motion on $S^3$]\label{ex:diffusion}
\emph{System.}  A spin-$\tfrac12$ subject to independent small random rotations with isotropic statistics,
$\langle\theta_i\theta_j\rangle=2D_{\mathrm r}\,t\,\delta_{ij}$ for the rotation vector (for example a fluctuating isotropic field in
the motional-narrowing regime).
\emph{Formula.}  The spinor performs Brownian motion on $S^3$.  Brownian motion as in \cite[Prop.~7.4]{PaperI} has generator
$\tfrac12\Delta$ on the unit sphere; since a rotation by $\theta$ is a step of length $\theta/2$ on $S^3$, the time of
Paper~I is $t_{\mathrm I}=D_{\mathrm r}t/2$.  The Laplacian on the unit $S^3$ has eigenvalue $-k(k+2)=-4j(j+1)$ on
harmonics of degree $k=2j$, so a degree-$k$ correlation decays as $e^{-j(j+1)D_{\mathrm r}t}$.  With $k=1$ \cite[Prop.~7.4]{PaperI},
\[
\mathbb E\,c(\psi_0,\psi_t)=\mathbb E\operatorname{Re}\langle\psi_0|U_t|\psi_0\rangle=e^{-3t_{\mathrm I}/2}=e^{-\frac34D_{\mathrm r}t},
\]
the phase-referenced contrast of an interferometer with the random rotation in one arm (as in
Example~\ref{ex:neutron}); with $k=2$, the Bloch vector (polarisation) decays as $e^{-4t_{\mathrm I}}=e^{-2D_{\mathrm r}t}$, so
the longitudinal and transverse relaxation times (towards zero polarisation) are equal,
$t_1=t_2=1/(2D_{\mathrm r})$.  The two rates are in the ratio $j(j+1)$ at $j=\tfrac12$ and $j=1$, $\tfrac34:2=3:8$, the
ratio of the quadratic Casimir values.  In standard quantum mechanics this follows from the expansion
of the $\SU(2)$ heat kernel in characters \cite{ex-Favro60}, and the zero-field equality of the
longitudinal and transverse rates under an isotropic fluctuating field is part of relaxation theory
\cite{ex-Abragam61}.  The integer-$j$ part of the same spectrum is Debye's law for rotational diffusion of a
rigid body, $\tau_j^{-1}=j(j+1)D_{\mathrm r}$ \cite{ex-Debye29,ex-Favro60}, whose ratio $\tau_1/\tau_2=3$ (dielectric versus NMR
or light-scattering correlation times) is the standard diagnostic of small-step diffusion.
\emph{Numbers.}  A Monte Carlo of $4\times10^4$ random walks with $D_{\mathrm r}=1$, $t=0.8$ gives contrast $0.549$
(formula $0.5488$) and polarisation $0.205$ (formula $0.2019$).
\emph{Status.}  (a): the decay rates are those of Brownian motion on $S^3$ \cite[Prop.~7.4]{PaperI}, which isotropic small-step
rotational noise produces.  The ratio $3:8$ is a standard consequence of spinor rotational diffusion;
every spinor theory gives it, so it does not test the axiom.  Noise-induced dephasing of neutron
interferometric contrast was studied in \cite{ex-Sulyok10}; we found no experiment comparing
interferometric contrast with polarisation under the same isotropic noise, so the ratio has not been
compared directly.
\end{example}

\subsection{Clock holonomy and geometric phase}

\begin{example}[Pancharatnam and Berry phases; (a) for the holonomy, (b) for the reading]\label{ex:pancharatnam}
\emph{System.}  Polarised light taken around a closed circuit on the Poincar\'e sphere in one arm of a
laser interferometer \cite{ex-Pancharatnam56,ex-BS88}; polarised neutrons in a helical field \cite{ex-BD87}; spin
flippers in a neutron interferometer \cite{ex-Wagh97}.
\emph{Measured quantity.}  The interferometric phase acquired by a state carried around a closed circuit
of rays.
\emph{Formula.}  For three rays with Bloch vectors $n_1,n_2,n_3$ and projections $P_1,P_2,P_3$,
Proposition~\ref{prop:threerays} gives
$B_{123}=\operatorname{Tr}(P_1P_2P_3)=\tfrac14\bigl(1+\sum_{i<j}n_i\!\cdot n_j+i\,n_1\!\cdot(n_2\times n_3)\bigr)$, of modulus
$(\beta_{12}\beta_{23}\beta_{31})^{1/2}$ and argument $+\Omega/2$, where $\Omega$ is the solid angle of the geodesic
triangle signed by $\operatorname{sign}\det(n_1,n_2,n_3)$ (positive when $n_1\to n_2\to n_3$ runs anticlockwise seen
from outside the sphere) \cite{ex-SB88,ex-MS93}.  By Proposition~\ref{prop:holonomy}, transport of the clock
(Hopf) connection around the geodesic triangle returns $x_1$ to $e^{-i\Omega/2}x_1$: $\arg\operatorname{Tr}(P_1P_2P_3)=+\Omega/2$, while the geometric phase
acquired by a spin-$\tfrac12$ state taken around the cycle $n_1\to n_2\to n_3\to n_1$ along geodesic arcs,
$\arg\langle\psi(0)|\psi(T)\rangle-\im\int_0^T\langle\psi|\dot\psi\rangle\,dt$ (Hermitian product antilinear in the first slot), is $-\Omega/2$.  Bhandari and Samuel found the measured phase equal in
magnitude to half the solid angle of the circuit \cite{ex-BS88}; Bitter and Dubbers measured the adiabatic
Berry phase \cite{ex-Berry84} of neutron spins in a helical field \cite{ex-BD87}, a continuous loop (the limit
of geodesic polygons), for which the Hopf holonomy is half the enclosed solid angle; it is not a three-point
invariant; Wagh \emph{et al.}\ obtained a pure
geometric phase shift of $\pi$ by reversing the current in one flipper \cite{ex-Wagh97}, which in this reading
is half the solid angle $2\pi$ of the lune between the two flip paths.
\emph{Relation to Paper~I.}  Proposition~\ref{prop:holonomy} and this example concern the Hopf (clock)
connection on the bundle $S^3\to S^2$ and the Bargmann invariant of the rays, the Hopf images; this is
what the experiments measure.  The triangle theorem \cite[Thm.~9.9]{PaperI} concerns Levi-Civita spin
transport along minimal arcs of $S^3$; its formula has the algebraic form of the Bargmann invariant
$\operatorname{Tr}(P_1P_2P_3)$ of three spin-$\tfrac12$ states, evaluated on the points $x,y,z$ themselves as unit vectors of
$\operatorname{span}(x,y,z)$, with $i$ replaced by $\nu\bar x$.  The two are different connections and different invariants, and
for points over given rays they differ in general: for $x=1$, $y=e^{i/2}$, $z=e^{j/2}$ the Levi-Civita excess is
$0.1302$, while the Hopf images of $x$ and $y$ coincide and the Bargmann invariant of the rays is real.
They share a half-angle law: the Hopf holonomy of a loop of rays multiplies by the phase of half the
enclosed solid angle, and the Levi-Civita holonomy of a loop bounding a region of area $A$ in a great
$2$-sphere is the rotation by $A$, with spin lift $\exp(\tfrac12A\,\nu\bar x_0)$, of spin angle $A$
\cite[Cor.~9.11(ii), Thm.~9.14]{PaperI}.  In the two conventions the half-angle laws carry opposite signs:
the companion's spin holonomy has the sign of $+\arg B_{123}$, the Hopf holonomy here is $e^{-i\Omega/2}$.  Both papers write $\arg\operatorname{Tr}(P_1P_2P_3)=+\Omega/2$ with the orientation
stated above.
\emph{Numbers.}  The octant triangle $n=(e_1,e_2,e_3)$: $B_{123}=(1+i)/4$, $\Omega=\pi/2$, $\arg B_{123}=\pi/4$, geometric phase of the cycle $-\pi/4$.
\emph{Status.}  (a) for the phase as the holonomy of the clock connection (Proposition~\ref{prop:holonomy}).
(b) for reading the measured phase as it: the phase is observed relative to a reference beam, the
composite reading of Example~\ref{ex:neutron}.  The helical-fibre experiment of Tomita and Chiao
\cite{ex-TC86} measures the photon helicity phase, a rotation equal to the full solid angle in momentum
space; it is a different instance of Berry's law and not of the half-angle formula.
\end{example}

\subsection{The clock reading: Born weight and time reversal}

\begin{example}[Malus's law; (b) under the clock reading]\label{ex:malus}
The transmission of linearly polarised light through an analyser at angle $\vartheta$ is $\cos^2\vartheta$; on the
Poincar\'e sphere the angle is $2\vartheta$, so $\cos^2\vartheta=\beta=\tfrac12(1+r\cdot r')$ (Theorem~\ref{thm:born}).  Malus's
law is the Born weight $\beta$ on $S^2$; the comparison $\kappa$ on $S^3$ appears only in a phase-referenced
experiment such as Example~\ref{ex:neutron}.  (b): that $\beta$ is the weight is Theorem~\ref{thm:readings}
under [L], [W], [Comp], [B], [N]; that it is a probability is [T].
\end{example}

\begin{example}[Kramers degeneracy; (a) for one spin, (b) for many]\label{ex:kramers}
Theorem~\ref{thm:wigner}(iv) gives the clock-reversing symmetry $\Theta=-i\sigma_yK$ with $\Theta^2=-1$ and
$\langle\psi|\Theta\psi\rangle=0$ for one spin-$\tfrac12$.  For $N$ spins with $\Theta_N=\Theta^{\otimes N}$, $\Theta_N^2=(-1)^N$, so every
level of a time-reversal-invariant Hamiltonian with an odd number of electrons is at least doubly
degenerate (Kramers' theorem \cite{ex-Kramers30}); in zero field these doublets are observed in the magnetic
resonance of odd-electron ions \cite{ex-AB70}.  A modern instance is the quantum spin Hall state in HgTe
quantum wells: for well width above $6.3$\,nm a residual conductance close to $2e^2/h$, destroyed by a
small magnetic field \cite{ex-Koenig07}; its protection uses $\Theta^2=-1$ (the band topology is additional
input).  (a) for one spin; (b) for $N$ spins, which uses the composite of C0, C1 and [Span]
and that time reversal acts factorwise.
\end{example}

\subsection{Composites}

\begin{example}[Exchange coupling and $\sqrt{\mathrm{SWAP}}$; (b) under C0, C1, {[Span]}]\label{ex:exchange}
\emph{System.}  Two electron spins in a double quantum dot \cite{ex-LD98,ex-Petta05}; pairs of ${}^{87}$Rb atoms in
double-well optical lattices \cite{ex-Anderlini07,ex-Trotzky08}.
\emph{Formula.}  The origin-free generators of a pair are $\alpha\mathbf1+JX$ with $X=\tfrac12(\mathbf1+\sigma\cdot\sigma)$ the
exchange (Theorem~\ref{thm:originfree}); this is the Heisenberg coupling
$J_{\mathrm{ex}}\,S_1\!\cdot S_2=\tfrac{J_{\mathrm{ex}}}2X-\tfrac{J_{\mathrm{ex}}}4$, so $J=J_{\mathrm{ex}}/2\hbar$.  At $Jt=\pi/4$, i.e.\ phase
$J_{\mathrm{ex}}t/\hbar=\pi/2$, the evolution is $e^{-i\pi/4}\sqrt X$ and maps $\lvert{\uparrow\downarrow}\rangle$ to
$(\lvert{\uparrow\downarrow}\rangle-i\lvert{\downarrow\uparrow}\rangle)/\sqrt2$, of concurrence $1$ (Corollary~\ref{cor:pair}).
\emph{Numbers.}  Petta \emph{et al.}: singlet return probability $\tfrac12\{1+\cos[J_{\mathrm{ex}}(\varepsilon)t_E/\hbar]\}$, SWAP at
$J_{\mathrm{ex}}t_E/\hbar=\pi$, a $180$\,ps $\sqrt{\mathrm{SWAP}}$, inhomogeneous dephasing time $10\pm1$\,ns, spin-echo
coherence time $1.2\,\mu$s \cite{ex-Petta05}; the $180$\,ps gate corresponds to $J_{\mathrm{ex}}/h=1.39$\,GHz,
$J_{\mathrm{ex}}=5.7\,\mu$eV (our conversion).  Anderlini \emph{et al.}: swap period $285(1)\,\mu$s, amplitude
$27(2)\%$, decay time above $10$\,ms \cite{ex-Anderlini07}, so $\sqrt{\mathrm{SWAP}}$ at a quarter period,
$\approx71\,\mu$s, and an exchange energy $\approx h\times3.5$\,kHz (our conversion).  Superexchange couplings from
$5$\,Hz to $1$\,kHz were observed in \cite{ex-Trotzky08}.
\emph{Status.}  (b): the tensor product of the composite is C0, C1 and [Span]; given it, the
classification of $\SU(2)$-invariant pair generators is Schur--Weyl duality.  In the dot experiment a
uniform field is present; it commutes with $X$ and leaves the $S$--$T_0$ dynamics origin-free, while
nuclear-field gradients do not.
\end{example}

\begin{example}[Real versus complex quantum theory; (b) under C0]\label{ex:real}
Renou \emph{et al.}\ showed that quantum theory over $\mathbb R$ with the real tensor product cannot reproduce
certain correlations of an entanglement-swapping network with two independent sources, although it
reproduces the outcomes of every multipartite experiment in which the parties may share arbitrary real
states \cite{Renou}.  Chen \emph{et al.}\ (superconducting qubits, swapping fidelity $0.952$) violate the real
bound $7.66$ by $43$ standard deviations \cite{ex-Chen22}; Li \emph{et al.}\ (photons, three parties, two
independent sources) by more than $4.5$ standard deviations \cite{ex-Li22}.  Relation to this paper: the
experiments test a composition rule, not a single-system reading.  For one system the clock reading
and the real four-level reading are both consistent (Remark~\ref{rem:ball}, Theorem~\ref{thm:readings}(iv)).
The real four-level reading $K=\{\pm1\}$, $w=c^2$, composes through $V\otimes_{\mathbb R}V$, which is real quantum
theory with the real tensor product; under the source-independence assumption of \cite{Renou} these
experiments therefore exclude it, independently of the binarity argument of
Theorem~\ref{thm:readings}(ii).  They are consistent with C0, which real quantum theory with the
real tensor product violates (commentary after Corollary~\ref{post:C0}).  Theorem~\ref{thm:oneclock} excludes
the other alternative, the real tensor product with two independent clocks, through C2 and a symmetric normalised weight; the
experiments do not test [Loc] or [O].
\end{example}

\subsection{Geometry of states}

\begin{example}[Partial polarisers as Lorentz boosts; (a) for identities shared by classical polarisation optics]\label{ex:boosts}
\emph{System.}  Partially polarised light; Stokes vector $(S_0,\mathbf S)$, a future-pointing vector of
$\mathbb R^{1,3}$ \cite{ex-HKN97}.  A diattenuator with intensity transmittances $t_{\max},t_{\min}$ has Mueller matrix
proportional to the pure boost of velocity $\lvert r\rvert=(t_{\max}-t_{\min})/(t_{\max}+t_{\min})$ (its
diattenuation) along its axis on the Poincar\'e sphere; retarders are rotations.
\emph{Formula.}  Theorem~\ref{thm:cone}(iii) and Theorem~\ref{thm:velocity}(ii): the filter $g=g^\dagger>0$ is the pure
boost, and two filters compose to a boost with the Einstein sum of velocities followed by a rotation,
visible in polarimetry as the retarder of the polar decomposition \cite{ex-LuChipman96}.  The degree of
polarisation from an unpolarised input is the velocity.
\emph{Numbers.}  Two diattenuators with $t_{\max}=1$, $t_{\min}=0.25$ ($\lvert r\rvert=0.6$, $\gamma=1.25$), axes at $0^\circ$ and
$45^\circ$ (orthogonal on the Poincar\'e sphere), unpolarised input: output degree of polarisation
$\sqrt{1-(1-\lvert r\rvert^2)^2}=0.76837$, $\gamma_{\mathrm{tot}}=\gamma^2=1.5625$; the residual Wigner rotation is about the
circular axis $S_3$ by $\arccos\bigl(2\gamma/(1+\gamma^2)\bigr)=12.680^\circ$ on the sphere, an optical rotation of the
azimuth by $6.340^\circ$.  These are standard polarimetry, computed here.
\emph{Status.}  (a) for the identities of Theorems~\ref{thm:cone}(iii) and \ref{thm:velocity}(ii), which the
coherency matrices of classical partially polarised light share with the states of this paper; the
example concerns classical light and tests no statement about quantum systems.  We found proposals
\cite{ex-HKN97} but no dedicated measurement of this Wigner rotation.
\end{example}

\begin{example}[Complementarity and $\gamma^{-1}=$ concurrence; (a) for the identity, (b) under C0--C2 for the reading]\label{ex:pvc}
\emph{System.}  A two-path interferometer whose path qubit is correlated with a second degree of freedom.
With predictability $\mathsf P$, visibility $\mathsf V$ \cite{ex-GY88,ex-Englert96} and concurrence $\mathsf C$
\cite{Wootters} of the joint pure state, $\mathsf P^2+\mathsf V^2+\mathsf C^2=1$ \cite{JB} (Theorem~\ref{thm:velocity}(iv)).  In
the velocity reading, $\lvert r\rvert^2=\mathsf P^2+\mathsf V^2$ and $\gamma(r)^{-1}=\mathsf C$; for classical fields the relation
between the degree of polarisation and $\mathsf V$, $\mathsf P$ is the polarisation coherence theorem \cite{ex-EQV17}.
\emph{Numbers.}  Qian, Vamivakas and Eberly measured $(\mathsf V,\mathsf D,\mathsf C)$ ($\mathsf D$ the which-path
distinguishability, equal to $\mathsf P$ here) for $13$ classical laser beams with path--polarisation
correlation; their tabulated $\mathsf V^2+\mathsf D^2+\mathsf C^2$ lies between $0.987$ and $1.000$ \cite{ex-QVE18}.  For
their state 6, $(\mathsf V,\mathsf D,\mathsf C)=(0.733,0.431,0.522)$: $\lvert r\rvert=0.8503$, $\gamma=1.900$, $\gamma^{-1}=0.526$ against the
measured $\mathsf C=0.522$.
\emph{Status.}  (a) for the identity on pure two-qubit states and on classical two-degree-of-freedom
fields; (b) for calling $\lvert r\rvert$ a velocity and $\mathsf C$ the concurrence of a composite, which uses
C0--C2.
\end{example}

\subsection{Fields and events}

\begin{example}[The electroweak mass relation and the hypercharges; (b) under F, F$_{\mathrm g}$(a) and G]\label{ex:ew}
\emph{Formula.}  Under Hypothesis~F the mass form has eigenvalue ratio $g^2:(g^2+g'^2)$ (Theorem~\ref{thm:gauge}(iv));
with F$_{\mathrm g}$(a) these are tree-level masses, $m_W=m_Z\cos\theta_W$ with $\rho=1$, in a model without a Higgs
boson (Section~\ref{sec:centre}); a form invariant on one side only
changes the ratio by the factor $\lambda_1/\lambda_3$ of Proposition~\ref{prop:berger}.
\emph{Numbers} (PDG 2026, electroweak review \cite{ex-PDG26}).  $m_Z=91.1879(20)$\,GeV;
$m_W=80.3625(77)$\,GeV (2026 combination without CDF; CDF II alone $80.432(16)$\,GeV, the value adjusted to the
CT18 parton distributions).  In the on-shell scheme $s_W^2:=1-m_W^2/m_Z^2$, so the relation holds by
definition; the masses give $0.22334(15)$, the fit $0.22342(9)$.  In $\overline{\mathrm{MS}}$, $\hat s_Z^2=0.23122(6)$
gives $m_Z\hat c_Z=79.954$\,GeV, below $m_W$ by $0.51\%$; the ratio $m_W^2/(m_Z^2\hat c_Z^2)=1.0103$ is the
radiative $\hat\rho=1.01016(9)$, dominated by the top--bottom splitting.  Physics beyond the Standard Model
enters as $\Delta\rho=\alpha(M_Z^2)\,T_{\mathrm{PT}}$ with the Peskin--Takeuchi parameter $T_{\mathrm{PT}}=0.021\pm0.055$
($U_{\mathrm{PT}}=0$) \cite{ex-PDG26}, i.e.\ $\Delta\rho=(1.6\pm4.3)\times10^{-4}$.  We use $\hat\alpha^{-1}(M_Z)=127.95$, the
$\overline{\mathrm{MS}}$ value of PDG 2024 \cite{PDG} that also enters Remark~\ref{rem:scales}; the on-shell value
$\alpha(M_Z^2)^{-1}\approx128.9$ \cite{ex-Steinhauser98} used by PDG 2026 gives the same digits.  The tree-level
statement of this paper is the $\rho=1$ part; the measured deviations are loop effects outside its scope.
The hypercharges of Theorem~\ref{thm:sm}(iv) ($Q_L:\tfrac16$, $u_R:\tfrac23$, $d_R:-\tfrac13$, $L:-\tfrac12$, $e_R:-1$,
$\nu_R:0$, with $Q=T_{3L}+Y$) coincide with the Standard Model assignment with a right-handed neutrino, and
all gauge and gravitational anomalies cancel per generation; $\sin^2\theta_W=\tfrac38$ under
Hypothesis~G(v) compares with $\hat s_Z^2=0.23122$ only after running (Remark~\ref{rem:scales}).
\emph{Status.}  (b): Hypothesis~F(i)--(iv) for the eigenvalue ratio, with F$_{\mathrm g}$(a) for the masses; under F and F$_{\mathrm g}$(a) alone the sector contains no scalar and is contradicted by the observed one (Section~\ref{sec:centre}); with F$_{\mathrm r}$ it is the Standard Model scalar sector (Section~\ref{sec:radial}); C0--C3, [Span] and Hypothesis~G(i)--(iv)
for the hypercharges; G(v) for $\tfrac38$.
\end{example}

\begin{example}[Positive cosmological constant; (c), with the kinematic-space literature as (a)]\label{ex:desitter}
Theorem~\ref{thm:desitter}(ii) states that the dichotomies of $M$ form $dS_4$ with $\Lambda=3$ in units of the
radius of $M$; the equation $\Lambda\ell^2=3$ for the de~Sitter radius $\ell$ is a pure number and fixes no physical
scale until a length is supplied.  The empirical side is $\Lambda>0$ \cite{ex-Riess98,ex-Perlmutter99}: Planck 2018
gives $H_0=67.36\pm0.54$\,km\,s$^{-1}$Mpc$^{-1}$ and $\Omega_\Lambda=0.6847\pm0.0073$ \cite{ex-Planck18}, hence
$\Lambda=3\Omega_\Lambda H_0^2/c^2=(1.089\pm0.021)\times10^{-52}\,\mathrm m^{-2}$, $\ell=\sqrt{3/\Lambda}=1.660\times10^{26}$\,m ($17.5$\,Gly),
against the Hubble radius $c/H_0=1.373\times10^{26}$\,m.  Identifying the radius of $M$ with $\ell$ would be an
additional hypothesis beyond [E], and with an ungauged back-reacting source no maximal solution carries the
full cap order (Proposition~\ref{prop:friedmann}).  The correspondence ``balls of a sphere $\leftrightarrow$ points of
de~Sitter space, nesting $\leftrightarrow$ causal order'' is used physically as the kinematic space of holography
\cite{CLMS15,dBHMN16,dBHHM16}; there it is an exact statement about the geometry of boundary regions, which
is the sense of Theorem~\ref{thm:desitter}.
\end{example}

\paragraph{Summary.}  Exact instances of theorems (a): Ramsey fringes (Proposition~\ref{prop:rebit};
Example~\ref{ex:ramsey}); stroboscopic Larmor precession (spinor recurrence, Theorem~\ref{thm:dyn}; conventional single-spin NMR and
ESR record Bloch components, the spinor recurrence appears in neutron and NMR interferometry;
Example~\ref{ex:larmor}); Grover search (reflection
recurrence; Example~\ref{ex:grover}); isotropic rotational diffusion (Brownian motion on $S^3$;
Example~\ref{ex:diffusion}); the geometric phase (clock holonomy, Proposition~\ref{prop:holonomy};
Example~\ref{ex:pancharatnam}); Kramers degeneracy of one spin (Example~\ref{ex:kramers}); partial polarisers
(classical light; Example~\ref{ex:boosts}); the complementarity identity (Example~\ref{ex:pvc}); the
state-independent fringe $\kappa(\psi,U\psi)=\cos^2(\alpha/4)$ of neutron $4\pi$ interferometry
(Example~\ref{ex:neutron}).  Instances under postulates or hypotheses (b): the reading of interferometric
fringes and phases as $\kappa$ and as holonomy (C0, C1, [Span]; Examples~\ref{ex:neutron},
\ref{ex:pancharatnam}); Malus's law (clock reading; Example~\ref{ex:malus}); Kramers degeneracy of many spins
(C0, C1, [Span]); the realisation of Grover search (C3); exchange-driven $\sqrt{\mathrm{SWAP}}$ (C0, C1, [Span];
Example~\ref{ex:exchange}); real-versus-complex network tests (C0; Example~\ref{ex:real}); the reading of
complementarity (C0--C2); the electroweak relations (F, F$_{\mathrm g}$(a), G; Example~\ref{ex:ew}; under F and F$_{\mathrm g}$(a) alone the sector contains no scalar and is contradicted by the observed one; with F$_{\mathrm r}$ it is the Standard Model scalar sector).  Analogy (c): $\Lambda>0$
(Example~\ref{ex:desitter}).  Standard consequences computed here that have not been compared directly
with a dedicated measurement: the ratio $3:8$ of spinor contrast to polarisation decay under isotropic
noise (Example~\ref{ex:diffusion}) and the Wigner rotation of two partial polarisers
(Example~\ref{ex:boosts}).  None of the examples distinguishes the inputs of this paper from standard
quantum mechanics (Section~\ref{sec:closed}).

%=====================================================================================
\section{Status and open problems}\label{sec:scope}

\begin{table}[h]
\caption{Status of the main statements.  The physical examples that instantiate them are listed in the
last column of Table~\ref{tab:answers} and tagged (a)/(b)/(c) in Section~\ref{sec:examples}.}
\label{tab:scope}
\footnotesize
\begin{tabular}{@{}p{0.55\textwidth}p{0.4\textwidth}@{}}
\toprule
Statement & Status\\
\midrule
Carrier $S^1$ or $S^3$; kernel; dimension $3$; compactness; positivity & theorem (Paper I)\\
Round-trip law $\Leftrightarrow$ two-readout determinacy & theorem (given $(\Omega1)$, $(\Omega2)$)\\
Linear space $\R^4$ & theorem\\
Complex structures commuting with $\cL$ $=$ clocks, all conjugate & theorem\\
{[Cl]} from the invisible total phase & theorem under [Ph] (Cor.~\ref{cor:clbridge}); [Ph] is an input; closure with invisible total phase alone does not suffice (Rem.~\ref{rem:clreal})\\
Closed register: singlets and marginals, constraint kernel, conditional dynamics with cross-cut source and error bounds, finite readouts, eight-site clock, restriction, finite modular flow and Gibbs criterion, entropy balances, finite-bath erasure & theorem under C0--C3, [Span], [T] (Thm.~\ref{thm:cusinglet}, Props.~\ref{prop:cukernel}, \ref{prop:cuspectrum}, Thms.~\ref{thm:cucond}, \ref{thm:cumodular}, Cor.~\ref{cor:cugibbs}, Props.~\ref{prop:cubalances}, \ref{prop:culandauer}); physical readings: [U], [A]; thermal time, area law, arrow, cosmology: programme; local equilibration: conjecture\\
States are clock orbits; weight $\beta$ (exponent $2$) & theorem under [L], [B], equivalently [L], [Cl], [NT] (Prop.~\ref{prop:clockB}); exponent $2$ with [W], [Comp], [N] (Thm.~\ref{thm:readings})\\
Real four-level and quaternionic one-level readings of the same $\R^4$ & theorem: composable but not binary; trivial\\
Born weight for one system; unitary/antiunitary symmetries; $\Theta^2=-1$ & theorem; probability reading: [T]\\
Schr\"odinger evolution (one system); uncertainty equality & theorem; $H$ free\\
Lorentz cones of signature $(1,4)$, $(1,3)$ & theorem\\
Complex composite with one clock (C0); probabilities of independent tests multiply (C1(a)) & theorem for clock-read factors under [Bil$_\cL$], [Span], [Loc], [O], [C2], [N] (Thm.~\ref{thm:oneclock}); C2 essential (Prop.~\ref{prop:productsalone}); excluded: superselection (Prop.~\ref{prop:superselection})\\
Tensor product; amplitudes multiply & theorem from the above and C1(b); or from C0, C1, [Span], and C2 for $N\ge3$ (product rule, Thm.~\ref{thm:product}, Cor.~\ref{cor:ampmult})\\
Entangled states; all unitaries are operations & theorem from C0--C3\\
Born weights for composites; their uniqueness & theorem from C0--C2, [Span]; uniqueness by Gleason among nonnegative functions of $\beta_N$ with (a)--(c), including marginal consistency (c)\\
Effects, POVMs, instruments, complete positivity, conditional states & theorem from C0--C3 and Convention T\\
Origin-free generators: clock, exchange; closed pair unique up to time unit & theorem; permutation-symmetric generator fixed up to time unit iff $N\in\{2,3\}$\\
Connected closed clock-invariant round-trip-closed sets are spheres of complex subspaces; all dimensions $1\le d\le2^N$; exchange sector $\mathrm{Sym}^N(\C^2)$ & theorem (Prop.~\ref{prop:rtclosed}; Cor.~\ref{cor:substructure} under C0--C2, [Span]; Prop.~\ref{prop:symsector} under C0, C1, [Span], C3); as systems: [Sub]\\
Skyrme term: area square, degree bound $6\pi^2v/e_{\mathrm S}$, scaling, Hopf charge and its bound, $\pi_1$ of configuration spaces, screening, radial unwinding & theorem under [F], [F$_{\mathrm S}$] (Thm.~\ref{thm:skbound}, Props.~\ref{prop:skhopf}, \ref{prop:skconfig}, \ref{prop:skscreen}, \ref{prop:skradial}, Thm.~\ref{thm:vk}); existence of minimisers and spin from rotations: cited; metastability: conjecture; particle readings: [P]\\
Finite memory: Gram extension and weak nesting, caps, capacity, reset heat, exact eraser criterion, gluing on chordal patterns, pointer records, redundancy, barycentre, holonomy bound & theorem (Thms.~\ref{thm:fmschur}, \ref{thm:fmeraser}, \ref{thm:fmglue}, \ref{thm:fmpointer}, \ref{thm:fmbary}, \ref{thm:fmberry}; Props.~\ref{prop:fmheat}, \ref{prop:fmredundancy}); window: [Gw]; geometry-to-entropy: [Ue]; access: [Ca]; broadcast: cited; small-phase control: conjecture (Conj.~\ref{conj:fmcontrol}); dynamical records: programme\\
Single outcomes & no covariant rule at ties; Kochen--Specker for composites (Prop.~\ref{prop:nooutcome}); theorem under [R]: Born at every stage $\Leftrightarrow$ martingale (Thm.~\ref{thm:sorigidity}), restart converse (Cor.~\ref{cor:sorestart}); no-signalling $\Rightarrow$ affine unread map (Prop.~\ref{prop:sochannels}); threshold and simplex laws (Thms.~\ref{thm:sosphere}, \ref{thm:sosimplex}); inverse-tick bound under [Tr] (Prop.~\ref{prop:soresolution}); invariant apparatus sampling: [Ap], Conj.~\ref{conj:somixing}; which dynamics is realised: open\\
Velocity space; $\gamma^{-1}=$ concurrence & theorem about states; (Q5) not answered\\
$dS_4$ of dichotomies; causal order; cosmological form; static patch & classical; theorem\\
Tetrad, torsion; uniqueness of the soldering & theorem, unique relative to the imposed equivariance\\
Observer $=$ binary alternative; static patch $=$ separating dichotomies & theorem; reading under [E]\\
Every invariant metric is de~Sitter (vacuum equations); order determines conformal class & theorem\\
Centre field: wave map, $\rho\propto R^{-2}$, $w=-\tfrac13$; linear stability in the expanding phase & theorem under Hypotheses E, F\\
Ungauged centre field: null-energy obstruction on the invariant metric; [E] forces the FLRW form; all Friedmann branches; no maximal solution carries the full cap order & theorem under [E], [F] (ungauged), [GR], $v\ne0$ (Props.~\ref{prop:nullobstruction}, \ref{prop:friedmann}); every solution is the restriction of a longer branch to an open conformal-time interval of length $\pi$ and is extendible; under F$_{\mathrm g}$ the centre field is pure gauge\\
Physical events are dichotomies & Hypothesis~E\\
Unit of length; Hypothesis~GR with sources; $G_Nv^2$ & open\\
Group $\U(2)$, residual $\U(1)$, mass-form eigenvalue ratio $g^2:(g^2+g'^2)$ & theorem under [F](i)--(iv)\\
$m_W=m_Z\cos\theta_W$ as tree-level masses; no scalar & under [F], [F$_{\mathrm g}$](a); contradicted by the observed scalar\\
Cone comparisons; degenerate Hessian; scalar doublet, $\rho=1$, $\kappa_V=\kappa_{2V}=1$, unitarity at tree level & theorem (Props.~\ref{prop:cone}, \ref{prop:conelinear}); under [F], [F$_{\mathrm g}$](a), [F$_{\mathrm r}$] (Prop.~\ref{prop:radial}); potential, $m_h$ free\\
Round $S^4$ target: vector masses, $\kappa_V=\sqrt{1-\xi}$, $\kappa_{2V}=1-2\xi$, Goldstone amplitudes $1/f^2$, Yukawa modifiers, two-term vacuum and mass, Hopf quotient, stabiliser obstruction & theorem under [F], [F$_{\mathrm g}$](a), [F$_{\mathrm H}$] (Thm.~\ref{thm:hbvector}, Cor.~\ref{cor:hbscatter}, Props.~\ref{prop:hbmass}, \ref{prop:hbhopf}, \ref{prop:hbobstruction}); Yukawas under [F$_{\mathrm{Hf}}$] (Prop.~\ref{prop:hbyuk}); microscopic origin: conjecture (Conj.~\ref{conj:hbmicro}); matching: programme\\
Stabiliser of $\Xi$, global form, hypercharges & theorem under [G](i)--(iv)\\
$\sin^2\theta_W=\tfrac38$ & under [G](v) (equal couplings), not realised by one-loop running\\
Number of structures; three generations; scales; couplings; masses & not derived\\
Single-system clauses of \cite{CRQO} & theorem after a clock (Thm.~\ref{thm:clauses})\\
Clauses \textsf{Q3}, \textsf{Q4} and R of \cite{CRQO} & theorem from C0--C3 and Convention T (Cor.~\ref{cor:clausesQ3})\\
\bottomrule
\end{tabular}
\end{table}

\begin{table}[h]
\caption{What could fail.  Classes: [M] mathematical consequence of a named hypothesis; [S] consequence of
standard physics given the named hypotheses; [D] measured input or fit; [X] a statement that could fail
against data or an explicit completion; [F] fails against data without an additional input.}
\label{tab:falsify}
\footnotesize
\begin{tabular}{@{}p{0.42\textwidth}p{0.53\textwidth}@{}}
\toprule
Statement & Class and status\\
\midrule
Kinematics: one system, composites, sub-structures of other dimension & [M] under the inputs of Sections~\ref{sec:one}--\ref{sec:many}; reproduces standard quantum mechanics, so no distinguishing prediction\\
Events: dimension four, Lorentzian, causal order by inclusion, static patches & [M] under [E]; no direct empirical test of nesting as causal order\\
$\Lambda>0$, $\Lambda\ell^2=3$ & [M] for the invariant metric; $\Lambda>0$ is [D]; no length calibration, hence no number\\
Ungauged centre field: $\rho\propto R^{-2}$, $p=-\rho/3$ (curvature-like), Friedmann branches, fixed-background spectrum & [S] under F, GR; it renormalises the curvature term, so the curvature radius measured geometrically and the curvature term of the expansion differ by the factor $1-4\pi G_Nv^2$, an [X] test that needs no value of $v$ in advance; no maximal solution with [E] (Props.~\ref{prop:nullobstruction}, \ref{prop:friedmann}); under F$_{\mathrm g}$ pure gauge\\
Mass-form eigenvalue ratio; $m_W=m_Z\cos\theta_W$, $\rho=1$, massless photon at tree level & [M] under F; masses [S] under F, F$_{\mathrm g}$(a); shared with the Standard Model; deviations are loop effects, [D]\\
Electroweak sector under F, F$_{\mathrm g}$(a): gauged $\OO(4)/\OO(3)$ sigma model, no scalar & [F]: contradicted by the observed scalar ($125.13\pm0.11$\,GeV, $\kappa_V=1.023\pm0.026$, $\kappa_{2V}=1.02\pm0.23$ \cite{ex-PDG26}); a further input is required, for instance a radial mode with a potential, which the data do not single out; not a statement about G, whose bidoublet has a radial mode\\
Electroweak sector under F, F$_{\mathrm g}$(a), F$_{\mathrm r}$: custodial Standard Model scalar sector & [S]: $\rho=1$, $\kappa_V=\kappa_{2V}=1$ at tree level, compatible with the quoted values; $m_h$ is [D] (fitted); no distinguishing statement\\
Electroweak sector under F, F$_{\mathrm g}$(a), F$_{\mathrm H}$: round $S^4$ target & [X]: $\kappa_{2V}=2\kappa_V^2-1$ with $\kappa_V\le1$ common to $W$ and $Z$, and Goldstone growth $s/f^2$; a significant deviation after radiative and matching corrections refutes F$_{\mathrm H}$; present bound from the quoted $\kappa_V$: $\xi\le0.039$, $f\ge1.24$\,TeV (one-sided 95\%, separate Gaussians, no covariance); $m_h$ is [D]\\
Pati--Salam group, global form, hypercharges, anomaly cancellation & [M] under G(i)--(iv); agree with the Standard Model; no breaking scale or heavy vectors supplied\\
$\sin^2\theta_W=\tfrac38$ (G(v)) or $\tfrac12$ as boundary values & [X]: with Standard Model running and no thresholds the colour condition fails; the required shifts are Remark~\ref{rem:couplings}; $\hat s_Z^2$ is [D]\\
Finite memory & [M]: extension, gluing and recovery criteria; [S]: the reset bound is the Landauer--Reeb--Wolf balance; physical record dynamics: programme\\
Single outcomes & [M]: martingale equivalence, threshold and simplex laws, inverse-tick bound under [Tr]; the Born statistics are those of standard quantum mechanics, so no example distinguishes; an objective-collapse completion would be [X] through recoherence tests (Example~\ref{ex:sorecoherence}), none is fixed here\\
Closed register and internal clocks & [M]: finite identities and bounds; Conjecture~\ref{conj:cueth} is a mathematical statement that computation at larger $n$ could refute; [S]: the conditional dynamics is that of standard quantum mechanics with a Page--Wootters reading; the physical readings [U], [A] are not tested here\\
Topological particles under F$_{\mathrm S}$ & [M]: bounds and topology; [X]: a particle identification (Hypothesis~P) would be testable once $v$, $e_{\mathrm S}$ and the quantisation character are fixed; none is fixed here\\
Three structures, generations, scales, masses, $g'/g$ & open; not derived\\
\bottomrule
\end{tabular}
\end{table}

\subsection{Audit of the assumptions}
Table~\ref{tab:audit} lists every assumption beyond Axiom~$\Omega$, where it is used, a model showing
that it is not implied by the others, and the weakest form that suffices.

\begin{table}[h]
\caption{Assumptions beyond Axiom $\Omega$ (the list of Section~\ref{sec:assumptions}).}
\label{tab:audit}
\footnotesize
\begin{tabular}{@{}p{0.13\textwidth}p{0.25\textwidth}p{0.3\textwidth}p{0.26\textwidth}@{}}
\toprule
Assumption & Used by & Not implied by the others & Weakest sufficient form\\
\midrule
{[NC]} & Secs.~\ref{sec:one}--\ref{sec:sm} & $S^1$ is a solution (rebit) & ---\\
{[L]} & Lemma~\ref{lem:normaliser}, Thm.~\ref{thm:readings} & $K$ a circle of left transports with $w=\lvert\ip xy\rvert^2$ for the left complex structure: binary, composable, normalised, not normalised by $\cL$ & left transports act on classes\\
{[W]} & Thm.~\ref{thm:readings}(i), (iii) & relative to [L]: $K=\{\pm1\}$ with $w=\beta$, $\cL$-invariant and not of this form (not binary); with [B] (a clock) no model with a continuous $\cL$-invariant weight & automatic for a clock (continuous $\cL$-invariant weights)\\
{[Cl]} & Prop.~\ref{prop:clockB} (with [NT], equivalent to [B] under [L]); a theorem under [Ph] (Cor.~\ref{cor:clbridge}) & relative to [L], [NT]: $K=\{\pm1\}$ (real four-level reading), not binary & a one-parameter group in $K$\\
{[Ph]} & Cor.~\ref{cor:clbridge} & the real reading: closed composites with invisible total sign, $K=\{\pm1\}$, no circle (Rem.~\ref{rem:clreal}) & (a) and (b) as stated\\
{[U]}, {[A]} & Sec.~\ref{sec:closedreg} (physical readings only) & the finite theorems hold without them & as stated\\
{[NT]} & Prop.~\ref{prop:clockB} & $K=N(T)$ or $\SU(2)$: [Cl] holds, no orthogonal classes & one orthogonal pair\\
{[B]} & Thm.~\ref{thm:readings}(ii) (under [L], given by [Cl] with [NT]) & not implied by [L], [W], [Comp], [N]: $K=\{\pm1\}$, $w=c^2$, composable, normalised, not binary & one class with exactly one orthogonal class\\
{[Comp]} & Thm.~\ref{thm:readings}(i), (iii) & $K=T$ with $\mathsf F(s)=3s^4-2s^6$: binary, normalised, $\mathsf F$ not multiplicative & as in Def.~\ref{def:reading}\\
{[N]} & Thm.~\ref{thm:readings}(iii), Thm.~\ref{thm:oneclock}(iii), (iv) & one system: a clock with $w=s_T^p$, $p\ne2$ (binary, composable); composites: $H_+\oplus H_-$ with $w=s_K^2$, all other inputs of Thm.~\ref{thm:oneclock} hold & over all maximal orthogonal families of classes\\
{[Bil]} & Thm.~\ref{thm:oneclock}(i), (iii) & no model constructed; on the second route C0 and C1(a) replace it (Thm.~\ref{thm:product}, no linearity) & unit values, multilinearity, equivariance for the clock torus\\
{[Bil$_\cL$]} & Thm.~\ref{thm:oneclock}(i), (iv) (form of $\Phi$) & no model satisfying all other inputs constructed; the $\C^3\oplus\C$ model after Prop.~\ref{prop:productsalone} is torus- but not $\cL$-equivariant and violates [N] & needed only for $\Phi=\sum a_\varepsilon P_\varepsilon\circ\otimes$\\
{[Loc]} & Thm.~\ref{thm:oneclock}(ii), (iii) & $K=\rho(T_1)$ on $V_1\otimes_\R V_2\cong\C^8$: all other inputs hold, phase of factor 2 observable & as stated\\
{[O]} & Thm.~\ref{thm:oneclock}(iii), (iv) & the asymmetric $w_{\mathrm{ss}}$ on $H_+\oplus H_-$ (Prop.~\ref{prop:superselection}) & symmetric; zero set $\{s_K=0\}$; $w(u,u)=1$ for (iv)\\
C0 & Sec.~\ref{sec:many}; theorem for clock-read factors under [Bil$_\cL$], [Span], [Loc], [O], [C2], [N] & $V\otimes_\R V$ with $T_1\times T_2$: weights multiply, relative phases unobservable, two-torus unobservable; excluded by [N] with a symmetric weight (Thm.~\ref{thm:oneclock}(iii)) & as an input (second route): equivalent to the complex Hilbert space with one clock it replaced; not to ``relative phases unobservable''\\
C1(a) & Thm.~\ref{thm:product}, Cor.~\ref{cor:ampmult} (second route only; a theorem by Thm.~\ref{thm:oneclock}(iv)) & $\mathcal P(p,q)=p$ (composite $=$ first factor) & as stated; no continuity\\
C1(b) & Thm.~\ref{thm:oneclock}(iv), Cor.~\ref{cor:ampmult} & $H_-$ (Thm.~\ref{thm:oneclock}); mixed type $[\psi\otimes\bar\phi]$ satisfies (a), not (b) & one triple with $\im B\ne0$ per factor\\
{[Span]} & Thm.~\ref{thm:oneclock}(i), (iv) (form, not one clock); Cor.~\ref{cor:ampmult} & $\mathcal P(p,q)=[\psi\otimes\phi\oplus0]$; $\Phi=x\otimes y\oplus0$ & definition of $\mathcal W$ and $V_N$\\
C2 & Thm.~\ref{thm:oneclock}(iii) (essential), Thm.~\ref{thm:comp}, Cor.~\ref{cor:ampmult} ($N\ge3$), Thms.~\ref{thm:bornN}, \ref{thm:instruments} & states $=\Sigma_N$ (Ex.~\ref{ex:mirror}); products only with $\mathcal E=\{+,-\}$ (Prop.~\ref{prop:productsalone}(i)) & for Thms.~\ref{thm:bornN}, \ref{thm:instruments}: all unit vectors are states\\
C3 & Prop.~\ref{prop:universal}, Thm.~\ref{thm:instruments}, Prop.~\ref{prop:pairsu3}(iv) & operations $=$ products of local unitaries & equivalent: the complex reflections $s_xs_{ix}$ and the clock (Prop.~\ref{prop:complexrefl})\\
{[Sub]} & Cor.~\ref{cor:substructure}, Props.~\ref{prop:symsector}, \ref{prop:subops} (as systems) & registers alone: only dimensions $2^N$ are systems & a restriction of preparations, tests and operations to $S(U)$\\
Conv.\ T & Thms.~\ref{thm:bornN}(ii), \ref{thm:instruments}; Prop.~\ref{prop:relative} & definitional (normalisation is [N] or, for composites, a theorem); other weights excluded by Thm.~\ref{thm:bornN}(iii) & probability reading of $\beta_N$ on complete tests\\
$w\ge0$, $w=f(\beta)$ & Thm.~\ref{thm:bornN}(iii) & --- & as stated (Gleason)\\
Hyp.\ E & readings in Sec.~\ref{sec:events}, Props.~\ref{prop:static}, \ref{prop:friedmann} & all theorems of Secs.~\ref{sec:geometry}, \ref{sec:eventspace} hold without it & as stated\\
Hyp.\ F & Thm.~\ref{thm:gauge}, Prop.~\ref{prop:berger}, Sec.~\ref{sec:centre} & a single structure & (i) for Thm.~\ref{thm:gauge}(i); (ii), (iii) for $\U(2)$; (iv) for the ratio of the mass-form eigenvalues\\
Hyp.\ G & Thm.~\ref{thm:sm}, Prop.~\ref{prop:angle38} & Hyp.~F alone gives $\U(2)$ & (i)--(iv) for Thm.~\ref{thm:sm}; (v) for $\tfrac38$; (iv) and the $(4,1,2)$ scalar are choices\\
Hyp.\ GR & Props.~\ref{prop:nullobstruction}, \ref{prop:friedmann} (vacuum: Prop.~\ref{prop:invmetric}) & the de~Sitter metric with $v\ne0$ is not a solution & field equations with a source; with [E] and the ungauged centre field, no maximal solution for $v\ne0$\\
Hyp.\ F$_{\mathrm S}$ & Sec.~\ref{sec:skyrme} & F without it (Derrick: no finite-size static texture) & a positive coefficient of the area square\\
Hyp.\ P & Sec.~\ref{sec:skyrme} & the topology equates the degree and the Hopf label (Def.~\ref{def:hopfcharge}) & as stated\\
Hyp.\ F$_{\mathrm H}$ & Sec.~\ref{sec:hopfbase} & F, F$_{\mathrm r}$ (flat target; also not forced by the pair: Prop.~\ref{prop:hbobstruction}) & a round $S^4$ target of radius $f$, gauged in $\SO(4)$\\
Hyp.\ F$_{\mathrm{Hf}}$ & Prop.~\ref{prop:hbyuk} & the embedding is a choice ($\mathbf4$, $\mathbf5$, $\mathbf{10}$, $\mathbf{14}$ give different $\kappa_f$) & as stated\\
Hyp.\ F$_{\mathrm r}$ & Sec.~\ref{sec:radial} & F with an $S^3$-valued configuration (no scalar) & values in $V\setminus\{0\}$; potential with a minimum on a sphere\\
Hyp.\ F$_{\mathrm g}$ & Sec.~\ref{sec:centre} (flat gauging, masses, no scalar), Rem.~\ref{rem:ew}, Ex.~\ref{ex:ew}, Table~\ref{tab:answers} (Q6, Q7) & F with a non-dynamical connection (Prop.~\ref{prop:friedmann}) & positive Yang--Mills term; flat covariantly constant vacuum\\
\bottomrule
\end{tabular}
\end{table}

Postulate~\ref{post:C3} is not derivable from C0--C2 and the
one-system theorems (Proposition~\ref{prop:universal}); its earlier form (all reflections are
operations) is stronger than needed.

\begin{problem}
Is there a principle, formulated in $(M,\hp{\cdot}{\cdot}{\cdot},\kappa)$, that fixes the ratio $g'/g$ of
Theorem~\ref{thm:gauge} (the missing step is stated after Proposition~\ref{prop:symmetric}), or the
permutation-symmetric generator for $N\ge4$?
\end{problem}
\begin{problem}
Derive Hypothesis~GR in the presence of sources, or another dynamics for the metric of $\cD$, from the
structure, consistently with Hypothesis~E (with an ungauged back-reacting centre field no maximal solution
carries the full cap order, Propositions~\ref{prop:nullobstruction}, \ref{prop:friedmann}); or derive Hypothesis~F$_{\mathrm g}$ and prove the coupled stability of its vacuum.
\end{problem}
\begin{problem}
Which record dynamics, if any, is realised: does an apparatus produce absorbing records and the
invariant auxiliary law of Theorems~\ref{thm:sosphere} and \ref{thm:sosimplex} (Conjecture~\ref{conj:somixing}), and does the
observer's filtration satisfy the martingale condition of Theorem~\ref{thm:sorigidity}?
\end{problem}
\begin{problem}
Find a principle, formulated in the comparison structure, that selects the potential of the radial
mode of Hypothesis~F$_{\mathrm r}$ (or, under G, of the bidoublet), and hence the scalar mass; the round-trip law on
the cone supplies no radial quadratic form (Proposition~\ref{prop:cone}(iii)).
\end{problem}
\begin{problem}
A binary reading is a clock reading for every weight (Theorem~\ref{thm:readings}(ii)), and on a clock every
continuous $\cL$-invariant weight has the form $\mathsf F\circ s_T$.  What remains are weights that are discontinuous
or not invariant under the left transports.  Can any of [Loc], [O], C2, [Cl] or [B] be derived from the axiom?
\end{problem}
\begin{problem}
Derive Hypothesis~G, or at least the number of structures (a pair and one more), from the axiom
(Remark~\ref{rem:hopf} states the missing principle); explain three generations; determine the
breaking scale and the failure of the strong coupling to meet the electroweak couplings at one loop
(Remark~\ref{rem:scales}).
\end{problem}
\noindent The later sections of \cite{CRQO} (a relational automaton, closure as a least fixed point,
and relational time) concern how histories are generated and are outside the scope of this paper.

%=====================================================================================
\section{Formalisation status and numerical checks}\label{sec:lean}
{\sloppy Nothing in this paper is machine-checked.  The formalisation targets, in increasing order of
difficulty (Certaine's lemma; the duplication lemma; the identities of Section~\ref{sec:one} in the
quaternion model; Theorem~\ref{thm:unc}; Lemma~\ref{lem:B0}; Theorem~\ref{thm:converse};
Theorem~\ref{thm:trd}; the Lie-algebra step of Theorem~\ref{thm:main}; the finite-dimensional statements
of Appendix~\ref{app:formal}), are preregistered at
\url{https://github.com/the-omega-institute/trureturing/issues/31739}.\par}

{\sloppy\emph{Library results.}  For Section~\ref{sec:outcomes}, source inspection of an existing Lean library (no build
for this paper) located declarations for finite parts: the rank-one weight pairing, weight sums,
conditional states and the unread ensemble (\texttt{Observer/Conditioning}); idempotence of record pinching and
the nonselective marginal; fresh-environment and repeated-record coherence decay; stabilisation of finite
contracting candidate sets (a stationary tail, without uniqueness); the qubit Bloch formulas; the
Fubini--Study angle triangle and record-time and record-count bounds (with the two-target speed premise);
the cosine fixed point; qubit pinching and its discarded mass; $R^2=-\mathbf1$ for skew-adjoint unitaries, Krylov growth
and permanent stability; and monitored first-return conservation.  No declaration was found for the
martingale equivalence, the threshold and simplex converses, the adaptive transcript bound or the Zeno
limit.  For Section~\ref{sec:finitemem} the same inspection located declarations for the Landauer bound, identity,
saturation and strictness from supplied balances; finite record pinching idempotence and the fixed-point
criterion; joint coherent reversal and the surviving-copy obstruction; unitary extension from equal Gram
matrices; orthogonal-complement projection identities; two-axis finite words and moving-axis invariance;
anchored and tree phase transport; a strictly increasing generation ledger; two- and three-anchor Gram
determinant formulas and Gram realisability; Kraus completion; equal effects with different two-step
joint laws; the axis-contraction commutator; the qubit Bargmann triple formula; a Bures--Wasserstein
barycentre from a dual certificate; and the separable-cone residual witness.  These are library results,
not checks of this paper's statements.\par}

{\sloppy Every identity stated above was checked numerically on random inputs by the accompanying
scripts (numpy; sympy optional): \texttt{check\_identities.py} (Section~\ref{sec:one}, Proposition~\ref{prop:avg}),
\texttt{check\_extensions.py} (Example~\ref{ex:mirror}, the Lie-algebra dimension in
Theorem~\ref{thm:comp}(ii) for $N=2$, Theorems~\ref{thm:originfree}, \ref{thm:velocity},
\ref{thm:gauge}(iv), Remark~\ref{rem:octonions}), \texttt{check\_events.py}
(Theorem~\ref{thm:desitter}(iv)(a,b)), \texttt{check\_examples.py} (the numbers of Section~\ref{sec:examples}),
and \texttt{check\_physics.py}, which covers all remaining statements and runs the further scripts
\texttt{check\_oneclock.py} (Theorem~\ref{thm:oneclock}), \texttt{check\_productsalone.py}
(Proposition~\ref{prop:productsalone}), \texttt{check\_clockB.py} (Proposition~\ref{prop:clockB}),
\texttt{check\_subsystems.py} (Section~\ref{sec:otherdim}), \texttt{check\_fields.py} (Sections~\ref{sec:fields}
and \ref{sec:sm}) and \texttt{check\_radial.py} (Section~\ref{sec:radial}, with own checks of the cone law, the
bilinear identity, the radial mass, the cancellation of the longitudinal growth and the fit to the
quoted data) and \texttt{check\_skyrme.py} (Section~\ref{sec:skyrme}, with own checks of the completed square, the
two normalisations, the trial integrals and binding, the Hopf coefficients, the compact centre field and
its conformal threshold $\sqrt2$, the barrier family and the Chern--Simons interpolation) and \texttt{check\_closedreg.py}
(Section~\ref{sec:closedreg}: twelve groups, including the singlet-chain data of Conjecture~\ref{conj:cueth}; own checks of the
Catalan dimensions, the swap sum, the constraint kernels, the eight-site clock and its code matrices, the
cross-edge example, the conditional equation and its bound, the modular and Gibbs identities, the
erasure balance and the fixed-point values) and \texttt{check\_hopfbase.py} (Section~\ref{sec:hopfbase}: eight groups;
own checks of the Hopf map and its factor two, the vector masses and modifiers, the Goldstone contact
amplitude, the Yukawa logarithmic derivatives, the two-term vacuum, mass and trilinear, the mass
counterexample, the stabiliser obstruction, the marginal example and the bounds on $\xi$ and $f$) and
\texttt{check\_outcomes.py} (Section~\ref{sec:outcomes}: eight groups; own checks of the two-tick counterexample,
the threshold laws and flows, the simplex race and the coarse block mass, the five-state chain, and
the cosine, spectral, basis and Zeno values) and \texttt{check\_finitemem.py} (Section~\ref{sec:finitemem}: nine
groups; own checks of the Schur slack and weak nesting, the reset balance with side information, the
eraser trace distance, the gluing witnesses, CHSH and the Peres--Mermin square, the redundancy bound,
the barycentre certificate, the latitude holonomy with zero work, the coverage bound, and the spiral
areas by independent quadrature).  Among them: Theorems~\ref{thm:trd}, \ref{thm:bornN}, \ref{thm:instruments}, \ref{thm:sm},
\ref{thm:product}; Propositions~\ref{prop:relative}, \ref{prop:solder}, \ref{prop:pairsu3}, \ref{prop:angle38},
\ref{prop:universal}, \ref{prop:cosmo}, \ref{prop:solderunique}, \ref{prop:centre}, \ref{prop:superselection} (with its
classification of weights and free split), \ref{prop:rtclosed}, \ref{prop:symsector}, \ref{prop:subops},
\ref{prop:nullobstruction}, \ref{prop:friedmann} (six cases) and \ref{prop:stability} (spectrum of $\mathcal J$ by
diagonalisation on spin blocks); Remarks~\ref{rem:trd}, \ref{rem:scales}, \ref{rem:couplings}; the corrections to
Examples~\ref{ex:neutron} and \ref{ex:larmor}; and the comparison of the Hopf holonomy with the Levi-Civita
triangle excess of Paper~I.  All scripts run from the paper's directory.  These are sanity checks, not
proofs.\par}

%=====================================================================================
\appendix
%=====================================================================================
\section{Relation to the closed relational reconstruction}\label{sec:crqo}

\subsection{The bottom-up route}
Clause names of \cite{CRQO} are set in sans-serif (\textsf{Q1}--\textsf{Q4}); (Q1)--(Q7) remain the questions of
Section~\ref{sec:intro}.  The manuscript \cite{CRQO} proceeds in layers.  \emph{B0}: a nonnegative self-weight $Q$ on a real
space whose binary cross-term $C_Q(x,y)=Q(x+y)-Q(x)-Q(y)$ obeys the merge law
$C_Q(x+y,z)=C_Q(x,z)+C_Q(y,z)$.  \emph{B1}: a common phase circle, $J^2=-1$ with
$Q((\cos\theta+\sin\theta J)x)=Q(x)$.  \emph{B}: binary positive relations
$M=\bigl(\begin{smallmatrix}a&c\\\bar c&b\end{smallmatrix}\bigr)\succeq0$ and their normalisation.
\emph{E1--E3}: for three directions, a positive real Gram block, a signed triple term, and one internal
$\U(1)$ fibre of minimal Chern class $\pm1$ whose triangle holonomy is fixed by the signed term.
\emph{S1--S4}: oriented directions on $S^2$, continuations preserving metric and orientation, retained
holonomy, and promotion of positive second-order residuals to new carrier directions.  \emph{\textsf{Q1}}: a
common Hermitian positive kernel with an orientation certificate.  \emph{\textsf{Q2}}: reversible
kernel-preserving continuation and zero Schur residual.  \emph{\textsf{Q3}}: complete readout invariant under
balanced re-encoding, and successor-kernel equality for instruments.  \emph{\textsf{Q4}}: product kernels for
independent systems, spanning the joint carrier, with complete joint effects.  \emph{R}: finite rank,
regularity, complete effect menu, ancilla-stable positivity.  Its main result is a conditional
representation theorem: models of these clauses are realised by finite-dimensional complex quantum
theory, uniquely up to unitary equivalence.  We now record which of its steps we have checked, then
compare clause by clause.

\subsection{Imported results}
The following statements are taken from \cite{CRQO}; the proofs are ours.

\begin{lemma}[quadratic merge, {\cite[\S4]{CRQO}}]\label{lem:B0}
Let $E$ be a real vector space and $Q\colon E\to[0,\infty)$ satisfy the merge law.  Then $B:=C_Q/2$ is
symmetric and biadditive and $Q(x)=B(x,x)$.  If moreover $J$ is real-linear and satisfies B1, then $B$ is
real bilinear,
$N=\{Q=0\}$ is a $J$-invariant subspace, and $K(x,y):=B(x,y)+iB(Jx,y)$ is a positive definite Hermitian
form on $E/N$, conjugate-linear in $x$ for the complex structure $J$.
\end{lemma}
\begin{proof}
$B$ is symmetric, and additive in the first argument by the merge law, hence biadditive.
$L(x):=Q(x)-B(x,x)$ is additive, since $Q(x+y)=Q(x)+Q(y)+2B(x,y)$.  Biadditivity gives
$B(x/n,x/n)=n^{-2}B(x,x)$ and $L(\pm x/n)=\pm L(x)/n$, so $0\le Q(\pm x/n)=n^{-2}B(x,x)\pm n^{-1}L(x)$,
i.e.\ $B(x,x)\ge n\lvert L(x)\rvert$ for all $n$; hence $L=0$ and $B(x,x)=Q(x)\ge0$.  The real
quadratic polynomial $q\mapsto Q(qx+y)$ is nonnegative on $\mathbb Q$, hence on $\R$, which gives
$B(x,y)^2\le Q(x)Q(y)$.  Under B1, $\cos\theta\,x=\tfrac12\bigl(u+u'\bigr)$ with
$u,u'=(\cos\theta\pm\sin\theta J)x$, and the parallelogram law gives $Q(\cos\theta\,x)\le\tfrac12(Q(u)+Q(u'))=Q(x)$;
so the additive map $t\mapsto B(tx,y)$ is bounded on $[-1,1]$ and therefore linear.  $Q(Jx)=Q(x)$ gives
$B(Jx,Jy)=B(x,y)$, hence $B(Jx,y)=-B(x,Jy)$, Hermitian symmetry of $K$, $K(Jx,y)=-iK(x,y)$, and
$J$-invariance of $N$, which is a subspace by Cauchy--Schwarz.
\end{proof}
In the present setting the hypotheses are theorems: $E=V$, $Q=\lVert\cdot\rVert_V^2$ is determined by
$\kappa$ (Proposition~\ref{prop:linear}), $J=J_T$, and $K=\ip{\cdot}{\cdot}_T$ is the formula of
Proposition~\ref{prop:clock}(ii).

\begin{proposition}[binary normalisation, {\cite[Thm.~4.2]{CRQO}}]\label{prop:binary}
For $M=\bigl(\begin{smallmatrix}a&c\\\bar c&b\end{smallmatrix}\bigr)$ Hermitian with $N=a+b>0$ put
$r=(2\re c,-2\im c,a-b)/N$.  Then $M/N=\tfrac12(\mathbf1+r\cdot\sigma)$, $1-\lvert r\rvert^2=4\det M/N^2$
and $\lvert r\rvert^2=2\tr(M/N)^2-1$; so $M\succeq0$ iff $\lvert r\rvert\le1$, and $\rk M=1$ iff
$\lvert r\rvert=1$.
\end{proposition}
\begin{proof}
The off-diagonal entry of $\tfrac12(\mathbf1+r\cdot\sigma)$ in the upper right is $\tfrac12(r_1-ir_2)$,
which fixes the sign of $r_2$.  The eigenvalues are $(1\pm\lvert r\rvert)/2$ because $(r\cdot\sigma)^2=\lvert r\rvert^2$;
the other identities follow by expansion.
\end{proof}
In \cite{CRQO} the second coordinate is printed as $+2\im c/N$ with $c$ the upper-right entry; with that
sign the first identity fails unless $\im c=0$ (\texttt{check\_physics.py}, 7b).  Here the proposition is
Theorem~\ref{thm:cone}(iii) in coordinates.

\begin{proposition}[stereographic lift, {\cite[Prop.~4.3]{CRQO}}]\label{prop:stereo}
For $w=x+iy\in\C$ let $\psi(w)=(1,w)/\sqrt{1+\lvert w\rvert^2}$.  Its Bloch vector is
$n(w)=(2x,2y,1-\lvert w\rvert^2)/(1+\lvert w\rvert^2)\in S^2$, and $w=(n_1+in_2)/(1+n_3)$ for $n_3\ne-1$; so
$\C\cup\{\infty\}\cong\C P^1\cong S^2$.  In the present setting $\psi(w)$ is a point of $M$ and $n(w)$
is its image in $P_T$ under the Hopf map of Theorem~\ref{thm:born}(i).
\end{proposition}
\begin{proof}
$4x^2+4y^2+(1-x^2-y^2)^2=(1+x^2+y^2)^2$; the Bloch components of $(\alpha,\beta)$ are
$(2\re\bar\alpha\beta,2\im\bar\alpha\beta,\lvert\alpha\rvert^2-\lvert\beta\rvert^2)$.
\end{proof}

\begin{proposition}[three rays, {\cite[Lemma~6.2]{CRQO}}]\label{prop:threerays}
For $x_1,x_2,x_3\in M$ put $r_i=r_{x_i}$, $a_{ij}=\ip{r_i}{r_j}$, $\chi=\ip{r_1}{r_2\times r_3}$, $G$ the
Gram matrix of the $r_i$, $P_i$ the projection onto $[x_i]$ and
$B_{123}=\ip{x_1}{x_2}_T\ip{x_2}{x_3}_T\ip{x_3}{x_1}_T$ (the Bargmann invariant).  Then
\[
\det G=\chi^2,\qquad B_{123}=\tr(P_1P_2P_3)=\tfrac14\bigl(1+a_{12}+a_{23}+a_{31}+i\chi\bigr),\qquad
16\,\beta_{12}\beta_{23}\beta_{31}=(1+a_{12}+a_{23}+a_{31})^2+\chi^2 .
\]
$B_{123}$ depends only on the rays.  It is invariant under $\cL$; clock reversal maps $\chi\mapsto-\chi$
and $B_{123}\mapsto\overline{B_{123}}$.  Hence the pairwise Born probabilities determine $\lvert\chi\rvert$,
and the sign of $\chi$ is fixed by the oriented clock.
\end{proposition}
\begin{proof}
$\det G=\chi^2$ is the Gram determinant of three vectors of $\R^3$.  The trace formula follows from
$(a\cdot\sigma)(b\cdot\sigma)=\ip{a}{b}+i(a\times b)\cdot\sigma$.  For the third identity expand
$16\prod\beta_{ij}=2\prod(1+a_{ij})$ and $\det G=1+2a_{12}a_{23}a_{31}-\sum a_{ij}^2$.  Under $\cL$ the
$r_i$ rotate; under $\Theta$, $r\mapsto-r$ and amplitudes are conjugated (Theorem~\ref{thm:wigner}).
\end{proof}

\begin{proposition}[the fibre and its holonomy]\label{prop:holonomy}
The distribution $H_x:=\{u\in T_xM:c(u,J_Tx)=0\}$ is a connection on the circle bundle $M\to P_T$,
whose fibres are the clock orbits; this bundle is the Hopf bundle, of Chern class $\pm1$.  Horizontal
transport along the geodesic of $P_T$ from $[x]$ to $[y]$ ($\beta\ne0$) ends at the representative $y'$ with
$\ip{x}{y'}_T>0$.  Transport around the geodesic triangle $[x_1],[x_2],[x_3]$ returns $x_1$ to
$\overline{B_{123}}/\lvert B_{123}\rvert\cdot x_1=e^{-i\Omega/2}x_1$, where $\arg B_{123}=+\Omega/2$ and $\Omega\in(-2\pi,2\pi)$ is the solid
angle of the geodesic triangle of Bloch vectors $r_1,r_2,r_3$, signed by $\operatorname{sign}\det(r_1,r_2,r_3)$
(positive when $r_1\to r_2\to r_3$ runs anticlockwise seen from outside the sphere).  For the octant
$r_i=e_i$, $B_{123}=(1+i)/4$ and $\Omega=\pi/2$.  So $\arg\operatorname{Tr}(P_1P_2P_3)=+\Omega/2$, while the geometric phase acquired by a spin-$\tfrac12$ state taken around
the cycle $r_1\to r_2\to r_3\to r_1$ along geodesic arcs, $\arg\langle\psi(0)|\psi(T)\rangle-\im\int_0^T\langle\psi|\dot\psi\rangle\,dt$ (Hermitian product antilinear in the first slot), is $-\Omega/2$.
Signed solid angles refer to the Bloch map of Proposition~\ref{prop:stereo}; the map $x\mapsto xi\bar x$ used in
Paper~I differs from it by an orientation-reversing isometry of $S^2$.  Paper~I writes
$\arg\operatorname{Tr}(P_1P_2P_3)=+\Omega/2$ for Bloch vectors with the same orientation convention; its triangle theorem
\cite[Thm.~9.9]{PaperI} concerns the Levi-Civita connection of $S^3$, a different connection
(Example~\ref{ex:pancharatnam}).
\end{proposition}
\begin{proof}
The great circle $\cos s\,x+\sin s\,u$ with $c(u,x)=c(u,J_Tx)=0$ has tangent orthogonal to $J_T$ of its
points at every $s$, and ends at $y'$ with $\ip{x}{y'}_T=\cos s>0$; it projects to a geodesic of $P_T$.
Composing three such steps gives $x_1'=e^{i\phi}x_1$ with $\ip{x_3'}{x_1'}_T>0$, so $e^{i\phi}$ is the phase
of $\overline{\ip{x_3'}{x_1}_T}$, which is the phase of $\overline{B_{123}}$ because the other two factors are
positive.  By Proposition~\ref{prop:threerays}, $4B_{123}=(1+\sum a_{ij})+i\chi$ with $\chi=\det(r_1,r_2,r_3)$, so $\arg B_{123}$ is
the angle with cosine proportional to $1+\sum a_{ij}$ and sine proportional to $\chi$; this is the formula
of van~Oosterom and Strackee \cite{VOS} for the signed $\Omega/2$; the Chern number of the Hopf bundle is
classical \cite{Hopf}.
\end{proof}

\begin{lemma}[Schur residual, {\cite[\S9]{CRQO}}]\label{lem:schur}
(i) For $\lvert a\rvert<1$ the Hermitian matrix with rows $(1,\bar a,\bar b)$, $(a,1,\bar a)$, $(b,a,1)$ is
positive semidefinite iff $\lvert b-a^2\rvert\le1-\lvert a\rvert^2$, and has rank two iff equality holds.
(ii) For any finite family in $M$, the Gram matrices of $\kappa$, $c$ and $\ip{\cdot}{\cdot}_T$ have
rank at most $5$, $4$ and $2$; relative to a subfamily spanning $V$, the Schur residual
$K(y,y)-K(y,X)K(X,X)^+K(X,y)$ of every further point vanishes.
\end{lemma}
\begin{proof}
(i) The Schur complement of the leading block is $(1-\lvert a\rvert^2)^{-1}\bigl((1-\lvert a\rvert^2)^2-\lvert b-a^2\rvert^2\bigr)$.
(ii) The residual is the squared distance from $y$ to the span of $X$ in the reproducing-kernel space,
whose dimensions are given by Theorem~\ref{thm:main} and Proposition~\ref{prop:clock}.
\end{proof}
For one system the carrier never needs enlargement.  For composites the product manifold $(S^3)^N$
does not span a comparison structure (Theorem~\ref{thm:obstruction}); its span $V_N$ is, by definition
([Span]), the tensor product of Corollary~\ref{cor:ampmult}, and Postulate~\ref{post:C2} fills $S(V_N)$ with the states that
\cite{CRQO} would call promoted residual directions.

\subsection{Single-system clauses as theorems}
We read the clauses of \cite{CRQO} with histories the points of $M$, directions their Bloch vectors
$r_x$, complex kernel $K=\ip{\cdot}{\cdot}_T$, and continuations the clock-preserving automorphisms
$\cL T\cong\U(2)$.

\begin{theorem}\label{thm:clauses}
Let $M=S^3$ carry a clock $T$.  In the reading above, clauses B0, B1, B, E1, E2, E3, S1--S4, \textsf{Q1}, \textsf{Q2} and
the finite-rank and $C^1$-time parts of R hold, as listed in Table~\ref{tab:clauses}.
\end{theorem}
\begin{proof}
B0, B1, \textsf{Q1}: Lemma~\ref{lem:B0} with $Q=\lVert\cdot\rVert_V^2$ and $J=J_T$, positivity of $\kappa$
(Theorem~\ref{thm:main}), and Proposition~\ref{prop:clock}; $J_T$ commutes with $\cL T$.  B:
Theorem~\ref{thm:cone}(iii) and Proposition~\ref{prop:binary}.  E1: $G$ is the Gram matrix of three
vectors of $\R^3$.  E2: Proposition~\ref{prop:threerays}.  E3, S3: Proposition~\ref{prop:holonomy}.  S1:
Theorem~\ref{thm:born}(i); the orientation of $P_T$ is the one for which $\chi$ has the sign of
Proposition~\ref{prop:threerays}.  S2: $\cL T$ acts on $P_T$ through $\SO(3)$
(Theorem~\ref{thm:wigner}(ii)), preserving the metric and $J_r(v)=r\times v$.  S4: at every parameter where $\psi$ and $\dot\psi$ are complex-linearly independent they span the
carrier $\C^2$, so the second-order residual vanishes there.  \textsf{Q2}: $\cL T$ is a group of
$K$-preserving maps acting unitarily (Theorem~\ref{thm:wigner}(iii)), and Lemma~\ref{lem:schur}(ii).  R:
the ranks are $5,4,2$; continuous one-parameter solutions are smooth (Theorem~\ref{thm:dyn}(ii)).
\end{proof}

\begin{corollary}\label{cor:clausesQ3}
Under C0--C3 and Convention~\ref{conv:T}, clauses \textsf{Q3} (complete readout,
successor-kernel equality, instruments, ancilla-stable complete positivity) and \textsf{Q4} (product kernel,
spanning, complete joint effects) hold as well, and so do the effect-menu and ancilla parts of R.
\end{corollary}
\begin{proof}
Theorems~\ref{thm:bornN} and \ref{thm:instruments} applied to one system and to composites;
Proposition~\ref{prop:avg} for the product kernel and spanning.  A Kraus map tensored with the identity
on an ancilla is again of Kraus form.
\end{proof}

In \cite[\S5.6]{CRQO} the residual $R=\lVert(\mathbf1-P_s)\ddot\psi\rVert^2$, with $P_s$ the projection onto
$\mathrm{span}_\C\{\psi,\dot\psi\}$, is stated to vanish identically for every curve in a fixed
two-dimensional carrier.  This requires $\psi,\dot\psi$ to be independent: the curve
$\psi(s)=e^{is}(1,s^2)/\sqrt{1+s^4}$ in $\C^2$ has $\dot\psi(0)=i\psi(0)$ and $R(0)=4$.  The reading of S4
used above includes this condition.

\begin{theorem}[a converse]\label{thm:converse}
Let $\mathcal H$ be a complex Hilbert space of dimension $2$.  For unit $x,y$ let $g_{x,y}\in\SU(\mathcal H)$ be
the unique element with $g_{x,y}y=x$.  Then $\hp{x}{y}{z}:=g_{x,y}z$ and
$\kappa(x,y):=\tfrac12(1+\re\ip{x}{y})$ make $S(\mathcal H)$ a comparison structure isomorphic to $S^3$,
with left transports $\SU(\mathcal H)$.  Multiplication by $e^{i\chi}$ is a clock, and its amplitude
$\ip{\cdot}{\cdot}_T$ is the given inner product.
\end{theorem}
\begin{proof}
$\SU(2)$ acts simply transitively on the unit sphere of $\C^2$: the unique element with first column
$v$ is $U(v)=\bigl(\begin{smallmatrix}v_0&-\bar v_1\\v_1&\bar v_0\end{smallmatrix}\bigr)$, so
$g_{x,y}=U(x)U(y)^\dagger$.  Identify $\mathcal H$ with $\HH$ by an orthonormal basis and
$(z_0,z_1)\mapsto z_0+jz_1$, with $\C$ acting on the right.  Left multiplication by a unit quaternion
$q$ is complex-linear, isometric and of determinant one (the determinant is a homomorphism from the
perfect group $S^3$ to $\U(1)$), so it equals $U(q)$; hence $g_{x,y}$ is left multiplication by $x\bar y$
and $\hp{x}{y}{z}=x\bar yz$.  With $\kappa=\frac12(1+\re\ip{x}{y})$ this is Example~\ref{ex:models}.
Multiplication by $e^{i\chi}$ is right multiplication, the right transport $z\mapsto\hp{z}{1}{e^{i\chi}}$, and
$\re\ip{x}{y}+i\re\ip{ix}{y}=\ip{x}{y}$.
\end{proof}

Theorem~\ref{thm:converse} applies to the Hilbert space produced by the representation theorem of
\cite{CRQO} when its continuations include all of $\SU(\mathcal H)$.  Two features block a converse
starting from the primitive clauses.  First, \cite{CRQO} allows the legal continuations to be a proper
subgroup of $\SU(2)$; a proper closed subgroup has dimension at most one, is not transitive on $S^3$,
and does not contain the left transports.  Second, clauses S and E concern rays and their phase
holonomy; the comparison $\kappa$ distinguishes $x$ from $-x$ and is defined only once the vector-level
kernel of \textsf{Q1} is available.  We do not claim an equivalence of the two axiom systems.

\begin{table}[h]
\caption{Clauses of \cite{CRQO} and their status under Axiom $\Omega$.  ``After clock'': derived once a
clock is chosen (Proposition~\ref{prop:clock}).}
\label{tab:clauses}
\footnotesize
\begin{tabular}{@{}p{0.07\textwidth}p{0.3\textwidth}p{0.22\textwidth}p{0.33\textwidth}@{}}
\toprule
Clause & Content in \cite{CRQO} & Status here & Argument\\
\midrule
B0 & quadratic merge law, $Q=B(x,x)$ & derived & $V$ is the kernel space of $c$; Lemma~\ref{lem:B0}\\
B1 & phase circle $J$, $J^2=-1$ & derived after clock & $J=J_T$; clocks exist and are conjugate (Prop.~\ref{prop:clock})\\
B & binary positive cone, Bloch ball, boundary $S^2$ & derived after clock & Thm.~\ref{thm:cone}(iii), Prop.~\ref{prop:binary}\\
E1 & positive $3\times3$ real Gram block & derived after clock & Bloch vectors lie in $\R^3$\\
E2 & signed triple term $\chi$ & derived after clock & Prop.~\ref{prop:threerays}: sign fixed by the oriented clock; mirror branch $=\Theta$-image\\
E3 & one $\U(1)$ fibre, $c_1=\pm1$, triangle holonomy & derived after clock & fibre $=$ clock orbit; $M\to P_T$ is the Hopf bundle (Prop.~\ref{prop:holonomy})\\
S1 & oriented directions on $S^2$ & derived after clock & Thm.~\ref{thm:born}(i)\\
S2 & continuations preserve metric and $J_r$ & derived after clock & Thm.~\ref{thm:wigner}(ii)\\
S3 & frame transport and holonomy retained & derived after clock & Prop.~\ref{prop:holonomy}\\
S4 & positive residuals enlarge the carrier & derived (one system) & carrier is $\C^2$; Lemma~\ref{lem:schur}\\
\textsf{Q1} & common Hermitian positive kernel; orientation certificate & real kernel derived; complex kernel after clock & Thm.~\ref{thm:main}, Prop.~\ref{prop:clock}\\
\textsf{Q2} & reversible continuation; zero Schur residual & derived after clock & Prop.~\ref{prop:aut}, Thms.~\ref{thm:wigner}, \ref{thm:dyn}; Lemma~\ref{lem:schur}\\
\textsf{Q3a} & complete readout $\Rightarrow$ effects, Born & derived from C0--C3 and Convention T & Thms.~\ref{thm:born}, \ref{thm:bornN}, \ref{thm:instruments}; Prop.~\ref{prop:effects}\\
\textsf{Q3b} & successor-kernel equality, instruments, complete positivity & derived from C0--C3 and Convention T & Thm.~\ref{thm:instruments}, Cor.~\ref{cor:clausesQ3}\\
\textsf{Q4} & product kernel; product histories span; joint effects complete & product kernel derived for clock-read factors from [Bil$_\cL$], [Span], [Loc], [O], C2, [N], C1(b) (Thm.~\ref{thm:oneclock}), or from C0, C1, [Span], and C2 for $N\ge3$ (Thm.~\ref{thm:product}, Cor.~\ref{cor:ampmult}); spanning is the definition of $V_N$; states and joint effects derived from C0--C3 & Cor.~\ref{cor:ampmult}, Prop.~\ref{prop:avg}, Thms.~\ref{thm:comp}, \ref{thm:instruments}\\
R & finite rank, regularity, effect menu, ancilla positivity & ranks and $C^1$ time derived; effect menu and ancillas derived from C0--C3; regularity of $\kappa$ derived \cite[Thm.~3.11]{PaperI} & Thms.~\ref{thm:main}, \ref{thm:dyn}, \ref{thm:instruments}\\
\bottomrule
\end{tabular}
\end{table}

\subsection{Complex, real and quaternionic}
The manuscript \cite{CRQO} selects the complex branch by declaration: its binary relation is defined with a complex
off-diagonal entry, and \textsf{Q1} requires an orientation certificate; real and quaternionic alternatives are
stated to lie outside the contract.  Here the selection is a sequence of theorems
(Remark~\ref{rem:complex}): the real qubit is the commutative solution $S^1$
(Proposition~\ref{prop:rebit}) and is excluded by noncommutativity; the noncommutative solution is a
quaternionic line with commuting left and right quaternionic actions, whose clock-free state space is
the five-dimensional spin factor and is not a qubit over any division algebra; and a clock, which
exists and is unique up to symmetry, makes it the complex qubit.  For composites
C0 (Theorem~\ref{thm:oneclock}) gives the complex tensor product (Corollary~\ref{cor:ampmult},
Proposition~\ref{prop:avg}).  What remains a
choice is the clock itself; all choices are conjugate.

\section{Statements ready for formalisation}\label{app:formal}
The following finite-dimensional statements are self-contained.  The notation follows the shape of a
Lean~4/Mathlib statement: hypotheses before the colon, conclusion after it.
\begin{enumerate}[label=(F\arabic*),leftmargin=2.6em]
\item \emph{Born weight as clock invariant.}  $z:\C$ $\vdash$
$(\forall\chi:\R,\ (\re(e^{-i\chi}z))^2\le\lvert z\rvert^2)\wedge(\exists\chi,\ (\re(e^{-i\chi}z))^2=\lvert z\rvert^2)
\wedge\frac1\pi\int_0^{2\pi}(\re(e^{-i\chi}z))^2d\chi=\lvert z\rvert^2$.
\item \emph{Three rays.}  $r_1,r_2,r_3:\R^3$, $\lVert r_i\rVert=1$ $\vdash$ $\det G=\chi^2$,
$\tr(P_1P_2P_3)=\frac14(1+a_{12}+a_{23}+a_{31}+i\chi)$ with $P_i=\frac12(\mathbf1+r_i\cdot\sigma)$, and
$2\prod(1+a_{ij})=(1+\sum a_{ij})^2+\chi^2$.
\item \emph{Uncertainty equality.}  $a,b,r:\R^3$, $\lVert r\rVert=1$ $\vdash$
$(\lVert a\rVert^2-\ip ar^2)(\lVert b\rVert^2-\ip br^2)=(\ip ab-\ip ar\ip br)^2+\ip{a\times b}{r}^2$.
\item \emph{Relative state.}  $\Psi:\C^m\otimes\C^n$, $a:\C^n$, $\phi:\C^m$ $\vdash$
$\ip{\phi\otimes a}{\Psi}=\ip{\phi}{\Psi_a}$ and $\sum_k\lvert\ip{\phi_k\otimes a}{\Psi}\rvert^2=\lVert\Psi_a\rVert^2$ for every
orthonormal basis $(\phi_k)$.
\item \emph{Stabiliser and hypercharges.}  In $\mathfrak{su}(4)\oplus\mathfrak{su}(2)$ acting on $\C^4\otimes\C^2$:
$\dim\{X:X(\omega\otimes v)=0\}=9$; $Y=T_{3R}+\frac12(B-L)$ lies in it, and the spectra of $Y$ on $S_L$ and
$S_R$ are $\{-\frac12,\frac16\}$ and $\{-1,-\frac13,0,\frac23\}$.
\item \emph{Mass form.}  $g,g',\lambda_1,\lambda_3:\R$ $\vdash$ the matrix of
$\lambda_1g^2((W^1)^2+(W^2)^2)+\lambda_3(gW^3-g'B)^2$ has eigenvectors $W^1,W^2$ (eigenvalue $\lambda_1g^2$),
$(0,0,g,-g')$ (eigenvalue $\lambda_3(g^2+g'^2)$) and $(0,0,g',g)$ (eigenvalue $0$).
\item \emph{Soldering isometry.}  $t,\tau:\R$, $v,w:\HH$, $\lVert v\rVert^2-t^2=1$, $t\tau=\re(\bar vw)$ $\vdash$
$\re(\overline{\bar nw_\perp})=0$ and $\lVert w\rVert^2-\tau^2=\lVert\bar nw_\perp\rVert^2-(\tau/\lVert v\rVert)^2$.
\item \emph{Product rule.}  $\mathcal P:P(\C^2)\times P(\C^2)\to P(\mathcal H)$ with
$\tau(\mathcal P(p,q),\mathcal P(p',q'))=\tau(p,p')\tau(q,q')$ $\vdash$ $\exists$ unitary $U$ onto a four-dimensional
subspace and types $\sigma,\sigma'$ with $\mathcal P([\psi],[\phi])=[U(\psi^\sigma\otimes\phi^{\sigma'})]$.  This one depends on the non-bijective
Wigner theorem \cite{Geher}.
\item \emph{Double counting and two families.}  $w$ symmetric and normalised on finite maximal orthogonal
families $A,B$ $\vdash$ $\lvert A\rvert=\lvert B\rvert$ (Lemma~\ref{lem:double}); in $H_+\oplus H_-\cong\C^4\oplus\C^4$ with componentwise
orthogonality, $\{e^\pm_i\}$ and $\{(e^+_i+e^-_i)/\sqrt2\}$ are maximal orthogonal families of sizes $8$ and $4$.
\item \emph{Shifts basis.}  The four product vectors of Proposition~\ref{prop:productsalone}(ii) are pairwise
orthogonal and no product vector of $(\C^2)^{\otimes3}$ is orthogonal to all of them.
\end{enumerate}


%=====================================================================================
\section*{Use of AI tools}
The authors used AI tools, namely Claude Fable 5.2 and Claude Opus 5.5 (via Claude Code) and Codex (gpt-6.2-sol), for literature
search, computation and double-checking.  All proofs can be checked by hand, and the computations can be
reproduced from the accompanying check scripts.  The authors take full responsibility for their
contributions.

\begin{thebibliography}{99}\footnotesize

\bibitem{PaperI} [Paper I] ``Ternions and minimal arcs: an origin-free calculus on the three-sphere'', companion paper, in preparation.

\bibitem{CRQO} [CRQO] ``The quantum origin of closed relational theory: a no-countermodel reconstruction from positive history, reversible continuation, complete readout, and composition'', unpublished manuscript, version 2.20, 2026.

\bibitem{AertsDaubechies} D.~Aerts and I.~Daubechies, Physical justification for using the tensor product to describe two quantum systems as one joint system, Helv.\ Phys.\ Acta \textbf{51}, 661--675 (1978).

\bibitem{AO53} A.~D.~Alexandrov and V.~V.~Ovchinnikova, Remarks on the foundations of special relativity, Vestnik Leningrad.\ Univ.\ \textbf{11}, 94--110 (1953) (in Russian). % UNVERIFIED: start page (94 or 95) seen only in a secondary source

\bibitem{AshtekarSchilling} A.~Ashtekar and T.~A.~Schilling, Geometrical formulation of quantum mechanics, in \emph{On Einstein's Path} (Springer, 1999), pp.~23--65, arXiv:gr-qc/9706069.

\bibitem{Adams} J.~F.~Adams, On the non-existence of elements of Hopf invariant one, Ann.\ of Math.\ \textbf{72}, 20--104 (1960).

\bibitem{Aaronson04} S.~Aaronson, Is quantum mechanics an island in theoryspace?, arXiv:quant-ph/0401062 (2004).

\bibitem{Adler95} S.~L.~Adler, \emph{Quaternionic Quantum Mechanics and Quantum Fields} (Oxford Univ.\ Press, 1995).

\bibitem{AJL} D.~Aharonov, V.~Jones, and Z.~Landau, A polynomial quantum algorithm for approximating the Jones polynomial, Algorithmica \textbf{55}, 395--421 (2009).

\bibitem{AB80} T.~Appelquist and C.~Bernard, Strongly interacting Higgs bosons, Phys.\ Rev.\ D \textbf{22}, 200--213 (1980).

\bibitem{Aronszajn} N.~Aronszajn, Theory of reproducing kernels, Trans.\ Amer.\ Math.\ Soc.\ \textbf{68}, 337--404 (1950).

\bibitem{BremnerEtAl} M.~J.~Bremner, C.~M.~Dawson, J.~L.~Dodd, A.~Gilchrist, A.~W.~Harrow, D.~Mortimer, M.~A.~Nielsen, and T.~J.~Osborne, Practical scheme for quantum computation with any two-qubit entangling gate, Phys.\ Rev.\ Lett.\ \textbf{89}, 247902 (2002).

\bibitem{Baez12} J.~C.~Baez, Division algebras and quantum theory, Found.\ Phys.\ \textbf{42}, 819--855 (2012).

\bibitem{BGW20} H.~Barnum, M.~A.~Graydon, and A.~Wilce, Composites and categories of Euclidean Jordan algebras, Quantum \textbf{4}, 359 (2020).

\bibitem{BruknerZeilinger09} \v{C}.~Brukner and A.~Zeilinger, Information invariance and quantum probabilities, Found.\ Phys.\ \textbf{39}, 677--689 (2009).

\bibitem{BaezHuerta} J.~C.~Baez and J.~Huerta, The algebra of grand unified theories, Bull.\ Amer.\ Math.\ Soc.\ \textbf{47}, 483--552 (2010).

\bibitem{BK12} W.~Bertram and M.~Kinyon, Torsors and ternary Moufang loops arising in projective geometry, arXiv:1206.2222 (2012).

\bibitem{Baez} J.~C.~Baez, The octonions, Bull.\ Amer.\ Math.\ Soc.\ \textbf{39}, 145--205 (2002).

\bibitem{Baer29} R.~Baer, Zur Einf\"uhrung des Scharbegriffs, J.\ Reine Angew.\ Math.\ \textbf{160}, 199--207 (1929).

\bibitem{Bell66} J.~S.~Bell, On the problem of hidden variables in quantum mechanics, Rev.\ Mod.\ Phys.\ \textbf{38}, 447--452 (1966).

\bibitem{BrodyHughston} D.~C.~Brody and L.~P.~Hughston, Geometric quantum mechanics, J.\ Geom.\ Phys.\ \textbf{38}, 19--53 (2001).

\bibitem{Busch03} P.~Busch, Quantum states and generalized observables: a simple proof of Gleason's theorem, Phys.\ Rev.\ Lett.\ \textbf{91}, 120403 (2003).

\bibitem{Bargmann} V.~Bargmann, Note on Wigner's theorem on symmetry operations, J.\ Math.\ Phys.\ \textbf{5}, 862--868 (1964).

\bibitem{BC} B.~A.~Bernevig and H.-D.~Chen, Geometry of the three-qubit state, entanglement and division algebras, J.\ Phys.\ A \textbf{36}, 8325--8339 (2003).

\bibitem{Certaine} J.~Certaine, The ternary operation $(abc)=ab^{-1}c$ of a group, Bull.\ Amer.\ Math.\ Soc.\ \textbf{49}, 869--877 (1943).

\bibitem{CMA21} G.~Carcassi, L.~Maccone, and C.~A.~Aidala, Four postulates of quantum mechanics are three, Phys.\ Rev.\ Lett.\ \textbf{126}, 110402 (2021).

\bibitem{CFMR04} C.~M.~Caves, C.~A.~Fuchs, K.~K.~Manne, and J.~M.~Renes, Gleason-type derivations of the quantum probability rule for generalized measurements, Found.\ Phys.\ \textbf{34}, 193--209 (2004).

\bibitem{CHKL97} X.~Chen, M.~Hindmarsh, M.~Kamionkowski, and A.~R.~Liddle, Instability of the one-texture universe, Phys.\ Rev.\ D \textbf{56}, 2051--2056 (1997).

\bibitem{CLMS15} B.~Czech, L.~Lamprou, S.~McCandlish, and J.~Sully, Integral geometry and holography, JHEP \textbf{10}, 175 (2015).

\bibitem{CU} J.-L.~Chen and A.~A.~Ungar, The Bloch gyrovector, Found.\ Phys.\ \textbf{32}, 531--565 (2002).

\bibitem{CDP} G.~Chiribella, G.~M.~D'Ariano, and P.~Perinotti, Informational derivation of quantum theory, Phys.\ Rev.\ A \textbf{84}, 012311 (2011).

\bibitem{Davis87} R.~L.~Davis, Texture: a cosmological topological defect, Phys.\ Rev.\ D \textbf{35}, 3705--3708 (1987).

\bibitem{Davis87b} R.~L.~Davis, Cosmic texture and the microwave background, Phys.\ Rev.\ D \textbf{36}, 997--999 (1987).

\bibitem{dBHMN16} J.~de~Boer, M.~P.~Heller, R.~C.~Myers, and Y.~Neiman, Holographic de~Sitter geometry from entanglement in conformal field theory, Phys.\ Rev.\ Lett.\ \textbf{116}, 061602 (2016).

\bibitem{dBHHM16} J.~de~Boer, F.~M.~Haehl, M.~P.~Heller, and R.~C.~Myers, Entanglement, holography and causal diamonds, JHEP \textbf{08}, 162 (2016).

\bibitem{DB} B.~Daki\'c and \v{C}.~Brukner, Quantum theory and beyond: is entanglement special?, in \emph{Deep Beauty}, edited by H.~Halvorson (Cambridge Univ.\ Press, 2011), pp.~365--392.

\bibitem{Davison08} T.~M.~K.~Davison, D'Alembert's functional equation on topological groups, Aequationes Math.\ \textbf{76}, 33--53 (2008).

\bibitem{Davison09} T.~M.~K.~Davison, D'Alembert's functional equation on topological monoids, Publ.\ Math.\ Debrecen \textbf{75}, 41--66 (2009).

\bibitem{Dixon} G.~M.~Dixon, \emph{Division Algebras: Octonions, Quaternions, Complex Numbers and the Algebraic Design of Physics} (Kluwer, 1994).

\bibitem{FLPS} S.~Friedland, C.-K.~Li, Y.-T.~Poon, and N.-S.~Sze, The automorphism group of separable states in quantum information theory, J.\ Math.\ Phys.\ \textbf{52}, 042203 (2011).

\bibitem{Dyson62} F.~J.~Dyson, The threefold way: algebraic structure of symmetry groups and ensembles in quantum mechanics, J.\ Math.\ Phys.\ \textbf{3}, 1199--1215 (1962).

\bibitem{FK} J.~Faraut and A.~Kor\'anyi, \emph{Analysis on Symmetric Cones} (Oxford Univ.\ Press, 1994).

\bibitem{Furey} C.~Furey, \emph{Standard model physics from an algebra?}, Ph.D.\ thesis, University of Waterloo (2015), arXiv:1611.09182.

\bibitem{Geher} G.~P.~Geh\'er, An elementary proof for the non-bijective version of Wigner's theorem, Phys.\ Lett.\ A \textbf{378}, 2054--2057 (2014).

\bibitem{GalleyMasanes17} T.~D.~Galley and Ll.~Masanes, Classification of all alternatives to the Born rule in terms of informational properties, Quantum \textbf{1}, 15 (2017).

\bibitem{GQW} H.~Georgi, H.~R.~Quinn, and S.~Weinberg, Hierarchy of interactions in unified gauge theories, Phys.\ Rev.\ Lett.\ \textbf{33}, 451--454 (1974).

\bibitem{GibbonsHawking} G.~W.~Gibbons and S.~W.~Hawking, Cosmological event horizons, thermodynamics, and particle creation, Phys.\ Rev.\ D \textbf{15}, 2738--2751 (1977).

\bibitem{Gleason57} A.~M.~Gleason, Measures on the closed subspaces of a Hilbert space, J.\ Math.\ Mech.\ \textbf{6}, 885--893 (1957).

\bibitem{GW} R.~Goodman and N.~R.~Wallach, \emph{Symmetry, Representations, and Invariants} (Springer, 2009).

\bibitem{GG} M.~G\"unaydin and F.~G\"ursey, Quark structure and octonions, J.\ Math.\ Phys.\ \textbf{14}, 1651--1667 (1973).

\bibitem{Barrett07} J.~Barrett, Information processing in generalized probabilistic theories, Phys.\ Rev.\ A \textbf{75}, 032304 (2007).

\bibitem{HancheOlsen83} H.~Hanche-Olsen, On the structure and tensor products of JC-algebras, Can.\ J.\ Math.\ \textbf{35}, 1059--1074 (1983).

\bibitem{HancheOlsen85} H.~Hanche-Olsen, JB-algebras with tensor products are C$^*$-algebras, in \emph{Operator Algebras and their Connections with Topology and Ergodic Theory}, edited by H.~Araki \emph{et al.}, Lecture Notes in Mathematics Vol.~1132 (Springer, Berlin, 1985).

\bibitem{McKagueMoscaGisin09} M.~McKague, M.~Mosca, and N.~Gisin, Simulating quantum systems using real Hilbert spaces, Phys.\ Rev.\ Lett.\ \textbf{102}, 020505 (2009).

\bibitem{ABW13} A.~Aleksandrova, V.~Borish, and W.~K.~Wootters, Real-vector-space quantum theory with a universal quantum bit, Phys.\ Rev.\ A \textbf{87}, 052106 (2013).

\bibitem{HardyWootters12} L.~Hardy and W.~K.~Wootters, Limited holism and real-vector-space quantum theory, Found.\ Phys.\ \textbf{42}, 454--473 (2012).

\bibitem{Bennett99} C.~H.~Bennett, D.~P.~DiVincenzo, T.~Mor, P.~W.~Shor, J.~A.~Smolin, and B.~M.~Terhal, Unextendible product bases and bound entanglement, Phys.\ Rev.\ Lett.\ \textbf{82}, 5385--5388 (1999).

\bibitem{DiVincenzo03} D.~P.~DiVincenzo, T.~Mor, P.~W.~Shor, J.~A.~Smolin, and B.~M.~Terhal, Unextendible product bases, uncompletable product bases and bound entanglement, Commun.\ Math.\ Phys.\ \textbf{238}, 379--410 (2003).

\bibitem{Hardy} L.~Hardy, Quantum theory from five reasonable axioms, arXiv:quant-ph/0101012 (2001).

\bibitem{HoehnWever17} P.~A.~H\"ohn and C.~S.~P.~Wever, Quantum theory from questions, Phys.\ Rev.\ A \textbf{95}, 012102 (2017).

\bibitem{HawkingEllis} S.~W.~Hawking and G.~F.~R.~Ellis, \emph{The Large Scale Structure of Space-Time} (Cambridge Univ.\ Press, 1973).

\bibitem{HJ} U.~Hertrich-Jeromin, \emph{Introduction to M\"obius Differential Geometry} (Cambridge Univ.\ Press, 2003).

\bibitem{Hopf} H.~Hopf, \"Uber die Abbildungen der dreidimensionalen Sph\"are auf die Kugelfl\"ache, Math.\ Ann.\ \textbf{104}, 637--665 (1931).

\bibitem{JB} M.~Jakob and J.~A.~Bergou, Quantitative complementarity relations in bipartite systems: entanglement as a physical reality, Opt.\ Commun.\ \textbf{283}, 827--830 (2010).

\bibitem{KT96} M.~Kamionkowski and N.~Toumbas, A low-density closed universe, Phys.\ Rev.\ Lett.\ \textbf{77}, 587 (1996).

\bibitem{HKM76} S.~W.~Hawking, A.~R.~King, and P.~J.~McCarthy, A new topology for curved space-time which incorporates the causal, differential, and conformal structures, J.\ Math.\ Phys.\ \textbf{17}, 174--181 (1976).

\bibitem{n6:skyrme} T.~H.~R.~Skyrme, A non-linear field theory, Proc.\ R.\ Soc.\ Lond.\ A \textbf{260}, 127--138 (1961).

\bibitem{n6:derrick} G.~H.~Derrick, Comments on nonlinear wave equations as models for elementary particles, J.\ Math.\ Phys.\ \textbf{5}, 1252--1254 (1964).

\bibitem{n6:faddeev} L.~D.~Faddeev, Some comments on the many-dimensional solitons, Lett.\ Math.\ Phys.\ \textbf{1}, 289--293 (1976).

\bibitem{n6:hatcher} A.~Hatcher, \emph{Algebraic Topology} (Cambridge University Press, 2002).

\bibitem{n6:esteban1986} M.~J.~Esteban, A direct variational approach to Skyrme's model for meson fields, Commun.\ Math.\ Phys.\ \textbf{105}, 571--591 (1986).

% author list as commonly cited; the publisher's abstract page lists M. J. Esteban
\bibitem{n6:esteban-muller} M.~J.~Esteban and S.~M\"uller, Sobolev maps with integer degree and applications to Skyrme's problem, Proc.\ R.\ Soc.\ Lond.\ A \textbf{436}, 197--201 (1992).

\bibitem{n6:esteban2004} M.~J.~Esteban, Existence of 3D Skyrmions, Commun.\ Math.\ Phys.\ \textbf{251}, 209--210 (2004) (erratum to \cite{n6:esteban1986}).

\bibitem{n6:hopf-original} H.~Hopf, \"Uber die Abbildungen der dreidimensionalen Sph\"are auf die Kugelfl\"ache, Math.\ Ann.\ \textbf{104}, 637--665 (1931).

\bibitem{n6:bundles} A.~Hatcher, \emph{Vector Bundles and K-Theory}, version 2.2 (2017).

\bibitem{n6:whitehead} J.~H.~C.~Whitehead, An expression of Hopf's invariant as an integral, Proc.\ Natl.\ Acad.\ Sci.\ USA \textbf{33}, 117--123 (1947).

\bibitem{n6:vk-original} A.~F.~Vakulenko and L.~V.~Kapitanski\u{\i}, Stability of solitons in $S^2$ in the nonlinear $\sigma$-model, Sov.\ Phys.\ Dokl.\ \textbf{24}, 433--434 (1979) [Dokl.\ Akad.\ Nauk SSSR \textbf{246}, 840 (1979)].

\bibitem{n6:faddeev-niemi} L.~Faddeev and A.~J.~Niemi, Stable knot-like structures in classical field theory, Nature \textbf{387}, 58--61 (1997).

\bibitem{n6:lin-yang} F.~Lin and Y.~Yang, Existence of energy minimizers as stable knotted solitons in the Faddeev model, Commun.\ Math.\ Phys.\ \textbf{249}, 273--303 (2004).

\bibitem{n6:lin-yang2} F.~Lin and Y.~Yang, Energy splitting, substantial inequality, and minimization for the Faddeev and Skyrme models, Commun.\ Math.\ Phys.\ \textbf{269}, 137--152 (2007).

\bibitem{n6:toda} H.~Toda, \emph{Composition Methods in Homotopy Groups of Spheres}, Annals of Mathematics Studies 49 (Princeton University Press, 1962).

\bibitem{n6:giulini} D.~Giulini, On the possibility of spinorial quantization in the Skyrme model, Mod.\ Phys.\ Lett.\ A \textbf{8}, 1917--1924 (1993).

\bibitem{n6:fr-original} D.~Finkelstein and J.~Rubinstein, Connection between spin, statistics, and kinks, J.\ Math.\ Phys.\ \textbf{9}, 1762--1779 (1968).

\bibitem{n6:krusch} S.~Krusch, Homotopy of rational maps and the quantization of Skyrmions, Ann.\ Phys.\ \textbf{304}, 103--127 (2003).

\bibitem{n6:witten-current} E.~Witten, Global aspects of current algebra, Nucl.\ Phys.\ B \textbf{223}, 422--432 (1983).

\bibitem{n6:witten-baryons} E.~Witten, Current algebra, baryons, and quark confinement, Nucl.\ Phys.\ B \textbf{223}, 433--444 (1983).

\bibitem{n6:axler} S.~Axler, P.~Bourdon, and W.~Ramey, \emph{Harmonic Function Theory}, 2nd ed.\ (Springer, 2001).

\bibitem{Manton87} N.~S.~Manton, Geometry of skyrmions, Commun.\ Math.\ Phys.\ \textbf{111}, 469--478 (1987).

\bibitem{n1-haar} A.~Haar, Der Massbegriff in der Theorie der kontinuierlichen Gruppen, Ann.\ of Math. \textbf{34}, 147--169 (1933).

\bibitem{n1-pw} D.~N.~Page and W.~K.~Wootters, Evolution without evolution: Dynamics described by stationary observables, Phys.\ Rev.\ D \textbf{27}, 2885--2892 (1983).

\bibitem{n1-wootters} W.~K.~Wootters, Time replaced by quantum correlations, Int.\ J.\ Theor.\ Phys. \textbf{23}, 701--711 (1984).

\bibitem{n1-glm} V.~Giovannetti, S.~Lloyd, and L.~Maccone, Quantum time, Phys.\ Rev.\ D \textbf{92}, 045033 (2015).

\bibitem{n1-hoehn} P.~A.~H\"ohn, A.~R.~H.~Smith, and M.~P.~E.~Lock, Trinity of relational quantum dynamics, Phys.\ Rev.\ D \textbf{104}, 066001 (2021), arXiv:1912.00033.

\bibitem{n1-smith} A.~R.~H.~Smith and M.~Ahmadi, Quantizing time: Interacting clocks and systems, Quantum \textbf{3}, 160 (2019).

\bibitem{n1-page} D.~N.~Page, Average entropy of a subsystem, Phys.\ Rev.\ Lett. \textbf{71}, 1291--1294 (1993).

\bibitem{n1-rt} S.~Ryu and T.~Takayanagi, Holographic derivation of entanglement entropy from AdS/CFT, Phys.\ Rev.\ Lett. \textbf{96}, 181602 (2006).

\bibitem{n1-takesaki} M.~Takesaki, \emph{Tomita's Theory of Modular Hilbert Algebras and Its Applications}, Lecture Notes in Mathematics \textbf{128} (Springer, Berlin, 1970).

\bibitem{n1-hhw} R.~Haag, N.~M.~Hugenholtz, and M.~Winnink, On the equilibrium states in quantum statistical mechanics, Commun.\ Math.\ Phys. \textbf{5}, 215--236 (1967).

\bibitem{n1-connes} A.~Connes and C.~Rovelli, Von Neumann algebra automorphisms and time-thermodynamics relation in generally covariant quantum theories, Class.\ Quantum Grav. \textbf{11}, 2899--2918 (1994).

\bibitem{n1-bw} J.~J.~Bisognano and E.~H.~Wichmann, On the duality condition for a Hermitian scalar field, J.\ Math.\ Phys. \textbf{16}, 985--1007 (1975).

\bibitem{n1-unruh} W.~G.~Unruh, Notes on black-hole evaporation, Phys.\ Rev.\ D \textbf{14}, 870--892 (1976).

\bibitem{n1-hawking} S.~W.~Hawking, Particle creation by black holes, Commun.\ Math.\ Phys. \textbf{43}, 199--220 (1975).

\bibitem{n1-blanco} D.~D.~Blanco, H.~Casini, L.-Y.~Hung, and R.~C.~Myers, Relative entropy and holography, J.\ High Energy Phys. \textbf{08}, 060 (2013).

\bibitem{n1-landauer} R.~Landauer, Irreversibility and heat generation in the computing process, IBM J.\ Res.\ Dev. \textbf{5}, 183--191 (1961).

\bibitem{n1-reeb} D.~Reeb and M.~M.~Wolf, An improved Landauer principle with finite-size corrections, New J.\ Phys. \textbf{16}, 103011 (2014).

\bibitem{n1-jacobson} T.~Jacobson, Thermodynamics of spacetime: The Einstein equation of state, Phys.\ Rev.\ Lett. \textbf{75}, 1260--1263 (1995).

\bibitem{n1-mo} Ll.~Masanes and J.~Oppenheim, A general derivation and quantification of the third law of thermodynamics, Nat.\ Commun. \textbf{8}, 14538 (2017).

\bibitem{n1-eth} M.~Srednicki, Chaos and quantum thermalization, Phys.\ Rev.\ E \textbf{50}, 888--901 (1994).

\bibitem{n7:ACP} K.~Agashe, R.~Contino, and A.~Pomarol, The minimal composite Higgs model, Nucl.\ Phys.\ B \textbf{719}, 165--187 (2005), arXiv:hep-ph/0412089.

\bibitem{n7:CPW} R.~Contino, L.~Da Rold, and A.~Pomarol, Light custodians in natural composite Higgs models, Phys.\ Rev.\ D \textbf{75}, 055014 (2007), arXiv:hep-ph/0612048.

\bibitem{n7:CGG} M.~Carena, L.~Da Rold, and E.~Pont\'on, Minimal composite Higgs models at the LHC, JHEP \textbf{06}, 159 (2014), arXiv:1402.2987.

\bibitem{n7:PW} G.~Panico and A.~Wulzer, The Composite Nambu--Goldstone Higgs, Lect.\ Notes Phys.\ \textbf{913}, 1--316 (2016), arXiv:1506.01961.

\bibitem{n7:SILH} G.~F.~Giudice, C.~Grojean, A.~Pomarol, and R.~Rattazzi, The strongly-interacting light Higgs, JHEP \textbf{06}, 045 (2007), arXiv:hep-ph/0703164.

\bibitem{n7:CW} S.~Coleman and E.~Weinberg, Radiative corrections as the origin of spontaneous symmetry breaking, Phys.\ Rev.\ D \textbf{7}, 1888--1910 (1973).

\bibitem{n7:PT} M.~E.~Peskin and T.~Takeuchi, Estimation of oblique electroweak corrections, Phys.\ Rev.\ D \textbf{46}, 381--409 (1992).

\bibitem{n7:BBRV} R.~Barbieri, B.~Bellazzini, V.~S.~Rychkov, and A.~Varagnolo, The Higgs boson from an extended symmetry, Phys.\ Rev.\ D \textbf{76}, 115008 (2007), arXiv:0706.0432.

\bibitem{n7:Connes} A.~Connes, Noncommutative geometry and reality, J.\ Math.\ Phys.\ \textbf{36}, 6194--6231 (1995).

\bibitem{n7:CC} A.~H.~Chamseddine and A.~Connes, The spectral action principle, Commun.\ Math.\ Phys.\ \textbf{186}, 731--750 (1997), arXiv:hep-th/9606001.

\bibitem{n7:CC12} A.~H.~Chamseddine and A.~Connes, Resilience of the spectral Standard Model, JHEP \textbf{09}, 104 (2012), arXiv:1208.1030.

\bibitem{n3:Gisin} N.~Gisin, Stochastic quantum dynamics and relativity, Helv.\ Phys.\ Acta \textbf{62}, 363--371 (1989).

\bibitem{n3:Simon} C.~Simon, V.~Bu\v{z}ek, and N.~Gisin, No-signaling condition and quantum dynamics, Phys.\ Rev.\ Lett.\ \textbf{87}, 170405 (2001).

\bibitem{n3:GisinPercival} N.~Gisin and I.~C.~Percival, The quantum-state diffusion model applied to open systems, J.\ Phys.\ A \textbf{25}, 5677--5691 (1992).

\bibitem{n3:Diosi} L.~Di\'osi, Continuous quantum measurement and It\^o formalism, Phys.\ Lett.\ A \textbf{129}, 419--423 (1988).

\bibitem{n3:Jacobs} K.~Jacobs and D.~A.~Steck, A straightforward introduction to continuous quantum measurement, Contemp.\ Phys.\ \textbf{47}, 279--303 (2006).

\bibitem{n3:Adler} S.~L.~Adler, D.~C.~Brody, T.~A.~Brun, and L.~P.~Hughston, Martingale models for quantum state reduction, J.\ Phys.\ A \textbf{34}, 8795--8820 (2001).

\bibitem{n3:BassiGhirardi} A.~Bassi and G.~C.~Ghirardi, Dynamical reduction models, Phys.\ Rep.\ \textbf{379}, 257--426 (2003).

\bibitem{n3:Bassi} A.~Bassi, K.~Lochan, S.~Satin, T.~P.~Singh, and H.~Ulbricht, Models of wave-function collapse, underlying theories, and experimental tests, Rev.\ Mod.\ Phys.\ \textbf{85}, 471--527 (2013).

\bibitem{n3:Schlosshauer} M.~Schlosshauer, Decoherence, the measurement problem, and interpretations of quantum mechanics, Rev.\ Mod.\ Phys.\ \textbf{76}, 1267--1305 (2005).

\bibitem{n3:Archimedes} Archimedes, On the sphere and cylinder, Book I, in \emph{The Works of Archimedes}, edited by T.~L.~Heath (Cambridge University Press, 1897).

\bibitem{n3:Aerts} D.~Aerts and M.~Sassoli de Bianchi, The extended Bloch representation of quantum mechanics and the hidden-measurement solution to the measurement problem, Ann.\ Phys.\ \textbf{351}, 975--1025 (2014).

\bibitem{n3:Zyczkowski} K.~\.Zyczkowski and H.-J.~Sommers, Induced measures in the space of mixed quantum states, J.\ Phys.\ A \textbf{34}, 7111--7125 (2001).

\bibitem{n3:vonMises} R.~von Mises and H.~Pollaczek-Geiringer, Praktische Verfahren der Gleichungsaufl\"osung, Z.\ Angew.\ Math.\ Mech.\ \textbf{9}, 58--77, 152--164 (1929).

\bibitem{n3:EckartYoung} C.~Eckart and G.~Young, The approximation of one matrix by another of lower rank, Psychometrika \textbf{1}, 211--218 (1936).

\bibitem{n3:Mirsky} L.~Mirsky, Symmetric gauge functions and unitarily invariant norms, Q.\ J.\ Math.\ \textbf{11}, 50--59 (1960).

\bibitem{n3:Zeno} B.~Misra and E.~C.~G.~Sudarshan, The Zeno's paradox in quantum theory, J.\ Math.\ Phys.\ \textbf{18}, 756--763 (1977).

\bibitem{n3:Attention} A.~Vaswani et al., Attention is all you need, Adv.\ Neural Inf.\ Process.\ Syst.\ \textbf{30}, 5998--6008 (2017).

\bibitem{n3:Hoel} E.~P.~Hoel, L.~Albantakis, and G.~Tononi, Quantifying causal emergence shows that macro can beat micro, Proc.\ Natl.\ Acad.\ Sci.\ U.S.A.\ \textbf{110}, 19790--19795 (2013).

\bibitem{n3:FR} D.~Frauchiger and R.~Renner, Quantum theory cannot consistently describe the use of itself, Nat.\ Commun.\ \textbf{9}, 3711 (2018).

\bibitem{n3:Bong} K.-W.~Bong et al., A strong no-go theorem on the Wigner's friend paradox, Nat.\ Phys.\ \textbf{16}, 1199--1205 (2020).

\bibitem{n3:Williams} D.~Williams, \emph{Probability with Martingales} (Cambridge University Press, 1991).

\bibitem{n2-schur} I.~Schur, \"Uber Potenzreihen, die im Innern des Einheitskreises beschr\"ankt sind, J.\ Reine Angew.\ Math. \textbf{147}, 205--232 (1917).

\bibitem{n2-albert} A.~Albert, Conditions for positive and nonnegative definiteness in terms of pseudoinverses, SIAM J.\ Appl.\ Math. \textbf{17}, 434--440 (1969).

\bibitem{n2-dirac} G.~A.~Dirac, On rigid circuit graphs, Abh.\ Math.\ Sem.\ Univ.\ Hamburg \textbf{25}, 71--76 (1961).

\bibitem{n2-grone} R.~Grone, C.~R.~Johnson, E.~M.~S\'a and H.~Wolkowicz, Positive definite completions of partial Hermitian matrices, Linear Algebra Appl.\ \textbf{58}, 109--124 (1984).

\bibitem{n2-lindblad} G.~Lindblad, Completely positive maps and entropy inequalities, Commun.\ Math.\ Phys. \textbf{40}, 147--151 (1975).

\bibitem{n2-uhlmann-dpi} A.~Uhlmann, Relative entropy and the Wigner--Yanase--Dyson--Lieb concavity in an interpolation theory, Commun.\ Math.\ Phys. \textbf{54}, 21--32 (1977).

\bibitem{n2-bennett} C.~H.~Bennett, The thermodynamics of computation---a review, Int.\ J.\ Theor.\ Phys. \textbf{21}, 905--940 (1982).

\bibitem{n2-scully} M.~O.~Scully and K.~Dr\"uhl, Quantum eraser: A proposed photon correlation experiment concerning observation and delayed choice in quantum mechanics, Phys.\ Rev.\ A \textbf{25}, 2208--2213 (1982).

\bibitem{n2-herzog} T.~J.~Herzog, P.~G.~Kwiat, H.~Weinfurter and A.~Zeilinger, Complementarity and the quantum eraser, Phys.\ Rev.\ Lett. \textbf{75}, 3034--3037 (1995).

\bibitem{n2-kim} Y.-H.~Kim, R.~Yu, S.~P.~Kulik, Y.~H.~Shih and M.~O.~Scully, Delayed choice quantum eraser, Phys.\ Rev.\ Lett. \textbf{84}, 1--5 (2000).

\bibitem{n2-chsh} J.~F.~Clauser, M.~A.~Horne, A.~Shimony and R.~A.~Holt, Proposed experiment to test local hidden-variable theories, Phys.\ Rev.\ Lett. \textbf{23}, 880--884 (1969).

\bibitem{n2-brukner} \v{C}.~Brukner, A no-go theorem for observer-independent facts, Entropy \textbf{20}, 350 (2018).

\bibitem{n2-zurek81} W.~H.~Zurek, Pointer basis of quantum apparatus: Into what mixture does the wave packet collapse?, Phys.\ Rev.\ D \textbf{24}, 1516--1525 (1981).

\bibitem{n2-zurek09} W.~H.~Zurek, Quantum Darwinism, Nature Phys.\ \textbf{5}, 181--188 (2009).

\bibitem{n2-blumekohout} R.~Blume-Kohout and W.~H.~Zurek, A simple example of quantum Darwinism: Redundant information storage in many-spin environments, Found.\ Phys. \textbf{35}, 1857--1876 (2005).

\bibitem{n2-bph} F.~G.~S.~L.~Brand\~ao, M.~Piani and P.~Horodecki, Generic emergence of classical features in quantum Darwinism, Nat.\ Commun. \textbf{6}, 7908 (2015), arXiv:1311.0258.

\bibitem{n2-horodecki} R.~Horodecki, J.~K. Korbicz and P.~Horodecki, Quantum origins of objectivity, Phys.\ Rev.\ A \textbf{91}, 032122 (2015).

\bibitem{n2-bjl} R.~Bhatia, T.~Jain and Y.~Lim, On the Bures--Wasserstein distance between positive definite matrices, Expo.\ Math. \textbf{37}, 165--191 (2019).

\bibitem{n2-pearl} J.~Pearl, Causal inference in statistics: An overview, Statist.\ Surv. \textbf{3}, 96--146 (2009).

\bibitem{n2-friston} K.~Friston, The free-energy principle: A unified brain theory?, Nat.\ Rev.\ Neurosci. \textbf{11}, 127--138 (2010).

\bibitem{n2-pancharatnam} S.~Pancharatnam, Generalized theory of interference, and its applications. Part I. Coherent pencils, Proc.\ Indian Acad.\ Sci. A \textbf{44}, 247--262 (1956).

\bibitem{n2-berry} M.~V.~Berry, Quantal phase factors accompanying adiabatic changes, Proc.\ R.\ Soc.\ Lond. A \textbf{392}, 45--57 (1984).

\bibitem{n2-simon} B.~Simon, Holonomy, the quantum adiabatic theorem, and Berry's phase, Phys.\ Rev.\ Lett. \textbf{51}, 2167--2170 (1983).

\bibitem{n2-sb} J.~Samuel and R.~Bhandari, General setting for Berry's phase, Phys.\ Rev.\ Lett. \textbf{60}, 2339--2342 (1988).

\bibitem{n2-docarmo} M.~P.~do Carmo, Differential Geometry of Curves and Surfaces, Prentice-Hall (1976), Chapter~4, \S~5.

\bibitem{n2-peres} A.~Peres, Incompatible results of quantum measurements, Phys.\ Lett.\ A \textbf{151}, 107--108 (1990).

\bibitem{n2-mermin} N.~D.~Mermin, Simple unified form for the major no-hidden-variables theorems, Phys.\ Rev.\ Lett. \textbf{65}, 3373--3376 (1990).

\bibitem{CLT74} J.~M.~Cornwall, D.~N.~Levin, and G.~Tiktopoulos, Derivation of gauge invariance from high-energy unitarity bounds on the S matrix, Phys.\ Rev.\ D \textbf{10}, 1145 (1974); Erratum \textbf{11}, 972 (1975).

\bibitem{LQT77} B.~W.~Lee, C.~Quigg, and H.~B.~Thacker, Weak interactions at very high energies: the role of the Higgs-boson mass, Phys.\ Rev.\ D \textbf{16}, 1519--1531 (1977).

\bibitem{Malament77} D.~B.~Malament, The class of continuous timelike curves determines the topology of spacetime, J.\ Math.\ Phys.\ \textbf{18}, 1399--1404 (1977).

\bibitem{Kibble79} T.~W.~B.~Kibble, Geometrization of quantum mechanics, Commun.\ Math.\ Phys.\ \textbf{65}, 189--201 (1979).

\bibitem{KS67} S.~Kochen and E.~P.~Specker, The problem of hidden variables in quantum mechanics, J.\ Math.\ Mech.\ \textbf{17}, 59--87 (1967).

\bibitem{Levichev87} A.~V.~Levichev, The causal structure of a Lorentzian manifold determines its conformal geometry, Dokl.\ Akad.\ Nauk SSSR \textbf{293}, 1301--1305 (1987); English translation: Soviet Math.\ Dokl.\ \textbf{35}, 452--455 (1987). % translation volume/pages from a secondary source

\bibitem{Lovelock} D.~Lovelock, The Einstein tensor and its generalizations, J.\ Math.\ Phys.\ \textbf{12}, 498--501 (1971).

\bibitem{Matolcsi75} T.~Matolcsi, Tensor product of Hilbert lattices and free orthodistributive product of orthomodular lattices, Acta Sci.\ Math.\ (Szeged) \textbf{37}, 263--272 (1975).

\bibitem{Mielnik68} B.~Mielnik, Geometry of quantum states, Commun.\ Math.\ Phys.\ \textbf{9}, 55--80 (1968).

\bibitem{MM} Ll.~Masanes and M.~P.~M\"uller, A derivation of quantum theory from physical requirements, New J.\ Phys.\ \textbf{13}, 063001 (2011).

\bibitem{MMAP} Ll.~Masanes, M.~P.~M\"uller, R.~Augusiak, and D.~P\'erez-Garc\'{\i}a, Existence of an information unit as a postulate of quantum theory, Proc.\ Natl.\ Acad.\ Sci.\ USA \textbf{110}, 16373--16377 (2013).

\bibitem{MMPA14} Ll.~Masanes, M.~P.~M\"uller, D.~P\'erez-Garc\'{\i}a, and R.~Augusiak, Entanglement and the three-dimensionality of the Bloch ball, J.\ Math.\ Phys.\ \textbf{55}, 122203 (2014).

\bibitem{MD} R.~Mosseri and R.~Dandoloff, Geometry of entangled states, Bloch spheres and Hopf fibrations, J.\ Phys.\ A \textbf{34}, 10243--10252 (2001).

\bibitem{MohapatraPati} R.~N.~Mohapatra and J.~C.~Pati, ``Natural'' left-right symmetry, Phys.\ Rev.\ D \textbf{11}, 2558--2561 (1975).

\bibitem{MuM} M.~P.~M\"uller and Ll.~Masanes, Three-dimensionality of space and the quantum bit: an information-theoretic approach, New J.\ Phys.\ \textbf{15}, 053040 (2013).

\bibitem{ONeill} B.~O'Neill, \emph{Semi-Riemannian Geometry with Applications to Relativity} (Academic Press, 1983).

\bibitem{Prufer24} H.~Pr\"ufer, Theorie der Abelschen Gruppen~I.\ Grundeigenschaften, Math.\ Z.\ \textbf{20}, 165--187 (1924).

\bibitem{PatiSalam} J.~C.~Pati and A.~Salam, Lepton number as the fourth ``color'', Phys.\ Rev.\ D \textbf{10}, 275--289 (1974).

\bibitem{PDG} S.~Navas \emph{et al.}\ (Particle Data Group), Review of particle physics, Phys.\ Rev.\ D \textbf{110}, 030001 (2024).

\bibitem{Renou} M.-O.~Renou, D.~Trillo, M.~Weilenmann, T.~P.~Le, A.~Tavakoli, N.~Gisin, A.~Ac\'{\i}n, and M.~Navascu\'es, Quantum theory based on real numbers can be experimentally falsified, Nature \textbf{600}, 625--629 (2021).

\bibitem{Robertson} H.~P.~Robertson, The uncertainty principle, Phys.\ Rev.\ \textbf{34}, 163--164 (1929).

\bibitem{Schrodinger30} E.~Schr\"odinger, Zum Heisenbergschen Unsch\"arfeprinzip, Sitzungsber.\ Preuss.\ Akad.\ Wiss., Phys.-Math.\ Kl.\ \textbf{19}, 296--303 (1930).

\bibitem{SenjanovicMohapatra} G.~Senjanovi\'c and R.~N.~Mohapatra, Exact left-right symmetry and spontaneous violation of parity, Phys.\ Rev.\ D \textbf{12}, 1502--1505 (1975).

\bibitem{SmithHarmonic} R.~T.~Smith, The second variation formula for harmonic mappings, Proc.\ Amer.\ Math.\ Soc.\ \textbf{47}, 229--236 (1975).

\bibitem{SSVZ} P.~Sikivie, L.~Susskind, M.~Voloshin, and V.~Zakharov, Isospin breaking in technicolor models, Nucl.\ Phys.\ B \textbf{173}, 189--207 (1980).

\bibitem{Stetkaer} H.~Stetk\ae r, \emph{Functional Equations on Groups} (World Scientific, 2013).

\bibitem{Stueckelberg} E.~C.~G.~Stueckelberg, Quantum theory in real Hilbert space, Helv.\ Phys.\ Acta \textbf{33}, 727--752 (1960).

\bibitem{VDD} F.~Verstraete, J.~Dehaene, and B.~De~Moor, Local filtering operations on two qubits, Phys.\ Rev.\ A \textbf{64}, 010101(R) (2001).

\bibitem{Veltman77} M.~Veltman, Limit on mass differences in the Weinberg model, Nucl.\ Phys.\ B \textbf{123}, 89--99 (1977).

\bibitem{vW} C.~F.~von~Weizs\"acker, \emph{Aufbau der Physik} (Hanser, Munich, 1985); English edition \emph{The Structure of Physics}, edited by T.~G\"ornitz and H.~Lyre (Springer, 2006).

\bibitem{VOS} A.~van~Oosterom and J.~Strackee, The solid angle of a plane triangle, IEEE Trans.\ Biomed.\ Eng.\ \textbf{BME-30}, 125--126 (1983).

\bibitem{Weinberg67} S.~Weinberg, A model of leptons, Phys.\ Rev.\ Lett.\ \textbf{19}, 1264--1266 (1967).

\bibitem{Wigner32} E.~Wigner, \"Uber die Operation der Zeitumkehr in der Quantenmechanik, Nachr.\ Ges.\ Wiss.\ G\"ottingen, Math.-Phys.\ Kl.\ 546--559 (1932). % UNVERIFIED: year and pages confirmed only from secondary sources.

\bibitem{Wootters12} W.~K.~Wootters, Entanglement sharing in real-vector-space quantum theory, Found.\ Phys.\ \textbf{42}, 19--28 (2012).

\bibitem{Wootters} W.~K.~Wootters, Entanglement of formation of an arbitrary state of two qubits, Phys.\ Rev.\ Lett.\ \textbf{80}, 2245--2248 (1998).

\bibitem{Xin80} Y.~L.~Xin, Some results on stable harmonic maps, Duke Math.\ J.\ \textbf{47}, 609--613 (1980).

\bibitem{Yamabe} H.~Yamabe, On an arcwise connected subgroup of a Lie group, Osaka Math.\ J.\ \textbf{2}, 13--14 (1950).

\bibitem{Zeeman} E.~C.~Zeeman, Causality implies the Lorentz group, J.\ Math.\ Phys.\ \textbf{5}, 490--493 (1964).


\bibitem{ex-Rauch75} H.~Rauch, A.~Zeilinger, G.~Badurek, A.~Wilfing, W.~Bauspiess, and U.~Bonse, Verification of coherent spinor rotation of fermions, Phys.\ Lett.\ A \textbf{54}, 425--427 (1975).

\bibitem{ex-WCOE75} S.~A.~Werner, R.~Colella, A.~W.~Overhauser, and C.~F.~Eagen, Observation of the phase shift of a neutron due to precession in a magnetic field, Phys.\ Rev.\ Lett.\ \textbf{35}, 1053--1055 (1975).

\bibitem{ex-Rauch78} H.~Rauch, A.~Wilfing, W.~Bauspiess, and U.~Bonse, Precise determination of the $4\pi$-periodicity factor of a spinor wave function, Z.\ Phys.\ B \textbf{29}, 281--284 (1978).

\bibitem{ex-Rauch88} H.~Rauch, Quantum measurements in time dependent neutron interferometry, Ann.\ Inst.\ Henri Poincar\'e A \textbf{49}, 355--368 (1988).

\bibitem{ex-AS67} Y.~Aharonov and L.~Susskind, Observability of the sign change of spinors under $2\pi$ rotations, Phys.\ Rev.\ \textbf{158}, 1237--1238 (1967).

\bibitem{ex-Bernstein67} H.~J.~Bernstein, Spin precession during interferometry of fermions and the phase factor associated with rotations through $2\pi$ radians, Phys.\ Rev.\ Lett.\ \textbf{18}, 1102--1103 (1967).

\bibitem{ex-CEMM98} R.~Cleve, A.~Ekert, C.~Macchiavello, and M.~Mosca, Quantum algorithms revisited, Proc.\ R.\ Soc.\ Lond.\ A \textbf{454}, 339--354 (1998).

\bibitem{ex-Ramsey50} N.~F.~Ramsey, A molecular beam resonance method with separated oscillating fields, Phys.\ Rev.\ \textbf{78}, 695--699 (1950).

\bibitem{ex-BdBBvO85} H.~Barkhuijsen, R.~de Beer, W.~M.~M.~J.~Bov\'ee, and D.~van Ormondt, Retrieval of frequencies, amplitudes, damping factors, and phases from time-domain signals using a linear least-squares procedure, J.\ Magn.\ Reson.\ \textbf{61}, 465--481 (1985).

\bibitem{ex-Prony1795} G.~R.~de Prony, Essai exp\'erimental et analytique sur les lois de la dilatabilit\'e de fluides \'elastiques et sur celles de la force expansive de la vapeur de l'alkool, \`a diff\'erentes temp\'eratures, J.\ \'Ecole Polytech.\ \textbf{1}, 24--76 (1795). % UNVERIFIED (see examples-refs.bib)

\bibitem{ex-CODATA22} P.~J.~Mohr, D.~B.~Newell, B.~N.~Taylor, and E.~Tiesinga, CODATA recommended values of the fundamental physical constants: 2022, Rev.\ Mod.\ Phys.\ \textbf{97}, 025002 (2025).

\bibitem{ex-Favro60} L.~D.~Favro, Theory of the rotational Brownian motion of a free rigid body, Phys.\ Rev.\ \textbf{119}, 53--62 (1960).

\bibitem{ex-Debye29} P.~Debye, \emph{Polar Molecules} (Chemical Catalog Company, New York, 1929).

\bibitem{ex-Pancharatnam56} S.~Pancharatnam, Generalized theory of interference, and its applications, Proc.\ Indian Acad.\ Sci.\ A \textbf{44}, 247--262 (1956).

\bibitem{ex-Berry84} M.~V.~Berry, Quantal phase factors accompanying adiabatic changes, Proc.\ R.\ Soc.\ Lond.\ A \textbf{392}, 45--57 (1984).

\bibitem{ex-BS88} R.~Bhandari and J.~Samuel, Observation of topological phase by use of a laser interferometer, Phys.\ Rev.\ Lett.\ \textbf{60}, 1211--1213 (1988).

\bibitem{ex-SB88} J.~Samuel and R.~Bhandari, General setting for Berry's phase, Phys.\ Rev.\ Lett.\ \textbf{60}, 2339--2342 (1988).

\bibitem{ex-MS93} N.~Mukunda and R.~Simon, Quantum kinematic approach to the geometric phase. I. General formalism, Ann.\ Phys.\ (N.Y.) \textbf{228}, 205--268 (1993).

\bibitem{ex-BD87} T.~Bitter and D.~Dubbers, Manifestation of Berry's topological phase in neutron spin rotation, Phys.\ Rev.\ Lett.\ \textbf{59}, 251--254 (1987).

\bibitem{ex-Wagh97} A.~G.~Wagh, V.~C.~Rakhecha, J.~Summhammer, G.~Badurek, H.~Weinfurter, B.~E.~Allman, H.~Kaiser, K.~Hamacher, D.~L.~Jacobson, and S.~A.~Werner, Experimental separation of geometric and dynamical phases using neutron interferometry, Phys.\ Rev.\ Lett.\ \textbf{78}, 755--759 (1997).

\bibitem{ex-TC86} A.~Tomita and R.~Y.~Chiao, Observation of Berry's topological phase by use of an optical fiber, Phys.\ Rev.\ Lett.\ \textbf{57}, 937--940 (1986).

\bibitem{ex-Kramers30} H.~A.~Kramers, Th\'eorie g\'en\'erale de la rotation paramagn\'etique dans les cristaux, Proc.\ Kon.\ Akad.\ Wetensch.\ Amsterdam \textbf{33}, 959--972 (1930).

\bibitem{ex-AB70} A.~Abragam and B.~Bleaney, \emph{Electron Paramagnetic Resonance of Transition Ions} (Clarendon Press, Oxford, 1970).

\bibitem{ex-Koenig07} M.~K\"onig, S.~Wiedmann, C.~Br\"une, A.~Roth, H.~Buhmann, L.~W.~Molenkamp, X.-L.~Qi, and S.-C.~Zhang, Quantum spin Hall insulator state in HgTe quantum wells, Science \textbf{318}, 766--770 (2007).

\bibitem{ex-LD98} D.~Loss and D.~P.~DiVincenzo, Quantum computation with quantum dots, Phys.\ Rev.\ A \textbf{57}, 120--126 (1998).

\bibitem{ex-Petta05} J.~R.~Petta, A.~C.~Johnson, J.~M.~Taylor, E.~A.~Laird, A.~Yacoby, M.~D.~Lukin, C.~M.~Marcus, M.~P.~Hanson, and A.~C.~Gossard, Coherent manipulation of coupled electron spins in semiconductor quantum dots, Science \textbf{309}, 2180--2184 (2005).

\bibitem{ex-Anderlini07} M.~Anderlini, P.~J.~Lee, B.~L.~Brown, J.~Sebby-Strabley, W.~D.~Phillips, and J.~V.~Porto, Controlled exchange interaction between pairs of neutral atoms in an optical lattice, Nature \textbf{448}, 452--456 (2007).

\bibitem{ex-Trotzky08} S.~Trotzky, P.~Cheinet, S.~F\"olling, M.~Feld, U.~Schnorrberger, A.~M.~Rey, A.~Polkovnikov, E.~A.~Demler, M.~D.~Lukin, and I.~Bloch, Time-resolved observation and control of superexchange interactions with ultracold atoms in optical lattices, Science \textbf{319}, 295--299 (2008).

\bibitem{ex-Chen22} M.-C.~Chen, C.~Wang, F.-M.~Liu, J.-W.~Wang, C.~Ying, Z.-X.~Shang, Y.~Wu, M.~Gong, H.~Deng, F.-T.~Liang, Q.~Zhang, C.-Z.~Peng \emph{et al.}, Ruling out real-valued standard formalism of quantum theory, Phys.\ Rev.\ Lett.\ \textbf{128}, 040403 (2022).

\bibitem{ex-Li22} Z.-D.~Li, Y.-L.~Mao, M.~Weilenmann, A.~Tavakoli, H.~Chen, L.~Feng, S.-J.~Yang, M.-O.~Renou, D.~Trillo, T.~P.~Le, N.~Gisin, A.~Ac\'{\i}n \emph{et al.}, Testing real quantum theory in an optical quantum network, Phys.\ Rev.\ Lett.\ \textbf{128}, 040402 (2022).

\bibitem{ex-Grover97} L.~K.~Grover, Quantum mechanics helps in searching for a needle in a haystack, Phys.\ Rev.\ Lett.\ \textbf{79}, 325--328 (1997).

\bibitem{ex-CGK98} I.~L.~Chuang, N.~Gershenfeld, and M.~Kubinec, Experimental implementation of fast quantum searching, Phys.\ Rev.\ Lett.\ \textbf{80}, 3408--3411 (1998).

\bibitem{ex-Figgatt17} C.~Figgatt, D.~Maslov, K.~A.~Landsman, N.~M.~Linke, S.~Debnath, and C.~Monroe, Complete 3-qubit Grover search on a programmable quantum computer, Nat.\ Commun.\ \textbf{8}, 1918 (2017).

\bibitem{ex-HKN97} D.~Han, Y.~S.~Kim, and M.~E.~Noz, Stokes parameters as a Minkowskian four-vector, Phys.\ Rev.\ E \textbf{56}, 6065--6076 (1997).

\bibitem{ex-LuChipman96} S.-Y.~Lu and R.~A.~Chipman, Interpretation of Mueller matrices based on polar decomposition, J.\ Opt.\ Soc.\ Am.\ A \textbf{13}, 1106 (1996).

\bibitem{ex-GY88} D.~M.~Greenberger and A.~Yasin, Simultaneous wave and particle knowledge in a neutron interferometer, Phys.\ Lett.\ A \textbf{128}, 391--394 (1988).

\bibitem{ex-Englert96} B.-G.~Englert, Fringe visibility and which-way information: an inequality, Phys.\ Rev.\ Lett.\ \textbf{77}, 2154--2157 (1996).

\bibitem{ex-QVE18} X.-F.~Qian, A.~N.~Vamivakas, and J.~H.~Eberly, Entanglement limits duality and vice versa, Optica \textbf{5}, 942--947 (2018).

\bibitem{ex-EQV17} J.~H.~Eberly, X.-F.~Qian, and A.~N.~Vamivakas, Polarization coherence theorem, Optica \textbf{4}, 1113 (2017).

\bibitem{ex-SVV77} M.~E.~Stoll, A.~J.~Vega, and R.~W.~Vaughan, Explicit demonstration of spinor character for a spin-$\tfrac12$ nucleus via NMR interferometry, Phys.\ Rev.\ A \textbf{16}, 1521 (1977).

\bibitem{ex-Sulyok10} G.~Sulyok, Y.~Hasegawa, J.~Klepp, H.~Lemmel, and H.~Rauch, Noise-induced dephasing in neutron interferometry, Phys.\ Rev.\ A \textbf{81}, 053609 (2010).

\bibitem{ex-Abragam61} A.~Abragam, \emph{The Principles of Nuclear Magnetism} (Clarendon Press, Oxford, 1961).

\bibitem{ex-PDG26} F.~Takahashi \emph{et al.} (Particle Data Group), Review of particle physics, Int.\ J.\ Mod.\ Phys.\ A \textbf{41}, 2630011 (2026).

\bibitem{ex-Steinhauser98} M.~Steinhauser, Leptonic contribution to the effective electromagnetic coupling constant up to three loops, Phys.\ Lett.\ B \textbf{429}, 158--161 (1998).

\bibitem{ex-Planck18} N.~Aghanim \emph{et al.} (Planck), Planck 2018 results. VI. Cosmological parameters, Astron.\ Astrophys.\ \textbf{641}, A6 (2020).

\bibitem{ex-Riess98} A.~G.~Riess \emph{et al.}, Observational evidence from supernovae for an accelerating universe and a cosmological constant, Astron.\ J.\ \textbf{116}, 1009--1038 (1998).

\bibitem{ex-Perlmutter99} S.~Perlmutter \emph{et al.}, Measurements of $\Omega$ and $\Lambda$ from 42 high-redshift supernovae, Astrophys.\ J.\ \textbf{517}, 565--586 (1999).

\end{thebibliography}

\end{document}
